 \documentclass[reqno,11pt]{preprint}
%%%%\usepackage{textcomp}
%%%%\usepackage[sb]{libertine}
%%%%\usepackage[varqu,varl]{inconsolata}% sans serif typewriter
%%%%\usepackage[libertine,bigdelims,vvarbb]{newtxmath} % bb from STIX
%%%%\usepackage[cal=boondoxo]{mathalfa} % mathcal
%%%%\useosf % osf for text, not math
%%%%\usepackage[supstfm=libertinesups,%
%%%%	supscaled=1.2,%
%%%%	raised=-.13em]{superiors}
%\usepackage[marginparwidth=1in]{geometry}
\usepackage[full]{textcomp}
\usepackage[osf]{newtxtext}
\usepackage{cabin}
\usepackage{enumerate}
\usepackage{enumitem}
\usepackage{bbm}
\usepackage{dsfont}


\usepackage[varqu,varl]{inconsolata}
\usepackage[cal=boondoxo]{mathalfa}

\usepackage{comment}
\usepackage{hyperref}
\usepackage{breakurl}
\usepackage{mathrsfs}
\usepackage{booktabs}
\usepackage{caption}

\usepackage{tikz}
\usetikzlibrary{decorations.pathreplacing,decorations.markings}
\usepackage{amsmath} 
\usepackage{mhequ} 
\usepackage{mhsymb} 
\usepackage{mhenvs} 
\usepackage{microtype}
\usepackage{amssymb} 
\usepackage{eufrak}

\usepackage{wasysym}
\usepackage{centernot}

%\usepackage{showlabels}



%\usepackage[notref,notcite]{showkeys}

\usepackage{oldgerm}

%\newtheorem{conjecture}{Conjecture}


\newcommand\varfrak[1]{\mathord{\text{\textgoth{#1}}}}

%\newtheorem{notation*}{Notation}

%%%%%%%%%\mathbb

\newcommand{\Exp}{\mathbf{E}}
\newcommand{\Prob}{\mathbf{P}}
\newcommand{\Var}{\mathrm{Var}}
\newcommand{\hffNn}{\hat{\ff}^{\eps,n}}
\newcommand{\fFN}{{\fF}^{N}}
\newcommand{\wN}{w^{\eps,n}}
\newcommand{\vN}{v^{\eps,n}}
\newcommand{\tvN}{\tilde v^{\eps,n}}
\newcommand{\VN}{\mathcal{V}^{\eps,n}}
\newcommand{\WN}{\mathcal{W}^{\eps,n}}
\newcommand{\Weps}{\mathcal{W}^{\eps}}


\renewcommand{\R}{\mathbb{R}}
\renewcommand{\C}{\mathbb{C}}
\renewcommand{\N}{\mathbb{N}}
\renewcommand{\Q}{\mathbb{Q}}
\newcommand{\T}{\mathbb{T}}
\renewcommand{\Z}{\mathbb{Z}}
\renewcommand{\E}{\mathbb{E}}
\renewcommand{\P}{\mathbb{P}}

\newcommand{\J}{\mathbb{J}}


%%%%%%%%%\mathbf

\newcommand{\BX}{\mathbf{X}}
\newcommand{\BM}{\mathbf{M}}
\newcommand{\BP}{\mathbf{P}}

%%%%%%%%%%\mathcal
%
\newcommand{\cA}{\mathcal{A}}
\newcommand{\cB}{\mathcal{B}}
\newcommand{\cC}{\mathcal{C}}
\newcommand{\cD}{\mathcal{D}}
\newcommand{\cF}{\mathcal{F}}
\newcommand{\cG}{\mathcal{G}}
\newcommand{\cH}{\mathcal{H}}
\newcommand{\cI}{\mathcal{I}}
\newcommand{\cJ}{\mathcal{J}}
\newcommand{\cK}{\mathcal{K}}
\newcommand{\cL}{\mathcal{L}}
\newcommand{\cM}{\mathcal{M}}
\newcommand{\cN}{\mathcal{N}}
\newcommand{\cO}{\mathcal{O}}
%\newcommand{\CP}{\mathcal{P}}
%\newcommand{\CQ}{\mathcal{Q}}
%\newcommand{\CR}{\mathcal{R}}
\newcommand{\cS}{\mathcal{S}}
%\newcommand{\CT}{\mathcal{T}}
\newcommand{\cU}{\mathcal{U}}
\newcommand{\cV}{\mathcal{V}}
\newcommand{\cW}{\mathcal{W}}

%%%%%%%%%\mathfrak

% mathfrak \FA, \FB, ...
% mathcal  \CA, \CB, ...

%\newcommand{\GH}{\mathcal{GH}}
%\newcommand{\FK}{\mathfrak{K}}
%\newcommand{\FI}{\mathfrak{I}}

%\newcommand{\sh}{\mathfrak{h}}
%\newcommand{\s}{\mathfrak{s}}
%\newcommand{\sw}{\mathfrak{w}}
%\newcommand{\sz}{\mathfrak{z}}


%%%%%%%%%\mathrm

\newcommand{\RH}{\mathrm{H}}
\newcommand{\RM}{\mathrm{M}}

\newcommand{\rb}{\mathrm{b}}
\newcommand{\rc}{\mathrm{c}}
\newcommand{\rl}{\mathrm{l}}
\newcommand{\rp}{\mathrm{p}}
\newcommand{\rpm}{\mathrm{pm}}
\newcommand{\rr}{\mathrm{r}}


\newcommand{\bc}{\mathrm{bc}}
\newcommand{\br}{\mathrm{b}}
\newcommand{\bw}{\mathrm{bw}}
\newcommand{\com}{\mathrm{c}}
\newcommand{\cl}{\mathrm{cl}}
\newcommand{\dd}{\mathrm{d}}      %%%%%% d
\newcommand{\dis}{\mathrm{dis}}
\newcommand{\eff}{\mathrm{eff}}
\newcommand{\fin}{\mathrm{fin}}
\newcommand{\Id}{\mathrm{Id}}
\newcommand{\loc}{\mathrm{loc}}
\newcommand{\pGH}{\mathrm{pGH}}
\newcommand{\pH}{\mathrm{pH}}
\newcommand{\pM}{\mathrm{pM}}
\newcommand{\pmH}{\mathrm{pmH}}
\newcommand{\Pro}{\mathrm{Pr}}
\newcommand{\Sk}{\mathrm{Sk}}
\newcommand{\Sp}{\mathrm{sp}}
\newcommand{\tot}{\mathrm{tot}}

%%%%%%%%%\mathscr

\newcommand{\SB}{\mathscr{B}}
\newcommand{\SD}{\mathscr{D}}
\newcommand{\SE}{\mathscr{E}}
\newcommand{\SF}{\mathscr{F}}
\newcommand{\SH}{\mathscr{H}}
\newcommand{\SL}{\mathscr{L}}
\newcommand{\proj}{\mathscr{P}}
\newcommand{\SP}{\mathscr{P}}
\newcommand{\SR}{\mathscr{R}}
\renewcommand{\SS}{\mathscr{S}}
\newcommand{\ST}{\mathscr{T}}
\newcommand{\SX}{\mathscr{X}}



%%%%%%%%%\mathsf

%%%%%%%%%%\mathfrak

\renewcommand{\fE}{\mathfrak{E}}
\newcommand{\fFNn}{\mathfrak{F}^{\eps,n}}


\newcommand{\fb}{\mathfrak{b}}
\newcommand{\fc}{\mathcal{c}}
\newcommand{\fd}{\mathcal{d}}

\newcommand{\ff}{\mathfrak{f}}
\newcommand{\ft}{\mathfrak{t}}
\newcommand{\fp}{\mathfrak{p}}
\newcommand{\fg}{\mathfrak{g}}
\newcommand{\fl}{\mathfrak{l}}
\newcommand{\fs}{\mathfrak{s}}
\newcommand{\fh}{\mathfrak{h}}
\newcommand{\fw}{\mathfrak{w}}
\newcommand{\nf}{\mathfrak{n}^\eps}


\newcommand{\ffNn}{\ff^{\eps,n}}
\newcommand{\tffNn}{\tilde{\ff}^{\eps,n}}
\newcommand{\fbNn}{\fb^{\eps,n}}
\newcommand{\tfbNn}{\tilde{\fb}^{\eps,n}}
\newcommand{\fhNn}{\fh^{\eps,n}}
\newcommand{\tfhNn}{\tilde{\fh}^{\eps,n}}
\newcommand{\sym}{\mathrm{sym}}

\newcommand{\cft}{\check{\ft}}
\newcommand{\tft}{\tilde{\ft}}


\newcommand{\bv}{\mathbf{v}}


%\newcommand{\FP}{\mathsf{P}}
%\newcommand{\FR}{\mathsf{R}}

\def\one{\mathrm{(I)}}
\def\two{\mathrm{(II)}}
\def\three{\mathrm{III}}
\def\A{\mathrm{(A)}}
\def\B{\mathrm{(B)}}

\newcommand{\1}{\mathds{1}}

\def\erf{\mathop{\mathrm{erf}}}
\def\restr{\mathord{\upharpoonright}}
\def\c{\mathfrak{c}}

\def\Wick#1{\mathopen{{:}}#1\mathclose{{:}}}
\newcommand{\da}{\downarrow}
\newcommand{\ua}{\uparrow}
\newcommand{\uda}{\uparrow\mathrel{\mspace{-1mu}}\downarrow}
\newcommand{\dotp}{\boldsymbol{\cdot}}
\newcommand{\vphi}{\varphi}
\def\sym{{\mathrm{sym}}}
\newcommand{\Di}{\mathrm{diag}}
\newcommand{\oD}{\mathrm{off}}
\newcommand{\oDi}{\mathrm{off_1}}
\newcommand{\ooDi}{\mathrm{off_2}}
\newcommand{\bulk}{\mathrm{bulk}}
\newcommand{\Dbulk}{\mathrm{D}_\bulk}
\newcommand{\SHE}{\mathrm{SHE}}
\newcommand{\nueff}{\nu_{\rm eff}}
\newcommand{\Mn}{\BM^{n,\eps}}


%%%%%%%%%%%%%%%%%%%%%%%% Notation for us


% General

\def\one{\mathrm{(I)}}
\def\two{\mathrm{(II)}}
\def\three{\mathrm{(III)}}
\def\four{\mathrm{(IV)}}

% Operators, functions, spaces

\newcommand{\Deltahalf}{A}
\newcommand{\half}{\frac{1}{2}}
\def\halff#1{\frac{#1}{2}}
\def\pmo{^{-1}}
\newcommand{\phalf}{^\half}
\def\ssq{\subseteq}

\def\hot{\otimes^n h}
\def\hotp{\otimes^{n+1} h}
\def\hotm{\otimes^{n-1} h}
\def\hotmm{\otimes^{n-2} h}

\newcommand{\thec}{c}			% the function c(.,.)
\def\indN#1{{	\J^\eps_{#1}	}}
\def\tindN#1{{	\tilde\J^\eps_{#1}	}}
\newcommand{\fock}{\Gamma L^2}	% Fock space
\def\fockprod#1#2{\langle #1 , #2 \rangle}
\def\sint{I}					% Stochastic integral
\def\wc{\SH}					% Wiener Chaos
\def\FLS#1{{\CF L^{#1}}}			% Space in which we show tightness
\def\CFLS#1{C_T\FLS{#1}}		% Space of continuous funtions with values in $\CF$
\def\CCFLS#1#2{{ C_T^{#1} \FLS{#2}}}	% Space of Hoelder cont. functions with values in $\CF$
\newcommand{\Log}{\mathrm{Log}}


\def\embed{\hookrightarrow}

\def\qvar#1{\langle #1 \rangle}	% Quadratic variation

% Noises

\def\wn{{	\xi	}}
\def\swn{{	\eta	}}
\def\filt{{ \CG	}}		%Filtration generated by $(\xi_t)_{t\ge 0}$.

% Generators

\newcommand{\gen}{\cL^\eps}
\def\genF#1#2{{	\cL^{\eps,#1}_{#2}	}}
\newcommand{\genn}{\cL^\eps_{2,n}}
\newcommand{\geneff}{\cL^\eff}
\newcommand{\gens}{\cL_0^{N}}
\newcommand{\gensy}{\cL_0}
\newcommand{\gensyt}{\widetilde{\cL_0}}
\newcommand{\gena}{\cA^{\eps}}
\def\genaF#1#2{{	\cA^{\eps,#1}_{#2}	}}
\newcommand{\genap}{\cA^{\eps}_+}
\newcommand{\genam}{\cA^{\eps}_-}
\newcommand{\genan}{\cA^{\eps}_{2,n}}
\newcommand{\genapn}{\cA^{\eps}_{+,2,n}}
\newcommand{\genamn}{\cA^{\eps}_{-,2,n}}
\newcommand{\cGN}{\cG^\eps}
\newcommand{\gensf}{\hat\cL_0^\eps}
\newcommand{\genaf}{\hat\cA^\eps}
\newcommand{\gensyf}{\mathfrak{L}_0}%Action in Fock space
\newcommand{\genapf}{\mathfrak{A}^\eps_+} %Action in Fock space
\newcommand{\genamf}{\mathfrak{A}^\eps_-} %Action in Fock space
\newcommand{\nonlin}{\cK^{N}}
\newcommand{\Nonlin}{\cN^{N}}
\newcommand{\inonlin}{\cB^\eps}
\newcommand{\inonlintilde}{\tilde\cB^\eps}
\newcommand{\nonlintilde}{\tilde\cK^\eps}
\newcommand{\nonlinint}{ \CF^\eps }
\newcommand{\nonlinlim}{ \CF }
\newcommand{\nonlinfock}{ \vphi^\eps }
\newcommand{\nonlinres}{ \psi^\eps }
\newcommand{\nonlinrestrunc}{ \psi^{N,M} }
\def\nonlinrestruncI#1{ \psi^{N,#1} }
\def\restrunc#1{ T^{N,#1} }
\newcommand{\Op}{\mathcal{H}^\eps}
\newcommand{\OpJ}{\mathcal{J}^\eps}

\newcommand{\NN}{\cN^\eps}

\newcommand{\bG}{\mathbf{G}}

\newcommand{\genapc}{A_+^\eps}

% Renormalization constants

\newcommand{\const}{\lambda_N}
\newcommand{\tconst}{\tilde\lambda_N}
\newcommand{\donst}{\nu_N}
\def\donsth{\donst\phalf}

% Notation for tightness

\def\energy{\CE^\eps}
\def\martf{M^\eps}			% Forward martingale
\def\martb{\bar M^\eps}	% Backward martingale
\def\pois{\Op}			% Solution to Poisson equation
\def\poisl{K^\eps}			% Solution linear Poisson
\newcommand{\martnl}{\cM^\eps}  		%martingale associated to the nonlinearity

\def\she{X^\eps} 	%solution of SHE with constant


\newcommand{\Ll}{\mathrm{L}^\eps}
\newcommand{\LB}{\mathrm{LB}}
\newcommand{\UB}{\mathrm{UB}}
\newcommand{\LlN}{\mathrm{L}^\eps}
\newcommand{\LBN}{\mathrm{LB}^\eps}
\newcommand{\UBN}{\mathrm{UB}^\eps}

\colorlet{darkblue}{blue!90!black}
\colorlet{darkred}{red!90!black}
\colorlet{darkgreen}{green!50!black}
\colorlet{darkyellow}{yellow!90!black}
\def\giuseppe#1{\comment[darkred]{GC: #1}}
\def\fabio#1{\comment[darkgreen]{FT: #1}}
\def\max#1{\comment[darkblue]{MG: #1}}


\long\def\giuseppeText#1{{\color{darkred}Giuseppe:\ #1}}
\long\def\fabioText#1{{\color{darkgreen}Fabio:\ #1}}
\long\def\maxText#1{{\color{darkblue}Max:\ #1}}



\title{Gaussian Fluctuations for the stochastic Burgers equation in dimension $d\geq 2$}


\begin{document}

\maketitle

\vspace{-2cm}

\noindent{ \bf Giuseppe Cannizzaro$^1$,  Massimiliano Gubinelli$^2$, Fabio Toninelli$^3$}
\newline

\noindent{\small $^1$University of Warwick, UK, %
    $^2$University of Oxford, UK, %
    $^3$Technical University of Vienna, Austria\\}

\noindent{\small\email{giuseppe.cannizzaro@warwick.ac.uk, 
gubinelli@maths.ox.ac.uk, \\
fabio.toninelli@tuwien.ac.at}}
\newline

\begin{abstract}
The goal of the present paper is to %(start to) fill in this gap and 
establish a framework which allows to rigorously determine the large-scale Gaussian  
fluctuations for a class of singular SPDEs at and above criticality, and therefore beyond 
the range of applicability of pathwise techniques, such as the theory of Regularity Structures. 
To this purpose, we focus on a $d$-dimensional generalization of the Stochastic Burgers equation (SBE) 
introduced in [H. van Beijeren, R. Kutner and H. Spohn, Excess noise for driven diffusive systems, PRL, 1985]. 
In both the critical $d=2$ and super-critical $d\geq 3$ cases, we show that the scaling limit of 
(the regularised) SBE is given by a stochastic heat equation with non-trivially renormalised coefficient, 
introducing a set of tools that we expect to be applicable more widely. 
For $d\ge3$  the scaling adopted is the classical diffusive one, while in $d=2$ it is the 
so-called {\it weak coupling} scaling which corresponds to tuning down the strength of the interaction 
in a scale-dependent way. 
%
%
%the super-critical case $d\ge 3$, we prove that the diffusively rescaled 
%
%
%obtain a scaling limit, in the form of a linear SPDE with non-trivially renormalized coefficients, when the regularization cut-off tends to zero. We obtain the analogous result in the critical dimension $d=2$ in the weak coupling regime, where the interaction strength tends suitably to zero together with the cut-off.
%{\color{red} }
\end{abstract}

\bigskip\noindent
{\it Key words and phrases.}
Stochastic Partial Differential Equations, critical and super-critical dimension, KPZ equation, weak coupling scaling, 
diffusion coefficient,  asymmetric simple exclusion 
\bigskip

\setcounter{tocdepth}{2} 
\tableofcontents

\section{Introduction}

The study of (singular) Stochastic Partial Differential Equations (SPDEs) has known tremendous advances 
in recent years. Thanks to the groundbreaking theory of regularity structures~\cite{Hai}, 
paracontrolled calculus~\cite{Para} or renormalization group techniques~\cite{Kup}, a local 
solution theory can be shown to hold for a large family of SPDEs in the so-called {\it sub-critical} 
dimensions (see~\cite{Hai} for the notion of {\it sub-criticality} in this context). 
In contrast, SPDEs in the {\it critical} and {\it super-critical} dimensions
%i.e. those at which the previous theories break down, 
are, from a mathematical viewpoint, still 
largely unexplored and poorly understood. 
The aim of the present paper is to %(start to) fill in this gap and 
establish a framework which allows to rigorously determine the large-scale (Gaussian) 
fluctuations for a class of SPDEs at and above criticality under (suitable) diffusive scaling. 

To present the ideas and the techniques in our analysis, 
we will focus on a specific example, the Stochastic Burgers Equation (SBE), which, in any dimension, 
can (formally) be written as 
\begin{equ}[e:BKS]
    \partial_t\eta = \tfrac{1}{2}\Delta\eta+\fw\cdot\nabla\eta^2 + \div \boldsymbol{\xi}
 \end{equ}
where the scalar field $\eta=\eta(x,t)$ depends on $t\ge 0$ and $x\in \R^d$ with $d\ge1$, 
$\fw\in \R^d$ is a fixed vector and $\boldsymbol{\xi}=(\xi_1,\dots,\xi_d)$ is a $d$-dimensional space-time white noise on 
$\R_+\times \R^d$, 
i.e. a centred $d$-dimensional Gaussian process whose covariance is formally given by 
$\Exp[\xi_i(t,x)\xi_j(s,y)]=\delta(t-s)\delta(x-y)\1_{i=j}$. 

The above equation was introduced by van Beijeren, Kutner and Spohn in
\cite{BKS85} as an approximation of the fluctuations of general driven diffusion systems with one conserved quantity, 
e.g. the Asymmetric Simple Exclusion Process (ASEP) on $\Z^d$. 
The nonlinear term in~\eqref{e:BKS} expresses the fact that the velocity at which the fluctuation density 
travels depends on its magnitude. In the context of ASEP, 
comes from the gradient of the current and is due to its asymmetry, so that in particular the vector $\fw$  
models how the ASEP under consideration is asymmetric in the different spatial directions. 
%In view of this derivation, it is not surprising that 
In the {\it sub-critical} dimension $d=1$,~\eqref{e:BKS} reduces to the (derivative of the)  
celebrated Kardar--Parisi--Zhang (KPZ) equation~\cite{Kardar} whose local solution theory is by now 
well-understood~\cite{KPZ,KPZreloaded,GonJara,GPuni} and whose large-scale statistics 
have been recently determined~\cite{QSkpz,Vir}. 
In higher dimensions the only available results are in the 
discrete setting. For $d\geq 3$, corresponding to the super-critical case, 
ASEP has been shown to be diffusive at large scales~\cite{EMY,LY} -- 
its fluctuations are Gaussian 
and given by the solution of a Stochastic Heat equation (SHE) with additive noise,
see~\cite{CLO,LOV} and~\cite{Komorowski2012} for a review. 
The critical dimension $d=2$ is more subtle. Only qualitative information is available and 
no scaling limit has been obtained thus far. The most important result is that of~\cite{Yau}, 
which states that two-dimensional ASEP is logarithmically superdiffusive, 
meaning that the diffusivity $D(t)$ satisfies $D(t)\sim (\log t)^{2/3}$ for $t$ large.\footnote{The diffusivity measures 
how the correlations of a process grow in space as a function of time so that, denoting by $\ell(t)$ 
the correlation length, it formally satisfies $\ell(t)\sim \sqrt{ t D(t)}$. }

As mentioned above, our focus is on the critical and super-critical dimensions, 
i.e. we will be studying~\eqref{e:BKS} (in the continuum) in $d\geq 2$. 
In both cases, we will prove that the large-scale behaviour is Gaussian and that the fluctuations 
evolve as a SHE with {\it renormalized coefficients} (see Theorem~\ref{thm:main}), 
but while for $d\ge3$ the scaling 
adopted is the classical diffusive scaling, for $d=2$ we will consider the so-called {\it weak coupling} 
scaling. This corresponds to tuning down the strength of the nonlinearity (morally, the ``norm'' of the vector $\fw$) 
at a suitable rate together with the scale parameter, see~\eqref{eq:lambdaeps}. 
The tuning is chosen so that we can ``tame'' the 
growth caused by the non-linear term without destroying its large-scale effects  
(see e.g.~\cite{CD,CSZ,Gu2020,CES,CETWeak} for a similar choice of scaling in the critical dimension for different equations). 
\medskip

Let us mention that the study of large-scale Gaussian fluctuations of space-time random fields 
with local non-linear random dynamics is  not restricted the context of SPDEs but
is a classical topic in probability theory and mathematical physics.
From the more general point of view of theoretical physics, large-scale Gaussian fluctuations are expected 
for equilibrium diffusive systems above their critical dimension 
and this belief is informally explained via (non-rigorous) Renormalization Group (RG) arguments 
and effective dynamical field theories. 

A variety of tools have been used to mathematically justify these heuristics in the {\it static} setting, i.e. 
for the equilibrium measure of local non-linear systems. RG methods have been applied 
to derive the scaling limit of unbounded spin systems (and other statistical mechanical models)  at and above their critical dimension $d\ge 4$  in~\cite{bauerschmidtRenormalisationGroupMethod2015a, bauerschmidtScalingLimitsCritical2014, bauerschmidtIntroductionRenormalisationGroup2019}\footnote{the authors further claim 
that ``Our emphasis here is on the critical dimension $d = 4$, which is more difficult than dimensions $d > 4$."}. 
The related problem of triviality of $\Phi^4_d$ theories for $d > 4$ was treated in~\cite{Aiz82, Fro82, fernandezRandomWalksCritical1992} using apriori bounds on correlation functions, an approach 
which ultimately lead to a proof in the more challenging critical dimension $d=4$, 
see~\cite{aizenmanMarginalTrivialityScaling2021a}. 
Another class of equilibrium models for which scaling limits have been successfully determined 
are equilibrium interface models, so called $\nabla \phi$ models in all dimensions~\cite{brydgesGradPhiPerturbations1990, naddafHomogenizationScalingLimit1997, Funaki19971,  giacominEquilibriumFluctuationsNabla2001a}, for which the Helffer-Sj\"ostrand representation of correlations together 
with homogeneization arguments was exploited. Analogous results have been proven, with very different methods, for \emph{discrete} $\nabla\phi$ models in dimension $d=2$ (in particular dimer models, both integrable ones \cite{Kenyon} and their non-integrable perturbations \cite{Giul}).

On the other hand, establishing large-scale Gaussian fluctuations 
for {\it dynamical} problems has proven to be challenging  
and we are not aware of general results in this direction.
The major issue seems to stem from the fact that there is currently no framework 
which allows to naturally identify the law of the limit process. 
A lot of attention has been recently devoted to the above mentioned KPZ equation in high dimension, 
whose expression in its most general form is as~\eqref{e:BKS} but 
with nonlinearity given by $\langle \nabla\eta, Q\nabla\eta\rangle$, for $Q$ 
a $d\times d$ deterministic symmetric matrix, 
and a real valued space-time white noise in place of $\div\boldsymbol{\xi}$. 
Most of the works focus on the case in which $Q$ is given by the identity matrix  
and the specific form of the resulting equation plays a crucial role. Indeed, one can use the Cole--Hopf 
transformation to linearise the SPDE and turn it into a {\it linear} multiplicative 
stochastic heat equation or the associated directed-polymer models. 
These latter models possess an ``explicit" representation (via, e.g., Feynman-Kac formula) 
and one can then leverage Malliavin calculus tools to derive the scaling limit 
(see~\cite{dunlapRandomHeatEquation2021, guEdwardsWilkinsonLimit2018} for 
the multiplicative stochastic heat equation and~\cite{Gu2018b, CCM1, LZ, CNN} for KPZ). 
In this context, let us  mention the recent striking results on the fluctuations of KPZ at the critical dimension $d=2$ 
in the weak coupling scaling (or the intermediate disorder regime) starting with \cite{CD,CSZ,Gu2020} (for KPZ)
and culminating with the characterisation of the stochastic heat flow~\cite{caravennaCritical2dStochastic2023}.
An alternative approach is that used in~\cite{magnenScalingLimitKPZ2018}, 
which is based on re-expressing the SPDE as a functional integral via the Martin--Siggia--Rose formalism 
(essentially a Girsanov transformation, which again is only possible in view of the specific form of the equation) 
and then leveraging constructive quantum field theory techniques and RG ideas to control 
the large scale fluctuations in the regime of small non-linearity. 

In the critical dimension $d=2$, for $Q$ diagonal and with diagonal entries $1$ and $-1$ (so-called 
Anisotropic KPZ or AKPZ), progress has been made in~\cite{CET}, 
where it is shown that the equation is logarithmically superdiffusive, and in~\cite{CETWeak}, 
where the large-scale behaviour in the weak coupling scaling is determined. In both articles, 
the tools adopted are closer to those in the present paper and we will comment more on them below. 

Let us also mention recent work in extending the homogeneization methods 
(originated in the work of Naddaf and Spencer~\cite{naddafHomogenizationScalingLimit1997}) to dynamical problems~\cite{armstrongQuantitativeHydrodynamicLimits2022, cardaliaguetScalingLimitsStochastic2022}.
\medskip

One of the main advantages of the approach taken in the present paper is that 
it moves away from the usual technique of Cole--Hopf transformation, which is {\it not applicable}, 
towards a set of tools which on the one hand makes a link with interacting particle systems, 
i.e. via the martingale problem and the associated infinite-dimensional generator, and on the other 
has better chances to apply more generally. A crucial tool for us is the fact that we have control 
over the invariant measure for~\eqref{e:BKS}, which significantly simplifies the already involved 
analysis of the generator. While this is another non-generic situation, there is a wide class of examples 
in which a similar control is available and to which our tools potentially directly apply, 
e.g. (super-)critical Stochastic Navier--Stokes~\cite{GT,CK,HZZ},  
diffusions in divergence-free random vector fields~\cite{CHT}, lattice gas models~\cite{LRY} and many others. 

Our results open the way to gradually attack more challenging models and we discuss some open problems below in Section~\ref{sec:openproblems}.

Before turning to the core of the proof, in the next section we will precisely state our main result 
and outline the ideas behind our arguments. 



\subsection{The equation and main result}

As written~\eqref{e:BKS} is meaningless as the noise is too irregular for the non-linearity to be well-defined. 
Nonetheless, as we are interested in the large-scale behaviour, 
we regularise it, so to have a well-defined field at the microscopic level, and 
our goal then becomes to control its fluctuations while zooming out at the correct scale. 
We choose the regularisation in such a way to retain a fundamental property of the solution, 
namely its (formal) invariant measure. 
This amounts to smoothening the quadratic term via a Fourier cut-off as follows
\begin{equ}[eq:toroN]
    \partial_t\eta =\frac{1}{2}\Delta\eta+ \fw\cdot\Pi_1\nabla(\Pi_1\eta)^2  + (-\Delta)^{\frac{1}{2}}\xi,
\end{equ}
where for $a>0$, $\Pi_a$ acts in Fourier space as
\begin{equ}[eq:FourPi]
\widehat{\Pi_a\eta}(k)=\widehat{\eta}(k)1_{|k|\leq a}\,.
\end{equ}
Without loss of generality, 
we also replaced $\div\boldsymbol{\xi}$ in~\eqref{e:BKS} with $(-\Delta)^{\frac{1}{2}}\xi$ for $\xi$ a space-time white noise 
(see~\eqref{e:fLapla} for the definition of the fractional Laplacian), since the two can be easily seen to have the same law. 

\begin{remark}\label{rem:pignolerie}
In principle, the results below can be shown to hold using the same techniques of the present paper, 
even if instead of the Fourier cut-off we used a more general mollifier. That said, 
to simplify some (minor but annoying) technical points, we will stick to $\Pi_1$ and further assume that 
  \begin{itemize}[noitemsep]
  \item for $d\ge3$, $1/\eps\in \N+1/2$ and that the norm in \eqref{eq:FourPi} is the sup-norm $|\cdot|_\infty$, 
  \item for $d=2$, the norm  in \eqref{eq:FourPi}  is the Euclidean norm $|\cdot|$ (no further assumption on $\eps$ is made).
  \end{itemize}
We refer the reader to Remark~\ref{rem:G} for some additional (technical) advantages of our choice. 
\end{remark}

In order to avoid technicalities related to infinite volume issues, 
we restrict~\eqref{eq:toroN} to 
the $d$-dimensional torus $\mathbb T^d_\eps$ of side-length $2\pi /\eps$. Here, $\eps$ is 
the ``scale'' parameter which will later be sent to $0$. 
Note that via the diffusive rescaling
\begin{equ}[e:scaling]
    \eta^\eps(x,t)\eqdef \eps^{-\frac d2}\eta\left(x/\eps,t/\eps^{2}\right)\,, 
\end{equ}
equation~\eqref{eq:toroN} becomes
\begin{equ}\label{e:BurgersScaled}
    \partial_t\eta^\eps = \frac{1}{2}\Delta\eta^\eps +\lambda_\eps \fw\cdot\Pi_{1/\eps}\nabla(\Pi_{1/\eps}\eta^\eps)^2  + (-\Delta)^{\frac{1}{2}}\xi
\end{equ}
where now the spatial variable takes values in the $d$-dimensional torus $\T^d\eqdef \T^d_1$ of side-length $2\pi$, 
and $\xi$ is a space-time white noise on $\R^+\times \mathbb T^d$. 
The coefficient $\lambda_\eps$ depends on the dimension 
and is defined as
\begin{equ}[eq:lambdaeps]
  \lambda_\eps\eqdef
  \begin{cases}
   \frac1{\sqrt{\log \eps^{-1}}} &\text{if $d=2$}\\
    \eps^{\frac{d}{2}-1} & \text{if $d\ge3$}.
  \end{cases}
\end{equ}
While in the super-critical case $d\geq 3$, $\lambda_\eps$ is directly determined by~\eqref{e:scaling}, 
for $d=2$,~\eqref{e:BKS} is formally scale invariant under the diffusive scaling  
(another reason why such dimension is critical) 
so that one would have $\lambda_\eps=1$. The choice made in~\eqref{eq:lambdaeps} 
corresponds to the so-called {\it weak coupling} scaling alluded to earlier and we will comment more on it 
after we state the main result. Before doing so, let us give the notion of solution for~\eqref{e:BurgersScaled} 
we will be working with. 

\begin{definition}\label{def:StatSol}
We say that $\eta^\eps$ is a stationary solution to~\eqref{e:BurgersScaled} with coupling constant $\lambda_\eps$ if 
$\eta^\eps$ solves~\eqref{e:BurgersScaled} and its initial condition is $\eta(0,\cdot)\eqdef\mu$, for $\mu$ 
a spatial white noise on $\T^d$, 
i.e. a mean zero Gaussian distribution on $\T^d$ such that 
\begin{equ}[e:CovS]
\E[\mu(\phi)\mu(\psi)]=\langle\phi,\psi\rangle_{L^2(\T^d)}\,,\qquad \text{for all $\phi,\,\psi\in L^2_0(\T^d)$}
\end{equ}
where $\langle\cdot,\cdot\rangle_{L^2(\T^d)}$ is the usual scalar product in $L^2(\T^d)$ and 
$L^2_0(\T^d)$ is the space of zero-average square integrable functions on $\T^d$. 
We will denote by $\P$ the law of $\mu$ and by $\E$ the corresponding expectation.
\end{definition}

The reason why the solution $\eta^\eps$ of~\eqref{e:BurgersScaled} started from a spatial white noise 
is called {\it stationary} is that $\P$ is its invariant measure 
and, as we will prove in Lemma \ref{lem:generalproperties} below,  this holds irrespective of the choice of $\fw,\eps$. 
Instead, the law $\mathbf P$ (with corresponding expectation $\mathbf E$) of a stationary 
solution to~\eqref{e:BurgersScaled} clearly depends on $\fw,\eps$. 

For $T>0$, let us denote by $C([0,T],\cS'(\T^d))$ the space of continuous 
functions with values in the space of tempered distributions $\cS'(\T^d)$ and by $\cF$ the Fourier transform 
on $\cS'(\T^d)$ (see~\eqref{e:FT} for a definition). 
We are now ready to state the main result of the paper. 

\begin{theorem}\label{thm:main}
Let $\fw\in\R^d\setminus\{0\}$ and $T>0$. 
For $\eps>0$, let $\lambda_\eps$ be defined according to \eqref{eq:lambdaeps} and 
$\eta^\eps$ be the stationary solution of \eqref{e:BurgersScaled} started from a spatial white noise $\mu$. 
Then, there exists a constant $D_\SHE>0$, depending only on the dimension $d$ and the norm $|\fw|$ of $\fw$, 
such that, as $\eps\to0$, $\eta^\eps$ converges in law in $C([0,T],\cS'(\T^d))$ to the unique stationary solution $\eta$ of
\begin{equ}[e:SHEintro]
\partial_t \eta =\tfrac12(\Delta + D_\SHE(\fw\cdot\nabla)^2)\eta + (-\Delta - D_\SHE(\fw\cdot\nabla)^2)^{1/2}\xi\,,\qquad \eta(0,\cdot)=\mu\,,
\end{equ}
where $(\fw\cdot\nabla)^2$ is the linear operator defined as $\cF((\fw\cdot\nabla)^2\phi)(k)\eqdef -(\fw\cdot k)^2\cF(\phi)(k)$, for $\phi\in \cS'(\T^d)$ and $k\in\Z^d$. 

In the case $d=2$, $D_\SHE$ is explicit and given by the formula
\begin{equ}
  \label{eq:DSHEd2}
  D_\SHE= \frac1{|\fw|^2}\left[\left(\frac{3|\fw|^2}{2\pi} +1\right)^{\frac{2}{3}}-1\right],
\end{equ}
$|\fw|$ being the Euclidean norm of $\fw$.
\end{theorem}

Before proceeding to the proof of the above theorem, let us make a few comments on 
its statement and on the choice of $\lambda_\eps$ in~\eqref{eq:lambdaeps}. 
Notice that in any dimension $d\geq 2$, the constant $D_\SHE$ appearing in the limiting 
equation is {\it strictly} positive. Hence, despite the presence of the vanishing factor 
$\lambda_\eps$ in front of the nonlinearity, not only the quadratic term does not go to $0$ 
but actually produces a new ``Laplacian'' and a new noise. Further, this new Laplacian (and noise), in a sense,  
``feels'' the small-scale behaviour of the field $\eta$ in~\eqref{eq:toroN} 
as it depends on the vector $\fw$, which in turn describes 
its microscopic dynamics. A similar phenomenon has already been observed for other 
critical and super critical SPDEs, even though we are not aware of examples 
in which a Laplacian of the form above was derived. 

We have already pointed out that in $d=2$~\eqref{e:BKS} is formally scale invariant. 
That said, the above mentioned result in~\cite{Yau} 
for ASEP suggests that also~\eqref{eq:toroN} is logarithmically superdiffusive and the additional 
diffusivity can only come from the non-linear term. What the previous theorem 
shows is that the choice in~\eqref{eq:lambdaeps} guarantees that 
the nonlinearity ultimately gives a non-vanishing order 1 contribution. 
What might appear puzzling is that, by translating the result of Yau to~\eqref{e:BurgersScaled}, 
the diffusivity $D^\eps(t)$ of the scaled process grows as $|\log \eps|^{2/3}$ so that one might be led 
to think that $\lambda_\eps$ should be chosen accordingly. While this is not 
the case, the expression in~\eqref{eq:DSHEd2} formally implies the result of Yau, as can 
be seen by taking $\fw=\fw_\eps$ such that $|\fw_\eps|\eqdef \lambda_\eps^{-1}$. 

Furthermore, a similar scaling has been considered 
in the critical dimension $d=2$ in the context of the KPZ 
equation~\cite{CD,CSZ,Gu2020, caravennaCritical2dStochastic2023}, 
of the Anisotropic KPZ equation~\cite{CES,CETWeak} and of the stochastic Navier--Stokes equation~\cite{CK}. 
What is interesting is that, even though these latter examples in principle have 
different large scale diffusivity, i.e., $D(t)\sim t^\beta$ for some $\beta>0$~\cite{Forrest} for KPZ, 
$D(t)\sim \sqrt{\log t}$ for AKPZ 
and Navier--Stokes ~\cite{CET,Adler}, they all display {\it non-trivial} 
(in that the non-linearity nontrivially contributes to the limit) {\it Gaussian
fluctuations} under the same weak coupling scaling $\lambda_\eps=|\log\eps|^{-1/2}$, thus suggesting some sort of universality for it. 
This is to be compared with the one-dimensional case, in which, setting $\lambda_\eps\eqdef\sqrt\eps$, 
one recovers the so-called {\it weakly asymmetric} scaling under which convergence of 
ASEP to the one-dimensional KPZ equation was first shown in~\cite{BG} and since then for a wide 
variety of models. 
\subsection{Open problems}\label{sec:openproblems}

Our results raise several interesting questions, as
\begin{itemize}[noitemsep]
\item obtain Gaussian fluctuations for~\eqref{e:BKS} in the case of more general non-linearity, 
i.e. $\fw\cdot\nabla \eta^2$ replaced by $\fw\cdot\nabla F(\eta)$, for instance for polynomial $F$. 
The main difficulty is to obtain operator estimates similar to those of Lemma \ref{l:GenaBound}, 
which though seem highly non-trivial.
\item Prove the analog of Theorem \ref{thm:main} (for $d=2$) for ASEP on $\Z^2$ in the same limit of weak asymmetry.
\item In dimension $d=2$, define $\tilde\eta^\eps(x,t)=\eps^{-1}\eta( x/\eps, t/(\eps^2\tau(\eps)))$ 
(to be compared with \eqref{e:scaling}) and find the right correction $\tau(\eps)$ to the diffusive scaling 
so that a non-trivial scaling limit as $\eps\to0$ exists. This corresponds to a \emph{strong coupling regime}. 
On the basis of \cite{Yau}, the natural guess is $\tau(\eps)=|\log \eps|^{2/3}$ but the identification of the limit 
is hard as the regularising properties of the Laplacian vanish. 
\end{itemize}

More challenging will be to move towards models for which the invariant measure is non-Gaussian or non-explicit. 
One guiding heuristic could be that the limiting Gaussian fluctuations and 
associated Gaussian stationary distribution could still provide a good  setting for the analysis.

%
%
%\newpage 
%Let us comment on~\eqref{eq:lambdaeps}. In the super-critical case $d\geq 3$, 
%we have no choice as $\lambda_\eps$ is determined by the diffusive scaling~\eqref{e:scaling}. 
%For $d=2$,~\eqref{e:BKS} is formally scale invariant under~\eqref{e:scaling} so that in principle 
%one would need to choose $\lambda_\eps=1$. That said, the above mentioned result in~\cite{Yau} 
%for ASEP suggests that also~\eqref{eq:toroN} is logarithmically superdiffusive and the additional 
%diffusivity cannot come but from the non-linear term. Our choice is made in such a way that 
%the nonlinearity ultimately gives a non-vanishing contribution of order $1$, as we will see in the 
%statement below, and is called {\it weak coupling} scaling. A similar scaling has been considered 
%in the critical dimension $d=2$ in the context of the KPZ 
%equation~\cite{CD,CSZ,Gu2020, caravennaCritical2dStochastic2023}, 
%of the Anisotropic KPZ equation~\cite{CES,CETWeak} and of the stochastic Navier--Stokes equation~\cite{CK}. 
%Let us also mention, that in $d=1$ setting $\lambda_\eps\eqdef\eps^{1/2}$, one recovers the 
%so-called {\it weakly asymmetric} scaling under which convergence of 
%ASEP to the one-dimensional KPZ equation was shown in~\cite{BG} and since then for a wide 
%variety of models~\cite{}. \giuseppe{Who do we cite? There might be countless examples...} 
%\medskip
%
%From now on, we will be working with the scaled equation~\eqref{e:BurgersScaled} started from the 
%stationary distribution. \giuseppeText{FROM HERE!}
%
%
%
%
%
%
%
%
%
%\newpage
%Since we are interested in the large scale behaviour, 
%we regularise it by restricting it to the $d$-dimensional torus $\T^d_\eps$ of side-length $2\pi /\eps$ 
%(the parameter $\eps$ represents the ``scale'' and we will be sending it to $0$)
%and introducing a Fourier 
%
%\giuseppeText{FROM HERE}
%We study an SPDE introduced by van Beijeren, Kutner and Spohn in
%\cite{BKS85} as a continuum limit of the fluctuation process
%of the Asymmetric Simple Exclusion Process on $\Z^d$. Formally, the
%equation is given by
%\begin{equation}
%    \partial_t\eta = \frac{1}{2}\Delta\eta+\fw\cdot\nabla\eta^2 + (-\Delta)^{\frac{1}{2}}\xi
%  \end{equation}
%  where $\eta=\eta(x,t),t\ge 0,x\in \R^d$ with $d\ge2$, $\fw\in \mathbb R^d$ is
%  a fixed vector and $\xi$ is a space-time white noise on
%  $\R^+\times \R^d$.  Actually, in \cite{BKS85} the noise is
%  written as $\mathrm{div}\xi$, i.e.  the divergence of a vector of
%  $d$ independent space-time white noises. On the other hand,
%  $\mathrm{div}\xi$ is easily seen to have the same law as
%  $(-\Delta)^{\frac{1}{2}}\xi$.
%
%  In order to make the equation well-posed, we regularize it by on one
%  hand restricting it to the $d$-dimensional torus $\mathbb T^d_\eps$
%  of side-length $2\pi /\eps$ ($\eps$ will tend to $0$ later) and, on the other hand by
%  introducing a cut-off in Fourier space in the non-linearity:
%\begin{equation}
% \label{eq:toroN}
%    \partial_t\eta =\frac{1}{2}\Delta\eta+ \fw\cdot\Pi_1\nabla(\Pi_1\eta)^2  + (-\Delta)^{\frac{1}{2}}\xi,
%\end{equation}
%where for $a>0$, $\Pi_a$ acts in Fourier space
%as
%\begin{eqnarray}
%  \label{eq:FourPi}
%\widehat{\Pi_a\eta}(k)=\widehat{\eta}(k)1_{|k|\leq a}  
%\end{eqnarray}
%and $\xi$
%is a space-time white noise on $\R^+\times \T^d_\eps$. We will specify later which norm $|\cdot|$ is used here.
%
%Note that via the rescaling
%\begin{equation}
%    \eta^\eps(x,t)\eqdef \eps^{-\frac d2}\eta\left({\eps^{-2}t},\eps^{-1}x\right),
%\end{equation}
%the equation \eqref{eq:toroN} becomes
%\begin{equation}\label{e:BurgersScaled}
%    \partial_t\eta^\eps = \frac{1}{2}\Delta\eta^\eps +\lambda_\eps \fw\cdot\Pi_{1/\eps}\nabla(\Pi_{1/\eps}\eta^\eps)^2  + (-\Delta)^{\frac{1}{2}}\xi
%\end{equation}
%on the torus $\mathbb T^d\eqdef \mathbb T^d_1$ of side-length $2\pi$,
%where $\lambda_\eps=\eps^{-1+\frac{d}{2}}$. Now $\xi$ is a space-time
%white noise on $\R^+\times \mathbb T^d$.  In the case $d=2$ we will
%let instead $\lambda_\eps=\frac1{|\log \eps|^{1/2}}.$ In this case we
%are not simply rescaling \eqref{eq:toroN}, we are looking at the weak
%non-linear regime, in the same spirit as the WASEP limit in interacting particle systems or the
%weak non-linearity/disorder regime of the 2D AKPZ or KPZ equation (REFERENZE).
%
%Summarizing, we will study the $\eps\to0$ limit of \eqref{e:BurgersScaled}, with
%\begin{equ}
%  \label{eq:lambdaeps}
%  \lambda_\eps=\left\{
%  \begin{array}{lll}
%   \frac1{\sqrt{|\log \eps|}} &\text{if}&d=2\\
%    \eps^{-1+\frac{d}{2}} & \text{if} & d\ge3.
%  \end{array}
%  \right.
%\end{equ}
%
%\begin{remark}
%  \label{rem:pignolerie}
%  In order to simplify some (minor but annoying) technical points, we will assume that:
%  \begin{itemize}
%  \item in the case $d\ge3$, $1/\eps\in \mathbb N+1/2$ and that the norm in \eqref{eq:FourPi} is the sup-norm $|\cdot|_\infty$;
%    
%  \item in the case $d=2$, the norm  in \eqref{eq:FourPi}  is the Euclidean norm $|\cdot|$ (no further assumption on $\eps$ is made).
%  \end{itemize}
%\end{remark}
%
%Let $\mathbb P$ denote the law of the white noise $\eta$ on
%$\mathbb T^d$, with zero spatial average. It is easy to check (see
%Lemma \ref{lem:generalproperties} below) that $\mathbb P$ is stationary,
%irrespective of the choice of $w,\eps$. We let  $\mathbf P$ (that depends on $w,\eps$) denote the law of the stationary process given by the solution of \eqref{e:BurgersScaled} with initial condition $\eta^\eps(0)\sim \mathbb P$.
%
%
%Here is our main result:
%\begin{theorem}
%  \label{thm:main}
%  Let $\eta^\eps$ be the stationary solution of \eqref{e:BurgersScaled}, with $\lambda_\eps$ as in \eqref{eq:lambdaeps}. If $w\ne 0$ there exists a constant $D_\SHE>0$ (that depends on the dimension $d$ and on the norm $|w|$ of $w$) such that, as $\eps\to0$, $\eta^\eps$ tends in distribution to the stationary  solution $\eta$ of
% \begin{equ}[e:SHE2]
%\partial_t \eta =\tfrac12(\Delta + D_\SHE(\fw\cdot\nabla)^2)\eta + (-\Delta - D_\SHE(\fw\cdot\nabla)^2)^{1/2}\xi\,,\qquad \eta(0,\cdot)\sim \mathbb P.
%\end{equ}
%% Where $\cD$ is the linear differential operator $\cD=\tfrac12D_\SHE(\fw\cdot \nabla)^2$,  whose Fourier multiplier is 
%%   \begin{equation}
%%     \label{e:cD}
%%     \hat\cD(k)=-\tfrac12(\fw\cdot k)^2 D_\SHE.
%%   \end{equation}
%In particular, for $d=2$ one has the explicit formula
%\begin{eqnarray}
%  \label{eq:DSHEd2}
%  D_\SHE= \frac1{|\fw|^2}\left[\left(\frac{3|\fw|^2}{2\pi} +1\right)^{\frac{2}{3}}-1\right],
%\end{eqnarray}
%with $|\fw|$ the Euclidean norm of $\fw$.
%\end{theorem}
%
%{\color{red} something on Yau? Here or before?}


\subsection{Idea of the proof }\label{sec:ideas}

The idea of the proof finds its roots in the approach detailed in~\cite{Komorowski2012} in the context of 
interacting particle systems. To summarise it and see how it translates to the present context, 
let us consider the weak formulation of~\eqref{e:BurgersScaled} started from a spatial white noise $\mu$,  
which, for a given test function $\phi$, reads 
\begin{equ}[e:WeakIntro]
\eta^\eps_t(\phi)-\mu(\phi)-\tfrac12\int_0^t \eta^\eps_s(\Delta\phi)\dd s-\int_0^t\cN^\eps_\phi(\eta^\eps_s)\dd s=M_t(\phi)
\end{equ}
where $\cN^\eps_\phi$ stands for the nonlinearity tested against $\phi$ and $M$ is the martingale 
associated to the space-time white noise $\xi$. Now, upon assuming the sequence $\{\eta^\eps\}_\eps$ 
to be tight in the space $C([0,T]\cS'(\T^d))$, we see that all the terms in the above expression 
converge (at least along subsequences) but we have no information concerning the nonlinearity. 
For Theorem~\ref{thm:main} to be true, we need this latter term to produce both a {\it dissipation} term 
of the form of $D_\SHE\gensy^\fw\eta$, where $\gensy^\fw\eqdef\frac12(\fw\cdot\nabla)^2$, and a {\it fluctuation} 
part which should encode the additional noise in~\eqref{e:SHEintro}. 
In other words, the problem is to identify a {\it fluctuation-dissipation} relation~\cite{LY}, i.e. to determine $D_\SHE>0$ 
and $\cW^\eps$ such that 
\begin{equ}[e:FD]
\cN^\eps_\phi(\eta^\eps)=-\gen \cW^\eps(\eta^\eps) + D_\SHE\gensy^\fw\eta^\eps(\phi)
\end{equ} 
where $\gen$ is the generator of $\eta^\eps$. Indeed, since by Dynkin's formula, there exists a martingale $\cM^\eps(\Weps)$ 
such that
\begin{equ}
\cW^\eps(\eta^\eps_t)-\cW^\eps(\mu)-\int_0^t \gen\cW^\eps(\eta^\eps_s)\dd s=\cM_t^\eps(\Weps), 
\end{equ}
given~\eqref{e:FD}, we could rewrite the nonlinear term in~\eqref{e:WeakIntro} as
\begin{equ}[e:NonlinExp]
\int_0^t\cN^\eps_\phi(\eta^\eps_s)\dd s=\int_0^t D_\SHE\gensy^\fw\eta_s^\eps(\phi)\dd s -\cM_t^\eps(\Weps)+o(1)\,,
\end{equ}
where $o(1)$ is a vanishing summand which contains the boundary terms. 
The advantage of the above is that we have expressed the nonlinearity in terms of a drift which captures  
the additional diffusivity and a martingale part which instead econdes the extra noise. 
At this point, tightness would ensure convergence of the former and we would be left to prove that the sequence 
$\cM^\eps(\Weps)$ converges to a martingale with the correct quadratic variation. 
\medskip

The (hardest) problem is clearly the derivation of the fluctuation-dissipation relation and 
this is the point from which our analysis departs from that in~\cite{Komorowski2012}. 
Even though, as we will see, it is enough to determine~\eqref{e:FD} approximately (i.e. that the difference 
of left and right hand side is small in a suitable sense), in the equation there are two unknowns, 
$D_\SHE\gensy^\fw$ which is {\it not} apriori given to us, and 
$\cW^\eps$ which, even if we were given the previous, is the solution of an infinite dimensional equation. 
To separate the two issues, we introduce a suitable truncation which removes the second 
summand from the right hand side of~\eqref{e:FD}, so that we can first solve for $\cW^\eps$ and 
then project back and determine $D_\SHE\gensy^\fw$. 
This is done in $d\geq 3$ and $d=2$ in very different ways. In the former case, we will control 
$\cW^\eps$ similarly to~\cite{LY}, as we can show it satisfies a graded sector condition 
in the spirit of~\cite[Section 2.7.4]{Komorowski2012}. Nonetheless, 
their functional analytic approach does not apply in our setting and in particular, for the identification 
of $D_\SHE$ and $\gensy^\fw$, we devise a new method which is detailed in Section~\ref{sec:Diff3}. 
For $d=2$ instead, we introduce a novel ansatz which is based on the idea that at large scales 
the generator of $\eta^\eps$ should approximate a {\it modulated version} of the generator of~\eqref{e:SHEintro}. 
This ansatz simplifies dramatically the (iterative) analysis performed in~\cite{CETWeak} and 
is heavily based on the Replacement Lemma in Section~\ref{sec:ReplLemma}. 


%
%At this point we are left to recover $D_\SHE\gensy^\fw$ by projecting $\cW^\eps$. 
%
%
%while in $d\geq3$ we control it a graded sector condition in the spirit of~\cite[Section 2.7.4]{Komorowski2012}, 
%for $d=2$ we devise a novel approximation which dramatically simplifies the approach taken in~\cite{CETWeak} 
%and is based on  


%\newpage 
%
%{\color{red}\begin{enumerate}
%\item mention the book~\cite{Komorowski2012} and say we are inspired by similar ideas. 
%\item Write the weak form of the equation~\eqref{e:Weak}.
%\item say that the difficult term is the nonlinearity and we need to determine a way to extract from it a drift (in the form 
%of the new Laplacian) and a noise term (in the form of a new martingale. 
%\item say that ideally we would need to prove a fluctuation-dissipation equation of the form $\gen \mathcal W = \cN_\phi + laplacian$, where the unknowns in the equation are 2: $\mathcal W$ and $laplacian$. Indeed if we were able to do that by Dynkin formula we could write~\eqref{e:DynkinCiao} which, 
%once plugged into the weak formulation would only require to prove that the boundary terms go to $0$ and the martingale converge. 
%\item the problem is that the solution of the fluctuation-dissipation is difficult. The way in which we proceed is that in 
%$d\geq 3$ we truncate and then take two limits. Here some sort of sector condition might be useful. In $d=2$ we prove something of the form 
%$\gen\approx \gensy^\fw G(\gensy)$ (maybe even more informal, ``we replace $\gen$ with a diagonal operator written 
%as...'') In spirit the approach is similar to that in~\cite{CETWeak} but we significantly improved it and a heuristic is provided 
%in Appendix~\ref{a:Heu}. 
%\end{enumerate}}

\subsection*{Organisation of the article}

The rest of this work is organised as follows. 
Below we introduce notations, conventions, function spaces and 
elements of Wiener space analysis which will be used throughout. 
Section~\ref{sec:EqandGen} is the bulk of the paper. In Section~\ref{sec:Burgers}, 
after recalling basic properties of~\eqref{e:BurgersScaled} 
and its generator, we discuss the two main technical tools - integrability estimates for $d\geq 3$ and the 
Replacement Lemma for $d=2$. In Section~\ref{sec:apriori} we determine apriori estimates on the solution to the 
generator equation and in the next two we identify $D_\SHE$ and $\gensy^\fw$, 
so to obtain the fluctuation-dissipation relation~\eqref{e:FD}. 
Section~\ref{sec:ProofofMain} is devoted to the proof of Theorem~\ref{thm:main} which, 
given the analysis in Section~\ref{sec:EqandGen}, essentially reduces to the control of the 
martingales $\cM^\eps(\Weps)$ in~\eqref{e:NonlinExp} which is carried out in Section~\ref{sec:martingales}. 
At last, Appendix~\ref{a:Approx} contains some technical steps necessary in the proof of the 
Replacement Lemma. 







%
%\newpage 
%we recall some results from previous works, 
%introduce the Burger's version of our equation, and prove some preliminary estimates. 
%Section~\ref{sec:Dbulk} is devoted to the proof of Theorem~\ref{thm:Dbulk}, 
%and, along the way, we derive a replacement lemma, 
%which basically allows one to replace the generator of~\eqref{e:AKPZ} 
%by the appropriately modulated (in Fourier space) generator of the stochastic heat equation. 
%In Sections~\ref{sec:truncated} and~\ref{s:theconvergence}, we prove Theorem~\ref{thm:Conv}. 
%First, we introduce the martingale problem and provide an outline of the proof. 
%Then, we carry out a thorough analysis of the generator $\gen$ of $h^N$ and identify the 
%class of observables needed to establish the statement, 
%whose proof is completed at the beginning Section~\ref{s:theconvergence}. 
%In the rest of the section, we show that the martingales associated to the observables 
%determined in the Section~\ref{sec:truncated} have quadratic variation which is asymptotically deterministic. 
%At last, in the appendix we obtain some technical results needed for the proof of the replacement lemma 
%(Appendix~\ref{a:Approx}), control of the quadratic variation (Appendix~\ref{app:S}) and provide 
%an example of the estimates to perform in Section~\ref{s:theconvergence} (Appendix~\ref{sec:example}). 

%\subsection{Motivations and related work}
%
%
%Large scale Gaussian fluctuations of space-time random fields with local non-linear random dynamics is a classical topic in probability theory. The more general approach to these problems comes from the fluctuation theory of Markov processes as  e.g.  described recently in the book~\cite{Komorowski2012}. From the point of view of theoretical physics, Gaussian large scale fluctuations are generally expected for equilibrium diffusive systems above their critical dimension and  it is heuristically  justified via renormalization group arguments and effective dynamical field theories. 
%
%The analysis of the large scale fluctuations of stochastic partial differential equations sits somehow between these two frameworks: from heuristic considerations it is generally expected that diffusive, non-linear SPDEs driven by (small scale regularisation of) space-time white noise should have Gaussian fluctuations in high dimensions. In particular this should hold for Kardar--Parisi--Zhang type equations for $d > 2$ and for reaction-diffusion models like the dynamic $\Phi^4_d$ model for $d > 4$. 
%
%Heuristically one would expect that renormalization group ideas should play an important role in establishing Gaussian fluctuations in high dimensions, both for the equilibrium measure and for the dynamics of local non-linear systems.  A  renormalization group methods has been applied to tackle the scaling limit of unbounded spin systems (and other statistical mechanical models)  in $d=4$ by Bauerschmidt, Brydges and Slade~\cite{bauerschmidtRenormalisationGroupMethod2015a, bauerschmidtScalingLimitsCritical2014, bauerschmidtIntroductionRenormalisationGroup2019}. In~\cite{bauerschmidtRenormalisationGroupMethod2015a} these authors claim also that ``Our emphasis here is on the critical dimension d = 4, which is more difficult than dimensions $d > 4$.". However to date we are not aware of any proof or sketch of strategy to tackle the $d>4$ case with RG methods. \max{Maybe I will ask again to Roland to clarify for me where is the problem. I believe there is an issue with the large field region since the non-linear couplings are going to zero too fast to provide coercivity.}
%
%The related problem of triviality of $\Phi^4_d$ theories for $d > 4$ has been attacked in the literature via techniques involving apriori bounds on correlation functions~\cite{Aiz82, Fro82, fernandezRandomWalksCritical1992} culminating in the recent proof of triviality for $\Phi^4_4$~\cite{aizenmanMarginalTrivialityScaling2021a}.  
%
%Another class of equilibrium models where scaling limits have been successfully addressed are equilibrium interface models, so called $\nabla \phi$ models~\cite{brydgesGradPhiPerturbations1990, naddafHomogenizationScalingLimit1997, funakiMotionMeanCurvature1997,  giacominEquilibriumFluctuationsNabla2001a, millerFluctuationsGinzburgLandauNabla2011}. Here the main approach is to use the Helffer-Sj\"ostrand representation of correlations together with homogeneisation arguments. The RG method of Bauerschmidt, Brydges and Slade has been also successfully applied in this context by Adams, Koteck\'y and M\"uller~\cite{adamsStrictConvexitySurface2016} to address the difficult situation of non-convex potential in the low temperature region.
%
%Establishing the existence of these high-dimensional Gaussian scaling limits for dynamical problems has proven to be quite difficult and we are not aware of general results in this direction. The major difficulty in obtaining Gaussian fluctuation results seems to have a framework with allow to identify quite naturally the law of the limit process. 
% A lot of attention has been devoted recently to the KPZ equation in high dimension, where the specific form of the equation allows to use the Cole--Hopf transformation to transform the problem in an analysis of the large scale behaviour of the stochastic heat equation or of the associated directed-polymer models. These latter models allow for ``explicit" solutions and one can leverage Malliavin calculus to tackle the scaling limit, see e.g. the works~\cite{dunlapRandomHeatEquation2021, guEdwardsWilkinsonLimit2018, dunlapFluctuationsSolutionsKPZ2019}. In this context, let us  mention the recent striking results on the fluctuations for KPZ at the critical dimension $d=2$~starting with \cite{CD} and ending up with the characterization of the stochastic heat flow~\cite{caravennaCritical2dStochastic2023}.
%
%An alternative approach is that used in~\cite{magnenScalingLimitKPZ2018}  and based on re-expressing the SPDE as a functional integral via the Martin--Sigga--Rose formalism (essentially a Girsanov transformation) and then leveraging constructive quantum field theory techniques and renormalization group ideas to control the large scale fluctuations and prove Gaussian fluctuations in the regime of small non-linearity. Let us also mention here recent work in extending the homogeneisation methods (originated in the work of Naddaf and Spencer~\cite{naddafHomogenizationScalingLimit1997}) to dynamical problems~\cite{armstrongQuantitativeHydrodynamicLimits2022, cardaliaguetScalingLimitsStochastic2022}.
%
%In our work we focus on a variation of the KPZ (or Burgers) equation for which we have a control on the invariant measure (which is not the case in general). While this is another non-generic situation, it moves away from the usual technique of Cole--Hopf transformation mentioned above, towards a set of techniques which make a link with interacting particle systems, i.e. via the martingale problem and the associated infinite-dimensional generator, and which have some hope to apply also more generally. The existence of an invariant measure simplifies drastically the already quite involved analysis of the generator. 






\subsection*{Notations, function spaces and Wiener space analysis}

We let $\T^d$ be the $d$-dimensional torus of side length $2\pi $. % If $N=1$ then we simply write $\T^2$ instead of $\T_N^2$.
We denote by $\{e_k\}_{k\in\Z^d}$ the Fourier basis defined via 
$e_k(x) \eqdef \frac{1}{(2\pi)^{d/2}} e^{i k \cdot x}$ which, for all $j,\,k\in\Z^d$, satisfies 
$\langle e_k, e_{-j}\rangle_{L^2(\T^d)}= \mathds{1}_{k=j}  $. 
% {\color{red} I think the following sentence can be omitted: The basis functions $e_k$ can be decomposed in their real and imaginary part, so that $e_k=a_k+i b_k$ 
% and the system $\{a_k\}_{k \in \Z_{\mathrm{diag}}^2} \sqcup \{b_k\}_{k \in \Z_{\mathrm{diag}}^2\setminus\{0\}}$ 
% forms a real valued orthogonal basis of $L^2(\T_L^2)$, where $\Z_{\mathrm{diag}}^2= \{(k_1,k_2)\in\Z^2:\, k_1\geq k_2\}$. }

The Fourier transform of a function $\phi\in L^2(\T^d)$ will be denoted by 
$\cF(\phi)$ or $\hat\phi$ and, for $k\in\Z^d$ is given by the formula
	\begin{equation}\label{e:FT}
	\cF(\varphi)(k) =\hat\varphi(k)\eqdef  \int_{\T^d} \varphi(x) e_{-k}(x)\dd x\,, 
	\end{equation}
so that in particular
\begin{equation}\label{e:FourierRep}
\varphi(x) = % \frac{1}{N^2}
\sum_{k\in\Z^d} \hat\varphi(k) e_k(x)\,,\qquad\text{for all $x\in\T^d$. }
\end{equation}
Let $\cS(\T^d)$ be the space of smooth functions on $\T^d$ and $\cS'(\T^d)$ 
the space of real-valued distributions given by the dual of $\cS(\T^d)$. 
For any $\eta\in\cS'(\T^d)$ and $k\in\Z^d$, 
we will denote its Fourier transform by 
\begin{equation}\label{e:complexPairing}
\hat \eta(k)\eqdef \eta(e_{-k})\,.%=\eta(a_k)-i \eta(b_k)
\end{equation} 
Note that $\overline{\hat\eta(k)}=\hat\eta(-k)$. Since the zero mode $\hat\eta(0)$ of the solution is automatically zero, 
we will only care about $\hat\eta(k)$ for $k\in\Z^d_0\eqdef\Z^d\setminus\{0\}$. 
Moreover, we recall that the Laplacian $\Delta$ on $\T^d$ has eigenfunctions $\{e_k\}_{k \in \Z^d}$ 
with eigenvalues $\{-|k|^2\,:\,k\in\Z^d\}$, so that the operator $(-\Delta)^{\frac12}$ in \eqref{eq:toroN} is defined
by its action on the basis elements 
	\begin{equation}\label{e:fLapla}
	(-\Delta)^{\frac12} e_k\eqdef |k|e_k\,,\qquad k\in\Z^d\,.
	\end{equation}	
        In particular, $(-\Delta)^{ \frac12}$ is an invertible linear
        bijection on distributions with null $0$-th Fourier mode.
We denote by $H^1(\T^d)$ the space of mean-zero functions $\phi$ such that the norm 
\begin{equ}
\|\phi\|_{H^1(\T^d)}^2\eqdef\|(-\Delta)^{\frac12}\phi\|_{L^2(\T^d)}^2\eqdef\sum_{k\in\Z^d_0} |k|^2 |\hat\phi(k)|^2
\end{equ} 
 is finite.  
\medskip 

Let $(\Omega,\cF,\P)$ be a complete probability space and 
$\eta$ be a mean-zero spatial white noise on the $d$-dimensional torus $\T^d$, i.e. 
$\eta$ is a Gaussian field with covariance
\begin{equ}\label{eq:spatial:white:noise}
\E[\swn(\vphi)\swn(\psi)]=\langle \vphi, \psi\rangle_{L^2(\T^d)}
\end{equ}
where $\varphi,\psi\in H\eqdef L^2_0(\T^d)$, 
the space of square-integrable functions with $0$ total mass, 
and $\langle\cdot,\cdot\rangle_{L^2(\T^d)}$ is the usual scalar product in $L^2(\T^d)$. 
For $n\in\N$, let $\SH_n$ be the {\it $n$-th homogeneous Wiener chaos}, i.e. the closed linear subspace of 
$L^2(\Omega)$ generated by the random variables $H_n(\eta(h))$, 
where $H_n$ is the $n$-th Hermite polynomial, and 
$h\in H$ has norm $1$. By~\cite[Theorem 1.1.1]{Nualart2006}, $\wc_n$ and $\wc_m$ are orthogonal whenever 
$m\neq n$ and $L^2(\Omega)=\bigoplus_{n}\SH_n$. 
Moreover, there exists a canonical contraction $\sint{}:\bigoplus_{n\ge 0} L^2(\T^{2n}) \to L^2(\Omega)$, 
which restricts to an isomorphism $\sint{}:\fock \to L^2(\Omega)$ on the Fock space $\fock:=\bigoplus_{n \ge 0} \fock_n$, 
where $\fock_n$ denotes the space $L_\sym^2(\T^{dn})$ of functions in $L^2_0(\T^{dn})$ which are symmetric
with respect to permutation of variables. The restriction $\sint_n$ of $\sint$ to $\fock_n$, 
is itself an isomorphism from $\fock_n$ to $\SH_n$
and, by~\cite[Theorem 1.1.2]{Nualart2006}, for every $F\in L^2(\Omega)$ there exists a family of kernels 
$(f_n)_{n\in\N}\in\fock$ such that $F=\sum_{n\geq 0} I_n(f_n)$ and 
\begin{equ}[e:Isometry]
\E[F^2]\eqdef\|F\|^2 = \sum_{n\geq 0} n!\|f_n\|_{L^2(\T^{dn})}^2
\end{equ}
and we take the right hand side as the definition of the scalar product on $\fock$, i.e. 
\begin{equ}[e:ScalarPFock]
 \langle f,g\rangle=\sum_{n\geq 0} \langle f_n,g_n\rangle \eqdef \sum_{n\geq 0} n!\langle f_n,g_n\rangle_{L^2(\T^{dn})}.
\end{equ}

\begin{remark}\label{rem:NotAbuse}
In view of the isometry between $\fock$ and $L^2(\Omega)$, throughout the paper, we will abuse notation 
and denote with the same symbol operators acting on $\fock$ and their composition 
with the isometry $I$, which is an operator acting on $L^2(\Omega)$. 
%More precisely, 
%if $L^2(\Omega)\ni F=\sum_{n\geq 0} I_n(f_n)$ and $\cO$ is an operator on $L^2(\Omega)$ with $\cO_n$ 
%defined on $\SH_n$, then we will write $\cO F=\sum_n \cO_n I_n(f_n)=\sum_n I_n(\cO_n f_n)$. 
\end{remark}

For $F(\eta)=f(\eta(\phi_1),\dots,\eta(\phi_n))$, $f\colon\R^n\to\R$ smooth and growing at most polynomially at infinity and 
$\phi_1,\dots,\phi_n\in L^2_0(\T^d)$, we define the Malliavin derivative $DF$ according to 
\begin{equ}[e:MalDer]
DF(\eta)(\cdot)\eqdef \sum_{i=1}^n \partial_i f(\eta(\phi_1),\dots,\eta(\phi_n))\phi_i(\cdot)
\end{equ}
and, for $k\in\Z_0^d$, denote its Fourier transform by 
\begin{equation}\label{e:Malliavin}
D_{k} F(v)\eqdef \cF(DF(v))(k)=\langle DF, e_{-k}\rangle\,.
\end{equation}
Notice that, for $p\in\Z$, the action of $D_p$ on $f\in\fock_n$ is 
\begin{equ}[e:MallDerFock]
\cF(D_pf)(k_{1:n-1})=n \hat f_n(k_{1:n-1},p)\,,
\end{equ}
so that in particular $D_p f\in\fock_{n-1}$. 
At last, we also recall the following integration by parts formula on Wiener space, i.e. 
\begin{equation}\label{eq:intbyparts}
\E[G D_kF]=\E[G\langle DF, e_{-k}\rangle]= \E[-F D_kG + FG\hat{\eta}(k)]\,. 
\end{equation}


Throughout the paper, we will write $a\lesssim b$ if there exists a constant $C>0$ such that $a\leq C b$. 
%and $a\sim b$ if $a\lesssim b$ and $b\lesssim a$. 
We will adopt the previous notations only in case in which 
the hidden constants do not depend on any quantity which is relevant for the result. 
When we write $\lesssim_T$ for some quantity $T$, it means that the constant $C$ implicit in the bound depends on $T$.


\section{The equation and its generator}\label{sec:EqandGen}

In this section, we first collect a number of preliminary properties of the solution $\eta^\eps$ to 
the regularised Burgers equation~\eqref{e:BurgersScaled} and its generator (Section~\ref{sec:Burgers}). 
Then, we carry out a detailed analysis of the generator equation (a.k.a. resolvent equation) and 
obtain the main estimates we will need in order to determine the large--scale behaviour of $\eta^\eps$. 


\subsection{Multidimensional Burgers generator}\label{sec:Burgers}

As in Section~\ref{sec:ideas}, let us first write~\eqref{e:BurgersScaled} in its weak formulation. 
For $\phi\in\cS(\T^d)$ and $t\geq 0$, it reads 
\begin{equ}[e:SPDEweak]
\eta^\eps_t(\phi)-\eta^\eps_0(\phi)=\int_0^t\eta^\eps_s(\Delta\phi)\dd s+\int_0^t\cN^\eps_\phi(\eta^\eps)\dd s +\int_0^t\xi(\dd s, (-\Delta)^{1/2}\phi)
\end{equ}
where $\eta^\eps_0$ is the initial condition, $\xi$ is a space-time white noise so that 
\begin{equ}
\xi(\dd t, (-\Delta)^{1/2}\phi)=\sum_{k\in\Z^d_0} |k|\hat\phi(k)\dd B_t(k)
\end{equ}
for $B(k)$'s complex valued Brownian motions satisfying 
$\overline{B(k)}= B(-k)$ and $\dd\langle B(k), B(\ell) \rangle_t =  \mathds{1}_{\{k+\ell=0\}}\, \dd t$, 
and $\NN_\phi(\eta)$ is the nonlinearity tested against $\phi$, i.e.
\begin{equs}[e:nonlin1]
\NN_\phi(\eta)&\eqdef \lambda_\eps \fw\cdot\Pi_{1/\eps}\nabla(\Pi_{1/\eps}\eta)^2(\phi)\\
&=\frac{\iota}{(2\pi)^{\frac{d}2}}\lambda_\eps\sum_{\ell,m}  \indN{\ell,m}\fw\cdot(\ell+m)\hat\phi(-\ell-m)\hat\eta(\ell)\hat\eta(m)
\end{equs}
with $\iota=\sqrt{-1}$ and 
\begin{equ}[eq:J]
    \indN{\ell, m}= \1\{0<\eps |\ell|_\infty\le 1, 0<\eps|m|_\infty\le 1,0<\eps|\ell+m|_\infty\le 1\}\,.
\end{equ}
In accordance with Remark \ref{rem:pignolerie}, for $d=2$ the sup-norm 
$|\cdot|_\infty$ in~\eqref{eq:J} is replaced by the Euclidean norm $|\cdot|$ instead.

From~\eqref{e:SPDEweak} follows that 
the Fourier modes $\hat \eta^\eps(k)$, $|k|\ge \eps^{-1}$ evolve like
independent (and $\eps$-independent) Ornstein-Uhlenbeck processes.

\begin{lemma}
\label{lem:generalproperties}
For every deterministic initial condition $\eta_0$, the solution $t\mapsto \eta^\eps_t$ 
of \eqref{e:SPDEweak} exists globally in time and is a strong Markov process. 
The generator $\gen$ of $\eta^\eps$ can be written as
$\gen=\gensy+\gena$ and the action of $\gensy$ and $\gena$ on smooth cylinder functions $F$
is given by 
 \begin{align}
 (\gensy F)(\eta) &\eqdef \frac{1}{2} \sum_{k \in \Z^d_0} |k|^2 (-\hat\eta(-k) D_k  +  D_{-k}D_k )F(\mu) \label{e:gens}\\
   (\gena F)(\eta) &\eqdef \frac\iota{(2
    \pi)^{d/2}}\lambda_\eps \sum_{m,\ell \in \Z_0^d} \indN{\ell,m} \fw\cdot (\ell+m) \hat \eta(m) \hat \eta(\ell) D_{-m-\ell} F(\eta). \label{e:gena}
 \end{align}
 Moreover, $\gensy$ is symmetric with respect to $\mathbb P$ while $\gena$ is skew-symmetric.
Finally, the law $\mathbb P$ of the average-zero space white noise is stationary.
\end{lemma}
\begin{proof}
Very similar arguments were provided in a number of references for equations which 
share features similar to those of~\eqref{e:SPDEweak}, e.g. \cite{CES, GubinelliJara2012} and, 
more comprehensively,~\cite{Gubinelli2016}, so 
we will limit ourselves to sketch some of the proofs. 

For the global in time existence of~\eqref{e:SPDEweak} for fixed $\eps$, we refer the reader to \cite[Prop. 3.4]{CES} 
or~\cite[Section 3]{GubinelliJara2012}, while the expressions \eqref{e:gens} and \eqref{e:gena} 
can be easily derived by It\^o's formula and the definition of Malliavin derivative in~\eqref{e:MalDer}.

To see that $\mathbb P$ is stationary, we start by noting that
stationarity for the linear equation with $w=0$ is well-known (and
easy to check), see e.g.~\cite[Section 2.3]{Gubinelli2016}, 
where it is further shown that $\gensy$ is symmetric with respect to $\P$. 
Hence, by~\cite{Ech}, it suffices to check that
$\E [\gena F] = 0$ for every cylinder function $F$. In fact, 
this follows immediately if we prove that $\gena$ is skew-symmetric, 
since, if so, $\E [\gena F]=-\E [F\gena 1]=0$ as the Malliavin derivative of the constant random variable $F=1$ is $0$. 
Now, to prove skew-symmetry for $\gena$, we apply Gaussian integration by parts~\eqref{eq:intbyparts} 
and get 
\begin{equs}
\E [G \gena F]=&-\E[F\gena G]\\
&+ \frac{\iota\lambda_\eps}{(2
 \pi)^{d/2}} \sum_{m,\ell \in \Z_0^d}\indN{\ell,m} \fw\cdot (\ell+m) \E \Big[\hat\eta(-\ell-m)\hat\eta(m)\hat\eta(\ell) F(\eta)G(\eta)\Big]\\
 =& -\mathbb E F \gena G +\E\Big[ \langle \cN^\eps_\phi(\eta),\eta\rangle_{L^2(\T^d)}  F(\eta)G(\eta)\Big]
\end{equs}
so that, since 
\begin{equ}[e:InvMeas]
\langle \cN^\eps_\phi(\eta),\eta\rangle_{L^2(\T^d)}=\lambda_\eps\langle \fw\cdot\nabla(\Pi_{1/\eps}\eta)^2, (\Pi_{1/\eps}\eta)\rangle=\frac{\lambda_\eps}{3}\langle \fw\cdot\nabla(\Pi_{1/\eps}\eta)^3, 1\rangle=0
\end{equ}
the proof is completed. 
\end{proof}

Next, we want to determine the action of $\gensy$ and $\gena$ on the Fock space $\fock$. 
To lighten notations, for a set of integers $I$, we will denote by $k_I$ 
the vector $(k_i)_{i\in I}$ and 
\begin{equ}[e:notations] 
|k_I|^2\eqdef \sum_{i\in I}|k_i|^2\qquad\text{and}\qquad k_{1:n}\eqdef k_{\{1:n\}}\,,
\end{equ}
where  $\{1:n\}\eqdef \{1,\dots,n\}$.

\begin{lemma}\label{lem:essid}
Let $\eps>0$, $\gensy$ and $\gena$ be the operators defined according to~\eqref{e:gens} and~\eqref{e:gena} 
respectively. Then, $\gena$ can be written as
  \begin{equ}[eq:split]
    \gena=\genap+\genam,\qquad\text{with}\qquad \genam=-(\genap)^*\,.
  \end{equ}
Moreover, for any $n\in\N$, the action of the operators $\gensy,\genam,\genap$ on $f\in\fock_n$ 
is given by
\begin{equs}
\cF(\gensy f) (k_{1:n})=&-\tfrac12|k_{1:n}|^2\hat f(k_{1:n})\\
\cF(\genap f) (k_{1:n+1})=&-\frac{2\iota}{(2\pi)^{d/2}} \frac{\lambda_\eps}{n+1}\times\label{e:gen}\\
&\times\sum_{1\le i<j\le n+1} \,[\fw\cdot (k_i+k_j)]\indN{k_i,k_j}\hat f(k_i+k_j,k_{\{1:n+1\}\setminus\{i,j\}})\\
\cF(\genam  f) (k_{1:n-1})=&\frac{2\iota}{(2\pi)^{d/2}} n\lambda_\eps\,\sum_{j=1}^{n-1}(\fw\cdot k_j) \sum_{\ell+m=k_j}\indN{\ell,m}\hat f(\ell,m,k_{\{1:n-1\}\setminus\{j\}}).
\end{equs}
and, if $n=1$, then $\genam f$ is identically $0$. 
\end{lemma}
\begin{proof}
The proof is identical to that of~\cite[Lemma 3.5]{CES}, to which we refer the interested reader. 
\end{proof}

In order to state the following lemma, we introduce the so-called {\it number operator}.

\begin{definition}\label{def:NoOp}
We define the {\it number operator} $\cN\colon \fock\to\fock$ to be the linear diagonal operator acting on $\psi=(\psi_n)_n\in\fock$ 
as $(\cN\psi)_n\eqdef n \psi_n$. 
\end{definition}
 
Our analysis will need some preliminary estimates on the asymmetric part of the generator, 
which corresponds to proving the so-called {\it graded sector condition} in~\cite[Section 2.7.4]{Komorowski2012}. 

\begin{lemma}\label{l:GenaBound}
For $d\ge2$ here exist a constant $C=C(d)>0$ such that for every $\psi\in\fock$ the following estimate holds
\begin{equ}\label{e:Abound}
\|(-\gensy)^{-\frac12}\gena_{\sigma} \psi\|\leq C \|\sqrt{\cN}(-\mathcal L_0)^\frac12 \psi\|\,,
\end{equ}
for $\sigma\in\{+,-\}$. 
In particular, for $\sigma\in\{+,-\}$, it  implies 
\begin{equ}\label{e:Tbound}
\|(-\gensy)^{-\frac12}\gena_{\sigma} (-\gensy)^{-\frac12} \psi\|\leq C \|\sqrt{\cN} \psi\|\,.
\end{equ}
\end{lemma}
\begin{proof}
Clearly, once we establish~\eqref{e:Abound},~\eqref{e:Tbound} follows immediately 
by choosing $\psi$ therein to be $(-\gensy)^{-1/2}\rho$ for $\rho\in\fock$.

For~\eqref{e:Abound}, we claim that we only need to prove that for any $\psi\in\fock_n$ and $\rho\in\fock_{n+1}$, we have
\begin{equ}[e:AprioriFirst]
|\langle\rho,\genap \psi\rangle|\lesssim \sqrt{n}\Big(\gamma\|(-\gensy)^{1/2}\psi\|^2+\frac{1}{\gamma} \|(-\gensy)^{1/2}\rho\|^2\Big)\,.
\end{equ}
We will show~\eqref{e:AprioriFirst} at the end. 
Assuming it holds, we first derive~\eqref{e:Abound} for $\genap$. 
The variational characterisation of the $(-\gensy)^{-1/2}\fock$-norm gives 
\begin{equs}
\|&(-\gensy)^{-1/2}\genap\psi\|^2=\sup_{\rho\in \fock_{n+1}}\left(2\langle \rho,\genap  \psi\rangle-
    \| (-\cL_0)^{\frac12}\rho\|^2\right)\\
&\leq \sup_{\rho\in \fock_{n+1}}\left(2C_0\sqrt{n}\Big(\gamma \|(-\gensy)^{1/2}\psi\|^2 + \frac{1}{\gamma}\|(-\gensy)^{1/2}\rho\|^2\Big)-
    \| (-\cL_0)^{\frac12}\rho\|^2\right)\\
    &\lesssim n\|(-\gensy)^{1/2}\psi\|^2=\|\sqrt{\cN} (-\gensy)^{1/2}\psi\|^2
\end{equs}
where in the first step we used~\eqref{e:AprioriFirst} (and $C_0$ is the universal 
constant implicit in that inequality) while in the last we chose $\gamma\eqdef 2C_0\sqrt{n}$.

For $\genam$ instead, we use that, by~\eqref{eq:split}, $\genam=-(\genap)^\ast$. Invoking
again the variational formula above, we deduce
\begin{equs}
\|&(-\gensy)^{-1/2}\genam\psi\|^2=\sup_{\rho\in \fock_{n-1}}\left(2\langle \rho,\genam  \psi\rangle-
    \| (-\cL_0)^{\frac12}\rho\|^2\right)\\
 &=\sup_{\rho\in \fock_{n-1}}\left(-2\langle \genap\rho, \psi\rangle-
    \| (-\cL_0)^{\frac12}\rho\|^2\right)\\
&\leq \sup_{\rho\in \fock_{n-1}}\left(2C_0\sqrt{n}\Big(\gamma \|(-\gensy)^{1/2}\rho\|^2 + \frac{1}{\gamma}\|(-\gensy)^{1/2}\psi\|^2\Big)-
    \| (-\cL_0)^{\frac12}\rho\|^2\right)\\
    &\lesssim n\|(-\gensy)^{1/2}\psi\|^2=\|\sqrt{\cN} (-\gensy)^{1/2}\psi\|^2
\end{equs}
where this time we chose $\gamma\eqdef 1/(2C_0\sqrt n)$.  
\medskip

It remains to prove~\eqref{e:AprioriFirst}. By definition of $\genap$, we have
\begin{equs}
|\langle\rho,\genap \psi\rangle|=&n!\frac{2}{(2\pi)^{\frac{d}2}} \lambda_\eps\Big|\sum_{k_{1:n+1}}\hat\rho(-k_{1:n+1})\times\\
&\times\sum_{1\le i<j\le n+1} [\fw\cdot (k_i+k_j)]\indN{k_i,k_j}\hat \psi(k_i+k_j,k_{\{1:n+1\}\setminus\{i,j\}})\Big|%(\ell_{1:p+1}, k_{1:n+2j\setminus\{i_1:i_{p+1}\}})\Big|\,.
\\
=&\frac{(n+1)! n}{(2\pi)^{\frac{d}2}}  \lambda_\eps\Big|\sum_{k_{1:n+1}}\hat\rho(-k_{1:n+1})[\fw\cdot(k_1+k_2)]\indN{k_1,k_2} \hat\psi(k_1+k_2, k_{3:n+1})\Big|\,.
\end{equs}
Let us look at the sum over $k_1,k_2$. This equals
\begin{equs}
&\Big|\sum_{k_{1:2}}\hat\rho(-k_{1:n+1})[\fw\cdot (k_1+k_2)]\indN{k_1,k_2} \hat\psi(k_1+k_2, k_{3:n+1})\Big|\\
&=\Big|\sum_{q} (\fw\cdot q)\hat\psi(q, k_{3:n+1})\sum_{k_1+k_2=q}\indN{k_1,k_2} \hat\rho(-k_{1:n+1})\Big|\\
&\lesssim \Big(\sum_q (\fw\cdot q)^2|\hat\psi(q, k_{3:n+1})|^2\Big)^{1/2}\Big(\sum_q \Big|\sum_{k_1+k_2=q}\indN{k_1,k_2} \hat\rho(-k_{1:n+1})\Big|^2\Big)^{1/2}\\
&=:A\times\one\label{e:Adef}
\end{equs}
We leave $A$ as it stands and focus on $\one$. Note that this can be bounded by 
\begin{equs}
\one&=\Big(\sum_q \Big|\sum_{k_1+k_2=q}\indN{k_1,k_2} \hat\rho(-k_{1:n+1})\Big|^2\Big)^{1/2}\\
&\leq \Big(\sum_q \sum_{k_1+k_2=q}|k_{1:2}|^2|\hat\rho(-k_{1:n+1})|^2\sum_{k_1+k_2=q}\frac{\indN{k_1,k_2}}{|k_{1:2}|^2}\Big)^{1/2}\\
&\lesssim \lambda_\eps^{-1} \Big(\sum_q \sum_{k_1+k_2=q}|k_{1:2}|^2|\hat\rho(k_{1:n+1})|^2\Big)^{1/2}=:\lambda_\eps^{-1} B\label{e:Bdef}
\end{equs}
where we used that for $d=2$, $\lambda_\eps^2=(\log\eps^{-2})^{-1}$ and 
\begin{equ}[e:CrucialBound2]
\lambda_\eps^2\sum_{k_1+k_2=q}\frac{\indN{k_1,k_2}}{|k_{1:2}|^2}\leq \frac{\eps^2}{\log\eps^{-2}}\sum_{\eps<|\eps k_1|\leq 1}\frac{1}{|\eps k_{1}|^2}\lesssim  \frac{1}{\log\eps^{-2}}\int_{|x|\in[\eps,1]}\frac{\dd x}{|x|^2}\lesssim 1
\end{equ}
while for $d\geq 3$, $\lambda_\eps^2=\eps^{d-2}$ and 
\begin{equs}[e:CrucialBound3+]
\eps^{d-2}\sum_{k_1+k_2=q}\frac{\indN{k_1,k_2}}{|k_{1:2}|^2}&\leq\eps^{d}\sum_{|\eps k_1|\leq 1}\frac{1}{|\eps k_{1}|^2}\lesssim \int_{|x|\leq 1}\frac{\dd x}{|x|^2}\lesssim 1\,.
\end{equs}
As a consequence, for any $\gamma>0$, we have 
\begin{equs}
|\langle\rho,\genap \psi\rangle|&\lesssim (n+1)! n\sum_{k_{3:n+1}}AB\leq (n+1)! n\Big(\frac{\gamma}{\sqrt n} \sum_{k_{3:n+1}} A^2 +\frac{\sqrt n}{\gamma}\sum_{k_{3:n+1}} B^2\Big)\\
&\leq (n+1)! n\Big(\frac{1}{n!}\frac{\gamma}{n^{3/2}}\|(-\gensy)^{1/2}\psi\|^2 +\frac{1}{(n+1)!}\frac{n^{-1/2}}{\gamma}\|(-\gensy)^{1/2}\rho\|^2\Big)\\
&\leq \sqrt n\gamma\|(-\gensy)^{1/2}\psi\|^2 +\frac{\sqrt n}{\gamma}\|(-\gensy)^{1/2}\rho\|^2\label{e:AlmostDoneBound}
\end{equs}
where we used that, by the definition of $A$ and $B$ in~\eqref{e:Adef} and~\eqref{e:Bdef} respectively, 
we have 
\begin{equ}
\sum_{k_{3:n+1}} A^2\lesssim\sum_{k_{1:n}}  |k_1|^2|\hat\psi(k_{1:n})|^2\leq\frac{n^{-1}}{n!} \|(-\gensy)^{1/2}\psi\|^2
\end{equ}
and similarly for $B$. Then,~\eqref{e:AprioriFirst} follows. % upon choosing $\alpha=1/2$.
\end{proof}

For $d=2$ and $d\geq 3$, we need some further but distinct estimates, which we separately give in the 
two sections below. 

\subsubsection{Integrability estimates: $d\geq 3$}

In the supercritical regime $d\geq 3$, in order to identify the limit process 
we will resort to use a sort of {\it uniform integrability condition}. 
This is better expressed in terms of modified $\genap$ and $\genam$ operators, 
which we now introduce. 
For $\delta\geq 0$, let $\cA^{\eps,\delta}_{\pm}$, be the operators 
whose action on $f\in\fock_n$ is given by 
\begin{equs}[e:Adelta]
\cF(\cA_+^{\eps,\delta} f) (k_{1:n+1})&= \eps^{d/2}\Big(\frac{2\eps^{-1}}{(2\pi)^{d/2}(n+1)}\Big)^{1+\delta}\\
&\times\sum_{1\le i<j\le n+1} \,|\fw\cdot(k_i+k_j)|^{1+\delta}\indN{k_i,k_j}\hat f(k_i+k_j,k_{\{1:n+1\}\setminus\{i,j\}})\\
\cF(\cA_-^{\eps,\delta}  f) (k_{1:n-1})&= \eps^{d/2}\Big(\frac{n\eps^{-1}}{(2\pi)^{d/2}}\Big)^{1+\delta}\\
&\times\sum_{j=1}^{n-1}|\fw\cdot k_j|^{1+\delta} \sum_{\ell+m=k_j}\indN{\ell,m}\hat f(\ell,m,k_{\{1:n-1\}\setminus\{j\}})\,.
\end{equs} 
%Let us point out that the operators $\cA^{\eps,\delta}_{\pm}$ have the same structure as $\gena_\pm$ 
%but in which all the quantities characterising the action of the  are raised  
In the next lemma, we establish suitable bounds on the previous operators, which 
are similar in flavour to those obtained in Lemma~\ref{l:GenaBound} for $\gena$. 

\begin{lemma}\label{l:GenaDBound}
Let $d\geq 3$. Then, there exists a $C=C(d)>0$ such that for any function $\psi\in\fock$ 
and $\delta\leq (d-2)/2$, the following estimate holds
\begin{equ}
\|(-\gensy)^{-\frac{1+\delta}{2}}\cA^{\eps,\delta}_{\sigma} \psi\|\leq C \|\cN^{1+\delta}(-\gensy)^{\frac{1+\delta}{2}} \psi\|\,,\label{e:Adeltabound}
\end{equ}
where $\sigma\in\{+,-\}$ and $\cN$ is the number operator given in Definition~\ref{def:NoOp}.
In particular, for $\sigma\in\{+,-\}$, the previous  implies
\begin{equs}
\|(-\gensy)^{-\frac{1+\delta}{2}}\cA^{\eps,\delta}_{\sigma}(-\gensy)^{-\frac{1+\delta}{2}} \psi\|&\leq C \|\cN ^{1+\delta}\psi\|\,.\label{e:Ubound}
\end{equs}
\end{lemma}
With respect to Lemma~\ref{l:GenaBound}, here we do not aim
  at the best possible estimate in terms of growth with respect to $\cN$,
  and the exponent $1+\delta$ is certainly not optimal. In fact, for
  the purposes of Section~\ref{sec:Diff3} (which is where the Lemma is used),
  the precise $\cN$-dependence is irrelevant.
\begin{proof}
As in the proof of Lemma~\ref{l:GenaBound}, it suffices to show~\eqref{e:Adeltabound}. 
We begin with $\cA^{\eps,\delta}_{+}$. We first apply Cauchy-Schwarz to the sums in $i,j$ and relabel the variables 
so that, for $\psi\in \Gamma L^2_{n}$, we obtain
\begin{equs}\label{e:H1A+}
    &\|(-\gensy)^{-\frac{1+\delta}{2}}\cA^{\eps,\delta}_{+} \hat\psi\|^2\\
    &\lesssim n! n^3\sum_{k_{1:n}}|k_1|^{2+2\delta} |\hat\psi(k_{1:n})|^2 \, \eps^{d-2(1+\delta)}\sum_{\ell+m=k_1}  \frac{\indN{\ell,m}}{(|\ell|^2+|m|^2+|k_{2:n}|^2)^{1+\delta}}\\
    &\lesssim n! n^2\sum_{k_{1:n}}(|k_{1:n}|^2)^{1+\delta} |\hat\psi(k_{1:n})|^2 \, \eps^{d-2(1+\delta)}\sum_{\ell+m=k_1}  \frac{\indN{\ell,m}}{(|\ell|^2+|m|^2)^{1+\delta}}\,.
  \end{equs}
On the other hand, arguing as in~\eqref{e:CrucialBound3+},
  \begin{equ}[eq:thesum]
    \eps^{d-2(1+\delta)}\sum_{\ell+m=k_1}  \frac{\indN{\ell,m}}{(|\ell|^2+|m|^2)^{1+\delta}}\leq\eps^{d}\sum_{|\eps\ell|\leq 1} \frac{}{|\eps\ell|^{2+2\delta}}\lesssim_d  \int_{|x|\leq 1} \frac{\dd x}{|x|^{2+2\delta}}
  \end{equ}
  with the right hand side independent of $k_1$, so that~\eqref{e:Adeltabound} follows provided that $\delta<(d-2)/2$. 

Next, we turn to $\cA^{\eps,\delta}_{-}$. From its expression in~\eqref{e:Adelta}, we have for $\psi\in \Gamma L^2_{n}$
  \begin{equs}
    \|&(-\gensy)^{-\frac{1+\delta}{2}}\cA^{\eps,\delta}_{-} \psi\|^2\\
    &\lesssim \eps^{d-2(1+\delta)} (n-1)! n^{4+2\delta}\sum_{k_{1:n-1}}\frac{|k_1|^{2(1+\delta)}}{|k_{1:n-1}|^{2(1+\delta)}}\Big(
    \sum_{\ell+m=k_1}  \indN{\ell,m} \hat\psi(\ell,m,k_{2:n-1})
    \Big)^2\\
    &= \eps^{d-2(1+\delta)} (n-1)! n^{4+2\delta}\sum_{k_{1:n-1}}\frac{| k_1|^{2(1+\delta)}}{|k_{1:n-1}|^{2(1+\delta)}}\times\\
    &\qquad\qquad\times\Big(
    \sum_{\ell+m=k_1}  \indN{\ell,m}\frac{(|\ell|^2+|m|^2)^{\frac{1+\delta}2}\hat\psi(\ell,m,k_{2:n-1})}{(|\ell|^2+|m|^2)^{\frac{1+\delta}2}}  \Big)^2\\
    &\le n! n^{3+2\delta}  \sum_{k_{1:n}}(|k_1|^2+|k_2|^2)^{1+\delta} |\hat\psi(k_{1:n+1})|^2\Big(\sum_{k_1+k_2=q}\frac{ \eps^{d-2(1+\delta)} \indN{\ell,m}}{(|\ell|^2+|m|^2)^{1+\delta}}\Big)\\
%    &\le n! n^{3+2\delta}\eps^{d-2(1+\delta)}  \sup_q\Big(\sum_{\ell+m=q}\frac{ \indN{\ell,m}}{(|\ell|^2+|m|^2)^{1+\delta}}\Big)\sum_{k_{1:n}}(|k_1|^2+|k_2|^2)^{1+\delta} |\hat\psi(k_{1:n+1})|^2\\
    &\lesssim_d n! n^{2+2\delta} \sum_{k_{1:n}}(|k_{1:n}|^2)^{1+\delta} |\hat\psi(k_{1:n})|^2= n^{2+2\delta} \|(-\cL_0)^{\frac{1+\delta}2}\psi\|^2
  \end{equs}
where again we used \eqref{eq:thesum}. The proof of~\eqref{e:Adeltabound} is complete. 
%
%We now turn to~\eqref{e:A+H-2}. Arguing as in~\eqref{e:H1A+}, we have 
%\begin{equs}
%\|(-\gensy)^{-\frac12}\genap \psi\|^2&\lesssim n! n^2\sum_{k_{1:n}}|k_{1:n}|^2 |\psi(k_{1:n})|^2 \, \lambda_\eps^2\sum_{\ell+m=k_1} \frac{ \indN{\ell,m}}{(|\ell|^2+|m|^2)^2},
%\end{equs}
%so that, since 
%\begin{equ}
% \lambda_\eps^2\sum_{\ell+m=k_1}  \indN{\ell,m}\frac1{(|\ell|^2+|m|^2)^2}\lesssim_d f(\eps)
%\end{equ}
%for $f$ as in~\eqref{e:f}, \eqref{e:A+H-2} follows. 
\end{proof}

\subsubsection{The replacement lemma: $d=2$}\label{sec:ReplLemma}

The reason why the estimates of Lemma~\ref{l:GenaDBound} were possible is that, when suitably scaled, 
the inverse Laplacian in dimension $d\geq 3$ is integrable close to the origin, which in 
our context means that the estimate~\eqref{eq:thesum} holds for $\delta$ sufficiently small but positive. 
In $d=2$, this is not the case. To overcome the lack of integrability, we devise an alternative route, 
which (morally) approximates the inverse of the generator with a linear 
and  {\it diagonal} (both in Fourier and in the chaos) operator, 
given by a non-trivial order $1$ perturbation of $\gensy$. 

To be more precise, we need to introduce some notation. For $\eps>0$, let $\Ll$ be the function 
defined on $[1/2,\infty)$ as 
\begin{equ}[e:L]
\Ll(x)\eqdef \lambda_\eps^2 \log \left(1+\frac{1}{\eps^2 x}\right)\,.
\end{equ}
%where $\fc$ is the constant given in~\eqref{}. 
Further, let $G$ be the function
\begin{equ}
  \label{eq:defG}
G(x)\eqdef \frac1{|\fw|^2}\left[\left(\frac{3|\fw|^2}{2\pi} x+1\right)^{\frac{2}{3}}-1\right],
\end{equ}
and $\cGN$ be the operator on $\fock$ given by 
\begin{equ}[e:G]
\cGN\eqdef G(\Ll(-\gensy))
\end{equ}
which means that for $\psi\in\fock_n$, $n\in\N$, the action of $\cGN$ on $\psi$ is 
\begin{equ}
\cF(\cGN\psi)(k_{1:n})\eqdef G(\Ll(\tfrac12|k_{1:n}|^2))\hat\psi(k_{1:n})\,,\qquad k_{1:n}\in\Z^{2n}_0\,.
\end{equ}
We are now ready to state the main result of this section, namely, the replacement lemma. 
Recall the definition of the scalar product $\langle\cdot,\cdot\rangle$ on $\fock$ given in~\eqref{e:ScalarPFock}.


\begin{lemma}[Replacement Lemma]\label{l:Replacement}
There exists a constant $C>0$ such that for every $\psi_1,\psi_2\in\fock$ we have 
\begin{equs}\label{e:Replacement}
|\langle \big[-\genam(-\gensy-\gensy^\fw\cGN)^{-1}&\genap+\gensy^\fw\cGN\big]\psi_1,\psi_2\rangle|\\
&\leq C\lambda_\eps^2 \|\cN(-\gensy)^{1/2}\psi_1\|\|\cN(-\gensy)^{1/2}\psi_2\|
\end{equs}
where $\cGN$ is the operator defined according to~\eqref{e:G}  and \eqref{eq:defG} and, for $n\in\N$, 
$\gensy^\fw$ is the operator acting on $\psi\in\fock_n$ as 
\begin{equ}[e:L0w]
\cF(\gensy^\fw\psi)(k_{1:n})=-\tfrac12(\fw\cdot k)^2_{1:n}\hat\psi(k_{1:n})\,,\qquad \text{for all $k_{1:n}\in\Z^{2n}_0$}
\end{equ}
and $(\fw\cdot k)_{1:n}^2\eqdef \sum_{i=1}^n(\fw\cdot k_i)^2$.
\end{lemma}
\begin{remark}\label{rem:G}
As it appears from the proof of Proposition~\ref{p:Approx} in the appendix, 
the form of the function $G$ in \eqref{e:L} is dictated by the
following fixed point equation 
\begin{equ}[e:IntegralG]
G(x)=\frac{1}{\pi}\int_0^x   \frac{\dd y}{\sqrt{1+|\fw|^2 G(y)}}\,.
\end{equ}
Now, if we had chosen a different regularisation for the nonlinearity, 
e.g. the $|\cdot|_\infty$-distance in~\eqref{eq:J} instead of the Euclidean one or a smooth regularisation 
instead of the Fourier cut-off in~\eqref{e:nonlin1}, then we would have not obtained an explicit $G$ 
but statement above would have remained true. 

Apart from working with an explicit function, the advantage of our choice lies in the fact that 
it allows us to make an explicit connection with the work of Yau~\cite{Yau} on the two-dimensional ASEP. 
Indeed, the exponent 
$2/3$ in~\eqref{eq:defG} precisely reflects that obtained therein and strongly suggests that, if $\lambda_\eps$ were 
taken to be a constant independent of $\eps$, then the diffusivity would diverge as $(\log t)^{2/3}$ for $t\to\infty$. 
\end{remark}
\begin{proof}
Notice first that, for any $n\in\N$, the operator $\genam(-\gensy-\gensy^\fw\cGN)^{-1}\genap$ maps $\fock_n$ 
into itself, so that, to establish~\eqref{e:Replacement}, it suffices to consider $\psi_1,\psi_2\in\fock_n$. Moreover, we can write
\begin{equ}[e:ScalProd]
  \langle \big[-\genam(-\gensy-\gensy^\fw\cGN)^{-1}\genap\big]\psi_1,\psi_2\rangle = \langle (-\gensy-\gensy^\fw\cGN)^{-1}\genap\psi_1,\genap\psi_2\rangle.
\end{equ}
Recalling the form of $\genap$ in~\eqref{e:gen}, we see that the scalar product at the right hand side can be split into a
{\it diagonal part}, corresponding to making the same choice for the indices $i,j$ in the two occurrences of $\genap$, 
an {\it off-diagonal part of type 1}, corresponding to choosing $i,j$ and $i',j'$ with one element in common, 
and {\it off-diagonal part of type 2}, corresponding to all remaining choices.  
(See for instance \cite[Lemma 3.6]{CET}, where a similar distinction was made.)

To write this in formulas, let $\cS^\eps \eqdef (-\gensy-\gensy^\fw\cGN)^{-1}$ and 
denote with $\sigma^\eps =(\sigma^\eps _n)_{n\ge1}$ its Fourier multiplier, i.e. for $\psi\in\Gamma L^2_n$ 
$\cF(\cS^\eps  \psi)(k_{1:n})=\sigma^\eps _n(k_{1:n})\hat \psi(k_{1:n})$ where 
\begin{equs}[e:sigman]
  \sigma^\eps _n(k_{1:n})&=\frac2{|k_{1:n}|^2+(\fw\cdot k)_{1:n}^2G(\Ll(\frac12|k_{1:n}|^2))}\,.
\end{equs}
Then, the scalar product in~\eqref{e:ScalProd} satisfies
\begin{equ}
	\langle \cS^\eps \genap\psi_1,\genap\psi_2\rangle=\langle \cS^\eps \genap\psi_1,\genap\psi_2\rangle_{\Di} +\sum_{i=1}^2 \langle \cS^\eps \genap\psi_1,\genap\psi_2\rangle_{\oD_i}
\end{equ}
where diagonal part, given by the first summand, is defined as 
\begin{equs}\label{e:Diag}
	\langle \cS^\eps \genap\psi_1,&\genap\psi_2\rangle_{\Di}\eqdef n!\, n \,\frac{4 \lambda_\eps^2}{(2\pi)^2} \times\\
				& \times\sum_{k_{1:n}}\tfrac12(\fw\cdot k_1)^2\overline{\hat\psi_1(k_{1:n})}{\hat\psi_2(k_{1:n})}\sum_{\ell+m=k_1}\sigma^\eps _{n+1}(\ell,m,k_{2:n})  \indN{\ell,m}
\end{equs}
while off-diagonal parts of type $1$ and $2$ are respectively given by 
\begin{equs}\label{e:OffDiag1}
	\langle &\cS^\eps \genap\psi_1,\genap\psi_2\rangle_{\oDi}\eqdef n! \,n\, c_{\oDi}(n)\sum_{k_{1:n+1}} \sigma^\eps _{n+1}(k_{1:n+1})     \indN{k_1,k_2}\indN{k_1,k_3}\times\\
	&\times(\fw\cdot (k_1+k_2))  (\fw\cdot (k_1+k_3))
			\overline{\hat{\psi_1}(k_1+k_2,k_3,k_{4:n+1})}{\hat\psi_2(k_1+k_3,k_2,k_{4:n+1})}\,,
\end{equs}
and
\begin{equs}\label{e:OffDiag2}
\langle &\cS^\eps \genap\psi_1,\genap\psi_2\rangle_{\ooDi}\eqdef n! \,n\, c_{\ooDi}(n)\sum_{k_{1:n+1}} \sigma^\eps _{n+1}(k_{1:n+1})     \indN{k_1,k_2}\indN{k_1,k_3}\times\\
			&\times (\fw\cdot (k_1+k_2))  (\fw\cdot (k_3+k_4))\overline{\hat \psi_1(k_1+k_2,k_{3:4},k_{5:n+1})}{\hat \psi_2(k_3+k_4,k_{1:2},k_{5:n+1})}
\end{equs}
where, for $i=1,\,2$, $c_{\oD_i}(n)$ is an explicit positive constant only depending on $n$ 
and such that $c_{\oD_i}(n)= O(n^{i+1})$.
\medskip

We are ready to control the left hand side of~\eqref{e:Replacement}. At first, we split all the diagonal parts 
from the off-diagonal 
\begin{equs}
|\langle &\big[-\genam(-\gensy-\gensy^\fw\cGN)^{-1}\genap+\gensy^\fw\cGN\big]\psi_1,\psi_2\rangle\label{e:Di}|\\
&\leq  \Big|\langle\cS^\eps \genap\psi_1, \genap\psi_2\rangle_{\Di}+ \langle \gensy^\fw\cGN\psi_1,\psi_2\rangle\Big|
%\qquad\qquad\qquad\qquad\qquad\qquad
+\sum_{i=1,2}|\langle\cS^\eps \genap\psi_1, \genap\psi_2\rangle_{\mathrm{off}_i}|\,.%\label{e:oDi}
\end{equs}
Let us begin by estimating the first summand. 
By~\eqref{e:G},~\eqref{e:L0w} and~\eqref{e:Diag}, we
see that
\begin{equs}
\Big|\langle&(-\gensy-\gensy^\fw\cGN)^{-1}\genap\psi_1, \genap\psi_2\rangle_{\Di}+ \langle \gensy^\fw\cGN\psi_1,\psi_2\rangle\Big|\\
&=n! n\Big|\sum_{k_{1:n}}\tfrac12(\fw\cdot k_1)^2\hat\psi_1(-k_{1:n})\hat\psi_2(k_{1:n})\Big(P^\eps(k_{1:n})- G(\Ll(\tfrac12|k_{1:n}|^2))\Big)\Big|
\end{equs}
where $P^\eps$ is defined as 
\begin{equ}[e:P1]
P^\eps(k_{1:n})\eqdef\frac{\lambda_\eps^2}{\pi^2} \sum_{\ell + m=k_1} \frac{\indN{\ell,m}}{\tilde\Gamma + \tilde\Gamma^wG(\Ll(\tilde\Gamma))}
\end{equ}
and, to shorten the notations, we wrote
\begin{equ}
\tilde\Gamma\eqdef\tfrac12(|\ell|^2+|m|^2+|k_{2:n}|^2)\,,\qquad\tilde\Gamma^w\eqdef \tfrac12((\fw\cdot\ell)^2+(\fw\cdot m)^2+(\fw\cdot k)^2_{2:n})\,.
\end{equ}
By Proposition~\ref{p:Approx}, 
we can upper bound the above by
\begin{equs}
&n! n\sum_{k_{1:n}}\tfrac12(\fw\cdot k_1)^2|\hat\psi_1(k_{1:n})||\hat\psi_2(k_{1:n})|\Big|P^\eps(k_{1:n})- G(\Ll(\tfrac12|k_{1:n}|^2))\Big|\\
&\lesssim \lambda_\eps^2 n! n \sum_{k_{1:n}}|k_1|^2|\hat\psi_1(k_{1:n})||\hat\psi_2(k_{1:n})|=\lambda_\eps^2 n! \sum_{k_{1:n}}|k_{1:n}|^2|\hat\psi_1(k_{1:n})||\hat\psi_2(k_{1:n})|\\
&\lesssim \lambda_\eps^2 \|(-\gensy)^{1/2}\psi_1\|\|(-\gensy)^{1/2}\psi_2\|\label{e:CSH1}
\end{equs}
so that this term satisfies~\eqref{e:Replacement}. 

We are now left to estimate the off diagonal terms. 
This is done by following {\it mutatis mutandis} the same steps as in the
proof of~\cite[Lemma 3.4]{CETWeak}, so we will only point out the necessary changes.
For the off-diagonal terms of type 1, we need to control
\begin{equs}[e:od1]
|\langle\cS^\eps \genap\psi_1, \genap\psi_2&\rangle_{\mathrm{off}_1}|\lesssim  n!c_{\oD_1}(n)\lambda_\eps^2\sum_{j_{1:3},k_{3:n}}|j_1+j_2|\, |j_1+j_3| \\
 & \times \overline{\hat \psi_1(j_1+j_2,j_3,k_{3:n})}
  \hat \psi_2(j_1+j_3,j_2,k_{3:n})\sigma^\eps _{n+1}(j_{1:3},k_{3:n})
\end{equs}
where $\sigma^\eps _{n+1}$, given in~\eqref{e:sigman}, is the Fourier multiplier of $\cS^\eps $ in $\fock_{n+1}$
and, as usual, we have bounded $|\fw\cdot k|\lesssim |k|$.
As in \cite[Eq. (3.21)]{CETWeak}, an application of Cauchy-Schwarz shows that \eqref{e:od1} is upper  bounded by
\begin{equ}[e:od12]
  \lambda_\eps^2c_{\oD_1}(n)\prod_{i=1}^2\Big(n!\sum_{k_{1:n}}|\hat\psi_i(k_{1:n})|^2 |k_1|^2 |k_2|\sum_{j_1+j_2=k_1}\frac1{|j_2|}\sigma^\eps _{n+1}(j_{1:2},k_{2:n})
  \Big)^{\frac12}\,.
\end{equ}
This is the same expression as in \cite[Eq. (3.21)]{CETWeak}, with $\sigma^\eps _{n+1}$ 
replacing $J^N$ there. Since $\sigma^\eps _{n+1}$ satisfies
\[
\sigma^\eps _{n+1}(j_{1:2},k_{2:n})\lesssim \frac1{|j_{1:2}|^2+|k_{2:n}|^2},
  \]
  which is the same estimate that $J^N$ satisfies, one can argue as in \cite{CETWeak} and deduce
  \begin{equ}[e:thesame]
\sum_{j_1+j_2=k_1}\frac1{|j_2|}\sigma^\eps _{n+1}(j_{1:2},k_{2:n})\lesssim \frac1{|k_{1:n}|}\le \frac1{|k_2|}.
\end{equ}
Now, since $c_{\oD_1(n)}\lesssim n^2$,~\eqref{e:od12} is upper bounded 
by a constant times $\lambda_\eps^2 \|\sqrt{\cN}(-\gensy)^{\frac12}\psi_1\|$ $\|\sqrt{\cN}(-\gensy)^{\frac12}\psi_2\|$.

As for the off-diagonal terms of type 2, we proceed similarly, 
recalling that, this time, $ c_{\oD_2}(n)\lesssim n^3$. Namely, 
applying Cauchy-Schwarz and exploiting once more~\eqref{e:thesame}, 
we bound them as
\begin{equs}
  |\langle&\cS^\eps \genap\psi_1, \genap\psi_2\rangle_{\mathrm{off}_2}|\\
  &\lesssim  \lambda_\eps^2c_{\oD_2}(n)\prod_{i=1}^2\Big(n!\sum_{k_{1:n}}|\hat\psi_i(k_{1:n})|^2 |k_1|^2 |k_2||k_3|\sum_{j_1+j_2=k_1}\frac{\sigma^\eps _{n+1}(j_{1:2},k_{2:n})}{|j_1||j_2|}
  \Big)^{\frac12}\\
  &\lesssim \lambda_\eps^2 n! c_{\oD_2}(n)\prod_{i=1}^2\Big(
    \sum_{k_{1:n}}|\hat\psi_i(k_{1:n})|^2\frac{|k_1| |k_2||k_3|}{|k_{1:n}|}
    \Big)^{1/2}\\
    &\lesssim \lambda_\eps^2 n^3 \prod_{i=1}^2\Big(
      n!    \sum_{k_{1:n}}|\hat\psi_i(k_{1:n})|^2|k_1|^2
      \Big)^{1/2}\lesssim \lambda_\eps^2 \prod_{i=1}^2 \|\mathcal N(-\gensy)^{1/2}\psi_i\|
    \end{equs}
    as desired.
\end{proof} 




\subsection{A priori estimates on the truncated generator equation}\label{sec:apriori}

The goal of this section is to derive a priori estimates on suitable elements of the 
$n$-th inhomogeneous Fock space, which will be crucial in characterising the limit $\eps\to 0$ 
of the solution to~\eqref{e:BurgersScaled}. The specific functions we need will depend on the dimension,  
as for $d=2$ we want to take full advantage of the Replacement Lemma~\ref{l:Replacement}. Let us begin with some definitions. 

Let $n\in\N$ and $P^n_2$ be the projection onto $n$-th inhomogeneous Fock space with the first 
chaos removed, i.e. onto $\bigoplus_{j=2}^n\Gamma L^2_j$. 
The truncated generator $\genn$, and the corresponding truncated operators $\genan$, $\cA_{2,n}^{\eps,+}$ 
and $\cA_{2,n}^{\eps,-}$ are respectively defined as 
\begin{equ}[e:genn]
\genn\eqdef \gensy +\genan\,,\quad \genan\eqdef P_2^n\gena P_2^n\quad \text{and}\quad\cA_{2,n}^{\eps,\sigma}\eqdef P_2^n\gena_\sigma P_2^n\,,
\end{equ}
for $\sigma\in\{+,-\}$. 
\medskip

Let $\psi\in\bigoplus_{j=2}^m\Gamma L^2_j$, $2\leq m\le n$. For $d\geq 2$, let $u^{\eps,n}$ be the unique solution 
of the {\it truncated generator equation} which is defined as 
\begin{equation}\label{eq:TruncGenEq}
-\genn u^{\eps,n}=\psi\,.
\end{equation}
In the specific case of $d=2$, we further consider $\tilde u^\eps$ which satisfies 
\begin{equ}[e:AnsatzNon]
\tilde u^{\eps}\eqdef (-\gensy-\gensy^\fw\cGN)^{-1}\genap \tilde u^{\eps} +(-\gensy-\gensy^\fw\cGN)^{-1}\psi
\end{equ}
where $\gensy^\fw$ and $\cGN$ are defined according to~\eqref{e:L0w} and~\eqref{e:G} respectively. 
Note that, even if in~\eqref{e:AnsatzNon} there is no truncation in the chaos, the equation admits a unique solution. 
Indeed, it is a triangular equation 
which can be explicitly solved starting from $\tilde u^{\eps}_1=0$, and then inductively 
setting 
\begin{equ}
\tilde u^{\eps}_{j}\eqdef (-\gensy-\gensy^\fw\cGN)^{-1}\genap \tilde u^{\eps}_{j-1} +(-\gensy-\gensy^\fw\cGN)^{-1}\psi_{j}\,,\qquad j\geq 2\,.
\end{equ}
That said, we will be concerned with the projection of $\tilde u^\eps$ onto $\oplus_{j=2}^n\fock_j$, that 
we denote by $\tilde u^{\eps,n}$ and can be easily seen to solve
\begin{equ}[e:Ansatz]
\tilde u^{\eps,n}\eqdef (-\gensy-\gensy^\fw\cGN)^{-1}\genap P_2^{n-1} \tilde u^{\eps,n} +(-\gensy-\gensy^\fw\cGN)^{-1}\psi\,.
\end{equ}
\medskip

Since the difficult term in~\eqref{e:BurgersScaled} is the quadratic one, we need to find a way to  
replace it with objects we are able to control. To do so, 
we will choose the input function $\psi$ 
at the right hand side of~\eqref{eq:TruncGenEq} and~\eqref{e:Ansatz}
to be the (kernel in $\fock$ of the) nonlinearity of~\eqref{e:BurgersScaled}. 

For a smooth function $\phi\in\cS(\T^d)$, recall that, as in~\eqref{e:nonlin1}, 
$\NN_\phi(\eta)$ is the nonlinearity tested against $\phi$. Now, notice that this equals
\begin{equ}[e:nonlin]
\NN_\phi(\eta)= \sum_{\ell,m}\frac{\iota}{(2\pi)^{\frac{d}2}}\lambda_\eps  \indN{\ell,m}\fw\cdot(\ell+m)\hat\phi(-\ell-m):\hat\eta(\ell)\hat\eta(m): = I_2(\nf_\phi)
\end{equ}
where $:\hat\eta(\ell)\hat\eta(m):\eqdef \hat\eta(\ell)\hat\eta(m)-\E[\hat\eta(\ell)\hat\eta(m)]$ 
denotes the Wick's product and we could replace the usual product with it since for $\ell=-m$ the summand vanishes. 
In the above, we defined $\nf_\phi$ to be the kernel of $\NN_\phi(\eta)$ in $\fock_2$. 

From now on, we consider a given test function $\phi\in\cS(\T^d)$ and we will study $\vN$ (that we distinguish from 
$u^{\eps,n}$ as the input function is fixed), 
which is the solution to the equation 
\begin{equ}[e:nonlinGenEq]
-\genn \vN=\nf_\phi\,.
\end{equ}
and for $d=2$, we consider also $\tvN$ which solves
\begin{equ}[e:nonlinAnsatz]
\tvN= (-\gensy-\gensy^\fw\cGN)^{-1}\genap P_2^{n-1} \tvN +(-\gensy-\gensy^\fw\cGN)^{-1}\nf_\phi\,.
\end{equ} 

We are ready to state the main result of this section. 

\begin{theorem}\label{thm:GenEqu}
Let $\phi\in\cS(\T^d)$ and $\nf_\phi$ be defined according to~\eqref{e:nonlin}. For $n\in\N$, let $\wN$ 
be the solution of the truncated generator equation $\vN$ in~\eqref{e:nonlinGenEq} if $d\geq 3$ or be $\tvN$
given by~\eqref{e:nonlinAnsatz} if $d=2$. Then, 
for any $k\in\N$, there exists a constant $C>0$, depending only on $k,m$ and the dimension $d$, 
and $\eps_d(n)>0$ such that for any $\eps<\eps_d(n)$ the following estimate holds
\begin{equ}
\|(-\gensy)^{-\frac12}(-\gen \wN-\nf_\phi+\genam \wN_2)\|\leq \frac{C}{n^{k}}\|(-\gensy)^{-\tfrac{1}{2}}\nf_\phi\|\,,\label{e:ApproxH1}%\\
%\|P_{>m}u^{\eps,n}\|&\leq C\sqrt{f(\eps)}\Big(1+\|(-\gensy)^{-\tfrac12}\psi\|\Big),\label{e:ApproxL2}
\end{equ}
and the latter norm at the right hand side satisfies 
\begin{equ}[e:nonlinUnifBound]
\|(-\gensy)^{-\tfrac{1}{2}}\nf_\phi\|\lesssim \|\phi\|_{H^1(\T^d)}
\end{equ}
uniformly over $\eps$. 
If $d\geq 3$, $\eps_d(n)$ can be taken to be $1$. 
%where $P^{>m}$ is the projection onto $\oplus_{j>m}\fock_j$. % and $f(\eps)$ is the function defined according to~\eqref{e:f}.  
\end{theorem}  


%At first, we notice that~\eqref{eq:TruncGenEq} can be written as a finite system of equation involving 
%the $n-1$ components of $u^{\eps,n}=(u^{\eps,n}_j)_{j=2,\dots,n}$, which reads 
%\begin{equation}\label{eq:truncgensys}
%\begin{cases}
%%-\gensy u^{\eps,n}_n-\genap u^{\eps,n}_{n-1}=0\\
%-\gensy u^{\eps,n}_{j}-\genap u^{\eps,n}_{j-1}-\1_{j\neq n}\genam u^{\eps,n}_{j+1}=0\,,& j=m+1,\dots, n\\
%-\gensy u^{\eps,n}_{j}-\1_{j\neq 2}\genap u^{\eps,n}_{j-1}-\genam u^{\eps,n}_{j+1}=\psi_{j}\,,& j=2,\dots, m\,.
%%\cdots\\
%%-\gensy u^{\eps,n}_{1}-\genam u^{\eps,n}_{2}=g_{1}\,.
%\end{cases}
%\end{equation}
%The system above can be solved by expressing the $j$-th coordinate of $u^{\eps,n}$ 
%in terms of the $j-1$-th starting from $u^{\eps,n}_n$ and using the operators $\Op_j$ 
%recursively defined as 
%\begin{equation}\label{e:Op}
%\begin{cases}
%\Op_1=0\,,\\
%\Op_j=-\genam (-\gensy + \Op_{j-1})^{-1}\genap.
%\end{cases}
%\end{equation}
%The operators above will play a crucial role in our analysis, and especially in the characterisation of the limiting 
%equation. We postpone our analysis to the next subsection, in which we will derive some of its 
%crucial properties. For now, it suffices to note that, for any $j$, $\Op_j$ is symmetric and positive. 
%
%We are ready to state the main result of this section, concerning the behaviour of the solution  
%of the truncated generator equation $u^{\eps,n}$. 

As a first step in the proof of the previous theorem, we derive weighted estimates 
on the solution to~\eqref{e:nonlinGenEq} which hold in {\it any} dimension $d\geq 2$. 
The argument we exploit follows closely 
that of \cite[Lemma 2.5]{LandimYau1997} and since the proof works for any choice of $\psi$, 
we formulate and prove the result in the more 
general case of~\eqref{eq:TruncGenEq}. 


\begin{proposition}\label{P:AprioriGeneratorEq}
Let $m\in\N$ and $\psi\in\bigoplus_{j=2}^m\Gamma L^2_j$. For $n\in\N$, $n\geq m$, let $u^{\eps,n}$ 
be the solution of the truncated generator equation in~\eqref{eq:TruncGenEq} in any dimension $d\geq 2$. 
Then, there exists a 
positive constant $C=C(m,k)$ independent of $n,\eps$ and $\psi$, such that
\begin{equation}\label{e:AprioriMain}
\|\mathcal N^k (-\gensy)^\frac12 u^{\eps,n}\|\leq C\|(-\gensy)^{-\frac12}\psi\|\,,
\end{equation}
where $\cN$ is the number operator in Definition~\ref{def:NoOp}. 
\end{proposition}
\begin{proof}
As $n$ and $\eps$ are fixed throughout the proof, we will write $u$ and $u_j,\,j=2,\dots,n$ in place of 
$u^{\eps,n}$ and $u^{\eps,n}_j$, respectively. 
Also, by convention, we let $u_{n+1}=0=u_1$. To denote constants which do not depend on $\eps,n$ or $\psi$ 
we will use $C$, and such $C$ might change from line to line. 
%
%As a first observation, we notice that, by orthogonality of Fock spaces with different indices, 
%the first component of $u$ is independent of the others, in the sense that on the one hand it solves the equation 
%\begin{equ}[e:u1]
%-\gensy u_1=\psi_1
%\end{equ}
%and on the other no other component of $u$ depends on $u_1$. 
%Further, a standard argument shows that~\eqref{e:u1} immediately implies 
%that~\eqref{e:AprioriMain} holds for $u_1$, hence we only need to consider $u_j$ for $j\geq 2$. 

Let us test the $j$-th component, $j\geq 2$, with $\gen u$, which, since by Lemma~\ref{lem:essid} 
$(\genap)^\ast=-\genam$, gives
  \begin{equs}
    \langle u_j,\gen u\rangle&= \langle u_j, \cL_0 u_j\rangle+\langle u_j,\genap u_{j-1}\rangle+\langle u_j,\genam u_{j+1}\rangle\\
    &=  \langle u_j, \cL_0 u_j\rangle- \Big[\langle u_{j+1},\genap u_j\rangle-\langle u_{j},\genap u_{j-1}\rangle\big].
  \end{equs}
Via a summation by parts, for any $a>0$, we have
  \begin{equs}
    \sum_{j=2}^n (a+j^{2k})\langle u_j,(-\cL_0) u_j\rangle=  &\sum_{j=2}^n (a+j^{2k})\langle u_j,(-\gen) u\rangle\\
    &+\sum_{j=2}^n (j^{2k}-(j-1)^{2k})\langle u_j,\genap u_{j-1}\rangle.
  \end{equs}
  Next, note that $-\gen u =\psi-\genap u_n-\genam u_2$ so that, by orthogonality of Fock spaces with different indices,
  \begin{equ}
    |\langle u_j,(-\gen) u\rangle|= |\langle u_j,\psi_j\rangle|\le \frac12 \langle u_j,(-\cL_0)u_j\rangle+\frac12\langle \psi_j,(-\cL_0)^{-1}\psi_j\rangle.
  \end{equ}
  As a consequence, there exists some finite strictly positive constant $C=C(a,m,k)$, 
  which might change from line to line,  such that
  \begin{equs}[eq:freddo]
\sum_{j=2}^n &(a+j^{2k})\langle u_j,(-\cL_0) u_j\rangle\\
&\le C\Big(\|(-\cL_0)^{-\frac12}\psi\|^2+\sum_{j=2}^n (j^{2k}-(j-1)^{2k})\langle u_j,\genap  u_{j-1}\rangle\Big)\\
&\leq C\Big(\|(-\cL_0)^{-\frac12}\psi\|^2+\sum_{j=2}^n j^{2k-1}\langle u_j,\genap  u_{j-1}\rangle \Big)   
  \end{equs}
where $ C= C(k)$ is a constant 
chosen in such a way that $(j^{2k}-(j-1)^{2k})\sqrt j\le  C j^{2k-\frac12}$. 
To handle the last sum, we bound
  \begin{equs}[e:Gammaaa]
    |\langle u_j,\genap  u_{j-1}\rangle |&\le \|(-\cL_0)^{\frac12} u_j\| \| (-\cL_0)^{-\frac12}\genap  u_{j-1}\|\\
    &\leq C\sqrt{j} \|(-\cL_0)^{\frac12} u_j\| \| (-\cL_0)^{\frac12} u_{j-1}\|\\
   &\leq C\sqrt{j}\Big(\frac12 \|(-\cL_0)^{\frac12} u_j\|^2+\frac12\| (-\cL_0)^{\frac12} u_{j-1}\|^2\Big)
  \end{equs}
where in the second bound we used~\eqref{e:Abound} 
in Lemma~\ref{l:GenaBound}. We now plug the result into~\eqref{eq:freddo} and rearrange the terms, 
so that we conclude 
\begin{equ}\label{eq:freddo2}
\sum_{j=1}^n j^{2k}\Big(1+\frac a{j^{2k}}-\frac{C}{\sqrt j}\Big)\langle u_j,(-\cL_0) u_j\rangle\le C(1+\|(-\cL_0)^{-\frac12}\psi\|^2)\,. \end{equ}
Choosing $a$ sufficiently large, 
in such a way that $(1+\frac a{j^{2k}}-\frac{C}{\sqrt j})\ge 1/2$ for every $j\ge2$,~\eqref{e:AprioriMain} follows.
%\medskip
%
%It remains to prove \eqref{eq:claimT0}. To do so, let $f\in \fock_n$ and $h\in \fock_{n+1}$. 
%Recalling the definition of $\genap$ in~\eqref{e:gen} 
%and using the symmetry of $f,h$, we have\giuseppe{Qua stiamo usando la versione non simmetrizzata. } 
%\begin{equ}
%  \langle h,\genap  f\rangle= (n+1)! n \iota\lambda_\eps\sum_{k_{1:n+1}}h(k_{1:n+1}) \, \fw\cdot(k_1+k_2)  \indN{k_1,k_2} f(k_1+k_2,k_{3:n+1}).
%\end{equ}
%Let us look at the sum in $k_1$ and $k_2$, ignoring for the moment the prefactors. We bound its absolute value as
%\begin{equs}
%\Big|  &\sum_{a} (\fw\cdot a)f(a,k_{3:n+1})\sum_{k_1+k_2=a}\indN{k_1,k_2} h(k_{1:n+1})
%\Big| \\
%&\le \sqrt{\sum_a|a|^2 f(a,k_{3:n+1})^2 }\sqrt{\sum_a \bigg(
%\sum_{k_1+k_2=a}\indN{k_1,k_2} h(k_{1:n+1})  \bigg)^2}.
%\end{equs}
%On the other hand,
%\begin{equ}
%\Big(
%\sum_{k_1+k_2=a}\indN{k_1,k_2} h(k_{1:n+1})  \Big)^2\le \sum_{k_1+k_2=a}|k_1|^2 h(k_{1:n+1})^2\times
%\sum_{k_1+k_2=a}\frac{\indN{k_1,k_2} }{|k_1|^2}
%\end{equ}
%and the latter sum is $O(\lambda_\eps^{-2})$ for every $d\ge2$.
%Altogether,
%\begin{equ}
%  | \langle h,\genap  f\rangle|\lesssim_d (n+1)! n \sum_{k_{3:n+1}}\sqrt{A B}\leq (n+1)! n \sum_{k_{3:n+1}}\Big(\frac A{\gamma\sqrt{n}}+B \gamma\sqrt n\Big)\,,
%\end{equ}
%where $A$ and $B$ are defined as 
%\begin{equ}
%  A\eqdef\sum_a |a|^2 f(a,k_{3:n+1})^2,\qquad B\eqdef\sum_{k_1,k_2}|k_1|^2h(k_{1:n+1})^2\,.
%\end{equ}
%and the last inequality holds for any $\gamma>0$ and $n>0$. 
%Changing the labels of the summed variables and using symmetry of the kernels we deduce that
%\begin{equs}[caffe]
%  | \langle h,\genap  f\rangle|&\lesssim \frac{\sqrt n}\gamma n!\sum_{k_{1:n}}\frac{|k_{1:n}|^2}2 f(k_{1:n})^2+
%  \gamma\sqrt n \, (n+1)! \sum_{k_{1:n+1}} \frac{|k_{1:n+1}|^2}2 h(k_{1:n+1})^2\\
%  &= \sqrt n \left(\frac{\|(-\cL_0)^{1/2}f \|^2}\gamma+\gamma \|(-\cL_0)^{1/2}h \|^2\right).
%\end{equs}
%Going back to  \eqref{eq:claimT0}, since $u_{j-1}\in \cH_{j-1}$,
%\begin{equs}
%  \| &(-\cL_0)^{-\frac12}\genap  u_{j-1}\|^2=\sup_{g\in \cH_j}\left(2\langle g,\genap  u_{j-1}\rangle-
%    \| (-\cL_0)^{\frac12}g\|^2\right)
%  \\
%  &\le \sup_{g\in \cH_j}\left(2C\sqrt j \frac{ \| (-\cL_0)^{\frac12}u_{j-1}\|^2}\gamma+2C\sqrt j\gamma\|(-\cL_0)^{\frac12}g\|^2 - \| (-\cL_0)^{\frac12}g\|^2\right)
%  \\
%  &\lesssim j \|(-\cL_0)^{\frac12}u_{j-1}\|^2
%\end{equs}
%where in the first inequality we applied \eqref{caffe} and then we chose $\gamma=1/(2C\sqrt j)$.
\end{proof}

In the next two lemmas, we consider the specific case of $d=2$. 
In the first one, we show that, for any $\psi$, $\tilde u^{\eps,n}$ defined according to~\eqref{e:Ansatz} 
can be estimated in terms of $\psi$. 

\begin{lemma}\label{l:aprioriAnsatz}
For $d=2$, let $\psi\in\bigoplus_{j=2}^m\Gamma L^2_j$, $n\in\N$ and $\tilde u^{\eps,n}$ be given by~\eqref{e:Ansatz}. 
Then, 
there exists a constant $C>0$ independent of $n,\eps$ and $\eps_2(n)>0$ such that for every $\eps<\eps_2(n)$ 
\begin{equ}[e:apriorivtilde]
\|(-\gensy)^{1/2}\tilde u^{\eps,n}\|^2
%+\|(-\gensy^\fw\cGN)^{1/2}u^{\eps,n}_n\|^2
\leq C \|(-\gensy)^{-1/2}\psi\|^2\,.
\end{equ}
\end{lemma}
\begin{proof}
At first, we study how the generator $\gen$ acts on $\tilde u^{\eps,n}$. We have
\begin{equs}[e:genAnsatz]
-\gen \tilde u^{\eps,n}=&(-\gensy-\gensy^\fw\cGN)\tilde u^{\eps,n}-\genap  \tilde u^{\eps,n}-\genam \tilde u^{\eps,n}+\gensy^\fw\cGN \tilde u^{\eps,n}\\
=&\psi -\genam (-\gensy-\gensy^\fw\cGN)^{-1}\psi -\genap \tilde u^{\eps,n}_n+\gensy^\fw\cGN \tilde u^{\eps,n}_n\\
&+\big[-\genam(-\gensy-\gensy^\fw\cGN)^{-1}\genap+\gensy^\fw\cGN\big]P_2^{n-1} \tilde u^{\eps,n}
\end{equs}
where in the last step we used~\eqref{e:Ansatz}. 
Now, $\genap \tilde u^{\eps,n}_n\in\fock_{n+1}$, so that, since $ \tilde u^{\eps,n}\in\oplus_{j=2}^n\fock_j$, 
the two are orthogonal. Therefore, by testing both sides by $ \tilde u^{\eps,n}$, we obtain
\begin{equs}
\|&(-\gensy)^{1/2} \tilde u^{\eps,n}\|^2\\
&=\langle  \tilde u^{\eps,n}, \psi-\genam (-\gensy-\gensy^\fw\cGN)^{-1}\psi\rangle+\langle  \tilde u^{\eps,n}_n, \gensy^\fw\cGN \tilde u^{\eps,n}_n\rangle\\
&\quad+\langle  \tilde u^{\eps,n}, \big[-\genam(-\gensy-\gensy^\fw\cGN)^{-1}\genap+\gensy^\fw\cGN\big]P_2^{n-1} \tilde u^{\eps,n}\rangle\\
&\leq \tfrac12\|(-\gensy)^{1/2} \tilde u^{\eps,n}\|^2+\tfrac12\|(-\gensy)^{-1/2}[\psi-\genam (-\gensy-\gensy^\fw\cGN)^{-1}\psi]\|^2 \\
&\quad- \|(-\gensy^\fw\cGN)^{1/2} \tilde u^{\eps,n}_n\|^2+C\lambda_\eps^2 n^2 \|(-\gensy)^{1/2} \tilde u^{\eps,n}\|^2\\
&\leq C'\|(-\gensy)^{-1/2}\psi\|^2+(\tfrac12 + Cn^2\lambda_\eps^2)\|(-\gensy)^{1/2} \tilde u^{\eps,n}\|^2\label{e:apriorivtilde1}
%=&-\|(-\gensy^\fw\cGN)^{1/2}\tvN_n\|^2\\
%&+\lambda_\eps^2 n^2 \|(-\gensy)^{1/2}\tvN\|^2
\end{equs}
where in the last step we neglected the negative term at the right hand side, we used~\eqref{e:Tbound} together with 
the positivity of $-\gensy^\fw\cGN$ and the fact that 
the operators $\gensy,\gensy^\fw$ and $\cGN$ commute to control the second term containing $\psi$, 
and the Replacement Lemma~\ref{l:Replacement} to bound the third summand. 
Now, we choose $\eps_2(n)$ in such a way that $C n^2\lambda_{\eps_2(n)}^2<1/2$. 
Therefore,~\eqref{e:apriorivtilde} follows upon rearranging the terms in~\eqref{e:apriorivtilde1}. 
\end{proof}

In the next lemma, we turn our attention to $\tvN$ given in~\eqref{e:nonlinAnsatz} and
prove that it satisfies the same weighted estimates as those in~\eqref{e:AprioriMain}. 
To do so, in view of Proposition~\ref{P:AprioriGeneratorEq} and the previous Lemma, 
it suffices to control the difference between $\tvN$ and the solution of the truncated generator equation $\vN$. % is small, so that in particular a bound analogous to~\eqref{e:AprioriMain} holds for $\vN$ as well. 

\begin{lemma}\label{l:AnsatzvsGenEq}
Let $d=2$. For $n\in\N$, let $\tvN$ be given by~\eqref{e:nonlinAnsatz}.
Then, for any $k\in\N$ there exists a constant $C>0$ and $\eps_2(n,k)>0$ such that for every $\eps<\eps_2(n,k)$ 
\begin{equ}[e:AprioriAnsatz]
\|\cN^k(-\gensy)^{1/2}\tvN\|\leq C\|(-\gensy)^{-1/2}\nf_\phi\|\,.
\end{equ}
\end{lemma}
\begin{proof}
Let $\vN$ be the solution of~\eqref{e:nonlinGenEq} in $d=2$. Notice that, trivially, 
\begin{equ}[e:TwoSum]
\|\cN^k(-\gensy)^{1/2}\tvN\|\leq n^k\|(-\gensy)^{1/2}(\tvN-\vN)\|+\|\cN^k(-\gensy)^{1/2}\vN\|\,.
\end{equ}
Now, in view of Proposition~\ref{P:AprioriGeneratorEq}, the second summand is bounded by (a constant times) 
$\|(-\gensy)^{-1/2}\nf_\phi\|$, so we only need to control 
the first. To do so, we claim that for any $j$ there exists $C>0$ such that for $\eps$ small enough, we have
\begin{equ}[e:AprioriDiff]
\|(-\gensy)^{1/2}(\tvN-\vN)\|^2\leq C\big( n^{-j} +n^2\lambda_\eps^2\big)\|(-\gensy)^{-1/2}\nf_\phi\|^2\,.
\end{equ}
Assuming the claim, the proof is concluded as it suffices to choose in~\eqref{e:AprioriDiff} 
$j=k$ and $\eps_2(n,k)>0$ such that $n^{k+2}\lambda_{\eps_2(n,k)}<1$, 
to ensure that also the first summand in~\eqref{e:TwoSum} is bounded. 
\medskip

\newcommand{\pes}{\psi_\eps^\sharp}
We now prove~\eqref{e:AprioriDiff}. 
To shorten the notation, set $\pes\eqdef (-\gensy-\gensy^\fw\cGN)^{-1}\nf_\phi$. 
We begin by evaluating $-\gen$ on $\tvN-\vN$. 
To do so, we exploit~\eqref{e:genAnsatz} and the fact that, as noted in the proof of Proposition~\ref{P:AprioriGeneratorEq}, 
since $\vN$ solves~\eqref{e:nonlinGenEq}, we have 
\begin{equ}%[e:genu]
-\gen \vN-\nf_\phi+\genam \vN_2%% =-\genn u^{\eps,n}-\genap u^{\eps,n}_n-\psi=
=  -\genap \vN_n\,.
\end{equ}
This means that 
\begin{equs}[e:genDiff]
-\gen (\tvN-\vN)=&-\genam (\pes-\vN_2)-\genap(\tvN_n-\vN_n)+\gensy^\fw\cGN \tvN_n\\
&+\big[-\genam(-\gensy-\gensy^\fw\cGN)^{-1}\genap+\gensy^\fw\cGN\big]P_2^{n-1} \tvN\,.
\end{equs}
Now, since $\pes\in\fock_2$, $\genam (\pes-\vN_2)\in\fock_1$ and 
$\genap(\vN_n-\tvN_n)\in\fock_{n+1}$, which implies that they are orthogonal 
to both $\tvN$ and $\vN$ as these belong to $\oplus_{j=2}^n\fock_j$. 
Hence, by testing both sides of~\eqref{e:genDiff} by $\tvN-\vN$, we obtain
\begin{equs}\label{e:aprioriDiff1}
\|(-\gensy)^{1/2}&(\tvN-\vN)\|^2=\langle \tvN-\vN, \gensy^\fw\cGN \tvN_n\rangle\\
&+\langle \tvN-\vN, \big[-\genam(-\gensy-\gensy^\fw\cGN)^{-1}\genap+\gensy^\fw\cGN\big]P_2^{n-1}\tvN\rangle\,.
\end{equs}
Let us analyse the two terms at the right hand side separately. 
For the first, note that the operator $\gensy^\fw\cGN$ is negative, so that 
\begin{equs}
\langle \tvN-\vN, \gensy^\fw\cGN\tvN_n\rangle&=\langle \vN_n, \gensy^\fw\cGN\tvN_n\rangle -\langle \tvN_n, (-\gensy^\fw\cGN)\tvN_n\rangle\\
&\leq \langle -\vN_n, (-\gensy^\fw\cGN)\tvN_n\rangle\\
&\lesssim \|(-\gensy)^{1/2}\vN_n\|\|(-\gensy)^{1/2}\tvN_n\|
\end{equs}
and in the last bound we used that $-\gensy^\fw\cGN\lesssim -\gensy$ and that $\gensy, \gensy^\fw$ and $\cGN$ commute. 
Now, the a priori estimates in Proposition~\ref{P:AprioriGeneratorEq} and 
Lemma~\ref{l:aprioriAnsatz} allow to upper bound the previous by
\begin{equs}
 n^{-k} \|\cN^k(-\gensy)^{1/2}\tvN\|\|(-\gensy)^{1/2}\vN_n\|\lesssim n^{-k} \|(-\gensy)^{-1/2}\nf_\phi\|^2\,.
\end{equs}
For the second term in~\eqref{e:aprioriDiff1}, we apply the Replacement Lemma~\ref{l:Replacement} 
first and the same a priori estimates as above, thus getting
\begin{equs}
\langle \tvN-\vN, &\big[\genam(-\gensy-\gensy^\fw\cGN)^{-1}\genap+\gensy^\fw\cGN\big]P_2^{n-1}\tvN\rangle\\
&\leq C \lambda_\eps^2n^2\|(-\gensy)^{1/2}(\tvN-\vN)\|\|(-\gensy)^{1/2}\tvN\|\\
&\leq \tfrac12C\lambda_\eps^2n^2\Big(\|(-\gensy)^{1/2}(\tvN-\vN)\|^2+\|(-\gensy)^{1/2}\tvN\|^2\Big)\\
&\leq \tfrac12C\lambda_\eps^2n^2\Big(\|(-\gensy)^{1/2}(\tvN-\vN)\|^2+\|(-\gensy)^{1/2}\nf_\phi\|^2\Big)\,,
\end{equs}
where the constant $C$ changed in the last two lines. Now, by using that for $\eps<\eps_2(n,k)$, 
$C n^2\lambda_\eps^2<1$, collecting the previous bounds and rearranging the terms,~\eqref{e:AprioriDiff} follows. 
\end{proof}

We are now ready for the proof of Theorem~\ref{thm:GenEqu}, which is an easy consequence 
of Lemma~\ref{l:GenaBound} and, for $d\geq 3$, Proposition~\ref{P:AprioriGeneratorEq}, while, 
for $d=2$ Lemmas~\ref{l:aprioriAnsatz} and~\ref{l:AnsatzvsGenEq}. 


\begin{proof}[of Theorem~\ref{thm:GenEqu}]
%Let us begin by~\eqref{e:ApproxH1}. 
Let us first treat the case of $d\geq 3$. As noted in the proof of Proposition~\ref{P:AprioriGeneratorEq}, 
since $\vN$ solves~\eqref{e:nonlinGenEq}, we have 
\begin{equ}[e:genu]
-\gen \vN-\nf_\phi+\genam \vN_2%% =-\genn \vN-\genap \vN_n-\psi=
=  -\genap \vN_n\,.
\end{equ}
By~\eqref{e:Abound}, for any $k\ge 1$, we have 
\begin{equ}[e:chaosdecay]
\|(-\gensy)^{-\tfrac12}\genap \vN_n\|\lesssim \sqrt{n}\|(-\gensy)^{\tfrac12}\vN_n\|\leq \frac1{n^{k}}\|\cN^{k+1}(-\gensy)^{\tfrac12} \vN\|\,.
\end{equ}
Hence,~\eqref{e:ApproxH1} follows directly by~\eqref{e:AprioriMain}. 
\medskip

We now turn to $d=2$, in which case $\tvN$ satisfies~\eqref{e:nonlinAnsatz}. 
Notice that, by the recursive definition of $\tvN$, $\tvN_1=0$ and therefore
\begin{equ}
\tvN_2=(-\gensy-\gensy^\fw\cGN)^{-1}\nf_\phi\,.
\end{equ}
Hence, upon replacing $\psi$ with $\nf_\phi$ in~\eqref{e:genAnsatz}, we see that 
\begin{equs}
-\gen \tvN-\nf_\phi+\genam \tvN_2=&-\gen \tvN-\nf_\phi+\genam (-\gensy-\gensy^\fw\cGN)^{-1}\nf_\phi\\
=&-\genap \tvN_n+\gensy^\fw\cGN \tvN_n\\
&+\big[-\genam(-\gensy-\gensy^\fw\cGN)^{-1}\genap+\gensy^\fw\cGN\big]P_2^{n-1} \tvN\,.
\end{equs} 
Now, for the first two terms we use Lemma~\ref{l:GenaBound} and $-\gensy^\fw\cGN\lesssim-\gensy$, 
so that we can bound them by 
\begin{equs}
\|(-\gensy)^{-1/2}[-\genap \tvN_n+\gensy^\fw\cGN \tvN_n]\|&\lesssim \sqrt{n}\|(-\gensy)^{1/2}\tvN_n\|\\
&\lesssim n^{-k}\|(-\gensy)^{-1/2}\nf_\phi\|
\end{equs}
where the last step follows by Lemma~\ref{l:AnsatzvsGenEq}. 
For the third term, we note that 
the $\fock$-norm satisfies a variational formulation, i.e.
\begin{equs}
\|&(-\gensy)^{-1/2}\big[-\genam(-\gensy-\gensy^\fw\cGN)^{-1}\genap+\gensy^\fw\cGN\big]P_2^{n-1} \tvN\|\\
&=\sup_{\|\psi\|=1}\langle\psi, (-\gensy)^{-1/2}\big[-\genam(-\gensy-\gensy^\fw\cGN)^{-1}\genap+\gensy^\fw\cGN\big]P_2^{n-1} \tvN\rangle\\
&=\sup_{\|\psi\|=1}\langle(-\gensy)^{-1/2}\psi, \big[-\genam(-\gensy-\gensy^\fw\cGN)^{-1}\genap+\gensy^\fw\cGN\big]P_2^{n-1} \tvN\rangle\\
&\lesssim C \lambda_\eps^2 \sup_{\|\psi\|=1}\|\psi\|\|\cN^2(-\gensy)^{1/2}\tvN\|=C \lambda_\eps^2\|\cN^2(-\gensy)^{1/2}\tvN\|\\
&\leq C \lambda_\eps^2n^{2}\|(-\gensy)^{-1/2}\nf_\phi\|\label{e:VarRepl}
\end{equs}
where in the third step we used the Replacement Lemma~\ref{l:Replacement} and in the last, 
Lemma~\ref{l:aprioriAnsatz}. By choosing $\eps_2(n,k)>0$ such that $\lambda_{\eps_2(n,k)}^2n^{2+k}<1$,~\eqref{e:ApproxH1} follows at once. 
\medskip

It remains to argue~\eqref{e:nonlinUnifBound}. For this, notice that for any $\phi$,~\eqref{e:nonlin} and~\eqref{e:gen} 
imply that
the following equality holds
\begin{equ}[e:nonlinA+]
\genap \phi=\nf_{\phi}\,.
\end{equ}
Hence, by Lemma~\ref{l:GenaBound} we have
\begin{equ}
\|(-\gensy)^{-1/2}\nf_\phi\|=\|(-\gensy)^{-1/2}\genap \phi\|\lesssim \|(-\gensy)^{1/2}\phi\|=\|\phi\|_{H^1(\T^d)}\,,
\end{equ}
so that the proof is concluded. 
\end{proof}



 
\subsection{The diffusivity}

The main advantage of Theorem~\ref{thm:GenEqu} is that it allows to replace the nonlinearity with 
the sum of two terms: one involves the generator of our process, and therefore encodes its fluctuations, 
while the other consists of $\genam$ applied to an element in the Fock space of degree $2$, 
so that overall it is of degree $1$. This suggests that this latter term should give us the additional 
diffusivity appearing in the limit~\eqref{e:SHEintro}. 
The goal of this section is to analyse this object and determine its behaviour. 

We begin our analysis with a lemma that provides a representation for $\genam\wN_2$, 
which will help us identify its limit. 

\begin{lemma}\label{l:A-}
Let $\nf_\phi$ be defined according to~\eqref{e:nonlin}. For $n\in\N$, let $\vN$ 
be the solution of the truncated generator equation in~\eqref{e:nonlinGenEq} and, for $d=2$, $\tvN$ be 
given by~\eqref{e:nonlinAnsatz} if $d=2$. 
Then, for $d\geq 3$, we have 
\begin{equ}[e:A-D3]
\genam\vN_2=-\sum_k \hat\phi(k)\|(-\gensy)^{1/2}\vN[-k]\|^2 e_{k}
\end{equ}
where, for $k\in\Z^d$, $\vN[k]$ is the solution to~\eqref{e:nonlinGenEq} with $\phi=e_{-k}$. 

For $d=2$ instead, we have 
\begin{equ}[e:A-D2]
\genam \tvN_2=\genam (-\gensy-\gensy^\fw\cGN)^{-1}\genap\phi\,.
\end{equ}
\end{lemma}
\begin{proof}
Recall~\eqref{e:nonlinA+}. Then, in $d=2$, by the recursive definition of $\tvN$, we have that $\tvN_2= (-\gensy-\gensy^\fw\cGN)^{-1}\nf_\phi$ 
which immediately gives~\eqref{e:A-D2}. 

We now show~\eqref{e:A-D3}. It is convenient 
to observe that $\nf_\phi$ can be written as
\begin{equ}[e:nonlinFourier]
\nf_\phi=\sum_k\hat\phi(k) \nf_{-k}
\end{equ}
where, to shorten the notation, we set $\nf_{-k}\eqdef \nf_{e_k}$. 
By linearity of $\genn$ and uniqueness for~\eqref{e:nonlinGenEq}, we have 
$\vN=\sum_k \hat\phi(k) \vN[-k]$, where, as in the statement, $\vN[k]$ is 
the solution of the truncated generator equation whose input is $\nf_k$. 
Hence, applying $\genam$ to $\vN_2$ we obtain
\begin{equ}
\genam \vN_2=\sum_k \hat\phi(k) \genam \vN_2[-k]\,,
\end{equ}
so that the proof is concluded once we verify that 
\begin{equ}[e:final]
\genam \vN_2[k]=-\|(-\gensy)^{1/2}\vN[k]\|^2 e_{-k}\,,\qquad k\in\Z^d_0\,.
\end{equ}
Since $\genam \vN_2[k]\in\fock_1$, we can expand it in the Fourier basis, i.e. 
\begin{equs}[e:TransInv]
\genam \vN_2[k]&=\sum_{j\in\Z^d_0} \langle \genam \vN_2[k], e_{-j} \rangle e_j=-\sum_{j\in\Z^d_0} \langle \vN_2[k], \genap e_{-j} \rangle e_j\\
&=- \langle  \vN_2[k], \genap e_{k} \rangle e_{-k}=-\langle \vN_2[k], \nf_{-k} \rangle e_{-k}\,,
\end{equs}
where the third equality is a consequence of the translation invariance of the equation 
(which can be easily seen to hold by verifying that, for any $i=1,\dots,d$, the momentum operator $M_i$ defined as 
$\cF(M_if)(k_{1:n})=(\sum_j k^i_j)\hat f(k_{1:n})$, $f\in\fock_n$, commutes with $\gensy, \genap$ and $\genam$), 
while the last, of~\eqref{e:nonlinA+}. 
Therefore, we can now exploit orthogonality of Fock spaces 
with different indices and the antisymmetry of $\gena$, to see that 
\begin{equs}
\langle \vN_2[k], \genap e_{k} \rangle&=\langle \vN_2[k], \nf_{-k} \rangle=\langle \vN[k], \nf_{-k} \rangle = \langle \vN[k], -\gen \vN[-k] \rangle \\
&=\langle \vN[k], -\gensy \vN[-k] \rangle=\|(-\gensy)^{1/2}\vN[k]\|^2\,,
\end{equs}
from which, by~\eqref{e:TransInv},~\eqref{e:final} and therefore the statement, follow at once. 
\end{proof}

From the previous statement and the Replacement Lemma~\ref{l:Replacement}, 
we can easily obtain the $\eps\to0$ (and $n\to\infty$) limit of $\genam\vN_2$ for $d=2$. 
This is summarised in the following theorem, which also shows that the $\fock$-norm 
of $\vN$ goes to $0$. 

\begin{theorem}\label{thm:MainConv2}
Let $d=2$. For $n\in\N$, $n\geq 2$, $\eps>0$ and $\phi\in\cS(\T^2)$, let $\tvN$ 
be given by~\eqref{e:nonlinAnsatz}. Then
\begin{equ}
\lim_{\eps\to 0} \|\genam\tvN_2-\cD\phi\|=0\,,\label{e:Diffusivity2}
\end{equ}
where $\cD$ is the operator introduced in Theorem~\ref{thm:main}, given by $\cD\eqdef D_{\SHE}\gensy^\fw$ 
for $\gensy^\fw$ defined according to~\eqref{e:L0w} and $D_\SHE=G(1)$, $G$ being the function given in~\eqref{eq:defG}. 
Further, we have  
\begin{equ}
\lim_{\eps\to0}\|\tvN\|=0\,.\label{e:L2norm2}
\end{equ}
\end{theorem}
\begin{proof}
We begin by showing~\eqref{e:Diffusivity2}. By~\eqref{e:A-D2} in Lemma~\ref{l:A-}, we have 
\begin{equs}
\|\genam \tvN_2-\cD\phi\|\leq&\|\big[\genam (-\gensy-\gensy^\fw\cGN)^{-1}\genap-\gensy^\fw\cGN\big]\phi\| +\|\big[\gensy^\fw\cGN-\cD\big]\phi\|\,.
\end{equs}
For the first term at the right hand side, we use the same variational principle as in~\eqref{e:VarRepl} so that 
\begin{equs}
\|\big[-\genam(-\gensy-&\gensy^\fw\cGN)^{-1}\genap+\gensy^\fw\cGN\big]\phi\|\\
&=\sup_{\|\psi\|=1}\langle\psi, \big[-\genam(-\gensy-\gensy^\fw\cGN)^{-1}\genap+\gensy^\fw\cGN\big] \phi\rangle\\
&\lesssim  \lambda_\eps^2 \sup_{\|\psi\|=1}\|\psi\|\|(-\gensy)\phi\|= \lambda_\eps^2\|\phi\|_{H^2(\T^2)}
\end{equs}
where in the last step we used the Replacement Lemma~\ref{l:Replacement} 
but in which we distributed the $-\gensy$ in~\eqref{e:CSH1} differently (note that, as $\phi\in\fock_1$ 
there are no off-diagonal terms). 
For the second, we use the explicit expression of $\gensy^\fw$, $\cGN$ and $\cD$, which give
\begin{equs}
\|\big[\gensy^\fw\cGN-\cD\big]\phi\|^2=\tfrac14\sum_k (\fw\cdot k)^2\Big[G(\Ll(\tfrac12|k|^2))- D_\SHE\Big]^2|\hat\phi(k)|^2\,.
\end{equs} 
Since $\phi$ is smooth we can take the limit in $\eps$ inside the sum and, provided we 
take $D_\SHE=G(1)$, we conclude that the right hand side goes to $0$. 
Hence, the proof of~\eqref{e:Diffusivity2} is completed. 
\medskip

Concerning~\eqref{e:L2norm2}, we note that for any $j\in\N$ and $\psi\in\fock_j$, 
\begin{equs}
\|(-\gensy-&\gensy^\fw\cGN)^{-1}\genap\psi\|^2\\
&\lesssim  j^2\lambda_\eps^2\sum_{k_{1:j+1}}\frac{\indN{k_1,k_2}(\fw\cdot (k_1+k_2))^2|\hat\psi(k_1+k_2, k_{3:j+1})|^2}{[\frac12 |k_{1:j+1}|^2+\frac12(\fw\cdot k)^2_{1:j+1}G(\Ll(\frac12|k_{1:j+1}|^2))]^2}\\
&\lesssim j^2\lambda_\eps^2\sum_{k_{1:j}}(\fw\cdot k_1)^2|\hat\psi(k_{1:j})|^2\sum_{\ell+m=k_1}\frac{\indN{\ell,m}}{ (|\ell|^2+|m|^2)^2}
\end{equs}
Now, the inner sum is clearly finite, and therefore we deduce 
%\begin{equs}
%\lambda_\eps^2\sum_{\ell+m=k_1}\frac{\indN{\ell,m}}{ (|\ell|^2+|m|^2)^2}&\lesssim\lambda_\eps^2\eps^4\sum_{\eps<|\eps\ell|<1}\frac{1}{ |\eps\ell|^4}\lesssim \lambda_\eps^2 \eps^2\int_{|x|\in[\eps,1]}\frac{\dd x}{|x|^4}\\
%&\lesssim\lambda_\eps^2\eps^2\int_\eps r^{-3}\dd r\lesssim \lambda_\eps^2\,.
%\end{equs}
%As a consequence, we obtain
\begin{equ}[e:L2control]
\|(-\gensy-\gensy^\fw\cGN)^{-1}\genap\psi\|^2\lesssim\lambda_\eps^2\|\cN (-\gensy)^{1/2}\psi\|^2\,.
\end{equ}
We apply the previous estimate to the definition of $\tvN$ together with~\eqref{e:nonlinA+}, so that 
we obtain
\begin{equ}[e:uffa]
\|\tvN\|^2\leq \|(-\gensy-\gensy^\fw\cGN)^{-1}\genap [\tvN +\phi]\|^2\lesssim \lambda_\eps^2\|\cN(-\gensy)^{1/2} [\tvN +\phi]\|^2
\end{equ}
where in the first step, we replaced the truncated operator $\cA^{\eps,+}_{2,n}$ with $\genap$. 
By the weighted a priori estimate on $\tvN$ in Lemma~\ref{l:AnsatzvsGenEq} 
and the bound~\eqref{e:nonlinUnifBound} on $\|(-\gensy)^{-1/2}\nf_\phi\|$, 
we see that the norm at the right hand side of~\eqref{e:uffa} is bounded uniformly in $\eps$. 
Now, since $\lambda_\eps$ goes to $0$ as $\eps\to 0$,~\eqref{e:L2norm2} follows. 
\end{proof}

We now turn to the case $d\geq 3$, whose analysis is more delicate. Our goal is to prove the following theorem. 


\begin{theorem}\label{thm:MainConv}
Let $d\geq 3$. For $n\in\N$, $n\geq 2$, $\eps>0$ and $\phi\in\cS(\T^d)$, let $\vN$ be the solution 
to~\eqref{e:nonlinGenEq}. Then, there exists a unique constant $D_\SHE>0$ independent of $\phi$ such that 
\begin{equ}[e:H1limit]
\lim_{n\to\infty}\lim_{\eps\to 0} \|(-\gensy)^{1/2}\vN\|^2=\langle\phi, (-\cD)\phi\rangle\,.
\end{equ}
where $\cD$ is the operator introduced in Theorem~\ref{thm:main}, given by $\cD\eqdef D_{\SHE}\gensy^\fw$ 
for $\gensy^\fw$ defined according to~\eqref{e:L0w}. 
As a consequence, 
\begin{equs}
%&\lim_{n\to\infty}\lim_{\eps\to 0} \|\vN\|=0\,,\label{e:L2limit}\\
\lim_{n\to\infty}\lim_{\eps\to 0} \|\genam\vN_2-\cD\phi\|=0 \,.\label{e:Diffusivity3}
\end{equs}
\end{theorem}

The proof of~\eqref{e:H1limit} requires to be able to pass to the limit 
for $\|(-\gensy)^{1/2}\vN\|$, first in $\eps$ and then in $n$. 
Since the first limit is the most delicate, we single it out in the statement below, which further provides 
essential information concerning the convergence of $\vN$ to $0$ in $\fock$, and postpone its proof 
to the next subsection. 

\begin{proposition}\label{p:epsLim}
In the setting of Theorem~\ref{thm:MainConv}, let $n$ be fixed. Then, there exists a unique constant $D^n>0$ 
such that 
\begin{equ}
\lim_{\eps\to 0} \|(-\gensy)^{1/2}\vN\|^2=\langle \phi, (-\cD^n)\phi\rangle \,,\label{e:nH1limit}
\end{equ}
where $\cD^n$ is the operator given by $\cD^n\eqdef D^n\gensy^\fw$ 
for $\gensy^\fw$ defined according to~\eqref{e:L0w}. 

Further, as $\eps\to 0$, $\vN$ converges to $0$ in $\fock$.
\end{proposition}

We now turn to the proof of~\eqref{e:H1limit} and consequently of Theorem~\ref{thm:MainConv}, for which, thanks to Proposition~\ref{p:epsLim} 
we only need to ensure that the limit in $n$ of the constants $D^n$ is unique and strictly positive.

\begin{proof}[of Theorem~\ref{thm:MainConv}]
Notice that, once~\eqref{e:H1limit} is established, by Lemma~\ref{l:A-} and smoothness of $\phi$,~\eqref{e:Diffusivity3} 
immediately follows. 

For~\eqref{e:H1limit}, thanks to Proposition~\ref{p:epsLim} it remains to argue uniqueness 
and that $D_\SHE$ is strictly positive. We start with uniqueness. 

Let $v^{\eps,m}$ and $v^{\eps,m+1}$ be 
the solutions to the truncated generator equation~\eqref{e:nonlinGenEq} with $n=m$ and $n=m+1$ 
respectively, and consider the difference $v^{\eps,m+1}-v^{\eps,m}$. Notice that 
\begin{equs}
-\gen (v^{\eps,m+1}-v^{\eps,m})=&-\gen_{2,m+1}v^{\eps,m+1}-\genap v_{m+1}^{\eps,m+1}-\genam v_2^{\eps,m+1}\\
&+\gen_{2,m}v^{\eps,m}+\genap v_m^{\eps,m}+\genam v_2^{\eps,m}\\
=&-\genap (v_{m+1}^{\eps,m+1}-v_m^{\eps,m})-\genam (v_{2}^{\eps,m+1}-v_2^{\eps,m})\,.
\end{equs}
We now test both sides by $v^{\eps,m+1}-v^{\eps,m}$ and, using the antisymmetry of $\gena$, obtain
\begin{equs}
\|(-\gensy)^{1/2}(v^{\eps,m+1}-v^{\eps,m})\|^2&=\langle v^{\eps,m+1}-v^{\eps,m}, -\genap (v_{m+1}^{\eps,m+1}-v_m^{\eps,m})\rangle\\
&=\langle v^{\eps,m+1}-v^{\eps,m}, \genap v_m^{\eps,m}\rangle\label{e:DiffBoundn}
\end{equs}
where the terms containing $\genam (v_{2}^{\eps,m+1}-v_2^{\eps,m})\in\fock_1$ and 
$\genap v_{m+1}^{\eps,m+1}\in\fock_{m+2}$ dropped out because of  
orthogonality of Fock spaces with different indices 
(recall that $v_{1}^{\eps,m+1}=v_1^{\eps,m}=0=v_{m+2}^{\eps,m+1}=v_{m+1}^{\eps,m}$). 
Now, the right hand side can be bounded above by 
\begin{equs}
\langle v^{\eps,m+1}-&v^{\eps,m}, \genap v_m^{\eps,m}\rangle\\
&\leq \tfrac12 \|(-\gensy)^{1/2}(v^{\eps,m+1}-v^{\eps,m})\|^2 +2\|(-\gensy)^{-1/2}\genap v_m^{\eps,m}\|^2\,.
\end{equs}
We can go back to~\eqref{e:DiffBoundn} 
and get 
\begin{equ}[e:Intermezzo]
\|(-\gensy)^{1/2}(v^{\eps,m+1}-v^{\eps,m})\|^2\lesssim \|(-\gensy)^{-1/2}\genap v_m^{\eps,m}\|^2\lesssim \|\cN(-\gensy)^{1/2}v_m^{\eps,m}\|^2
\end{equ}
where the last step follows by Lemma~\ref{l:GenaBound}. To control the last term, 
we apply Proposition~\ref{P:AprioriGeneratorEq}, which, for any $j>1$, gives 
\begin{equs}
\|\cN(-\gensy)^{1/2}v_m^{\eps,m}\|^2\leq m^{-2j}\|\cN^{j}(-\gensy)^{1/2}v_m^{\eps,m}\|^2\lesssim m^{-2j} \|(-\gensy)^{-1/2}\nf_\phi\|^2\,.
\end{equs}
Since $\|(-\gensy)^{-1/2}\nf_\phi\|^2$ is bounded uniformly in $\eps$ by~\eqref{e:nonlinUnifBound}, by applying the triangular inequality at the left hand side of~\eqref{e:Intermezzo} we get
\begin{equ}
\left|\|(-\gensy)^{1/2}v^{\eps,m+1}\|-\|(-\gensy)^{1/2}v^{\eps,m}\|\right|^2\lesssim m^{-2j}\,,
\end{equ}
which, multiplying both sides by $(\|(-\gensy)^{1/2}v^{\eps,m+1}\|+
\|(-\gensy)^{1/2}v^{\eps,m}\|)^2$ and 
invoking~\eqref{e:AprioriMain} once more, implies
\begin{equs}
\Big|\|(-\gensy)^{1/2}&v^{\eps,m+1}\|^2-\|(-\gensy)^{1/2}v^{\eps,m})\|^2\Big|^2\\
&\lesssim m^{-2j}\Big(\|(-\gensy)^{1/2}v^{\eps,m+1}\|+\|(-\gensy)^{1/2}v^{\eps,m})\|\Big)^2\lesssim m^{-2j}\,.
\end{equs}
By Proposition~\ref{p:epsLim}, for any $n$, the limit as $\eps\to 0$ of $\|(-\gensy)^{1/2}\vN\|^2$ exists, is unique 
and given by 
$\langle \phi, (-\cD^n)\phi\rangle=\tfrac12D^n\sum_k (\fw\cdot k)^2 |\hat\phi(k)|^2$,
hence we can take the $\eps$-limit in the above inequality and obtain
\begin{equ}
\left|D^{m+1}-D^m\right|\sum_k (\fw\cdot k)^2 |\hat\phi(k)|^2\lesssim m^{-j}\,,
\end{equ}
which holds for any $\phi$ (in particular for at least one which is not concentrated on $k$'s orthogonal to $w$). 
Therefore, the sequence $\{D^n\}_n$ is Cauchy, thus it converges to a unique limit.
\medskip

We now show that $D_\SHE$ is strictly positive. To do so, we consider once more
the solution $\vN$ of~\eqref{e:nonlinGenEq}. Notice that thanks to the antisymmetry of $\genan$ in~\eqref{e:genn}
\begin{equ}
\|(-\gensy)^{1/2}\vN\|^2=\langle \vN, (-\genn) \vN\rangle=\langle \nf_\phi, (-\genn)^{-1}\nf_\phi\rangle\,.
\end{equ}
Thanks to~\cite[Theorem 4.1]{Komorowski2012}, the term at the 
right hand side can be written as 
%for a matrix $M$, the equality $([M^{-1}]_s)^{-1}=M^\ast  (M_s^{-1}) M$ always holds, where the subscript $s$ denotes 
%the symmetric part. Hence, we have 
\begin{equs}
  \label{e:princvar}
\langle \nf_\phi, &(-\genn)^{-1}\nf_\phi\rangle\\%=\sup_{\psi\in\fock} \left\{2\langle \psi, \nf_\phi\rangle -\langle \psi, ([(-\genn)^{-1}]_s)^{-1}\psi\rangle\right\}\\
%&=\sup_{\psi\in\fock} \left\{2\langle \psi, \nf_\phi\rangle -\langle \psi, (-\genn)^\ast(-\genn)_s^{-1}(-\genn)\psi\rangle\right\}\\
&=\sup_{\psi\in\fock} \left\{2\langle \psi, \nf_\phi\rangle -\langle \psi, (-\gensy)\psi\rangle-\langle \psi, (-\genan)^\ast(-\gensy)^{-1}(-\genan)\psi\rangle\right\}\\
&\geq \sup_{\psi\in\fock_2} \left\{2\langle \psi, \nf_\phi\rangle -\langle \psi, (-\gensy)\psi\rangle-\langle \genap\psi, (-\gensy)^{-1}\genap\psi\rangle\right\}\,,
\end{equs}
where in the last step we restricted the supremum to $\fock_2$, thus making it smaller, and exploited 
Lemma~\ref{lem:essid} and the truncation $P^n_2$ in the definition of $\genan$ in~\eqref{e:genn} to 
simplify the last term. 
%\fabioText{ ma scusa, perche' non affermiamo direttamente la terza linea citandola da QUastel salmhofer etc? senza ridimostrarla?}
%where, in the last equality, we used that the symmetric and antisymmetric part of $\genn$ are respectively given by 
%$\gensy$ and $\genan$. We lower bound the previous by restricting to $\psi\in\fock_2$, so that the projection $P^n_2$ 
%has no effect \fabioText{non mi pare vero che non abbia effetto. Se non ci fosse, sotto avremmo anche il termine $\langle psi A_+ (L0)^{-1}A_-\psi\rangle$} and we get
%\begin{equs}
%\langle \nf_\phi, &(-\genn)^{-1}\nf_\phi\rangle\\
%&\geq \sup_{\psi\in\fock_2} \left\{2\langle \psi, \nf_\phi\rangle -\langle \psi, (-\gensy)\psi\rangle-\langle \psi, {\color{blue}(-\genap)^\ast}(-\gensy)^{-1}(-\genap)\psi\rangle\right\}\\
%&=\sup_{\psi\in\fock_2} \left\{2\langle \psi, \nf_\phi\rangle -\langle \psi, (-\gensy)\psi\rangle-\langle \genap\psi, (-\gensy)^{-1}\genap\psi\rangle\right\}\,.
%\end{equs}
According to Lemma~\ref{e:Abound}, for $\psi\in\fock_2$, there exists a constant $C>0$ such that 
\begin{equ}
\langle \genap\psi, (-\gensy)^{-1}\genap\psi\rangle=\|(-\gensy)^{-1/2}\genap \psi\|^2\leq C\|(-\gensy)^{1/2} \psi\|^2
\end{equ}
which implies a further lower bound of the form 
\begin{equs}
\langle \nf_\phi, (-\genn)^{-1}\nf_\phi\rangle&\geq\sup_{\psi\in\fock_2} \left\{2\langle \psi, \nf_\phi\rangle -\langle \psi, (-\gensy)\psi\rangle-C\langle \psi, (-\gensy)\psi\rangle\right\}\\
%% &=\sup_{\psi\in\fock_2} \left\{2\langle \psi, \nf_\phi\rangle -\langle \psi, (-(1+C)\gensy)\psi\rangle\right\}\\
&=\frac{1}{1+C}\|(-\gensy)^{-1/2}\nf_\phi\|^2\,.
\end{equs}
We are left with estimating $\|(-\gensy)^{-1/2}\nf_\phi\|^2$ from below. Thanks to~\eqref{e:nonlin} and~\eqref{e:gen}, 
we have an explicit expression for the latter, which reads
\begin{equs}
  \label{e:whichreads}
\|(-\gensy)^{-1/2}\nf_\phi\|^2&=\frac{2}{(2\pi)^{d}}\lambda_\eps^2\sum_{\ell,m} \indN{\ell,m}\frac{(\fw\cdot(\ell+m))^2}{|\ell|^2+|m|^2}|\hat\phi(\ell+m)|^2\\
&=\frac{2}{(2\pi)^{d}}\sum_{k}(\fw\cdot k)^2 |\hat\phi(k)|^2 \left(\lambda_\eps^2\sum_{\ell+m=k}\frac{\indN{\ell,m}}{|\ell|^2+|m|^2}\right)\,.
\end{equs}
We can now take the limit for $\eps\to 0$. Since $\phi$ is smooth, by dominated convergence, we can pass 
the limit in $\eps$ inside the sum in $k$ and observe that the quantity in parenthesis converges 
to a strictly positive constant $\bar C\eqdef\int_{|x|_\infty\leq 1}\dd x/(2|x|^2)$. 

By the first part of the statement, we know that the limit as $\eps\to 0$ first and $n\to\infty$ after of 
$\|(-\gensy)^{1/2}\vN\|$ exists, is unique and given by $\langle \phi, -\cD\phi\rangle$. In conclusion, we 
have shown that 
\begin{equs}
D_\SHE \sum_{k}(\fw\cdot k)^2 &|\hat\phi(k)|^2=\langle \phi, -\cD\phi\rangle=\lim_n\lim_\eps \|(-\gensy)^{1/2}\vN\|^2\\
&\geq \frac{1}{1+C}\lim_\eps \|(-\gensy)^{-1/2}\nf_\phi\|^2\geq \frac{4 \bar C}{1+C}\sum_{k}(\fw\cdot k)^2 |\hat\phi(k)|^2\,.
\end{equs}
%
%
%Summarising, we have shown that there exists a strictly positive finite constant $C>0$ such that for any $\eps$ and $n$,
%\begin{equ}[e:almostFinal]
%\|(-\gensy)^{1/2}\vN\|^2\geq \frac{1}{1+C}\sum_{k}(w\cdot k)^2 |\hat\phi(k)|^2 \left(\lambda_\eps^2\sum_{\ell+m=k}\frac{\indN{\ell,m}}{|\ell|^2+|m|^2}\right)\,.
%\end{equ}
%At this point, we want to pass to the limit as $\eps\to 0$ first and $n\to\infty$ after. By the first part of the proof, 
%we know that the limit of the left hand side exists, is unique and given by $\langle \phi, -\cD\phi\rangle$. 
%The sum in the right-hand side instead is independent of $n$ and, since $\phi$ is smooth, 
%by dominated convergence its $\eps$-limit is given by 
%\begin{equ}
%\sum_{k}(w\cdot k)^2 |\hat\phi(k)|^2 \left(\lim_{\eps\to 0}\lambda_\eps^2\sum_{\ell+m=k}\frac{\indN{\ell,m}}{|\ell|^2+|m|^2}\right)= \bar C\sum_{k}(w\cdot k)^2 |\hat\phi(k)|^2
%\end{equ}
%where $\bar C=\int_{|x|_\infty\leq 1}\dd x/(2|x|^2)$. In conclusion, 
%we are allowed to take the limits at both sides of~\eqref{e:almostFinal}, and we obtain
%\begin{equ}
%D_\SHE \sum_{k}(w\cdot k)^2 |\hat\phi(k)|^2=\langle \phi, -\cD\phi\rangle\geq \frac{4 \bar C}{1+C}\sum_{k}(w\cdot k)^2 |\hat\phi(k)|^2\,.
%\end{equ}
Since the previous holds for any $\phi$, $D_\SHE$ is strictly positive. % and the proof of the theorem is completed.
\end{proof}

%\begin{proof}[of Proposition~\ref{p:epsLim}]
%\giuseppeText{MOVE THE LAST SECTION HERE.}
 
\subsection{Proof of Proposition~\ref{p:epsLim}}\label{sec:Diff3}

The proof of Proposition~\ref{p:epsLim} is the most delicate of the present paper. 
Since in $d\geq 3$ we have no replacement lemma (and we do not even expect it to hold), we need to resort to 
a different strategy. Our proof is centred around a suitable expansion of the norms of $\vN$ (see e.g.~\eqref{e:defDn} below) 
which involves several subsequent applications of the operators $\genap$, $\genam$ interposed 
by inverses of $-\gensy$. In order to see why such an expansion is meaningful, 
it is convenient to write it terms of operators which are skew-Hermitian and bounded in $\fock$ (uniformly in $\eps$), 
and this is the first thing we will do. 

Let $n\in\N$, $n\geq 2$, be fixed throughout the section. Let $\eps>0$ and 
$\cA^\eps_{2,n}$ be given as in \eqref{e:genn}. 
We introduce the operators $T^\eps_{2,n}$ and for $\sigma\in\{+,-\}$, $T^{\eps,\sigma}, T^{\eps,\sigma}_{2,n}$, 
according to
\begin{equs}[eq:papiro]
  T^\eps_{2,n}&\eqdef (-\cL_0)^{-1/2}(\cA^\eps_{2,n})(-\cL_0)^{-1/2}\,,\\
  T^{\eps,\sigma}&\eqdef (-\cL_0)^{-1/2}(\cA^\eps_{\sigma})(-\cL_0)^{-1/2}\,,\\
  T^{\eps,\sigma}_{2,n}&\eqdef (-\cL_0)^{-1/2}(\cA^{\eps,\sigma}_{2,n})(-\cL_0)^{-1/2}\,.
\end{equs}
Note that, by Lemma~\ref{l:GenaBound} and in particular~\eqref{e:Tbound}, 
each of the operators above is uniformly bounded in $\fock$ in $\eps$ (the bound depends on $n$, 
but this is irrelevant as $n$ is fixed) and, by Lemma~\ref{lem:essid}, it is easy to see that $T^\eps_{2,n}$ 
is skew-Hermitian. 

To have a notation for multiple applications of the operators $T^{\eps,\sigma}$, $\sigma\in\{+,-\}$, 
we set, for $m\ge1$ and $a\ge 0$, $\Pi^{(n)}_{a,m}$ to be the set
of simple random walk paths $p=(p_0,\dots,p_a)$ consisting of $a$ steps, starting at height $1$, i.e. $p_0=1$,
ending at height $m$, i.e. $p_a=m$, that do not reach height $n+1$ and that can have
height $1$ only at the endpoints.  
Given $p\in \Pi^{(n)}_{a,m}$ we let
$\mathcal T^\eps_p$ be defined according to 
\begin{equ}[e:TTT]
\ST^\eps_p\eqdef  T^{\eps,\sigma_a}T^{\eps,\sigma_{a-1}}\dots T^{\eps,\sigma_1}
\end{equ}
where $\sigma_i=+$, if $p_{i}-p_{i-1}=1$, and $\sigma_i=-$, if $p_{i}-p_{i-1}=-1$. 
In other words, there is a factor $T^{\eps,+}$
(resp. $T^{\eps,-}$) whenever the path takes a step up
(resp. down)\footnote{For instance, for $a=4$ the path
  $p=(1,2,3,2,3)\in \Pi^{(n)}_{4,3}, n\ge 4$ corresponds to $\mathcal T^\eps_p=T^{\eps,+}T^{\eps,-}T^{\eps,+}T^{\eps,+}$.}.

We enclose the most technical aspect of the proof of Proposition~\ref{p:epsLim} 
in Proposition~\ref{p:limite} below. Before giving the proof of the latter, 
we will show how it implies the former. 

\begin{proposition}
  \label{p:limite}
In the notations introduced above, for every $a>0,m\ge1$ and $p\in \Pi^{(n)}_{a,m}$, the limit
\begin{equation}
  \label{e:cp}
  \lim_{\eps\to0} \langle e_{-k},\mathcal T^\eps_pe_k\rangle= c(p)\frac{(\fw\cdot k)^2}{|k|^2}
\end{equation}
exists, with $c(p)$ depending only on the path $p$.  Moreover, $c(p)=0$ unless $a\in 2\N, m=1$. 
Furthermore, for every $a>0, m\ge2$, $q_1,\dots,q_m\in \mathbb Z^d_0$, $m'\geq 1$ and 
$p\in \Pi^{(n)}_{a,m'}$, one has
\begin{equation}
\label{e:nullo}  \lim_{\eps\to0} \langle e_{-q_1}\otimes\dots\otimes e_{-q_m},\mathcal T^\eps_p e_k\rangle= 0\,.
\end{equation}
\end{proposition}

\begin{proof}[of Proposition \ref{p:epsLim} given Proposition \ref{p:limite}]
  We start with \eqref{e:nH1limit}.
Let $\vN$ be the solution of~\eqref{e:nonlinGenEq}.  As shown in the proof of Lemma~\ref{l:A-}, we have that
$\vN=\sum_k\hat\phi(k)\vN[-k]$, where, for $j\in\Z^d_0$, $\vN[j]$
solves the truncated generator equation~\eqref{e:nonlinGenEq} with $\phi=e_{-j}$. 
Hence,
recalling~\eqref{e:nonlinA+} and~\eqref{e:nonlinFourier}, the object
whose limit we want to determine equals
\begin{equs}
\langle \vN,(-\gensy)\vN\rangle&=\langle \vN,(-\genn)\vN\rangle=\langle \vN,\nf_\phi\rangle\\
&=\sum_{k,j}\hat\phi(k)\hat\phi(-j)\langle \vN[-k],\nf_{j}\rangle\\
&=\sum_k |\hat\phi(k)|^2\langle \vN[-k],\nf_{k}\rangle =\sum_k |\hat\phi(k)|^2\|(-\gensy)^{1/2} \vN[-k]\|^2
\end{equs} 
where, as in~\eqref{e:TransInv}, we used translation invariance to remove summands in which $j\neq -k$. 
Each of the squared norms at the right hand side are uniformly (in $\eps$) bounded by a quantity that grows at most 
quadratically in $k$. In fact, by Proposition \ref{P:AprioriGeneratorEq}, each squared norm is bounded by a constant times
$ \|(-\gensy)^{-1/2}\nf_k\|^2$. In turn,
by \eqref{e:whichreads} we have 
\begin{equ}
  \|(-\gensy)^{-1/2}\nf_k\|^2
=2(\fw\cdot k)^2 \left(\lambda_\eps^2\sum_{\ell+m=k}\frac{\indN{\ell,m}}{|\ell|^2+|m|^2}\right)\,
\end{equ}
and the sum is uniformly bounded by a constant independent of $\eps$ and $k$.
Since $\phi$ is smooth,  to prove~\eqref{e:nH1limit} it is enough (by dominated convergence) 
to show that there exists a unique $D^n$ such that for all $k$
\begin{equ}\label{e:nH1limitk}
  \lim_{\eps\to 0} \|(-\gensy)^{1/2}v^{\eps,n}[-k]\|^2
  = \lim_{\eps\to 0} \langle \nf_k,v^{\eps,n}[-k]\rangle
  =\tfrac12D^n (\fw\cdot k)^2\,. 
\end{equ}

Recall the definition of the operators $T^\eps_{2,n}$ and $T^{\eps,\sigma}$, $\sigma\in\{+,-\}$ in~\eqref{eq:papiro}. 
Note that, thanks to \eqref{e:nonlinA+} and \eqref{e:nonlinGenEq},
\begin{equ}
  \label{e:ven-k}
  v^{\eps,n}[-k]=(-\genn)^{-1}\cA^\eps_+e_k=\frac{|k|}{\sqrt2}(-\gensy)^{-1/2}(I-T^\eps_{2,n})^{-1}T^{\eps,+}e_k\,,
\end{equ}
where in the last equality we further used~\eqref{e:gen}. Therefore, 
\begin{equ}
  \label{e:integrs}
  \langle \nf_k,v^{\eps,n}[-k]\rangle=  \frac{|k|^2}2\langle T^{\eps,+}e_{-k},(I-T^\eps_{2,n})^{-1}T^{\eps,+}e_k\rangle\,.
\end{equ}
Note that $T^\eps_{2,n}$ is skew-Hermitian and, by Lemma \ref{l:GenaBound} and in particular~\eqref{e:Tbound}, 
it is bounded uniformly in $\eps$ (for $n$ fixed).
Then, we write the inverse operator above as
\begin{equs}[e:ss]
 \frac2{|k|^2} \langle \nf_k,v^{\eps,n}[-k]\rangle&= \langle  T^{\eps,+}e_{-k},
  \Big(\int_0^\infty e^{-s} e^{s T^\eps_{2,n}}\dd s\Big)\,T^{\eps,+}e_k\rangle\\
  &=\int_0^\infty e^{-s} \langle T^{\eps,+}e_{-k},e^{s T^\eps_{2,n}}T^{\eps,+}e_k\rangle \dd s%\\
%  &=\int_0^\infty e^{-s}
% \langle T^{\eps,+}e_{-k},  \sum_{a\ge0}\frac{s^a}{a!}(T^\eps_{2,n})^aT^{\eps,+}e_k\rangle\dd s%\\
% &=\int_0^\infty e^{-s}
%\sum_{a\ge0}\frac{s^a}{a!} \langle T^{\eps,+}e_{-k},  (T^\eps_{2,n})^aT^{\eps,+}e_k\rangle\dd s
\end{equs}
where in the second step we applied Fubini, which is allowed as the operator 
$e^{s T^\eps_{2,n}}$ has norm at most $1$, since the spectrum of $T^\eps_{2,n}$ is purely imaginary. 
For the same reason, once we take the $\eps$-limit at both sides we can apply dominated convergence 
and pass the limit inside the integral in $s$, i.e. 
\begin{equs}\label{e:defDn}
\lim_{\eps\to0} \frac2{|k|^2} \langle \nf_k,v^{\eps,n}[-k]\rangle&= \int_0^\infty e^{-s} \lim_{\eps\to0}\langle T^{\eps,+}e_{-k},e^{s T^\eps_{2,n}}T^{\eps,+}e_k\rangle \dd s\\
&=\int_0^\infty e^{-s} \lim_{\eps\to0}\langle T^{\eps,+}e_{-k},\sum_{b\ge0}\frac{s^b}{b!}(T^\eps_{2,n})^bT^{\eps,+}e_k\rangle \dd s\\
%&=\int_0^\infty e^{-s}
%\lim_{\eps\to0}\sum_{a\ge0}\frac{s^a}{a!} \langle T^{\eps,+}e_{-k},  (T^\eps_{2,n})^aT^{\eps,+}e_k\rangle\dd s\\
&=\int_0^\infty e^{-s}\sum_{b\ge0}\frac{s^b}{b!}\lim_{\eps\to0} \langle T^{\eps,+}e_{-k},  (T^\eps_{2,n})^bT^{\eps,+}e_k\rangle\dd s
\end{equs}
where in the last step we could exchange the limit with both the summation over $b$ and the scalar product 
because, for fixed $n,s$, the norm of $s T^\eps_n$ is uniformly bounded (recall that after the first, in the 
subsequent steps $s$ can be treated as fixed). 
Therefore, we are left to determine, for every $b\geq 0$, the limit of
\begin{equ}
\langle T^{\eps,+}e_{-k},  (T^\eps_{2,n})^bT^{\eps,+}e_k\rangle=-\langle e_{-k},  T^{\eps,-}(T^\eps_{2,n})^bT^{\eps,+}e_k\rangle\,.
\end{equ}
Let us begin with a few remarks. First
of all, since $T^\eps_{2,n}$ is skew-Hermitian, the scalar product is zero unless $b$ is even, hence we set $a\in\N$ 
to be such that $b=2a$. 
Next, we expand 
\begin{equs}\label{e:expansion}
\langle e_{-k},  T^{\eps,-}(T^\eps_{2,n})^{2a}T^{\eps,+}e_k\rangle&=\langle e_{-k},  T^{\eps,-}(T^{\eps,-}_{2,n}+T^{\eps,+}_{2,n})^{2a}T^{\eps,+}e_k\rangle \\
&=\sum_{p\in\widetilde\Pi_{2a+2,1}}\langle e_{-k},  T^{\eps,\sigma_{2a+2}}T^{\eps,\sigma_{2a+1}}_{2,n}\dots T^{\eps,\sigma_2}_{2,n}T^{\eps,\sigma_1}e_k\rangle
\end{equs}
where $\widetilde\Pi_{2a+2,1}$ is the set of all walks $p$ which start at time $1$ from $1$, i.e. $p_0=1$, 
ends at time $2a+2$ at $1$, i.e. $p_{2a+2}=1$, have increments of size $1$ and are such that 
$p_1-p_0=1=p_{2a+1}-p_{2a+2}$. If either of the previous 
conditions were violated we would be taking the scalar product of quantities in Fock spaces with 
different indices, which would be $0$. In the above, 
for $p\in\widetilde\Pi_{2a+2,1}$ and $i=1,\dots, 2a+2$, 
$\sigma_i=+$ if $p_{i}-p_{i-1}=1$ and $\sigma_i=-$ if $p_{i}-p_{i-1}=-1$. Note that the last condition on 
the $p$'s ensures that $\sigma_1=1=-\sigma_{2a+2}$ (so that the $T^{\eps,-}$ and $T^{\eps,+}$ in the line 
above are present). 

Because of the
projector $P^n_2$ entering the definition of $T^{\eps,\sigma}_{2,n}$,
if a path $p$ reaches height $n+1$ or height $1$ (except at the
endpoints), then the corresponding scalar product is annihilated. 
Hence, we can replace the index set of the sum at the right hand side of~\eqref{e:expansion} to 
$\Pi^{(n)}_{2a+2,1}$, and get
\begin{equs}%[e:ClaimWalks]
\langle e_{-k},  T^{\eps,-}(T^\eps_{2,n})^{2a}T^{\eps,+}e_k\rangle&=\sum_{p\in\Pi^{(n)}_{2a+2,1}}\langle e_{-k},  T^{\eps,\sigma_{2a+1}}T^{\eps,\sigma_{2a}}_{2,n}\dots T^{\eps,\sigma_2}_{2,n}T^{\eps,\sigma_1}e_k\rangle\\
&=\sum_{p\in\Pi^{(n)}_{2a+2,1}}\langle e_{-k},  T^{\eps,\sigma_{2a+1}}T^{\eps,\sigma_{2a}}\dots T^{\eps,\sigma_2}T^{\eps,\sigma_1}e_k\rangle\\
&=\sum_{p\in\Pi^{(n)}_{2a+2,1}}\langle e_{-k},  \ST^\eps_p e_k\rangle\,,
\end{equs}
where in the second step replaced all instances of $T^{\eps,\sigma}_{2,n}$ with $T^{\eps,\sigma}$, since 
for $p\in\Pi^{(n)}_{2a+2,1}$ the projection $P^n_2$ acts as the identity operator. 

At last, getting back to~\eqref{e:defDn} and applying~\eqref{e:cp} in Proposition~\ref{p:limite}, 
we obtain
\begin{equs}
\lim_{\eps\to0} \langle \nf_k,v^{\eps,n}[-k]\rangle&=\frac{|k|^2}{2}\int_0^\infty e^{-s}\sum_{a\ge0}\frac{s^{2a}}{(2a)!}\lim_{\eps\to0} \sum_{p\in\Pi^{(n)}_{2a+2,1}}\langle e_{-k},  \ST^\eps_p e_k\rangle\dd s\\
&=\frac{|k|^2}{2}\int_0^\infty e^{-s}\sum_{a\ge0}\frac{s^{2a}}{(2a)!} \sum_{p\in\Pi^{(n)}_{2a+2,1}}\lim_{\eps\to0}\langle e_{-k},  \ST^\eps_p e_k\rangle\dd s\\
&=\frac{(\fw\cdot k)^2}{2}\int_0^\infty e^{-s}\sum_{a\ge0}\frac{s^{2a}}{(2a)!} \sum_{p\in\Pi^{(n)}_{2a+2,1}}c(p)\dd s\\
&=:\frac12D^n (\fw\cdot k)^2
\end{equs}
where the last constant $D^n$ is uniquely defined by the expression before, so that~\eqref{e:nH1limit} follows. 
\medskip
%
%
%
%\newpage
%into products of $2a$ operators $T^{\eps,\sigma}_{2,n}$, with an
%additional $T^{\eps,-}$ (resp. $T^{\eps,+}$) on the left
%(resp. right), we get $2^{2a}$ such products that we identify with
%simple-random walk paths length $2a+2$ steps, starting from height $1$
%at time $0$ zero and decreasing (resp. increasing) by $1$ whenever a
%factor $T^{\eps,-}_{2,n}$ or $T^{\eps,-}$ (resp.  a factor
%$T^{\eps,+}_{2,n}$ or $T^{\eps,+}$) is encountered, reading the
%product from left to right.
%
%
%
%\newpage 
%
%
%
%Assume now that for every  $k\in \mathbb Z^d_0$, $a\ge0, n\ge 2$ the limit
%\begin{multline}
%  \label{esiste}
%  \lim_{\eps\to0}\langle T^{\eps,+}e_{-k}, (T^\eps_{2,n})^aT^{\eps,+}e_k\rangle=  -\lim_{\eps\to0}\langle e_{-k}, T^{\eps,-}(T^\eps_{2,n})^a T^{\eps,+}e_k\rangle\\= d(n,a)\frac{(w\cdot k)^2}{|k|^2}
%\end{multline}
%exists  and has the form as in the r.h.s., with $d(n,a)$ depending only on $n,a$.
%If this is the case, then
%\begin{multline}
%  \label{e:defDn}
%  \frac{(w\cdot k)^2}{|k|^2}D^n\eqdef \frac{(w\cdot k)^2}{|k|^2}\int_0^\infty e^{-s}\sum_{a\ge0}d(n,a)\frac{s^a}{a!}\dd s \\
%  = \int_0^\infty e^{-s}\sum_{a\ge0}\frac{s^a}{a!} \lim_{\eps\to0}\langle T^{\eps,+}e_{-k}, (T^\eps_{2,n})^aT^{\eps,+}e_k\rangle\dd s\\
%  = \int_0^\infty e^{-s}\lim_{\eps\to0}
% \langle T^{\eps,+}e_{-k},  \sum_{a\ge0}\frac{s^a}{a!}(T^\eps_{2,n})^aT^{\eps,+}e_k\rangle\dd s\\
%    = \int_0^\infty e^{-s}\lim_{\eps\to0}
% \langle T^{\eps,+}e_{-k}, e^{s T^\eps_{2,n}}T^{\eps,+}e_k\rangle\dd s.
%\end{multline}
%In the second equality we have exchanged the limit $\eps\to0$, the
%summation over $a$ and the scalar product; this is allowed because,
%for fixed $n,s$, the norm of $s T^\eps_n$ is uniformly bounded. Since
%the exponential $e^{s T^\eps_{2,n}}$ has norm at most $1$, by
%dominated convergence we exchange the limit $\eps\to0$ and the average
%with respect to the probability measure $\exp(-s)ds$.  Recalling
%\eqref{e:ss}, we see then that \eqref{e:nH1limitk} holds, with $D^n$
%defined as in the first line of \eqref{e:defDn}, provided that we
%prove \eqref{esiste}.  Note that in the definition of $D^n$ we do not
%interchange the summation and the integral w.r.t. $s$, because we have
%no good control of the growth in $a$ of $d(n,a)$, uniformly in $n$.
%
%
%
%%
%\newpage
%In order to prove \eqref{esiste}, we start with a few remarks. First
%of all, since $T^\eps_{2,n}$ is skew-Hermitian, the scalar product in
%\eqref{esiste} is zero unless $a$ is even (i.e., \eqref{esiste} holds
%with $d(n,a)=0$ if $a$ is odd).  Expanding
%\[T^{\eps,-}(T^\eps_{2,n})^{2a}T^{\eps,+}=T^{\eps,-}(T^{\eps,-}_{2,n}+T^{\eps,+}_{2,n})^{2a}T^{\eps,+},a\ge0\]
%into products of $2a$ operators $T^{\eps,\sigma}_{2,n}$, with an
%additional $T^{\eps,-}$ (resp. $T^{\eps,+}$) on the left
%(resp. right), we get $2^{2a}$ such products that we identify with
%simple-random walk paths length $2a+2$ steps, starting from height $1$
%at time $0$ zero and decreasing (resp. increasing) by $1$ whenever a
%factor $T^{\eps,-}_{2,n}$ or $T^{\eps,-}$ (resp.  a factor
%$T^{\eps,+}_{2,n}$ or $T^{\eps,+}$) is encountered, reading the
%product from left to right. The height of the path at the starting
%point is taken to be $1$ because $e_k$ is the kernel of a random
%variable in the first Wiener chaos.  However, only some of these paths
%give a non-zero contribution when they are taken in the scalar product
%in \eqref{esiste}. First of all, the paths must contain an equal
%number of up-steps and down-steps, i.e. they must terminate at height
%1. This guarantees that the two vectors we are taking the scalar
%product of, are both in the first chaos. Secondly, because of the
%projector $P^n_2$ entering the definition of $T^{\eps,\sigma}_{2,n}$,
%the path cannot reach height $n+1$ or height $1$ (except at the
%endpoints).  In other words, the only paths that can give a non-zero
%contribution in the scalar products are those in $ \Pi^{(n)}_{2a+2,1}$
%for some $a\ge0$.
%
%
%Given $p\in \Pi^{(n)}_{2a+2,m}$, we let $\mathcal T^{\eps,(n)}_p$ be
%the product of $a$ operators $T^{\eps,\pm}_{2,n}$ (with an additional
%$T^{\eps,-}$ on the left and an additional $T^{\eps,+}$, in the correct
%order corresponding to the path $p$. Going back to \eqref{esiste}, by
%the previous discussion one has
%\begin{equation}
%  \label{e:restrict}
%  \langle e_{-k}, T^{\eps,-}(T^\eps_{2,n})^{2a} T^{\eps,+}e_k\rangle=
%  \sum_{p\in \Pi^{(n)}_{2a+2,1}} \langle e_{-k},\mathcal T^{\eps,(n)}_p e_k
%    \rangle.
%\end{equation}
%On the other hand, we can replace the operators
%$T^{\eps,\sigma}_{2,n}$ entering the definition of $\mathcal
%T^{\eps,(n)}_p$ by $T^{\eps,\sigma}$, because for $p\in
%\Pi^{(n)}_{2a+2,1}$ the projectors $P^n_2$ act as the identity
%operator. In other words, we can replace $\mathcal T^{\eps,(n)}_p$ by
%$\mathcal T^\eps_p$: the only effect of the projector $P^n_2$ is then to
%restrict, in \eqref{e:restrict}, the set of admissible paths to $\Pi^{(n)}_{2a+2,1}$. Finally, \eqref{e:cp} gives the
%existence of the $\eps\to0$ limit of $\langle e_{-k},\mathcal
%T^{\eps}_p e_k \rangle$ and \eqref{esiste} (and therefore
%\eqref{e:nH1limit}) follows.
%\newpage 

We are left to show that $v^{\eps,n}$ converges to zero in $\fock$. Let us begin by writing
\begin{equ}
   \|v^{\eps,n}\|^2= \sum_{m=2}^n \sum_{k_{1:m}}|v^{\eps,n}_m(k_{1:m})|^2,
 \end{equ}
 with $v^{\eps,n}_m$ the component of $v^{\eps,n}$ in the $m^{th}$ Wiener chaos.
Let $N>0$. We split the norm as
\begin{equs}
  \|v^{\eps,n}\|^2 &= \sum_{m=2}^n\sum_{|k_{1:m}|>N}|v^{\eps,n}_m(k_{1:m})|^2+ \sum_{m=2}^n\sum_{|k_{1:m}|\le N}|v^{\eps,n}_m(k_{1:m})|^2
  \\
  &\le
  \frac2{N^2}\sum_{m=2}^n\sum_{|k_{1:m}|>N}\frac{|k_{1:m}|^2}2|v^{\eps,n}_m(k_{1:m})|^2+ \sum_{m=2}^n\sum_{|k_{1:m}|\le N}|v^{\eps,n}_m(k_{1:m})|^2\\
  &\le  \frac2{N^2}\|(-\cL_0)^{1/2}v^{\eps,n}\|^2+ \sum_{m=2}^n\sum_{|k_{1:m}|\le N}|v^{\eps,n}_m(k_{1:m})|^2\,.
\end{equs}
By the just proved \eqref{e:nH1limit}, if we first take the limit $\eps\to 0$ and then $N\to\infty$, 
we deduce that the convergence of $v^{\eps,n}$ to zero in $\Gamma L^2$ follows provided that
\begin{equation}
  \label{e:componentem2}
  \lim_{\eps\to0}|v^{\eps,n}_m(k_{1:m})|^2= \lim_{\eps\to0}|\langle e_{k_1}\otimes\dots\otimes e_{k_m},v^{\eps,n}\rangle|^2=0
\end{equation}
for every fixed $2\le m\le n$ and any given $k_{1:m}$. Since $v^{\eps,n}=\sum_{k}\hat \phi(k)v^{\eps,n}[-k]$ and $\phi$ is smooth, it is enough to show that
\begin{equ}
  \lim_{\eps\to0}\langle e_{k_1}\otimes\dots\otimes e_{k_m},v^{\eps,n}[-k]\rangle=0
\end{equ}
for every fixed $k$. To this purpose, by \eqref{e:ven-k} and, following steps similar to those after \eqref{e:integrs}, 
we come to
\begin{equs}
  \lim_{\eps\to0}&\langle e_{k_1}\otimes\dots\otimes e_{k_m},v^{\eps,n}[-k]\rangle\\
  &=\frac{|k|}{|k_{1:m}|}\int_0^\infty e^{-s}\sum_{b\ge0}\frac{s^b}{b!}\lim_{\eps\to0}
  \langle e_{k_1}\otimes\dots\otimes e_{k_m},(T^{\eps,+}_{2,n}+T^{\eps,-}_{2,n})^b T^{\eps,+}e_k\rangle\\
  &=\frac{|k|}{|k_{1:m}|}\int_0^\infty e^{-s}\sum_{b\ge0}\frac{s^b}{b!}\sum_{p\in\Pi^{(n)}_{b,m}}\lim_{\eps\to0}
  \langle e_{k_1}\otimes\dots\otimes e_{k_m},\ST^\eps_pe_k\rangle\,.
\end{equs}
Finally, \eqref{e:nullo} implies that the right-hand side of the previous expression is zero, thereby proving \eqref{e:componentem2}.
\end{proof}

We now turn to the proof of Proposition~\ref{p:limite}, for which we will need a preliminary lemma. 
Given $p\in \Pi^{(n)}_{a,m}$, $\mathcal T^\eps_pe_k$  belongs to the $m^{th}$ Fock space 
and we write the Fourier transform of its kernel in the form
\begin{equ}\label{eq:formkernel}
f^\eps(k_{1:m},k)= \eps^{\frac d2(m-1)}F^\eps_p(\eps k_{1:m};\eps k)\,,
\end{equ}
in which we singled out the scaling factor $\eps^{\frac d2(m-1)}$. 
We derive an expression for $F^\eps_p$ in terms of polynomials, whose form and main properties are summarised 
in the following lemma.

\begin{lemma}\label{lemma:kernels}
Given $p\in \Pi^{(n)}_{a,m}$ and $a\ge 1$, let $M=(1+a-m)/2\,(\ge 0)$ be the number of $T^{\eps,-}$ in the product 
$\mathcal T^\eps_p$.
The function $F_p^\eps(x_{1:m};x)$ in \eqref{eq:formkernel} % is zero if $x_1+\dots+x_m\ne x$.  If $x_1+\dots+x_m\ne x$ and $M>0$, then it 
  can be written as
  \begin{equs}[eq:formkernelG]
    F_p^\eps(x_{1:m};x)&=\1_{x_1+\dots+x_m= x}\frac{(\fw\cdot x)}{|x|} \frac{\iota^a}{|x_{1:m}|}\times\\
    &\times \sum_{g\in\mathcal G[p]}\eps^{d M}\sum_{\substack{y_{1:M}\in (\eps \Z^d_0)^M:\\|y_i|_
        \infty\le 1, i\le M}} \frac{P_g(y_{1:M};x_{1:m})}{Q_g(y_{1:M};x_{1:m})}I_g(y_{1:M};x_{1:m})\,,
  \end{equs}
  where
  \begin{enumerate}[noitemsep]
  \item [(i)] $\mathcal G[p]$ is a finite set whose cardinality depends only on the path $p$,
  
  \item  [(ii)] $I_g$ is a product of  indicator functions, each  imposing that certain linear combinations of its arguments have $|\cdot|_\infty$ norm in $(0,1]$,
    
  \item [(iii)] $P_g$ (resp. $Q_p$) is a homogenous  real-valued polynomial of degree $a-1$ (resp. $2(a-1)$) 
  in $y_{1:M},x_{1:m}$. Further, $Q_g\ge0$ and it does not vanish on the support of $I_g$,
    
  \item [(iv)] if $m=1$, then $Q_g(y_{1:M};0)>0$ for almost every $y_{1:M}$.
    %If the left-most operator in $\mathcal T^\eps_p$ is $T^{
%\eps,-}$, then $Q_g(y_{1:M};0)>0$ whenever $y_{i}\ne 0$ for all $i\le M$;
    
  %% \item $A_g(x_{1:m})$ is a non-identically vanishing linear function of $x_{1:m}$ and if $m=1$, then $A_g(x)=(\fw\cdot x)$;
    
  
  \end{enumerate}
  In particular, if $m=1$ then there exists a polynomial $\tilde P_g$ of degree $a-2$ such that 
  \begin{equs}[eq:casem1]
    \langle e_{-k_1}\,,\mathcal T^\eps_pe_k\rangle&=\iota^a \1_{k_1= k}\frac{(\fw\cdot k)^2}{|k|^2}\\
    &\times \sum_{g\in\mathcal G[p]}
     \eps^{d M}\sum_{\substack{y_{1:M}\in (\eps \Z^d_0)^M:\\|y_i|_\infty\le 1, i\le M}} \frac{\tilde P_g(y_{1:M};\eps k)}{Q_g(y_{1:M};\eps k)}I_g(y_{1:M};\eps k)\,.
  \end{equs}
  %where $\tilde P_g$ is a polynomial of degree $a-2$.
\end{lemma}
% The case $M=0$ has been left out because it is a bit special (and much easier, see later).
We emphasise that the functions $P_g, Q_g, I_g$ and $\tilde P_g$ are {\it $\eps$-independent}. 
The only $\eps$-dependence in $F^\eps_p$ comes from the fact that the Riemann sums runs over $\eps \Z^d_0$. 


\begin{proof}%[of Lemma \ref{lemma:kernels}]
The statement is proven via induction on $a$. In order to verify condition {\it (iv)}, we will 
simultaneously show that the $Q_g$'s further satisfy
\begin{enumerate}
\item[$(iii')$] $Q_g(y_{1:M};x_{1:m})$ is the product of homogeneous quadratic polynomials $q_r$, $1\leq r\leq a-1$, 
each being the sum of squared norms of linear combinations of its arguments, and, for every $1\leq r\leq a-1$, 
$q_r\neq |x_1+\dots+x_m|^2$. 
\end{enumerate}

From now on, since we are in dimension $d\geq3$, we will not write $\lambda_\eps$ but its value $\eps^{d/2-1}$ 
so that the scaling exponents become more transparent. 

We begin our induction with $a=1$. In this case, if $p\in  \Pi^{(n)}_{a,m}$, 
then necessarily $m=2,M=0$ and $\mathcal T^\eps_p=T^{\eps,+}$. 
Then, from~\eqref{eq:papiro} and~\eqref{e:gen} with $n=1$ (see also \eqref{eq:T+} below), one sees that
\begin{equ}\label{eq:casoa=1}
    F^\eps_p(x_{1:2};x)= \1_{x_1+x_2= x}\frac{(\fw\cdot x)}{|x|}\frac{\iota}{|x_{1:2}|}\frac{-4}{(2\pi)^{d/2}}\mathbb J^1_{x_1,x_2}
\end{equ}
because the kernel of $e_k$ equals $\1_{\cdot =k}$. This expression is of the correct form \eqref{eq:formkernelG}, 
with $\cG[p]=\emptyset$, $P_g\equiv-4/(2\pi)^{d/2}$, $Q_g\equiv 1$, $I_g\equiv \mathbb J^1_{x_1,x_2}$, 
so that $(i)-(iii)$ and $(iii')$ are clearly satisfied. 

  
As for the induction step, we assume the validity of the statement for a given $a\ge2$ and take $p\in\Pi^{(n)}_{a+1,m}$. 
Write $\mathcal T^\eps_p=T^{\eps,\sigma}\mathcal T_{p'}^\eps$ for some $p'\in\Pi^{(n)}_{a,m-\sigma_{a-1}}$. 
Consider first the case $\sigma=+$, so that $\sigma_{a-1}=1$. 
Recalling the definition of $T^{\eps,+}$ and~\eqref{eq:formkernel}, 
we have 
 \begin{equs}
 (T^{\eps,+}f^\eps_{p'})(k_{1:m};k)&=-\frac{1}{(2\pi)^{d/2}}\frac{4\eps^{d/2-1+d/2(m-2)}}{m}\frac\iota{|k_{1:m}|}\times\\
\times
\sum_{1\le i<j\le m}&\frac{\fw\cdot(k_i+k_j) \indN{k_i,k_j}}{\sqrt{|k_i+k_j|^2+|k_{1:m\setminus\{i,j\}}|^2}} F^\eps_{p'}(\eps k_i+\eps k_j,\eps k_{1:m\setminus\{i,j\}};\eps k)\\
&=\eps^{d/2(m-1)}F^\eps_{p}(\eps k_{1:m};\eps k), \label{eq:T+}
\end{equs} 
with
\begin{equs}\label{eq:F+}
F^\eps_{p}(x_{1:m};x)=&\frac\iota{|x_{1:m}|}\sum_{1\le i<j\le m}
\frac{-4}{m(2\pi)^{d/2}}\times\\
&\times \frac{\fw\cdot (x_i+x_j)}{\sqrt{|x_i+x_j|^2+|x_{1:m\setminus\{i,j\}}|^2}}\mathbb J^1_{x_i, x_j}F^\eps_{p'}(x_i+x_j, x_{1:m\setminus\{i,j\}};x).
\end{equs}
By induction, $F_{p'}$ has the form in~\eqref{eq:formkernelG}, 
and therefore so does $F_p$ once we set 
$\cG[p]=\{(g',(i,j)),g'\in \cG[p'],1\le i<j\le m\}$ and, for $g=(g',(i,j))\in \cG[p]$,
\begin{equs}
P_g(y_{1:M};x_{1:m})&=\frac{-4}{m(2\pi)^{d/2}} (\fw\cdot (x_i+x_j))P_{g'}(y_{1:M};x_i+x_j,x_{1:m\setminus\{i,j\}}),\\
Q_g(y_{1:M};x_{1:m})&=(| x_i+x_j|^2+|x_{1:m\setminus\{i,j\}}|^2)Q_{g'}(y_{1:M};x_i+x_j,x_{1:m\setminus\{i,j\}})\\
    I_g(y_{1:M};x_{1:m})&= \mathbb J^1_{x_i,x_j}I_{g'}(y_{1:M};x_i+x_j,x_{1:m\setminus\{i,j\}})\,.\label{e:casosale}
\end{equs}
Similarly, if $\sigma=-$ (so that $\sigma_{a-1}=-1$), by~\eqref{e:gen} one has 
\begin{equs}
  (T^\eps_-f^\eps_{p'})(k_{1:m};k) =&\frac{\eps^{d/2-1}}{(2\pi)^{d/2}}\frac{4\iota (m+1)}{|k_{1:m}|}
  \sum_{j=1}^{m}(\fw\cdot k_j)\\
  &\times\sum_{y,z\in \mathbb
    Z^d:y+z=k_j}\indN{y,z}\frac{\eps^{dm/2}F^\eps_{p'}(\eps y,\eps z,
    \eps k_{\{1:m\}\setminus\{j\}};\eps k)}{\sqrt{|y|^2+|z|^2+|k_{\{1:m\}\setminus\{j\}}|^2}}
  \\
  =&\eps^{d/2(m-1)}F^\eps_{p}(\eps k_{1:m};\eps k),  
\end{equs}
where in this case
\begin{equs}  \label{eq:F-}
F^\eps_{p}&(x_{1:m};x)= \frac{4\iota (m+1)}{(2\pi)^{d/2}}\frac1{|x_{1:m}|} \sum_{j=1}^{m}(\fw\cdot x_j)\1_{0<|x_j|_\infty\le 1}\\
&\times\sum_{y\in \eps \mathbb Z^d_0: |y|_\infty
    \le 1}\eps^d\1_{0<|y-x_j|_\infty\le 1}\frac{F^\eps_{p'}(y,x_j-y, x_{\{1:m\}\setminus\{j\}};x)}{\sqrt{|y|^2+|x_j-y|^2+|x_{\{1:m\}\setminus\{j\}}|^2}}.
\end{equs}
% In this case, either
  % $\mathcal T^\eps_p= T^{\eps,+}T^{\eps,+}$ or
  % $\mathcal T^\eps_p= T^{\eps,-}T^{\eps,+}$. In the former case, $p=\{1,2,3\}\in\Pi^{(n)}_{2,3}$ and $M=0$. 
  % Using \eqref{eq:F+} and \eqref{eq:casoa=1} for the kernel of
  % $T^{\eps,+}e_k$, we see that
  % \begin{multline}
  %   F^\eps_p(x_{1:3};x)=\1_{x_1+x_2+x_3=x}\iota \frac{w\cdot x}{|x|} \frac{2\iota}3
  %   \sum_{1\le i<j\le 3}\frac{w\cdot(x_i+x_j)}{|x_{1:3}|}\\\times\mathbb J^1_{x_i,x_j}\mathbb J^1_{x_i+x_j,x_{\{1:3\}\setminus\{i,j\}}}\frac1{|x_i+x_j|^2+|x_{\{1:3\}\setminus\{i,j\}}|^2}
  % \end{multline}
  % which is of the form \eqref{eq:formkernelG}, with $\mathcal G[p]$ the set of pairs $g=(i,j),1\le i<j\le 3$,
  % \begin{multline}
  %    P_g(\emptyset;x_{1:3})=-\frac23w\cdot(x_i+x_j),\;  \\Q_g(\emptyset;x_{1:3})=|x_i+x_j|^2+|x_{\{1:3\}\setminus\{i,j\}}|^2\ge0,\\
  % I_g(\emptyset;x_{1:3})=\mathbb J^1_{x_i,x_j}\mathbb J^1_{x_i+x_j,x_{\{1:3\}\setminus\{i,j\}}}.
  % \end{multline}
  % Similarly,  in the case $\mathcal T^\eps_p= T^{\eps,-}T^{\eps,+}$, where $p=\{1,2,1\}\in\Pi^{(n)}_{2,1}$, $M=1$, we leave it to the reader to check that \eqref{eq:formkernelG} hold with $\mathcal G[p]$ containing a single element $g$ and
  % \begin{multline}
  %     P_g(y_1;x_1)=-2w\cdot x_1, \\ Q_g(y_1;x_1)=|y_1|^2+|x_1-y_1|^2\ge0,\; I_g(y_1;x_1)=
  %   \mathbb J^1_{y,x_1-y}.
  % \end{multline}
% Note that, as claimed, $Q_g(y;0)>0$ unless $y=0$.
Again, as $F^\eps_{p'}$ is of the form \eqref{eq:formkernelG} by the induction hypothesis, 
the same holds for $F^\eps_{p}$, where this time we need to set 
$\cG[p]=\{(g',j),g'\in \mathcal G[p'],1\le j\le m\}$ and, for $g=(g',j)\in \cG[p]$,
  \begin{equs}
     P_g(y_{1:M};x_{1:m})&=\frac{4 (m+1)}{(2\pi)^{d/2}} (\fw\cdot x_j)P_{g'},\\
    Q_g(y_{1:M};x_{1:m})&=(|y_M|^2+|x_j-y_M|^2+|x_{1:m\setminus\{j\}}|^2)Q_{g'}\\
    I_g(y_{1:M};x_{1:m})&=\1_{0<|x_j|_\infty\le 1,0<|y_M-x_j|\le 1}I_{g'}\,,\label{e:casoscende}
  \end{equs}
where the omitted variables of the functions $P_{g'}$, $Q_{g'}$ and $I_{g'}$ are 
$y_{1:M-1};y_M,x_j-y_M,x_{1:m\setminus \{j\}}$. 

Now, for both $\sigma=+,-$,~\eqref{e:casosale} and~\eqref{e:casoscende} clearly ensure that 
if $\cG[p']$, $P_{g'}$, $Q_{g'}$ and $I_{g'}$ satisfy conditions $(i)-(iii)$ and $(iii')$ 
so do $\cG[p]$, $P_{g}$, $Q_{g}$ and $I_{g}$. 
As for property $(iv)$, note that by $(iii')$, $Q_g$ can be written as the product of quadratic polynomials 
$q_r$, $1\leq r\leq a-1$, and if $q_r$ depends on at least one of the $y_i,\,i\le M$, then 
necessarily $q_r(y_{1:M};0)\ne0$ for almost all $y_{1:M}$. 
The only potentially problematic case 
arises when one of the $q_r$ is such that $q_r(y_{1:M};x)=|x|^2$ but 
this is ruled out by $(iii')$ for $m=1$. 

At last, to verify~\eqref{eq:casem1}, we note that \eqref{eq:formkernel} and \eqref{eq:formkernelG} 
give 
\begin{equs}\label{eq:plugged}
    \langle e_{-k_1}\mathcal T^\eps_pe_k\rangle=& \iota^a\1_{k_1= k}\frac{(\fw\cdot k)}{|
      k|} \frac{1}{\eps|k|}\\
   &\times \sum_{g\in\mathcal G[p]}\eps^{d M}\sum_{\substack{y_{1:M}\in (\eps \Z^d_0)^M:\\|y_i|\le 1, i\le M}} \frac{P_g(y_{1:M};\eps k)}{Q_g(y_{1:M};\eps k)}I_g(y_{1:M};\eps k)\,.
\end{equs}
But, since $m=1$, the left-most operator in the product $\mathcal T^\eps_p$ is $T^{\eps,-}$, 
so that $\mathcal T^\eps_pe_k=T^{\eps,-}\mathcal T^\eps_{p'}e_k$ and $\mathcal T^\eps_{p'}e_k\in\fock_2$.
Then, from \eqref{e:casoscende} we see that
\begin{equ}
    P_g(y_{1:M};\eps k)= \frac{4}{(2\pi)^{d/2}} (\fw\cdot (\eps k))P_{g'}(y_{1:M-1};y_M,\eps k-y_M)
\end{equ}
which, once plugged into \eqref{eq:plugged}, reduces to \eqref{eq:casem1}, upon defining
\begin{equ}
\tilde P_g(y_{1:M};\eps k)\eqdef \frac{8}{(2\pi)^{d/2}} P_{g'}(y_{1:M-1};y_M,\eps k-y_M)\,.
\end{equ}
\end{proof}

\begin{proof}[of Proposition \ref{p:limite}]

We start with the proof of \eqref{e:cp}. If $a$ is odd or $m>1$,
$\mathcal T^\eps_pe_k$ belongs to a homogeneous Wiener chaos of index
different from $1$ and the scalar product is identically
zero. Therefore, we assume that $p\in\Pi^{(n)}_{2a,1},a>0$, so that $M=a$. In this
case, we want to show that the sum in \eqref{eq:casem1} has a limit as
$\eps\to0$ and that the limit depends only on $p$ (in particular, not
on $k$).  Formally, the $\eps\to0$ limit of  \eqref{eq:casem1} is
\begin{equ}
\iota^a \1_{k_1= k}\frac{(\fw\cdot k)^2}{|k|^2} 2^{ad}\int \mu(\dd y_{1:a}) \sum_{g\in \mathcal G[p]}\frac{P_g(y_{1:a};0)}{Q_g(y_{1:a};0)}I_g(y_{1:a};0)
\end{equ}
where\footnote{note that we added the volume factor $2^{ad}$ 
to pass from the Lebesgue measure to the uniform one. The scope of this change of measure will soon become apparent.} $\mu$ is the uniform probability measure on $[-1,1]^{ad}$ 
and  the integrand is well-defined and finite (except possibly on a zero-measure set) thanks to the  properties of $Q_g$ stated in Lemma \ref{lemma:kernels}.
It remains to prove that the sum indeed converges to the integral, and that the integral is finite. 

At first, we want to write the Riemann-sum in~\eqref{eq:casem1} as an integral with respect to the 
uniform probability measure on $[-1,1]^{ad}$. 
This is quite standard as it suffices to take a piece-wise constant extension of the polynomials 
$P_g$ and $Q_g$, but since the quotient displays singularities, we provide the details below. 

Recall that we assume $1/\eps\in\mathbb N+1/2$. Then, the box
$[-1,1]^d$ contains exactly $(2/\eps)^d$ points of $\eps
\Z^d$ and it can be written
as the union of cubes $C(y^*)$ of side $\eps$, centered at each of the
points $y^*\in \eps\Z^d\cap[-1,1]^d$. Given $y\in [-1,1]^d$, we let
$y^*(y)$ denote the (unique, up to Lebesgue-measure zero sets) $y^*$ such that $y\in
C(y^*)$.  Note that we can rewrite the sum in \eqref{eq:casem1} as
\begin{equ}
  \label{eq:sumint}
2^{ad} \int \mu(\dd y_{1:a})\sum_{g\in \mathcal G[p]}
  \frac{P^{(\eps)}_g(y_{1:a};\eps k)}{Q^{(\eps)}_g(y_{1:a};\eps k)}I^{(\eps)}_g(y_{1:a};\eps k),
\end{equ}
where $\mu$ is the uniform probability measure on 
$[-1,1]^{ad}$, $P^{(\eps)}_g$ is given by
\begin{equ}
P^{(\eps)}_g(y_{1:a};\eps k)\eqdef P_g(y^*(y_1),\dots,y^*(y_a))\1_{y_{i}\not\in[-\eps/2,\eps/2]^d, i\le a}\,,
\end{equ} 
and similarly $Q^{(\eps)}_g,I^{(\eps)}_g$, and the indicator function has been added 
because in \eqref{eq:casem1} the sums are restricted to $y_i\ne0$.
Note that, by construction, the integrand is piecewise constant. 
\medskip

At this point, existence, finiteness of the $\eps\to0$ limit of \eqref{eq:sumint} and its independence
w.r.t. $k$ follow if we can prove uniform integrability of the
integrand.
For this, we will introduce operators $U^{\eps,\pm}$ analogous to the $T^{\eps,\pm}$ but 
with $\cA^{\eps,\delta}$ in~\eqref{e:Adelta} replacing $\gena$, and their composition $\cU^{\eps}_p$ 
which is going to give us the same expression as in~\eqref{eq:sumint} but in 
which the integrand is raised to the power $1+\delta$ for $\delta$ small but positive. 
Upon proving that the latter is bounded, the claimed uniform integrability will follow. 
Let us now provide the details. 

Let $\delta\in(0,(d-2)/2)$ be independent of $\eps$ and define the operators $U^{\eps,\pm}$ in Fock space as
$U^{\eps,\pm} \eqdef (-\gensy)^{-(1+\delta)/2}$ $\cA^{\eps,\delta}_\pm (-\gensy)^{-(1+\delta)/2}$. 
More explicitly, $U^{\eps,+}$ acts on $f\in\fock_n$, $n\ge2$, as
 \begin{equs}\label{eq:U+}
&U^{\eps,+}f(k_{1:n+1})=\eps^{d/2}\left(\frac{4}{(2\pi)^{d/2}(n+1)}\frac1{|\eps k_{1:n+1}|}\right)^{1+\delta}\\
&\times\sum_{1\le i<j\le n+1}\frac{|\fw\cdot(\eps k_i+\eps k_j)|^{1+\delta}\,\indN{k_i,k_j}}{\Big(|\eps(k_i+k_j)|^2+|\eps k_{1:n+1\setminus\{i,j\}}|^2\Big)^{(1+\delta)/2}} f(k_i+k_j,k_{1:n+1\setminus\{i,j\}})
\end{equs}
while $U^{\eps,-}$ as
\begin{equs}
  &U^{\eps,-}f(k_{1:n-1}) =\eps^{d/2}\left(\frac{4 n}{(2\pi)^{d/2}|\eps k_{1:n-1}|}\right)^{1+\delta}\\
  &\times\sum_{j=1}^{n-1}|\fw\cdot (\eps k_j)|^{1+\delta}\sum_{\ell+m=k_j}\indN{\ell,m}\frac{f(\ell,m,
    k_{\{1:n-1\}\setminus\{j\}})}{\Big(|\eps \ell|^2+|\eps m|^2+|\eps k_{\{1:n-1\}\setminus\{j\}}|^2\Big)^{(1+\delta)/2}}\,.
\end{equs}
Also, define the  product of $U^{\eps,\pm}$ operators $\mathcal U^\eps_p$ corresponding to the path $p$, 
in analogy with $\mathcal T^\eps_p$, just with $T^{\eps,\pm }$
replaced by $U^{\eps,\pm}$. 
Then, the same argument as in Lemma~\ref{lemma:kernels} together with the derivation of~\eqref{eq:sumint}, 
shows that 
\begin{equs}\label{e:stanco}
  \langle e_{-k}, \mathcal U^\eps_p e_k\rangle=&2^{da}\frac{|\fw\cdot k|^{2(1+\delta)}}{|k|^{2(1+\delta)}}\\
 &\times \int\mu(\dd y_{1:a}) \sum_{g\in \mathcal G[p]}
  \left|\frac{P^{(\eps)}_g(y_{1:a},\eps k)}{Q^{(\eps)}_g(y_{1:a},\eps k)}\right|^{1+\delta}I^{(\eps)}_g(y_{1:a},\eps k)\,.
\end{equs}
Note that, as anticipated, the integral at the right hand side of~\eqref{e:stanco} is the same 
as that in~\eqref{eq:sumint}, but in which the modulus of the integrand is raised to the power $1+\delta$. 
Now, for any $\delta\in(0,(d-2)/2)$, the quantity at the left hand side is bounded uniformly in $\eps$ 
(but the bound might depend on $m$ and $a$, which is though irrelevant as they are fixed)
since $\mathcal U^\eps_p$ is just a composition of 
$U^{\eps,\pm} = (-\gensy)^{-(1+\delta)/2}$ $\cA^{\eps,\delta}_\pm (-\gensy)^{-(1+\delta)/2}$.  
In turn, these operators are uniformly bounded on $\fock$ by~\eqref{e:Ubound}, which 
 means that so is the integral at the right hand side of~\eqref{e:stanco}.  
Hence, the uniform integrability of the integrand in~\eqref{eq:sumint} is proven and~\eqref{e:cp} follows at once. 
\medskip

Next, we move to the proof of \eqref{e:nullo}.  First of all, the
scalar product is zero unless $q_1+\dots+q_m=k$ (by translation
invariance) and $p\in \Pi^{(n)}_{a,m'}$ with $m'=m$ (by
orthogonality of Wiener chaoses). We assume therefore that these
conditions are satisfied.  If $p\in\Pi^{(n)}_{a,m}$, the kernel of
$\mathcal T^\eps_p e_k$ is of the form \eqref{eq:formkernel}, with
$F_p^\eps$ as in Lemma \ref{lemma:kernels}. The scalar product in
\eqref{e:nullo} is then
\[
\langle e_{-q_1}\otimes\dots\otimes e_{-q_m},\mathcal T^\eps_p e_k\rangle=\eps^{d/2(m-1)}F_p^\eps(\eps q_{1:m};\eps k).
  \]
with $q_1+\dots+q_m= k$. 
By Lemma \ref{lemma:kernels} and arguing as in the derivation of~\eqref{eq:sumint}, we get
\begin{equ}\label{e:vaillant}
  \langle e_{-q_1}\otimes\dots\otimes e_{-q_m},\mathcal T^\eps_p e_k\rangle=2^{dM} \frac{(\fw\cdot k)}{|k|}\int \mu(\dd y_{1:M}) W^\eps_p(y_{1:M},\eps q_{1:m})\,,
  \end{equ}
where
\begin{equ}
W^\eps_p(\ell_{1:M},\eps q_{1:m})\eqdef \eps^{d/2(m-1)}  \sum_{g\in \mathcal G[p]}\frac{\iota^a}{|q_{1:m}|}
  \frac{P^{(\eps)}_g(y_{1:M},\eps q_{1:m})}{Q^{(\eps)}_g(y_{1:M},\eps  q_{1:m})}I^{(\eps)}_g(y_{1:M},\eps  q_{1:m}).
\end{equ}
Let us stress that the main difference w.r.t. \eqref{eq:sumint} is the presence of the
factor $ \eps^{d/2(m-1)}$, which  tends to zero  as
$\eps\to0$, because $m\ge2$, and this is the fact we want to leverage. 
For this, we will consider once again the operators $U^{\eps,\pm}$ and their composition $\cU^\eps$. 
Let $\delta\in(0,(d-2)/2)$. 
Proceeding as in~\eqref{e:stanco}, we see that 
\begin{equs}\label{e:stanco2}
&\langle e_{-q_1}\otimes\dots\otimes e_{-q_m},\mathcal U^\eps_p e_k\rangle=   2^{dM} \frac{|\fw\cdot k|^{1+\delta}}{|k|^{1+\delta}}\times\\
&\times\int\mu(\dd y_{1:M})\bigg[ \eps^{d/2(m-1)} \sum_{g\in \mathcal
      G[p]}\left|\frac{1}{|q_{1:m}|}
      \frac{P^{(\eps)}_g(y_{1:M},\eps
        q_{1:m})}{Q^{(\eps)}_g(y_{1:M},\eps
          q_{1:m})}\right|^{1+\delta}I^{(\eps)}_g(y_{1:M},\eps
    q_{1:m})\bigg].
%   \eps^{d/2(m-1)}|B(0,1)|^M \\\times
% \mathbf E \sum_{g\in \mathcal G[p]}
%   \left|\frac{P_g(\ell_{1:M},\eps q_1,\dots,\eps q_m)}{\sqrt{Q_g(\ell_{1:M},\eps  q_1,\dots,\eps q_m)}}\right|^{1+\delta}I_g(\ell_{1:M},\eps  q_1,\dots,\eps q_m)
\end{equs}
Now, notice that, Jensen's 
inequality (which holds as $\mu$ is a probability measure) applied at the right hand side of~\eqref{e:vaillant} 
gives
\begin{equs}
|&\langle e_{-q_1}\otimes\dots\otimes e_{-q_m},\mathcal T^\eps_p e_k\rangle|^{1+\delta}\\
&\lesssim \frac{|\fw\cdot k|^{1+\delta}}{|k|^{1+\delta}} \int \mu(\dd y_{1:M}) |W^\eps_p(y_{1:M},\eps q_{1:m})|^{1+\delta}\\
&\lesssim \eps^{d/2(m-1)\delta} \frac{|\fw\cdot k|^{1+\delta}}{|k|^{1+\delta}}\\
&\times\int\mu(\dd y_{1:M})\bigg[ \eps^{d/2(m-1)} \sum_{g\in \mathcal
      G[p]}\left|\frac{1}{|q_{1:m}|}
      \frac{P^{(\eps)}_g(y_{1:M},\eps
        q_{1:m})}{Q^{(\eps)}_g(y_{1:M},\eps
          q_{1:m})}\right|^{1+\delta}I^{(\eps)}_g(y_{1:M},\eps
    q_{1:m})\bigg]\\
    &\lesssim \eps^{d/2(m-1)\delta}\langle e_{-q_1}\otimes\dots\otimes e_{-q_m},\mathcal U^\eps_p e_k\rangle\lesssim \eps^{d/2(m-1)\delta}
\end{equs}
where we used~\eqref{e:stanco2} and, in the last step, that, 
since by Lemma~\ref{l:GenaBound}, the operators $U^{\eps,\pm}$ are uniformly bounded (in $\eps$) in $\fock$, 
so is the scalar product $\langle e_{-q_1}\otimes\dots\otimes e_{-q_m},\mathcal U^\eps_p e_k\rangle$. 
Hence, \eqref{e:nullo} follows at once and the proof of Proposition~\ref{p:limite} is completed. 
\end{proof}



\section{Convergence of Burgers to renormalised SHE}\label{sec:ProofofMain}

Thanks to the results in the previous section, we have all the ingredients we need in order to 
prove Theorem~\ref{thm:main}. 
To do so, we will first show that the sequence $\eta^\eps$ is tight and then verify that every limit 
point solves the stationary stochastic heat equation (SHE) 
\begin{equ}[e:SHE2]
\partial_t \eta =\tfrac12(\Delta + D_\SHE(\fw\cdot\nabla)^2)\eta + (-\Delta - D_\SHE(\fw\cdot\nabla)^2)^{1/2}\xi\,,\qquad \eta(0,\cdot)=\mu
\end{equ}
%\fabio{ check se pure il D storto si becca l'un mezzo. Per ora ho assunto di si e quindi ho modificato delle formule sotto. }
where $\mu$ is a  zero-average space white noise and 
$D_\SHE>0$ is the constant identified in Theorem~\ref{thm:MainConv2} for $d=2$ and 
Theorem~\ref{thm:MainConv} for $d\geq 3$.   


\subsection{Tightness}

The proof of tightness follows a by now standard route (see
e.g.~\cite{GubinelliJara2012, GT, CK, CES}) which exploits Mitoma's
criterion~\cite{Mitoma}, Kolmogorov's continuity theorem and, more
importantly, the so-called It\^o trick, introduced
in~\cite{GubinelliJara2012} and since then exploited in a variety of
contexts, see~\cite{Gubinelli2016} for a pedagogical introduction.
For the reader's convenience, we recall below the statement of the
latter and refer the reader to~\cite[Lemma 4.1]{CES}  for a proof.

\begin{lemma}[It\^o trick]\label{l:Ito}
Let $d\geq 2$, $\eta^\eps$ be the stationary solution to \eqref{e:BurgersScaled} with $\lambda_\eps$ 
given as in~\eqref{eq:lambdaeps}. For any
$p \ge 2$, $T > 0$ and $F\in L^2(\Omega)$ with finite chaos expansion, i.e. $F\in\oplus_{j=1}^n\cH_j$ for some $n\in\N$, 
there exists a constant $C=C(p,n)>0$ such that
  \begin{equ}[e:Itotrick]
    \mathbf E\left[\sup_{t\in[0,T]}\left|\int_0^t F(\eta^\eps_s)
      \dd s
      \right|^p
      \right]^{1/p}\leq C T^{1/2}\|(-\gensy)^{-1/2}F\|\,.
  \end{equ}
\end{lemma}

We now move to the proof of tightness. 

\begin{proposition}\label{p:Tight}
The sequence $\{\eta^\eps\}_\eps$ is tight in $C([0,T],\cS'(\T^d))$. 
\end{proposition}
\begin{proof}
By Mitoma's criterion~\cite{Mitoma}, the sequence $\{\eta^\eps\}_\eps$ is tight in $C([0,T],\cS'(\T^d))$ if and only if 
$\{\eta^\eps(\phi)\}_\eps$ is tight in $C([0,T],\R)$ for all $\phi\in\cS(\T^d)$. Therefore, we fix $\phi\in\cS(\T^d)$ and 
consider the process $\{\eta^\eps_t(\phi)\}_{t\in[0,T]}$. In view of~\eqref{e:SPDEweak}, it suffices to show that 
each of the terms at the right hand side is tight. This is clear for the last, as it is independent of $\eps$, 
while for the first and second, Lemma~\ref{l:Ito} implies that, for every $0\leq t\leq T$ and $p>2$, we have 
\begin{equ}[e:Lapla]
\Exp\left[\left(\int_0^t\eta^\eps_s(\Delta\phi)\dd s\right)^p\right]^{\frac1{p}}\lesssim t^{1/2}\|(-\gensy)^{-1/2}(\Delta\phi)\|=t^{1/2}\|\phi\|_{H^1(\T^d)}
\end{equ} 
and 
\begin{equs}[e:NonlinT]
\Exp\left[\left(\int_0^t\cN^\eps_\phi(\eta^\eps)\dd s\right)^p\right]^{\frac1{p}}\lesssim t^{1/2}\|(-\gensy)^{-1/2}\nf_\phi\|\lesssim t^{1/2}\|\phi\|_{H^1(\T^d)}
\end{equs}
where in the last step we used~\eqref{e:nonlinUnifBound}. 
Since the process $\{\eta^\eps_t(\phi)\}_{t\in[0,T]}$ is Markov,~\eqref{e:Lapla} and~\eqref{e:NonlinT} together with 
Kolmogorov's criterion ensure that the sequence $\{\eta^\eps(\phi)\}_\eps$ is tight in $C([0,T],\R)$, 
so that the proof of the statement is concluded. 
\end{proof}



\subsection{Convergence}

In order to identify the limit equation, we will prove that any limit point satisfies 
the martingale problem associated to~\eqref{e:SHE2} which we now recall. 

\begin{definition}\label{def:MPSHE}
Let $T>0$, $\Omega=C([0,T], \cS'(\T^d))$ and $\CB$ the canonical Borel $\sigma$-algebra on $C([0,T], \cS'(\T^d))$. 
Let $\mu$ be a zero-average space white noise on $\T^d$.
We say that a probability measure $\BP$ on $(\Omega,\CG)$ {\it solves the martingale problem for 
$\geneff\eqdef  \gensy+\cD= \frac12(\Delta+D_\SHE (\fw\cdot \nabla)^2)$
with initial distribution $\mu$}, if for all $\phi\in\cS(\T^d)$, the canonical process $\eta$ under $\BP$ is such that 
\begin{equs}[e:Mart]
\BM_t(\phi)&\eqdef \eta_t(\phi) - \mu(\phi) - \int_0^t \geneff\eta_s(\phi)\dd s\\
\Gamma_t(\phi)&\eqdef(\BM_t(\phi))^2-t\, \|(-\Delta-D_\SHE (\fw\cdot \nabla)^2)^{1/2}\phi\|^2_{L^2(\T^2)}
\end{equs}
are local martingales. 
\end{definition}

In the next theorem, we state well-posedness for the previous martingale problem. 
The proof is classical (see for example~\cite[Theorem D.1]{MW}) therefore we omit it. 

\begin{theorem}\label{thm:WellPosedMP}
The martingale problem for $\geneff$ with initial distribution $\mu$ in Definition~\ref{def:MPSHE} has a unique solution 
and uniquely characterises the law of the solution to~\eqref{e:SHE2} on $C([0,T], \cS'(\T^d))$.
%, i.e. there exists a unique probability 
%measure $\P$ on $C([0,T], \cD'(\T^2))$ such that the processes in~\eqref{e:Mart1} and~\eqref{e:Mart2} are 
%local martingales. 
\end{theorem}

We are now ready to prove Theorem~\ref{thm:main}. 

\begin{proof}[of Theorem~\ref{thm:main}]
By Proposition~\ref{p:Tight}, we know that the sequence $\{\eta^\eps\}_\eps$ is tight in $C([0,T],\cS'(\T^d))$, 
hence it converges along subsequences. If we prove that any limit point is a solution of the martingale 
problem in Definition~\ref{def:MPSHE} then the statement follows by Theorem~\ref{thm:WellPosedMP}. 

Let $\eta\in C([0,T],\cS'(\T^d))$ be a limit point. To verify that $\eta$ satisfies the martingale problem, 
we need to show that for every given $\phi\in\cS(\T^d)$ the processes
$\BM(\phi)$ and $\Gamma(\phi)$, defined according to~\eqref{e:Mart}, are local martingales. 
Notice that, by definition of quadratic variation, $\Gamma(\phi)$ is a local martingale 
if and only if $\BM(\phi)$ has quadratic variation given by $t\mapsto t\|(-\Delta-D_\SHE(\fw\cdot\nabla)^2)^{1/2}\phi\|^2_{L^2(\T^2)}$, 
and we will prove this latter statement. 
Before delving into the details, we proceed as in Section~\ref{sec:ideas} and 
write the equation for $\eta^\eps$ in such a way that 
we can identify the elements which are relevant to the limit. 

Let us fix a smooth test function $\phi\in\cS(\T^d)$. Recall the weak formulation of~\eqref{e:BurgersScaled} 
given in~\eqref{e:SPDEweak}, i.e. 
\begin{equ}[e:Weak]
\eta^\eps_t(\phi)-\mu(\phi)-\tfrac12\int_0^t \eta^\eps_s(\Delta\phi)\dd s-\int_0^t\cN^\eps_\phi(\eta^\eps_s)\dd s=M_t(\phi)
\end{equ}
where $M_t(\phi)$ is the martingale 
\begin{equ}[e:M1]
M_t(\phi)=\int_0^t \xi(\dd s,(-\Delta)^{\frac12}\phi)=\sum_k |k| \hat\phi(k) B_t(k)=\int_0^t\sum_k |k| D_k \eta^\eps_s(\phi) \dd B_s(k)
\end{equ} 
and in the last step, we used the notion of Malliavin derivative in~\eqref{e:Malliavin}. 

In order to describe the term containing the nonlinearity, we consider the random variable $\WN\in L^2(\P)$, 
for $n\in\N$, which has kernel $\wN$ given by the solution $\vN$ of the truncated generator equation~\eqref{e:nonlinGenEq} 
if $d\geq 3$ and by $\tvN$ as in~\eqref{e:nonlinAnsatz} if $d=2$. 
By It\^o's formula applied to $\WN$, we have 
\begin{equ}[e:DynkinCiao]
\WN(\eta^\eps_t)-\WN(\mu)-\int_0^t \gen\WN(\eta^\eps_s)\dd s=\cM_t^\eps(\WN)
\end{equ}
where $\cM^\eps(\WN)$ is the martingale given by 
\begin{equ}[eq:M2]
\cM_t^\eps(\WN)=\int_0^t\sum_k|k|[D_k\WN](\eta^\eps_s)dB_s(k)\,.
\end{equ}
We now go back to~\eqref{e:Weak}. In view of Theorem~\ref{thm:GenEqu}, we rewrite it as 
\begin{equ}[e:etaWeak]
\eta^\eps_t(\phi)-\mu(\phi)-\tfrac12\int_0^t \eta^\eps_s\Big((\Delta+\cD)\phi\Big)\dd s +R^{\eps,n}(t)=M_t(\phi) +\cM_t^\eps(\WN)
\end{equ}
where $R^{\eps,n}(t)\eqdef\one+\two+\three$ and 
\begin{equs}
\one&\eqdef \WN(\eta^\eps_t)-\WN(\mu)\\
\two&\eqdef -\int_0^t \Big(\gen \WN+\cN^\eps_\phi-\genam \WN_2\Big)(\eta^\eps_s)\dd s\\
\three&\eqdef -\int_0^t \Big(\genam \WN_2(\eta^\eps_s) -\eta^\eps_s(\cD\phi)\Big)\dd s\,.
\end{equs}
Notice that
\begin{equs}\label{e:rem1}
\Exp[\one^2]&\lesssim \Exp[\WN(\eta_t)^2]+\Exp[\WN(\mu)^2]=2\E[\WN(\eta)^2]=2\|\wN\|^2\,,\\
\Exp[\three^2]&\lesssim t^2\Exp\Big[\frac1{t}\int_0^t \Big(\genam \WN_2(\eta^\eps_s) -\eta^\eps_s(\cD\phi)\Big)^2\dd s\Big]\\
&= t^2 \|\genam \wN_2 -\cD\phi\|^2\label{e:rem3}
\end{equs}
where the second follows by Jensen's inequality,
and, by the It\^o trick, Lemma~\ref{l:Ito}, we have 
\begin{equs}
\Exp[\two^2]&\lesssim t \|(-\gensy)^{-\frac12} (\gen \wN+\nf_\phi-\genam \wN_2)\|^2\,.\label{e:rem2}
%\Exp[\three^2]&\lesssim t \|(-\gensy)^{-\frac12} (\genam \wN_2 -\cD\phi)\|^2\,.\label{e:rem3}
\end{equs}
The right hand sides of~\eqref{e:rem1},~\eqref{e:rem3} and~\eqref{e:rem2} all converge to $0$ in the double 
limit as $\eps\to 0$ first and $n\to\infty$ then,~\eqref{e:rem2} in view of Theorem~\ref{thm:GenEqu}, 
while~\eqref{e:rem1} and~\eqref{e:rem3} by Proposition~\ref{p:epsLim} and~\eqref{e:Diffusivity3} in Theorem~\ref{thm:MainConv} for $d\geq 3$, and by Theorem~\ref{thm:MainConv2} for $d=2$. 
Hence, $R^{\eps,n}$ converges to $0$ in $L^2(\Omega)$ and consequently, almost surely (at least along subsequences). 

It follows that the left hand side of~\eqref{e:etaWeak} converges to 
\begin{equ}
\eta_t(\phi)-\mu(\phi)-\tfrac12\int_0^t \eta_s\Big((\Delta+D_\SHE(\fw\cdot\nabla)^2)\phi\Big)\dd s
\end{equ}
which by definition, equals $\BM_t(\phi)$. Hence, we can determine a common subsequence 
such that the convergence holds for any $t$ in a countable dense set of $[0,T]$.  
Since tightness holds in $C([0,T],\R)$, $\BM(\phi)$ is the almost sure 
pointwise (in time) limit of the sequence of martingales $\{M_t(\phi) +\cM^\eps_t(\WN)\}_\eps$. 
Hence, it suffices to show that the sequence $\{M_t(\phi) +\cM^\eps_t(\WN)\colon t\in[0,T]\}_\eps$ converges to a 
martingale whose quadratic variation is $t\mapsto t\|(-\Delta-D_\SHE(\fw\cdot\nabla)^2)^{1/2}\phi\|^2_{L^2(\T^2)}$. 
This in turn is the content of Theorem~\ref{thm:MartConv} in the next section, 
hence the proof of Theorem~\ref{thm:main} is concluded. 
%
%
%
%
%
%
%
%
%
%
%
%\newpage$\lim_{n\to\infty}\lim_{\eps\to0}$, $\one$ in $L^2(\P)$ by Proposition~\ref{p:epsLim}, $\two$ and $\three$ 
%by the It\^o's trick Lemma~\ref{l:Ito} 
%and Theorem~\ref{thm:GenEqu} and Corollary~\ref{c:effective}, respectively. 
%Therefore, the only quantity we need to control
%is the martingales $\{M_t(k) +\cM_t^\eps(\VN)\}_{k\in\Z^d_0}$, to which the next section is devoted. Namely, we need to verify that {\color{blue} these martingales
%  converge to  space-time white noises
%with the correct covariance, see Theorem \ref{thm:MartConv} below. Once this is proven, the proof of Theorem ... will follow immediately from Theorem \ref{thm:WellPosedMP}, see ... WHERE WE WILL SHORTLY CONCLUDE THE PROOF OF THE MAIN THEOREM}. \giuseppe{This section can be improved, in particular the complete proof could be given, with a statement. FT: A me pare ok gia' come e' scritto ora. piuttosto, non e' molto visibile dove si conclude la prova del teorema principale del lavoro.}
%
%
%
%
%
%
\end{proof}

























%
%
%\newpage
%\fabioText{Give here the usual martingale characterisation of the limit equation. Hopefully we can take $\phi$ as $C^\infty$.}
%
%
%To identify the limit equation we study for smooth $\varphi$:
%\begin{equation}
%  \label{eq:integrata}
%\eta^\eps_t(\varphi)-\eta^\eps_0(\varphi)-\int_0^t\eta^\eps_s(\Delta\varphi)ds-\int_0^t \nf_\phi(\eta^\eps_s)ds=M_t(\varphi),
%\end{equation}
%where $\nf_\phi$ is the non-linearity, i.e.
%\begin{align}
%\nf_\phi(\eta)&=\label{eq:non-linearity}
%\sum_{\ell,m}\lambda_\eps  \indN{\ell,m}w\cdot(\ell+m)\hat \varphi(\ell+m)\eta(\ell)\eta(m)\\\nonumber
%&=\sum_{\ell,m}\widehat{\mathcal A_+\varphi}(\ell,m):\eta(\ell)\eta(m):
%=\mathcal A_+^\eps\eta(\varphi),
%\end{align}
%\fabioText{C'e' l'eterna confusione tra A che agisce in Fock e A che agisce in L2.}
%
%and
%\begin{equation}\label{eq:M1}
%  M_t(\varphi)=\int_0^t \xi_s((-\Delta)^{\frac12}\phi)ds= 
%    \int_0^t\sum_k|k|\hat \varphi(k)dB_s^k.
%\end{equation}
%
%\subsection{Replacing the non-linearity}
%For $g\in\Gamma L^2_1\simeq \mathcal H_1$, let $v^{\eps,n}$ be the solution of
%\begin{equation}%\label{e:nonlinGenEq}
%-\gen_nv^{\eps,n}=\nf_\phi+g
%\end{equation}
%Then, by Dynkin's formula
%\begin{equation}
%v^{\eps,n}(\eta^\eps_t)-\vN(\eta^\eps_0)-\int_0^t\gen \vN(\eta^\eps_s)ds=\mathcal M_t^\eps(\vN),
%\end{equation}
%where
%\begin{equation}
%\mathcal M^\eps_t(v)=\int_0^t\sum_k|k|D_kv(\eta^\eps_s)dB_s^k.
%\end{equation}
%
%Then by (\ref{e:nonlinGenEq}) we get
%\begin{equation}
%-\int_0^t\gen \vN(\eta^\eps_s)ds=\int_0^t\nf_\phi(\eta^\eps_s)ds-\int_0^tg(\eta^\eps_s)ds-\int_0^t\genap\vN_n(\eta^\eps_s)ds
%\end{equation}
%which in particular implies
%\begin{align}\label{eq:nonlinearityIntegral}
%\int_0^t\nf_\phi(\eta^\eps_s)ds=
%&\int_0^tg(\eta^\eps_s)ds+\mathcal M^\eps_t(\vN)\\\nonumber
%+&\vN(\eta^\eps_0)-\vN(\eta^\eps_t)
%+\int_0^t\genap\vN_n(\eta^\eps_s)ds.
%\end{align}
%Define the two error terms
%\begin{equation}
%  R^{\eps,n}_1(t)\eqdef \vN(\eta^\eps_0)-\vN(\eta^\eps_t),\qquad R^{\eps,n}_2(t)\eqdef
% \int_0^t\genap\vN_n(\eta^\eps_s)ds \,.
%\end{equation}
%
%By  \eqref{eq:integrata} and \eqref{eq:nonlinearityIntegral} we obtain:
%\begin{equation}\label{e:etaWeak}
%  \eta^\eps_t(\varphi)-\eta^\eps_0(\varphi)-\int_0^t\eta^\eps_s(\Delta\varphi)ds+\int_0^tg(\eta^\eps_s)ds
%= M_t(\varphi)+\mathcal M^\eps_t(\vN)+  R^\eps_1(t)+  R^\eps_2(t).
%\end{equation}
%The aim is to show that if one chooses 
%$g=g_\phi^n\in\mathcal H_1$ (independent of $\eps$ but depending on $n$) in \eqref{e:nonlinGenEq} as follows:
%\begin{eqnarray}
%  \label{eq:g}
%  \widehat{g^n_\phi}(k):= \mathcal D_n(k)\hat \phi(k)
%\end{eqnarray}
%\fabioText{presumably one of the k will be -k}
%with $\mathcal D_n(k)$ as in Theorem  \ref{th:H}, 
%then the following three facts hold:
%
%\begin{enumerate}
%\item The error terms $ R^{\eps,n}_1(t)$ and $ R^{\eps,n}_2(t)$ converge to $0$ \fabioText{in L2?}, as $\eps\rightarrow0$ first and $n\rightarrow\infty$ then.
%  
%  
%\item The average of the quadratic variation of the martingales satisfies
%  \begin{eqnarray}
%    \label{eq:avM}
%    \lim_{n\to\infty}\lim_{\eps\to0}  \mathbb E \langle M_\cdot(\phi) + \mathcal M_\cdot^\eps(\vN)\rangle_t= t
%    \langle \phi,(-\mathcal L_0+\mathcal D)\phi\rangle.
%  \end{eqnarray}
%
%\item The quadratic variation of the martingales concentrates:
%\begin{equation}
%\lim_{n\to\infty}\lim_{\eps\to0}{\rm Var}\left(\langle M_\cdot(\phi) + \mathcal M_\cdot^\eps(\vN)\rangle_t\right)=0.
%\end{equation}
%\end{enumerate}
%
%\fabioText{Conclude using the martingale characterization}
%
%\subsection{The error terms}
%To show that $R^{\eps,n}_1(t)$ goes to $0$ it suffices to show that, after the choice \eqref{eq:g} for $g$ is made, 
%\begin{equation}\label{e:L2}
%\|\vN\|\rightarrow 0\text{ as }\eps\to 0, n\rightarrow\infty.
%\end{equation}
%By Lemma \ref{L:higherChaos}, since the r.h.s. of \eqref{e:nonlinGenEq} belongs to the second
%chaos,
%we have
%\begin{eqnarray}
%\|P_{>3}\vN\|\lesssim \sqrt{f(\eps)}\Big(1+\|(-\cL_0)^{-\frac12}(\nf_\phi+g^n_\phi)\|\Big).
%\end{eqnarray}
%Because of the upper bound in \eqref{eq:boundsDn}, we have $\|(-\cL_0)^{-\frac12}g^n_\phi\|\lesssim \|\phi\|$ and
%since $\nf_\phi=\cA^\eps_+\phi$ we have by Lemma \ref{l:GenaBound}
%$\|(-\cL_0)^{-\frac12}\nf_\phi\|\lesssim \|(-\cL_0)^{\frac12}\varphi\|$. Altogether,
%\begin{eqnarray}
%  \|P_{>3}\vN\|\lesssim_\phi \sqrt{f(\eps)}
%\end{eqnarray}
%which tends to zero (uniformly in $n$) as $\eps\to0$.
%As for the components of $\vN$ in the first and  second chaos, note that by solving
% the system of equation \eqref{eq:truncgensys} (with $g$ replaced by $\nf_\phi+g^n_\phi=\cA^\eps_+\phi+g^n_\phi$) one finds 
%\begin{equation}
%\vN_2=(-\mathcal L_0+\Op_{n})^{-1}\genap(\varphi+\vN_1)
%\end{equation}
%and
%\begin{equation}
%  \label{eq:eqv1}
%-\mathcal L_0\vN_1-\mathcal A_-^\eps \vN_2 =g^n_\phi,
%\end{equation}
%and the latter gives
%\begin{equation}
%(-\mathcal L_0+\Op_{n})\vN_1=-\Op_{n}\varphi+g^n_\phi.
%\end{equation}
%Thanks to the choice \eqref{eq:g} for $g=g^n_\phi$ we see that $\|\vN_1\|\to0$.
%\fabioText{non so se bisogna giustificare meglio}
%Finally, \begin{equation}
%\|\vN_2\|\leq\|(-\mathcal L_0+\Op_{n})^{-1}\genap\varphi\|+\|(-\mathcal L_0+\Op_{n})^{-1}\genap\vN_1\|.
%\end{equation}
%In order to estimate the r.h.s., we proceed as in  the proof of Lemma
%\ref{L:higherChaos} and we  drop the positive operator $\Op_{n}$.
%Then, the first norm in the r.h.s. is bounded by
%\[
%\sqrt{f(\eps)}\|(-\cL_0)^{\frac12}\phi\|,
%\]
%and the second one by
%\[
%\sqrt{f(\eps)}\|(-\cL_0)^{\frac12}\vN_1\|\lesssim \sqrt{f(\eps)}\|(-\cL_0)^{-\frac12}(g_\phi-\Op_{n})\phi\|.
%\]
%Both quantities vanish as $\eps\to0$.
%
%As for $R_2^{\eps,n}(t)$, by Itô's trick we have
%\begin{align}
%  \label{eq:stimaR2}
%\mathbb E\left(\sup_{t\leq T}|R^{\eps,n}_2(t)|^2\right)^{\frac{1}{2}}
%&=\mathbb{E}\left(\sup_{t\leq T}\left|\int_0^t\genap\vN_n(\eta^\eps_s)ds\right|^2\right)^{\frac{1}{2}}\\
%&\lesssim t^{\tfrac12}\|(-\mathcal L_0)^{-\tfrac12}\vN_n\|\\
%&\lesssim t^{\tfrac12}n^{-k}\|(-\mathcal L_0)^{-\tfrac12}(\nf_\phi+g^n_\phi)\|= t^{\tfrac12}n^{-k}\|(-\mathcal L_0)^{-\tfrac12}(\cA_+^\eps\phi+g^n_\phi)\|.
%\end{align}
%where the last step is due to Proposition \ref{P:AprioriGeneratorEq}.
%
%Choosing $k>1$ and using Lemma \ref{l:GenaBound}, we see that the l.h.s. of \eqref{eq:stimaR2} tends to zero as $n\to\infty$, uniformly w.r.t. $\eps$.
%%\levi{Expand Itô trick step?}

\subsection{The Martingales}\label{sec:martingales}

In this section, we want to prove that for any $\phi\in\cS(\T^d)$, 
the martingale part at the right hand side of~\eqref{e:etaWeak} converges to a Brownian motion 
with a deterministic quadratic variation which can be suitably expressed in terms of the operators $\gensy$ and $\cD$. 
To do so, recall first that by the definition of $M_t(\phi)$ and $\mathcal M_t^\eps(\WN)$ 
in~\eqref{e:M1} and \eqref{eq:M2} respectively, we have 
\begin{equ}[e:Martingales]
\Mn_t(\phi)\eqdef  M_t(\phi)+\mathcal M_t^\eps(\wN)= \int_0^t \sum_k |k|D_k\Big(\eta_s^\eps(\phi)+\WN(\eta^\eps_s)\Big)\dd B_s(k)\,,
\end{equ}
where $\WN\in L^2(\P)$ is the random variable whose kernel $\wN\in\fock$ is either $\vN$, the solution 
of the truncated generator equation in~\eqref{e:nonlinGenEq}, for $d\geq 3$ 
or $\tvN$ given by~\eqref{e:nonlinAnsatz} if $d=2$. 
The quadratic variation of $\Mn_\cdot(\phi)$ is 
\begin{equ}[e:QV]
  \langle \Mn_t(\phi)\rangle_t=\int_0^t F^{\eps,n}(\eta_s)\dd s\,.
\end{equ}
for $F^{\eps,n}(\eta)$ the random variable defined as 
\begin{equ}[e:IntQV]
F^{\eps,n}(\eta)=\sum_k|k|^2|D_k(\eta(\phi)+\WN(\eta))|^2\,.
\end{equ}
Our goal is to prove the following theorem. 

\begin{theorem}\label{thm:MartConv}
For $\phi\in\cS(\T^d)$, $\eps>0$ and $n\in\N$, let $\WN$ be the random variable whose kernel 
$\wN\in\fock$ is either $\vN$, the solution 
of the truncated generator equation in~\eqref{e:nonlinGenEq}, for $d\geq 3$, 
or $\tvN$ given by~\eqref{e:nonlinAnsatz} if $d=2$. 
Then, in the limit as $\eps\to$ first and $n\to\infty$ after, 
the sequence of martingales $\{\Mn(\phi)\}_{n,\eps}$ defined in~\eqref{e:Martingales} 
converges in law to a Brownian motion whose quadratic variation is given by 
\begin{equ}[e:QVlimit]
t\quad\mapsto\quad 2t\langle\phi,(-\gensy-\cD)\phi\rangle=t\|(-\Delta-D_\SHE(\fw\cdot\nabla)^2)^{1/2}\phi\|^2
\end{equ}
where the operator $\cD$ was defined in Theorem~\ref{thm:MainConv} for $d\geq 3$ 
and in Theorem~\ref{thm:MainConv2} if $d=2$. 
\end{theorem}

The proof of the previous statement is a direct consequence of the following propositions. In the first, 
we show that the expectation of~\eqref{e:QV} converges to~\eqref{e:QVlimit}. 

\begin{proposition}\label{p:MeanQV}
In the same setting as in Theorem~\ref{thm:MartConv}, let $\Mn_t(\phi)$ be 
the martingale in~\eqref{e:Martingales}. Then, for any $t>0$
\begin{equ}[e:limQV]
\lim_{n\to\infty}\lim_{\eps\to0}\Exp\langle \Mn_{\cdot}(\phi)\rangle_t=2t\langle\phi,(-\gensy-\cD)\phi\rangle\,.
\end{equ}
\end{proposition}
\begin{proof}
By~\eqref{e:QV} and stationarity, we have 
\begin{equ}
\Exp  \langle \Mn_{\cdot}(\phi)\rangle_t=\int_0^t \Exp[F^{\eps,n}(\eta_s)]\dd t =t \E[F^{\eps,n}(\eta)]\,.
\end{equ}
so that we only need to focus on the latter expectation. To compute it, we notice first that for any random variable 
$G\in L^2(\P)$, Gaussian integration by parts~\eqref{eq:intbyparts} gives
\begin{equ}
\E[G(\eta) D_{-k} D_{k} G(\eta) ]=\E[ G(\eta) \hat\eta(-k)D_{k} G(\eta)]-\E\left[|D_{k} G(\eta)|^2 \right]
\end{equ}
so that, by~\eqref{e:gens}, we get 
\begin{equs}
\E\left[G(\eta) (-\gensy)G(\eta)\right]&=\frac12\sum_k |k|^2 \E\left[G(\eta) \left(\eta_{-k}D_{k} G(\eta)-D_{-k} D_{k} G(\eta)\right)\right]\\
&=\frac12\sum_k |k|^2\E\left[|D_{k} G(\eta)|^2 \right]\,.
\end{equs}
As a consequence, we see that
\begin{equ}[eq:123]
\tfrac12\E[F^{\eps,n}(\eta)]=
\|(-\cL_0)^{\frac12}(\phi+\wN)\|^2= \|(-\cL_0)^{\frac12}\phi\|^2+\|(-\cL_0)^{\frac12}\wN\|^2\,,
\end{equ}
where the last step holds as $\wN\in\oplus_{j=2}^n\fock_j$, while $\phi\in\fock_1$, 
so that in particular the two are orthogonal. 
For $d\geq3$, the result then follows immediately by Theorem~\ref{thm:MainConv}. 

For $d=2$ (recall that $\wN=\tvN$ in~\eqref{e:nonlinAnsatz}), by~\eqref{e:genAnsatz}, we can write
\begin{equs}
\|(-\cL_0)^{\frac12}\tvN\|^2=&\langle\tvN,(-\gen)\tvN\rangle\\
=&\langle \tvN,\nf_\phi\rangle +\langle \tvN_n, \gensy^\fw\cGN \tvN_n\rangle\\
&+\langle \tvN, \big[-\genam(-\gensy-\gensy^\fw\cGN)^{-1}\genap+\gensy^\fw\cGN\big]P_2^{n-1}\tvN\rangle
\end{equs}
where we removed $-\genam (-\gensy-\gensy^\fw\cGN)^{-1}\nf_\phi$ and $\genap\tvN_n$ as they are 
orthogonal to $\tvN$. Now, the second term at the right hand side only involves the 
$n$-th coordinate of $\tvN$ so that it goes to $0$ as $n\to\infty$ by the weighted a priori 
estimates in Lemma~\ref{l:AnsatzvsGenEq}. 
The last term can be bounded via the replacement Lemma~\ref{l:Replacement}, 
and after controlling $\|(-\gensy)^{1/2}\tvN\|$ with Lemma~\ref{l:aprioriAnsatz} first and~\eqref{e:nonlinUnifBound} then, 
we can conclude that, for any $n$ fixed, it converges to $0$ as $\eps\to 0$. 
Hence, it remains to consider the first term for which, by~\eqref{e:nonlinA+}, we have
\begin{equ}
\langle \tvN,\nf_\phi\rangle=\langle \tvN_2,\genap\phi\rangle=\langle-\genam \tvN_2,\phi\rangle,
\end{equ}
and the right hand side converges to $\langle\phi,(-\cD)\phi\rangle$ 
by~\eqref{e:Diffusivity2} in Theorem~\ref{thm:MainConv2}. 
%
%
%
%The second term at the right hand side can be written as
%\begin{equs}[eq:lunga]
%  \|(-\cL_0)^{\frac12}\vN\|^2&= \langle\vN,(-\cL_0)\vN\rangle= \langle\vN,(-\cL^\eps)\vN\rangle\\
%  &= \langle\vN,\nf_\phi+g^n_\phi\rangle= \langle \vN_2,\nf_\phi\rangle+\langle\vN_1,g^n_\phi\rangle\\
%  &= \langle \vN_2,\genap\phi\rangle+\langle\vN_1,g^n_\phi\rangle=\langle-\genam\vN_2,\phi\rangle+\langle\vN_1,g^n_\phi\rangle\,,
%\end{equs}
%where we used the antisymmetry of $\genan$, the fact that $\vN$ solves~\eqref{e:nonlinGenEq}, orthogonality 
%of different chaoses and that, by definition of $\genap$, $\nf_\phi=\genap\phi$. 
%Since $\vN$ solves~\eqref{e:nonlinGenEq}, we further have
%\begin{equ}
%-\gensy \vN_1-\genam\vN_2=g^n_\phi
%\end{equ}
%which implies that the right hand side of~\eqref{eq:lunga} equals
%\begin{equ}
%\langle g^n_\phi+\gensy\vN_1,\phi\rangle+\langle\vN_1,g^n_\phi\rangle=\langle g^n_\phi,\phi\rangle+\langle\vN_1,\gensy\phi+g^n_\phi\rangle\,.
%\end{equ}
%Collecting the previous equalities, we have shown that the expectation of $F^{\eps,n}$ in~\eqref{eq:123} 
%is given by 
%\begin{equ}
%\E[F^{\eps,n}(\eta)]=\langle \phi, g_n^\phi\rangle +\langle \vN_1, -\gensy\phi +g^n_\phi\rangle\,.
%\end{equ}
%By applying Cauchy-Schwarz at the second term at the right hand side, we obtain
%\[
%\langle \vN_1, -\gensy\phi +g^n_\phi\rangle\lesssim_\phi \|\vN_1\|
%\]
%and the norm of $\vN$ goes to $0$ in the limit as $\eps\to0$ thanks to Proposition~\ref{}\giuseppe{Convergence to $0$ in $L^2$}. 
%Since we also have that $g^n_\phi=\cD_n\phi$ 
%converges to $\cD\phi$ by Proposition~\ref{}\giuseppe{Convergence of $\cD_n$}, the statement follows at once. 
\end{proof}  

It remains to show that the quadratic variation of $\Mn_\cdot(\phi)$ concentrates around its mean. 
Since we always take the limit 
in $\eps$ before that in $n$, from now on we will fix $n$ and prove that, as $\eps\to0$, 
the variance of $\langle\Mn_\cdot(\phi)\rangle$ converges to $0$. 
To do so, we will actually prove a stronger statement, namely that the sequence of random 
variables $F^{\eps,n}$ in~\eqref{e:IntQV} is asymptotically deterministic. 
%Let us introduce that random variable 
%\begin{equ}[e:IntQV]
%F^{\eps,n}(\eta)=\sum_k|k|^2|D_k(\eta(\phi)+\vN(\eta))|^2\,,
%\end{equ}
%which is nothing but the integrand in~\eqref{e:QV}. In the next proposition we show that 
%as $\eps\to0$ the variance of $F^{\eps,n}$ vanishes. 

\begin{proposition}\label{p:VanVar}
Let $n\in\N$ be fixed and for $\eps>0$, let $F^{\eps,n}(\eta)$ be the random variable defined according to~\eqref{e:IntQV}. 
Then, 
\begin{equ}[e:VanVar]
\lim_{\eps\to 0} {\rm Var}(F^{\eps,n}(\eta))=0\,.
\end{equ}
\end{proposition}
\begin{proof}
Notice that, by Poincaré inequality for i.i.d. Gaussian measures, we have 
\begin{equs}
{\rm Var}(F^{\eps,n}(\eta))\leq \E[\|D F^{\eps,n}(\eta)\|_{L^2(\T^d)}^2]&=\E\left[\sum_{p\in\Z^d_0}|D_p F^{\eps,n}(\eta)|^2\right]\\
&=\sum_{p\in\Z^d_0}\|D_p F^{\eps,n}\|^2\,,
\end{equs}
where the last step follows by monotone convergence theorem. 
Therefore, if we were able to show that the expectation of the norm of the Malliavin derivative of $F^{\eps,n}(\eta)$ 
vanished, we would be done. In what follows, we will indeed prove that, for any $p\in\Z^d_0$ fixed, 
$\|D_p F^{\eps,n}\|$ goes to $0$, but exchanging the limit in $\eps$ and the sum in $p\in\Z^d_0$ 
would require a much stronger control. To bypass the issue, 
we introduce a martingale approximation $\{M^{\eps,n}_\ell\}_\ell$ of $F^{\eps,n}$, and 
to prove~\eqref{e:VanVar}, we show that for $\ell$ fixed and $\eps\to0$, $M^{\eps,n}_\ell$ converges 
to a constant, and that the limit in $L^2$ of the difference between $M^{\eps,n}_\ell$ and $F^{\eps,n}$ vanishes 
in the double limit, first in $\eps$ and then in $\ell\to\infty$. We will now detail the idea presented above. 

For $\ell\in\N$, denote by $\cF_\ell$ the $\sigma$-algebra generated by $\{\hat\eta(k)\colon|k|\leq \ell\}$ and 
define 
\begin{equ}
M^{\eps,n}_\ell\eqdef \E[F^{\eps,n}|\cF_\ell]\,.
\end{equ}
Clearly, $\{M^{\eps,n}_\ell\}_\ell$ is a martingale in $\ell$ and for all $\ell$
\begin{equ}[e:ExpMartell]
\E[M^{\eps,n}_\ell]= \E[F^{\eps,n}]\,.
\end{equ}
We now bound the variance of $F^{\eps,n}$ by 
\begin{equ}
{\rm Var}(F^{\eps,n})=\E[(F^{\eps,n}-\E[F^{\eps,n}])^2]\lesssim \E[(F^{\eps,n}-M^{\eps,n}_\ell)^2]+\E[(M^{\eps,n}_\ell-\E[M^{\eps,n}_\ell])^2]
\end{equ}
where we used~\eqref{e:ExpMartell}. 
Notice that upon proving 
\begin{equs}
&\lim_{\ell\to\infty}\limsup_{\eps\to 0} \E[(F^{\eps,n}-M^{\eps,n}_\ell)^2]=0\label{one}\\
&\lim_{\eps\to 0} \E[(M^{\eps,n}_\ell-\E[M^{\eps,n}_\ell])^2]=0\label{two}
\end{equs}
the conclusion follows, and this is what we will show next. 
\medskip

\noindent{\bf Proof of \eqref{one}} At first we want to ensure that, as $\eps\to 0$, at least along subsequences, 
$F^{\eps,n}$ and $M^{\eps,n}_\ell$ converge to square integrable random variables. 

Now, $F^{\eps,n}$ lives in the $n$-th inhomogeneous chaos, i.e. in $\oplus_{j\leq n} \SH_j$, hence 
Gaussian hypercontractivity~\cite[Theorem 5.10]{Jan} implies that, for any $p\geq 1$ there exists a constant $C=C(n,p)>0$ 
such that 
\begin{equ}
\E[(F^{\eps,n})^p]^{\frac1{p}}\leq C\, \E[F^{\eps,n}]\,.
\end{equ}
By~\eqref{eq:123} and Proposition~\ref{P:AprioriGeneratorEq} for $d\geq3$ and Lemma~\ref{l:aprioriAnsatz} for $d=2$, 
the expectation is bounded by 
\begin{equs}
\E[F^{\eps,n}]&=2\|(-\cL_0)^{\frac12}(\phi+\wN)\|^2= 2\|\phi\|^2_{H^1(\T^d)}+2\|(-\cL_0)^{\frac12}\wN\|^2\\
&\lesssim \|\phi\|^2_{H^1(\T^d)}+\|(-\cL_0)^{-\frac12}\nf_\phi\|^2\lesssim \|\phi\|^2_{H^1(\T^d)}
\end{equs}
where in the last step we used~\eqref{e:nonlinUnifBound}. 

Hence, uniformly in $\eps$, $F^{\eps,n}$  has bounded moments 
of all orders, which ensures that the sequence $\{F^{\eps,n}\}_\eps$ is tight. 
As a consequence, along subsequences (that, abusing the notation, will still be denoted by $\{F^{\eps,n}\}_\eps$) 
$\{F^{\eps,n}\}_\eps$ can be embedded in a common probability space and there exists $F^n$ 
such that $\{F^{\eps,n}\}_\eps$ converges 
almost surely and in $L^2(\P)$ to $F^n$, so that in particular
\begin{equ}[e:L2F]
\lim_{\eps\to0}\E[(F^{\eps,n}- F^n)^2]=0\,.
\end{equ}
Taking the same subsequence (in $\eps$) for $M^{\eps,n}_\ell$, by Jensen, we immediately have that 
$M^{\eps,n}_\ell$ has uniformly bounded moments in $\ell$ and $\eps$, and that, along this subsequence,
\begin{equ}
\lim_{\eps\to 0}\E[(M^{\eps,n}_\ell-M^n_\ell)^2]=0
\end{equ}
where $M^n_\ell$ is given by $M^n_\ell =\E[F^n|\cF_\ell]$. Notice that, by definition $\{M^n_\ell\}_\ell$ is again a 
square-integrable martingale. Hence, by Martingale convergence theorem, we conclude that, as $\ell\to\infty$ 
$M_\ell^n$ converges in $L^2(\P)$ to $\E[F^n|\cF_\infty]=F^n$. Summarising, we showed that
\begin{equ}
\lim_{\ell\to\infty}\limsup_{\eps\to 0} \E[(F^{\eps,n}- M^{\eps,n}_\ell)^2]\lesssim \lim_{\ell\to\infty} \E[(F^{n}- M^{n}_\ell)^2]=0
\end{equ}
along the above-mentioned subsequence in $\eps$. However, we also get \eqref{one} along \emph{any}
  sequence $\{\eps_k\}$ converging to zero, simply because (by
  exactly the same argument as above) for any such sequence one can
  extract a subsequence for which the convergence holds.
\medskip

\noindent{\bf Proof of \eqref{two}} Poincar\'e's inequality implies 
\begin{equ}
\E[(M^{\eps,n}_\ell-\E[M^{\eps,n}_\ell])^2]\leq \E[\|D M^{\eps,n}_\ell\|_{L^2(\T^d)}^2]=\sum_{|p|\leq\ell}\|D_p M^{\eps,n}_\ell\|^2\,,
\end{equ}
where the last equality follows by the fact that $M^{\eps,n}_\ell$ is a random variable that only depends on 
$\{\hat\eta(k)\colon|k|\leq \ell\}$ and therefore, for $|p|>\ell$, $D_p M^{\eps,n}_\ell=0$. Since $\ell$ is fixed, 
in order to prove that the right hand side above converges to $0$, it suffices to show that, 
for every $p$ fixed, $\lim_{\eps\to0}\|D_p M^{\eps,n}_\ell\|$ vanishes.
Note that 
\begin{equs}
\|D_p M^{\eps,n}_\ell\|&=\E[|D_p \E[F^{\eps,n}|\cF_\ell]|^2]^{1/2}\\
&=\E[| \E[D_pF^{\eps,n}|\cF_\ell]|^2]^{1/2}\leq \E[| D_pF^{\eps,n}|^2]^{1/2}=\|D_p F^{\eps,n}\|
\end{equs}
so that we can focus on the latter. By~\eqref{e:IntQV}, the $p$-th Fourier mode of the Malliavin 
derivative of $F^{\eps,n}$ is given by
\begin{equs}
D_p F^{\eps,n}(\eta)=&\sum_{k} |k|^2D_p|D_k(\eta(\phi)+\WN(\eta))|^2\\
=&\sum_{k} |k|^2\Big[D_{k}(\eta(\phi)+\WN(\eta))D_pD_{-k}(\eta(\phi)+\WN(\eta))\\
&\qquad\qquad+D_pD_k(\eta(\phi)+\WN(\eta))D_{-k}(\eta(\phi)+\WN(\eta))\Big]
\end{equs}
so that applying Cauchy-Schwarz in both $\E$ and the sum in $k$, we have 
\begin{equs}
\|D_p F^{\eps,n}\|^2&\lesssim \E\left[\sum_{k} |k|^2|D_k(\eta(\phi)+\WN(\eta))|^2\right] \E\left[\sum_{k} |k|^2|D_kD_p\WN(\eta)|^2\right]\\
&=4\|(-\gensy)^{1/2}(\phi+\wN)\|^2\|(-\gensy)^{1/2}D_p\wN\|^2\,,
\end{equs}
where we used the fact that $D_pD_k\eta(\phi)=0$ for any $p,k\in\Z^d_0$. Now, the first factor at the right hand side is 
uniformly bounded in $\eps$ (see the first part of the proof of~\eqref{one}). For the second, we need some 
more refined estimates, which will be given in the next subsection. More specifically, in Proposition~\ref{p:BoundH1Mall}, 
we will show that 
\begin{equ}
\lim_{\eps\to0}\|(-\gensy)^{1/2}D_p\wN\|^2=0
\end{equ}
and therefore the proof of~\eqref{two} is complete. 
\end{proof}

Thanks to the previous propositions, the proof of Theorem~\ref{thm:MartConv} is straightforward. 

\begin{proof}[of Theorem~\ref{thm:MartConv}]
$\{\Mn_{\cdot}(\phi)\}_\eps$ is a sequence of continuous martingales which has uniformly bounded (in $\eps$) 
moments of all orders. 
Indeed, by Burkholder-Davis Gundy inequality, we have 
\begin{equ}
\Exp[\sup_{s\leq t} | \Mn_{s}(\phi)|^p]\lesssim_{p} \Exp\langle \Mn_{\cdot}(\phi)\rangle_t^{p/2}\lesssim_{t} \E[|F^{\eps,n}(\eta)|^{p/2}]%\lesssim_p\E[F^{\eps,n}(\eta)]
\end{equ}
where we used the definition of the quadratic variation of $\Mn_\cdot(\phi)$ in~\eqref{e:QV} and Jensen's inequality. 
Now, in the proof of Proposition~\ref{p:VanVar}, we showed that the right hand side is uniformly bounded in $\eps$, 
from which the claim follows. 
According to~\cite[Theorem 7.1.4]{EK}, to complete the proof it suffices to show that, for any $t>0$ the sequence of 
quadratic variations $\{\langle \Mn_{\cdot}(\phi)\rangle_t\}_{\eps,n}$ converges to the deterministic linear function of $t$ 
given by~\eqref{e:QVlimit}. 
Let us abbreviate $Q(\phi)\eqdef \langle\phi,(-\gensy-\cD)\phi\rangle$. 
Since
\begin{equs}
&\Exp\left[\Big(\langle \Mn_{\cdot}(\phi)\rangle_t-2tQ(\phi)\Big)^2\right]\\
&\lesssim {\rm Var}(\langle \Mn_{\cdot}(\phi)\rangle_t)+\Big(\E[\langle \Mn_{\cdot}(\phi)\rangle_t]-2tQ(\phi)\Big)^2\\
&= \E\left[\Big(\int_0^t (F^{\eps,n}(\eta_s)-\E[F^{\eps,n}(\eta_s)])\dd s\Big)^2\right]+\Big(\E\langle \Mn_{\cdot}(\phi)\rangle_t-2tQ(\phi)\Big)^2\\
&\lesssim_t {\rm Var}(F^{\eps,n}(\eta)) +\Big(\E\langle \Mn_{\cdot}(\phi)\rangle_t-2tQ(\phi)\Big)^2\,,
\end{equs}
by taking the limit in $\eps$ first, Proposition~\ref{p:VanVar} implies that the first summand goes to $0$, 
and taking the one in $n$ after also the second summand vanishes, this time by Proposition~\ref{p:MeanQV}.
\end{proof}


\subsubsection{Bounds on the Malliavin derivative}

The aim of this section is to show that, for any given $p$, $\|(-\gensy)^\half D_p \wN\|$ converges 
to $0$ as $\eps\to 0$, where $\wN\in\fock$ is either the solution $\vN$ 
of the truncated generator equation in~\eqref{e:nonlinGenEq} for $d\geq 3$ 
or $\tvN$ given by~\eqref{e:nonlinAnsatz} for $d=2$. 
The proof boils down to obtain a priori estimates on the Malliavin derivative of $\wN$  
for which we need suitable bounds on the commutators 
between $\gensy$, $\gensy^\fw\cGN$, $\genap$ and $\genam$, and $D_p$, where we recall that, 
for two operators $A,B$ the commutator between $A$ and $B$ is defined as $[A,B]\eqdef AB-BA$. 
In the next lemma, we identify an explicit expression for the action of the commutators above 
on functions in the Fock space. 

\begin{lemma}
Let $f\in\fock_n$ and $\genap,\,\genam$ be the operators defined in~\eqref{e:gen} and $\cS$ 
an operator which, for every $n$, maps $\fock_n$ to itself and whose action 
on $\phi\in\fock_n$ is given in terms of the symmetric Fourier multiplier $\sigma_n\colon\Z_0^{dn}\to\R$ as
\begin{equ}
\cF(\cS\phi)(k_{1:n})=\sigma_n(k_{1:n})\hat \phi(k_{1:n})\,.
\end{equ}
For any $p\in\Z_0^d$,
the commutators $[\cS, D_p],\, [\genap, D_p],\, [\genam, D_p]$ act on $f$ as
\begin{align}
 &\cF([\cS, D_p]f)(k_{1:n-1})=(\sigma_{n-1}(k_{1:n-1})-\sigma_n(k_{1:n-1},p))\cF(D_pf)(k_{1:n-1})\label{e:CommL0}\\
 &\cF([\genap, D_p]f)(k_{1:n})=-2\iota\lambda_\eps\sum_{i=1}^n \,[\fw\cdot(k_i+p)]\indN{k_i,p}\hat f(k_i+p,k_{\{1:n\}\setminus\{i\}})\label{e:CommA+}\\
 &\cF([\genam, D_p]f)(k_{1:n-2})=-\iota n(n-1)\lambda_\eps\,(\fw\cdot p) \sum_{\ell+m=p}\indN{\ell,m}\hat f(\ell,m,k_{1:n-2})\,.\label{e:CommA-}
\end{align}
\end{lemma}
\begin{proof}
Recall how the Malliavin derivative $D_p$ acts on $\fock$ in~\eqref{e:MallDerFock}. 
Then, 
\begin{equs}
\cF([\cS, D_p]f)(k_{1:n-1})&=n(\sigma_{n-1}(k_{1:n-1})-\sigma_n(k_{1:n-1},p))f(k_{1:n-1},p)\\
&=(\sigma_{n-1}(k_{1:n-1})-\sigma_n(k_{1:n-1},p))D_pf(k_{1:n-1})\,.
\end{equs}
We now turn to the commutator of $\genap$ and $D_p$. We have 
\begin{equs}
\cF(D_p\genap f)(k_{1:n},p)=&2\iota\lambda_\eps\sum_{1\le i<j\le n} \,[\fw\cdot(k_i+k_j)]\indN{k_i,k_j} \hat f(k_i+k_j,k_{\{1:n\}\setminus\{i,j\}},p)\\
&+2\iota\lambda_\eps\sum_{i=1}^n \,[\fw\cdot(k_i+p)]\indN{k_i,p} \hat f(k_i+p,k_{\{1:n\}\setminus\{i\}})\\
=\cF(\genap D_p f)&(k_{1:n})+2\iota\lambda_\eps\sum_{i=1}^n \,[\fw\cdot(k_i+p)]\indN{k_i,p} \hat f(k_i+p,k_{\{1:n\}\setminus\{i\}})
\end{equs}
from which~\eqref{e:CommA+} follows. At last, for the commutator of $\genam$ and $D_p$, note that
\begin{equs}
\cF(D_p\genam  f) (k_{1:n-2})=&\iota n(n-1)\lambda_\eps\,\sum_{j=1}^{n-2}(\fw\cdot k_j) \sum_{\ell+m=k_j}\indN{\ell,m} \hat f(\ell,m,k_{\{1:n-2\}\setminus\{j\}}, p)\\
&+\iota n(n-1)\lambda_\eps\,(\fw\cdot p) \sum_{\ell+m=p}\indN{\ell,m} \hat f(\ell,m,k_{1:n-2})\\
=\cF(\genam D_p f)&(k_{1:n-2})+\iota n(n-1)\lambda_\eps\,(\fw\cdot p) \sum_{\ell+m=p}\indN{\ell,m} \hat f(\ell,m,k_{1:n-2})
\end{equs}
which proves~\eqref{e:CommA-}. 
\end{proof}

We now determine bounds 
on the commutators in the previous lemma. 

\begin{lemma}\label{l:CommEstimates}
Let $n\geq 1$, $ f\in\fock_n$ and $p\in\Z^d_0$. Then, there exists a constant $C>0$ independent of $n,\eps$ and $p$ 
for which  following estimates hold
\begin{align}
&\| \cN[\genap, D_p] f\|\leq C \lambda_\eps\|\cN^{3/2}(-\gensy)^{1/2} f\|\,,\label{e:CommA+estimate}\\
&\| \cN[\genam, D_p] f\|\leq C |\fw\cdot p|\|\cN^{5/2}(-\gensy)^{1/2} f\|\,,\label{e:CommA-estimate}\,,
\end{align}
and
\begin{equs}[e:CommL0w-estimate]
\| (-\gensy)^{-1/2} [-\gensy-&\gensy^\fw\cGN, D_p] f\|^2\\
&\leq C|p|^4\Big(\|\cN f\|^2+\lambda_\eps^4 \|\cN(-\gensy)^{1/2}f\|^2\Big)\,.
\end{equs}
\end{lemma}
\begin{proof}
We begin with~\eqref{e:CommA+estimate}. By~\eqref{e:CommA+}, we have 
\begin{equs}
\|\cN [\genap, D_p] f\|^2&=4n^2\lambda_\eps^2 \sum_{k_{1:n}}\Big|\sum_{i=1}^n \,[\fw\cdot(k_i+p)]\indN{k_i,p} \hat f(k_i+p,k_{\{1:n\}\setminus\{i\}})\Big|^2\\
&\leq 4n^3 \lambda_\eps^2 \sum_{k_{1:n}}\sum_{i=1}^n[\fw\cdot(k_i+p)]^2 | \hat f(k_i+p,k_{\{1:n\}\setminus\{i\}})|^2\\
&\lesssim n^3\lambda_\eps^2 \sum_{k_{1:n}}\frac12 |k_{1:n}|^2 | f(k_{1:n})|^2=\lambda_\eps^2\|\cN^{3/2}(-\gensy)^{1/2} f\|^2\,.
\end{equs}
For~\eqref{e:CommA-estimate} instead, we have 
\begin{align*}
&\| \cN[\genam, D_p] f\|^2=n^4(n-1)^2(\fw\cdot p)^2\sum_{k_{1:n-2}}\left|\lambda_\eps\sum_{\ell+m=p}\indN{\ell,m} \hat f(\ell,m,k_{1:n-2})\right|^2\\
&\leq  n^4(n-1)^2(\fw\cdot p)^2\left(\lambda_\eps^2\sum_{\ell}\frac{\indN{\ell,p-\ell}}{|\ell|^2+|p-\ell|^2}\right)\times\\
&\qquad\qquad\qquad\qquad\qquad\qquad\times\sum_{k_{1:n-1}}(|k_1|^2+|p-k_1|^2)| \hat f(k_1,p-k_1,k_{2:n-1})|^2
\end{align*}
where we simply applied Cauchy-Schwarz. Now, the second sum at the right hand side is bounded uniformly in $p$ 
and $\eps$, hence the previous is bounded above by %\giuseppe{This estimate is not really optimal but we do not need more}
\begin{equ}
n^4(n-1)^2(\fw\cdot p)^2\sum_{k_{1:n}}|k_1|^2| f(k_{1:n})|^2\lesssim(\fw\cdot p)^2\|\cN^{5/2}(-\gensy)^{1/2} f\|^2\,.
\end{equ}
At last, we consider~\eqref{e:CommL0w-estimate}. By~\eqref{e:CommL0}, we have 
\begin{equs}
\|&(-\gensy)^{-1/2}[-\gensy-\gensy^\fw\cGN, D_p\big]f\|^2\\
&=\tfrac{n^2}4\sum_{k_{1:n-1}}\frac{1}{|k_{1:n-1}|^2}| f(k_{1:n-1},p)|^2\Big[|p|^2+ (\fw\cdot p)^2 G(\Ll(\tfrac12 |k_{1:n-1}|^2+\tfrac12|p|^2))\\
&\qquad+(\fw\cdot k)^2_{1:n-1}\Big(G(\Ll(\tfrac12 |k_{1:n-1}|^2+\tfrac12|p|^2))-G(\Ll(\tfrac12(|k_{1:n-1}|^2)))\Big)\Big]^2\\
&\lesssim n^2\sum_{k_{1:n-1}} | f(k_{1:n-1},p)|^2\times \\
&\qquad\times \Big[|p|^4+ |k_{1:n}|^2\Big(G(\Ll(\tfrac12 |k_{1:n-1}|^2+\tfrac12|p|^2))-G(\Ll(\tfrac12(|k_{1:n-1}|^2)))\Big)^2\Big]
\end{equs}
where we used that $G$ in~\eqref{eq:defG} is locally bounded and $\Ll$ 
in~\eqref{e:L} is bounded uniformly in $\eps$ for $x\geq 1/2$. 
Now, to control the last term, we use that the function $x\mapsto G(\Ll(x))$ has first derivative uniformly bounded 
by a constant times $\lambda_\eps^2$ so that~\eqref{e:CommL0w-estimate} follows at once. 
%
%
% we obtain 
%\begin{equ}
%\|(-\gensy)^{-1/2}[-\gensy-\gensy^\fw\cGN, D_p\big]f\|^2\lesssim |p|^4\Big(\|\cN f\|^2+\lambda_\eps^4 \|\cN(-\gensy)^{1/2}f\|^2\Big)
%\end{equ}
%so that the proof of the lemma is concluded. 
\end{proof}

The following is the main result of this section, which completes the proof of Proposition~\ref{p:VanVar}. 

\begin{proposition}\label{p:BoundH1Mall}
For $n\in\N$ and $\eps>0$, let $\wN\in\fock$ be either the solution $\vN$  
of the truncated generator equation in~\eqref{e:nonlinGenEq} if $d\geq 3$ 
or $\tvN$ given by~\eqref{e:nonlinAnsatz} if $d=2$. 
Then, for any $p\in\Z^d_0$, there exists a constant $C=C(p)>0$ such that 
\begin{equ}[e:BoundH1Mall]
\lim_{\eps\to 0}\|(-\gensy)^\half D_p \wN\|^2=0\,.
\end{equ}
\end{proposition}
\begin{proof}
Let us separately consider the cases $d=2$ and $d\geq 3$, starting with the former. 

For $d=2$, note that by~\eqref{e:nonlinAnsatz} we have 
\begin{equs}
(-\gensy-\gensy^\fw\cGN) D_p\tvN &= [-\gensy-\gensy^\fw\cGN, D_p]\tvN+D_p(-\gensy-\gensy^\fw\cGN)\tvN\\
&= [-\gensy-\gensy^\fw\cGN, D_p]\tvN +D_p\genap P_2^{n-1}\tvN +D_p\nf_\phi\\
&=\genap P_2^{n-1}D_p\tvN + \psi^\sharp\label{e:NewShape}
\end{equs}
where $\psi^\sharp$ is given by 
\begin{equ}[e:psisharp]
\psi^\sharp\eqdef [-\gensy-\gensy^\fw\cGN, D_p]\tvN+[\genap, D_p]\tvN +D_p \tvN_2+D_p\nf_\phi\,.
\end{equ}
Equation~\eqref{e:NewShape} implies that $D_p\tvN$ solves~\eqref{e:Ansatz} 
in which $\psi$ is replaced by $\psi^\sharp$. Hence, by Lemma~\ref{l:aprioriAnsatz}, 
it satisfies the bound
\begin{equ}
\|(-\gensy)^{1/2}D_p\tvN\|\lesssim \|(-\gensy)^{-1/2}\psi^\sharp\|\,.
\end{equ}
This means that~\eqref{e:BoundH1Mall} follows if we prove that the $(-\gensy)^{-1/2}\fock$-norm 
of each of the summands in~\eqref{e:psisharp} vanishes. 
For last three terms we upperbound $(-\gensy)^{-1/2}$ by $1$. 
Then, for the fourth we use that by~\eqref{e:MallDerFock} and~\eqref{e:nonlin}, we have
\begin{equ}[e:Dpnf]
\|D_p\nf_\phi\|^2\leq \lambda_\eps^2\sum_\ell \indN{\ell,p}(\fw\cdot(\ell+p))^2|\hat\phi(\ell+p)|^2\lesssim \lambda_\eps^2\|\phi\|_{H^1(\T)}^2\,.
\end{equ}
For the third, we proceed as 
\begin{equ}[e:Dpv2]
\|D_p \tvN_2\|\lesssim\|(-\gensy-\gensy^\fw\cGN)^{-1}\nf_\phi\|=\|(-\gensy-\gensy^\fw\cGN)^{-1}\genap\phi\|\leq \lambda_\eps\|\phi\|_{H^1(\T)}^2
\end{equ} 
where we used~\eqref{e:nonlinA+} and~\eqref{e:L2control}. 
For the second term, we exploit~\eqref{e:CommA+estimate}, which gives
\begin{equs}
\|[\genap, D_p]\tvN\|\lesssim \lambda_\eps \|\cN^{5/2}(-\gensy)^{1/2}\tvN\|\lesssim \lambda_\eps\|(-\gensy)^{-1/2}\nf_\phi\|
\end{equs}
where in the last step we used~\eqref{e:AprioriAnsatz}. 
Hence, invoking also~\eqref{e:nonlinUnifBound}, the last three terms in~\eqref{e:psisharp} converge to $0$, 
since so does $\lambda_\eps$. 
For the first, by~\eqref{e:CommL0w-estimate},  we have 
\begin{equs}
\| (-\gensy)^{-1/2} [-\gensy-&\gensy^\fw\cGN, D_p] \tvN\|^2\\
&\lesssim n^2\| \tvN\|^2+\lambda_\eps^4 \|\cN(-\gensy)^{1/2}\tvN\|^2
\end{equs}
Now, the $\fock$-norm of $\tvN$ converges to $0$ as $\eps\to0$ by Theorem~\ref{thm:MainConv2}, 
while as argued above, $\|\cN(-\gensy)^{1/2}\tvN\|$ is uniformly bounded in both $n$ and $\eps$, 
so that the presence of the $\lambda_\eps$ at the second summand ensures it converges to $0$ as well. 
\medskip

We now turn to the case $d\geq 3$. 
Since $\vN$ solves~\eqref{e:nonlinGenEq}, it satisfies
\begin{equs}[e:MallD]
-\gen D_p \vN &= [-\gen, D_p]\vN -D_p\gen \vN\\
&= [-\gen, D_p]\vN +D_p\nf_\phi -D_p\genap \vN_n-D_p\genam \vN_2\,.
\end{equs}
We test both sides by $D_p \vN$, so that, using the antisymmetry of $\gena$, we get
\begin{equs}
\|&(-\gensy)^\half D_p \vN\|^2\\
&=\langle[-\gen, D_p]\vN+D_p\nf_\phi, D_p \vN\rangle -\langle D_p\genap \vN_n+D_p\genam \vN_2, D_p\vN\rangle\\
&=\langle[-\gen, D_p]\vN +D_p\nf_\phi, D_p \vN\rangle \\
&\leq  \langle[-\genap, D_p]\vN+[-\genam, D_p]\vN+D_p\nf_\phi, D_p \vN\rangle \label{e:normDp}
\end{equs}
where, in the second step, we used the orthogonality of Fock spaces with different indices, while in the last that, 
by~\eqref{e:CommL0}, we have
\begin{equ}
\langle[-\gensy, D_p]\vN, D_p \vN\rangle=-\frac12 |p|^2  \|D_p \vN\|^2\leq 0\,.
\end{equ}
At this point we apply Cauchy-Schwarz to the right hand side of~\eqref{e:normDp}, so that 
\begin{equs}[e:normDp1]
\|&(-\gensy)^\half D_p \vN\|^2 \\
&\leq \big(\|\cN [-\genap, D_p]\vN\|+\|\cN [-\genam, D_p]\vN\|+\|\cN D_p\nf_\phi\|\big)\|\cN^{-1}D_p\vN\|\\
&\lesssim \Big(\lambda_\eps \|\cN^{3/2}(-\gensy)^{1/2}\vN\| \\
&\qquad\qquad+ |\fw\cdot p| \|\cN^{5/2}(-\gensy)^{1/2}\vN\|+\|\cN D_p\nf_\phi\|\Big)\|\cN^{-1}D_p\vN\|\\
&\lesssim \Big(\|(-\gensy)^{-1/2}\nf_\phi\|+\| D_p\nf_\phi\|\Big)\|\vN\|
\end{equs}
where in the second step we applied Lemma~\ref{l:CommEstimates}, while in the last 
the a priori estimate in Proposition~\ref{P:AprioriGeneratorEq} and the fact that 
by~\eqref{e:MallDerFock}, we have $\|\cN^{-1}D_p\vN\|\leq \|\vN\|$. 
Now, since the term containing the norm of $\nf_\phi$ is bounded uniformly in $\eps$ 
and, by Proposition~\ref{p:epsLim}, $\vN$ converges to $0$ in $\fock$,~\eqref{e:BoundH1Mall} follows. 
\end{proof}




%\section{Uniqueness in $n$}\label{sec:limit}
%
%We would like to understand the limit
%\begin{equation}
%    \lim_{n,N\to\infty}\langle e_k,H_n^\eps e_j\rangle.
%\end{equation}
%The goal is to first take $\eps$ to $0$ and then $n$ to infinity.
%In this section we will show that the limit in $n$ exists for fixed $\eps$ and actually that this sequence is Cauchy uniformly w.r.t. $\eps$, which guarantees that we can still take the limit in $n$ after taking the limit in $\eps$, if the limit in $\eps$ exists.
%
%We first deal with the case $k\neq -j$ which is a simply argument using translation invariance before taking any limit.
%\begin{lemma}
%For any $k\neq -j$
%\begin{equation}
%    \langle e_k,H_n^\eps e_j\rangle=0.
%\end{equation}
%\end{lemma}
%\begin{proof}
%For $i=1,\dots,d$ let  $P^{(i)}$ be the operator acting on an element $\psi$ of the $n$-th chaos by
%\begin{equation}
%    P^{(i)}\psi(k_{1:n})=\sum_{j=1}^nk_j^{(i)}\psi(k_{1:n}),
%  \end{equation}
%  where $k^{(i)}_j$ is the $i$-th component of $k_j$. Note that this (apart for a missing $\iota=\sqrt{-1
%   }$ in front) is just the generator of the translations in direction $i$.  It is easily checked that $P^{(i)}$ commutes with 
%% Note that this operator commutes with
%  $\mathcal A^\eps_\pm,$ (this is due to the translation invariance of the equation), and that $(P^{(i)})^*=-P^{(i)} $ (recall that $\langle e_j,e_k\rangle= {\bf 1}_{k+j=0}$).
%% \begin{align}
%%     P^{(i)}\mathcal{A}_+\psi(k_{1:n+1})&=\sum_{j=1}^{n+1}k_j^{(i)}in\epsilon_Nw\cdot( k_1+k_2)\psi(k_1+k_2,k_{3:n+1})\\
%%     &=in\epsilon_Nw\cdot( k_1+k_2)(k_1^{(i)}+k_2^{(i)}+\sum_{j=3}^{n+1}k_j^{(i)})\psi(k_1+k_2,k_{3:n+1})\\
%%     &=\mathcal{A}_+P^{(i)}\psi(k_{1:n+1}).
%% \end{align}
%Since $P^{(i)}$ also clearly commutes with $\mathcal L_0$, it also commutes with $\Op_1$ and therefore by induction with $\Op_n$ for all $n$.
%
%We deduce:
%\begin{align}
%    \langle e_k,\Op_ne_j\rangle
%    &=\langle e_k,\Op_nP^{(i)}\frac{1}{j^{(i)}}e_j\rangle
%    =\langle e_k,P^{(i)}\Op_n\frac{1}{j^{(i)}}e_j\rangle\\
%    &=-\langle P^{(i)}e_k,H_n\frac{1}{j_i}e_j\rangle=-\frac{k^{(i)}}{j^{(i)}}\langle e_k,H_ne_j\rangle,
%\end{align}
%from which the statement follows.
%\end{proof}
%
%We deduce that $e_k$ is an eigenvector of $\Op_n$, with some eigenvalue that we call $\mathcal D_n^\eps(k)$. We need to prove 
%that  $\cD_n^\eps(k)$ is a Cauchy sequence w.r.t. $n$, uniformly in $\eps$.
%We use the following estimate:
%\begin{lemma}
%  \label{lem:cauchy}
%For any $m\ge0$ there exists a constant $C$ such that for any $n,\eps$ and $\psi$ in the first chaos we have that
%\begin{equation}
%\|(-\mathcal L_0)^{-\tfrac12}(\Op_n-\Op_{n+1})(-\mathcal L_0+\Op_n)^{-1}\psi\|\leq Cn^{-m}(1+\|(-\mathcal L_0)^{-\tfrac12}\psi\|).
%\end{equation}
%\end{lemma}
%Applying this to $\varphi=(-\mathcal L_0+\Op_n)e_k$ and recalling that
%\begin{eqnarray}
%  \label{eq:unifBH}
%H_n^\eps\le H_2^\eps\le C (-\cL_0)
%\end{eqnarray}
% for some universal constant $C$, 
%\fabioText{THis has to be written somewhere}
%we see that
%\[
%|\cD^\eps_n(k)-\cD^\eps_{n+1}(k)|\le C(k) n^{-m}
%\]
%for some constant $C(k)<\infty$, uniformly in $\eps,n$. If $m>2$, this implies the desired (uniform) Cauchy property.
%\begin{proof}[of Lemma \ref{lem:cauchy}]
%  \fabioText{I have to review this}
%For an arbitrary $\psi$ in the first chaos and $j\in\N$, let $u^{\eps,n}$ be the solution of the truncated resolvent equation
%\begin{equation}
%-\mathcal L_n^\eps u^{\eps,n}=\psi\,.
%\end{equation}
%For lightness of notation, in this proof we write just $u^n$ and drop the $\eps$ dependence.
%Now, let $n\in\N$ and consider the difference $u^{n+1}-u^{n}$. Notice that 
%\begin{equs}
%-\gen (u^{n+1}-u^{n})&=-\gen_{n+1} u^{n+1}-\genap u^{n+1}_{n+1} +\gen_{n} u^{n}+\genap u^{n}_{n}\\
%&=-\genap(u^{n+1}_{n+1}-u^{n}_{n})\,.
%\end{equs}
%Upon testing both sides of the above equation by $u^{n+1}-u^{n}$, we get 
%\begin{equs}
%\|(-\gensy)^{1/2} (u^{n+1}-u^{n})\|^2&=\langle u^{n+1}-u^{n}, -\genap(u^{n+1}_{n+1}-u^{n}_{n})\rangle=\langle u^{n+1}-u^{n}, \genap u^{n}_{n}\rangle\\
%&\le \delta^{-1}\|(-\gensy)^{-1/2}\genap u^{n}_{n}\|^2 + \delta \|(-\gensy)^{1/2} (u^{n+1}-u^{n})\|^2
%\end{equs}
%where the last equality holds for any $\delta>0$. Upon choosing $\delta$ small enough and exploiting~\eqref{e:chaosdecay} 
%with $m\geq 1$ and Proposition \ref{P:AprioriGeneratorEq}, we obtain
%\begin{equ}
%\|(-\gensy)^{1/2} (u^{n+1}-u^{n})\|^2\lesssim_{\delta}\|(-\gensy)^{-1/2}\genap u^{n}_{n}\|^2\lesssim_{\delta,m} n^{-m}(1+\|(-\mathcal L_0)^{-\tfrac12}\psi\|)
%\end{equ}
%uniformly over $\eps$. Now, by definition of the scalar product~\eqref{} and~\eqref{e:firstchaos}, we have 
%\begin{equs}
%\|(-\gensy)^{1/2} (u^{n+1}-u^{n})\|&\geq \|(-\gensy)^{1/2} (u_1^{n+1}-u_1^{n})\|\\
%&=\|(-\gensy)^{1/2}\Big((-\mathcal L_0+\Op_{n+1})^{-1}-(-\mathcal L_0+\Op_n)^{-1}\Big)\psi\|\\
%&=\|(-\gensy)^{1/2}(-\mathcal L_0+\Op_{n+1})^{-1}\Big(\Op_n-\Op_{n+1}\Big)(-\mathcal L_0+\Op_n)^{-1}\psi\|\,.
%\end{equs}
%To conclude, it suffices to note that by monotonicity $\Op_{n+1}\leq H_1\leq -C\gensy$, so that the result follows by 
%an easy estimate. \fabioText{Isn't it simply that H is positive? anyway for the psi we have above,
%  which is an eigenvector, this expression is much simpler}
%\end{proof}
%

%\section{Uniqueness in $\eps$}
%\giuseppeText{MOVE EARLIER}
%Here we prove that for each $k\in \Z^d$ and $n\ge1$, $\cD^\eps_n(k)$ has a limit $\mathcal D_n(k)$ as $\eps\to0$.
%Moreover, the limit is of the form
%\begin{eqnarray}
%  \label{eq:oftherform}
%  \cD_n(k)=|k|^2 f_n((w\cdot k)^2/|k|^2)
%\end{eqnarray}
%for some function $f_n:\mathbb R^+\mapsto \R^+$ (the fact that $f_n$ is non-negative follows simply from the fact that the operators $\Op_n$ are).
%We start by recalling that
%\begin{eqnarray}\label{eq:rapporto}
%  \frac1{|k|^2+\cD^\eps_n(k)}= \langle e_{-k},(-
%  \cL_0+\Op_n)^{-1}e_k\rangle= \langle e_{-k},(-\cL^\eps_n)^{-1}e_k\rangle
%\end{eqnarray}
%where the second equality follows from \eqref{eq:truncgensys} when $g=e_k$ is concentrated on the single Fourier mode $k$.
%

%\section{NOTES: General non-linearity}
%
%I just want to write down some notes on how we could address a more general nonlinearity. 
%The (scaled) equation we want to look at is
%\begin{equ}\label{e:GenBurgersScaled}
%\partial_t\eta^\eps = \frac{1}{2}\Delta\eta^\eps +\eps^{-\frac{d}{2}-1} w\cdot\Pi_{1/\eps}\nabla F(\eps^{d/2}\Pi_{1/\eps}\eta^\eps) + (-\Delta)^{\frac{1}{2}}\xi
%\end{equ}
%where $F\colon \R\to\R$ is (for now) a generic function. 
%Since the Hermite polynomials form a basis for the space $L^2(\nu)$ for $\nu(\dd x)\eqdef\gamma(x)\dd x$ and 
%$\gamma$ the standard Gaussian density, $F$ can be written as 
%\begin{equ}[e:F]
%F(x)=\sum_{m\geq 0} c_m(F) H_m(x)\,,\qquad x\in\R
%\end{equ}
%provided $\E[F(U)^2]<\infty$, where $U$ is a standard Gaussian random variable. 
%In the above, $H_m$ is the standard $m$-th  Hermite polynomial and  
%the coefficients $c_m(F)$ are given by 
%\begin{equ}
%c_m(F)\eqdef \frac{1}{m!}\E[F(U) H_m(U)]=\frac{1}{m!}\int_\R F(x) H_m(x) \gamma(x)\dd x\,.
%\end{equ}
%
%Following~\cite{GPHQ}, when evaluated at $\eps^{d/2}\Pi_{1/\eps}\eta^\eps(x)$, 
%$F$ is
%\begin{equ}[e:Fexpeta]
%F(\eps^{d/2}\Pi_{1/\eps}\eta^\eps(x))=\sum_{m\geq 0} c_n(F) \eps^{\frac{d}{2}m}\sum_{k_{1:m}\colon |k_i|\leq \eps^{-1}}\frac{e_{k_1+\dots+k_m}(x)}{(2\pi)^{\frac{d}{2}(m-1)}}\Wick{\prod_{i=1}^m\hat\eta(k_i)}\,.
%\end{equ}
%Without loss of generality, we can assume $c_1(F)=0$. Indeed, if $c_1(F)\neq 0$, we consider the process
%\begin{equ}
%\tilde\eta^\eps(t,x)\eqdef \eta^\eps(t, x-\eps^{-1}c_1(F) t w)
%\end{equ}
%which solves
%\begin{equ}
%\partial_t\tilde\eta^\eps = \frac{1}{2}\Delta\tilde\eta^\eps +\eps^{-\frac{d}{2}-1} w\cdot\Pi_{1/\eps}\nabla \tilde F(\eps^{d/2}\Pi_{1/\eps}\tilde\eta^\eps) + (-\Delta)^{\frac{1}{2}}\xi
%\end{equ}
%for $\tilde F$ given by 
%\begin{equ}
%\tilde F(x)\eqdef F(x)-c_1(F) x
%\end{equ}
%and $\tilde F$ is clearly such that $c_1(\tilde F)=0$. 
%
%\subsection{Tightness}
%
%At first, we want to verify under what conditions on $F$ we (can hope to) have tightness. Let $\NN_\phi$ 
%denote the nonlinearity in~\eqref{e:GenBurgersScaled} tested against a smooth test function $\phi$, i.e. 
%\begin{equs}[e:NonlinGen]
%\NN_\phi(\eta)&\eqdef \eps^{-\frac{d}{2}-1} w\cdot\Pi_{1/\eps}\nabla F(\eps^{d/2}\Pi_{1/\eps}\eta^\eps) (\phi)\\
%&=\sum_{m\geq 0} c_m(F)\iota \eps^{\frac{d}{2}(m-1)-1}\sum_{k_{1:m}}\indN{k_{1:m}}\frac{(w\cdot \overline{k_{1:m}})\hat\phi(-\overline{k_{1:m}})}{(2\pi)^{\frac{d}{2}(m-1)}}\Wick{\prod_{i=1}^m\hat\eta(k_i)}\\
%&=\sum_{m\geq 1} c_m(F)  I_m\left(\iota\frac{\eps^{\frac{d}{2}(m-1)-1}}{(2\pi)^{\frac{d}{2}(m-1)}}\sum_{k_{1:m}}\indN{k_{1:m}}(w\cdot \overline{k_{1:m}})\hat\phi(-\overline{k_{1:m}})e_{-k_{1:m}}\right)
%\end{equs}
%where $I_m$ was introduced in Section~\ref{}, we introduced the notations $\overline{k_{1:m}}\eqdef k_1+\dots +k_m$, 
%$e_{k_{1:m}}\eqdef e_{k_1}\otimes\dots\otimes e_{k_m}$ and denoted by $\indN{k_{1:m}}$ 
%the indicator 
%\begin{equ}[e:Indicator]
%\indN{k_{1:m}}\eqdef \1\{0<|k_i|_\infty\le \eps^{-1}, 0<|\overline{k_{1:n}}|_\infty\le \eps^{-1}\colon i=1,\dots, m \}.
%\end{equ}
%In what follows, we will write $\NN_k(\eta)\eqdef \NN_{e_{-k}}(\eta)$ for the $k$-th Fourier component of $\NN(\eta)$. 
%By the It\^o trick (actually Kipnis-Varadhan), we now that 
%\begin{equ}
%\Exp\left[\sup_{t\leq T}\left(\int_0^t\NN_\phi(\eta^\eps_s)\dd s\right)^2\right]\lesssim T \|(-\gensy)^{-1/2}\NN_\phi\|^2\,.
%\end{equ}
%Thanks to~\eqref{e:gen},~\eqref{e:Fexpeta} and orthogonality of Fock spaces with different indices, 
%the term at the right hand side can be written explicitly and equals
%\begin{equs}[e:tightness]
%\|(-\gensy)^{-1/2}\NN_\phi\|^2&=\sum_{m\geq 2} m! c_m(F)^2\eps^{d(m-1)-2}\sum_{k_{1:m}}(w\cdot (\overline{k_{1:m}}))^2\frac{|\hat\phi(\overline{k_{1:m}})|^2}{|k_{1:m}|^2}\indN{k_{1:m}}\\
%&=\sum_{m\geq 2} n! c_n(F)^2\sum_{k}(w\cdot k)^2|\hat\phi(k)|^2\left(\eps^{d(m-1)}\sum_{\overline{k_{1:m}}=k}\frac{\indN{k_{1:m}}}{|\eps k_{1:m}|^2}\right)\,.
%\end{equs}
%Notice that, as $\eps\to 0$ the quantity in parenthesis at the right hand side converges to a finite constant 
%independent of $k$ {\color{red}but depending on $m$} - let it be $\bar C(m)$, i.e. 
%\begin{equ}[e:TightConst]
%\bar C(m)\eqdef \int\frac{\mathbb J^1_{x_{1:m-1}}\dd x_{1:m-1}}{|x_{1:m-1}|^2+|\overline{x_{1:m-1}}|^2}\,,
%\end{equ}
%where $\overline{x_{1:m-1}}\eqdef \sum_{i=1}^{m-1} x_i$. 
%Hence, we obtain 
%\begin{equs}
%\lim_{\eps\to 0}\|&(-\gensy)^{-1/2}\NN_\phi\|^2= \left(\sum_{k}(w\cdot k)^2|\hat\phi(k)|^2\right)\sum_{m\geq 2} m! \bar C(m) c_m(F)^2\,.
%\end{equs}
%It follows that, in order to get tightness and be sure that the previous equalities hold, we need $F$ to satisfy
%\begin{equs}[e:assF]
%\sum_{m\geq 2} m!  \bar C(m) c_m(F)^2<\infty\,,
%\end{equs}
%where $\bar C(m)$ is given in~\eqref{e:TightConst}. 
%Notice that, it is quite important to have some sort of understanding of how $\bar C(m)$ depends on $m$ as this will tell us 
%what is the assumption on $F$...
%
%\begin{remark}
%\giuseppeText{Let me point out that~\eqref{e:assF} is not enough for tightness. 
%Indeed, it is not sufficient to bound the second moment, since in the nonliinearity  
%a martingale term is hidden, so that the power of $T$ cannot be improved. 
%To control higher moments, we would need to use Gaussian hypercontractivity, 
%which in turn will introduce some growth in the chaos. We will probably need something of the form
%\begin{equ}
%\sum_{m\geq 2}  m! C_p(m) \bar C(m) c_m(F)^2<\infty
%\end{equ} 
%where $C_p(m)$ is the hypercontractivity constant which depends on $m$. }
%%I need to check what is 
%%the best constant, but anyway we can bound the previous by 
%%\begin{equ}
%%\sum_{n\geq 2} C_p(n) n! c_n(F)^2\lesssim \E F^{(j(p))}(U)^2]
%%\end{equ} }
%%where $F^{(\ell)}$ is the $\ell$-th derivative of $F$ and $j(p)$ is an integer whose value depend on $p$ 
%%via $C_p$. 
%\end{remark}
%
%\subsection{Analysis of the generator}
%
%Arguing as usual, it can be shown that the generator $\gen$ for~\eqref{e:GenBurgersScaled} can 
%be written as the sum of $\gensy$ and $\gena$ whose action on a cylinder random variable $G$ is given by 
%\begin{align}
%(\gensy G)(\mu) &\eqdef \frac{1}{2} \sum_{k \in \Z^d} |k|^2 (-\mu_{-k} D_k  +  D_{-k}D_k )G(\mu) \label{e:gensGen}\\
%(\gena G)(\mu) &\eqdef \sum_{k \in \Z_0^d} \NN_k(\mu) D_{-k} G(\mu)\,. \label{e:genaGen}
%\end{align}
%Since the analysis of $\gensy$ is always the same, we will only focus on $\gena$ and study 
%how it acts on the Fock space $\fock$. 
%Notice first that, by~\eqref{e:NonlinGen} we can rewrite~\eqref{e:genaGen} as
%\begin{equs}
%&(\gena G)(\mu)=\sum_{k \in \Z_0} \cN^\eps_k(\mu) D_{-k} G(\mu)\\
%&=\sum_{k \in \Z_0}\sum_{m\geq 2} c_m(F)  I_m\left(\iota\frac{\eps^{\frac{d}{2}(m-1)-1}}{(2\pi)^{\frac{d}{2}(m-1)}}\sum_{k_{1:m}}\indN{k_{1:m}}(w\cdot \overline{k_{1:m}})\hat\phi(-\overline{k_{1:m}})e_{-k_{1:n}}\right)D_{-k} G(\mu)\\
%&=\sum_{m\geq 2} c_m(F) \sum_{k \in \Z_0} I_m\left(\iota\frac{\eps^{\frac{d}{2}(m-1)-1}}{(2\pi)^{\frac{d}{2}(m-1)}}\sum_{k_{1:m}}\indN{k_{1:m}}(w\cdot \overline{k_{1:m}})\hat\phi(-\overline{k_{1:m}})e_{-k_{1:n}}\right)D_{-k} G(\mu)\\
%&=:\sum_{m\geq 2} c_m(F) (\genaF{m}{} G)(\eta)\,.\label{e:GenDecomp}
%\end{equs}
%Hence, we need to study the operators $\genaF{m}{}$ and analyse their action on $\fock$. 
%
%\subsubsection{The operators $\genaF{m}{}$}
%
%In order to analyse how $\genaF{m}{}$ acts on $I_n(K)$, by polarisation, it suffices to consider the case of 
%$K=\otimes^n h$, $h\in L^2(\T^d)$, i.e. $K$ given by the tensor product of $n$ copies of $h$. 
%To do so notice that, by~\eqref{e:genaGen} we need to consider $D_{-k}I_n(h)$, which, in view of~\eqref{e:Der} 
%is given by 
%\begin{equ}
%D_{-k} I_n(K)=n I_{n-1}(\otimes^{n-1} h) \hat h(-k)\,.
%\end{equ}
%Hence, we obtain
%\begin{equs}\label{e:A}
%&\genaF{m}{} I_n(K)\\%=\sum_{k \in \Z_0^d} \NN_k(\mu) D_{-k} I_n(K)=n I_{n-1}(\otimes^{n-1} h) \sum_{k\in\Z_0^d}\NN_k(\mu)\hat h(-k)\\
%&=\iota\frac{\eps^{\frac{d}{2}(m-1)-1}}{(2\pi)^{\frac{d}{2}(m-1)}} n I_{n-1}(\otimes^{n-1} h) I_m\left(\sum_{k_{1:m}}\indN{k_{1:m}}(w\cdot \overline{k_{1:m}})\hat h(-\overline{k_{1:m}})e_{-k_{1:m}}\right)\,.
%\end{equs}
%In order to evaluate the product of $I_{n-1}$ and $I_m$, we apply~\eqref{e:Mult}, which gives
%\begin{equs}[e:prodI]
%&I_{n-1}(\otimes^{n-1} h) I_m\left(\sum_{k_{1:m}}\indN{k_{1:m}}(w\cdot \overline{k_{1:m}})\hat h(-\overline{k_{1:m}})e_{-k_{1:m}}\right)\\
%&=\sum_{p=0}^{ m} \1_{p\leq n-1}  p! \binom{n-1}{p}\binom{m}{p} \times\\
%&\quad\times I_{n-1+m-2p}\left(\sum_{k_{1:m}}\indN{k_{1:m}}(w\cdot \overline{k_{1:m}})\hat h(-\overline{k_{1:m}})\prod_{i=1}^p \hat h(k_i)e_{-k_{p+1:m}}\otimes^{n-1-p} h\right)\,.
%\end{equs}
%As a first observation, we point out that the summand corresponding to $p=m$ is always $0$. 
%Indeed, 
%\begin{equs}
%I_{n-1-m}&\left(\sum_{k_{1:m}}\indN{k_{1:m}}(w\cdot \overline{k_{1:m}})\hat h(-\overline{k_{1:m}})\prod_{i=1}^m \hat h(k_i)\otimes^{n-1-m} h\right)\\
%&=\left(\sum_{k_{1:m}}\indN{k_{1:m}}(w\cdot \overline{k_{1:m}})\hat h(-\overline{k_{1:m}})\prod_{i=1}^m \hat h(k_i)\right)I_{n-1-m}(\otimes^{n-1-m} h)\\
%&=\langle h, w\cdot\Pi_{1/\eps}\nabla (\Pi_{1/\eps}h)^m\rangle_{L^2(\T^d)} I_{n-1-m}(\otimes^{n-1-m} h)
%\end{equs}
%and the scalar product is $0$ since
%\begin{equs}
%\langle h, w\cdot\Pi_{1/\eps}\nabla (\Pi_{1/\eps}h)^m\rangle_{L^2(\T^d)}&=-\langle w\cdot\nabla \Pi_{1/\eps}h,  (\Pi_{1/\eps}h)^m\rangle_{L^2(\T^d)}\\
%&=-\frac{1}{m+1}\langle w\cdot\nabla (\Pi_{1/\eps}h)^{m+1}, 1 \rangle_{L^2(\T^d)}=0\,.
%\end{equs}
%For the other terms, we want to understand what is the action of $\genaF{m}{}$ on $K$, 
%but to do so, solely for notational convenience, we further distinguish 
%the cases $m$ even or odd. 
%
%\subsection*{$m$ odd}
%
%If $m$ is odd, we write 
%\begin{equs}
%\genaF{m}{} I_n(K)=&\sum_{p=0}^{(m-1)/2-1} I_{n-1+m-2p}\left(\genaF{m}{m-1-2p} K\right) + I_{n}\left(\genaF{m}{0} K\right) \\
%&+ \sum_{p=(m-1)/2+1}^{m-1} I_{n-1+m-2p}\left(\genaF{m}{m-1-2p} K\right)
%\end{equs}
%where $\genaF{m}{m-1-2p} K\in\fock_{n+m-1-2p}$ if $p\leq n-1$ and $0$ otherwise. 
%The previous expression means that, if $\psi\in\fock$, the $\gena\psi$ can be written as 
%\begin{equ}
%\genaF{m}{}  \psi=\genaF{m}{+} \psi +\genaF{m}{-} \psi  + \genaF{m}{0}  \psi 
%\end{equ}
%where further
%\begin{equ}
%\genaF{m}{+} \psi \eqdef \sum_{j=1}^{(m-1)/2} \genaF{m}{2j}   \psi\,,\qquad\text{and}\qquad\genaF{m}{-} \psi \eqdef \sum_{j=1}^{(m-1)/2} \genaF{m}{-2j}  \psi\,.
%\end{equ}
%Then, the operators above can immediately be seen to satisfy the following properties
%\begin{itemize}[noitemsep]
%\item for every $j=1,\dots, (m-1)/2$ and $n\geq (m+1-2j)/2$, 
%	\begin{equ}
%	\genaF{m}{2j} \colon \fock_n\to\fock_{n+2j},,\quad\genaF{m}{-2j} \colon \fock_n\to\fock_{n-2j}\,,\quad\genaF{m}{0} \colon \fock_n\to\fock_{n}
%	\end{equ}
%	so that in particular $\genaF{m}{+} $ increases the Fock space index, 
%	$\genaF{m}{-} $ decreases it, while $\genaF{m}{0} $ leaves it invariant, 
%\item by the previous, orthogonality of Fock spaces with different indices and since $\genaF{m}{}$ is antisymmetric, 
%	for every $j$ we have
%	\begin{equ}[e:AdjointsGen]
%	(\genaF{m}{2j} )^\ast =-\genaF{m}{-2j} \qquad\text{and}\qquad (\genaF{m}{0} )^\ast=-\genaF{m}{0} 
%	\end{equ}
%	which in particular implies that $(\genaF{m}{+} )^\ast =-\genaF{m}{-} $
%\end{itemize}
%We now want to obtain an expression for the operators above. 
%In view of~\eqref{e:AdjointsGen}, it suffices to look at $\genaF{m}{2j}$, $j=0,\dots,(m-1)/2$. 
%
%Let $j=-(m-1)/2,\dots,(m-1)/2$ be as above and $p\eqdef (m-1-2j)/2$. 
%Thanks to~\eqref{e:A} and~\eqref{e:prodI}, the Fourier transform of 
%$\genaF{m}{2j} \psi$, for $\psi\in\fock_n$ with $n\geq (m+1-2j)/2$ is 
%\begin{equs}
%&\genaF{m}{2j} \psi(k_{1:n+2j})=\iota\frac{\eps^{\frac{d}{2}(m-1)-1}}{(2\pi)^{\frac{d}{2}(m-1)}} n p! \binom{n-1}{p}\binom{m}{p}\\
%&\times \sum_{\ell_{1:p}}\indN{\ell_{1:p},k_{1:m-p}}(w\cdot (\overline{\ell_{1:p}}+\overline{k_{1:m-p}})) \hat\psi(-\overline{\ell_{1:p}}-\overline{k_{1:m-p}}, \ell_{1:p}, k_{m-p+1:n+2j})
%\end{equs}
%which still needs to be symmetrised. Before doing so, we want to rewrite the sum in a way which 
%makes it more transparent. 
%Note that changing variables by setting $q\eqdef -(\overline{\ell_{1:p}}+\overline{k_{1:m-p}})$, 
%we get
%\begin{equs}
%&\sum_{\ell_{1:p}}\indN{\ell_{1:p},k_{1:m-p}}(w\cdot (\overline{\ell_{1:p}}+\overline{k_{1:m-p}})) \hat\psi(-\overline{\ell_{1:p}}-\overline{k_{1:m-p}}, \ell_{1:p}, k_{m-p+1:n+2j})\\
%&=-\sum_{q+\overline{\ell_{1:p}}=-\overline{k_{1:m-p}}}\indN{\ell_{1:p},k_{1:m-p}}(w\cdot q) \hat\psi(q, \ell_{1:p}, k_{m-p+1:n+2j})\\
%&=-\sum_{\overline{\ell_{1:p+1}}=-\overline{k_{1:m-p}}}\tindN{\ell_{1:p+1}}\tindN{k_{1:m-p}}(w\cdot \ell_{p+1}) \hat\psi(\ell_{1:p+1}, k_{m-p+1:n+2j})\\
%&=-\frac{1}{p+1}\sum_{i=1}^{p+1}\sum_{\overline{\ell_{1:p+1}}=-\overline{k_{1:m-p}}}\tindN{\ell_{1:p+1}}\tindN{k_{1:m-p}}(w\cdot \ell_i) \hat\psi(\ell_{1:p+1}, k_{m-p+1:n+2j})\\
%&=-\frac{1}{p+1}\sum_{\overline{\ell_{1:p+1}}=-\overline{k_{1:m-p}}}\tindN{\ell_{1:p+1}}\tindN{k_{1:m-p}}(w\cdot \overline{\ell_{1:p+1}}) \hat\psi(\ell_{1:p+1}, k_{m-p+1:n+2j})\\
%&=\tindN{k_{1:m-p}}\frac{(w\cdot \overline{k_{1:m-p}}) }{p+1}\sum_{\overline{\ell_{1:p+1}}=-\overline{k_{1:m-p}}}\tindN{\ell_{1:p+1}} \hat\psi(\ell_{1:p+1}, k_{m-p+1:n+2j})\,,
%\end{equs}
%where to shorten the notation, we wrote
%\begin{equ}[e:tindN]
%\tindN{a_{1:r}}\eqdef \1\{1<|a_i|<\eps^{-1}\colon i=1,\dots,r\}
%\end{equ}
%so that the only difference with respect to $\indN{a_{1:r}}$ is that there is no condition on the sum of the variables $a_{1:r}$. 
%As a consequence, the unsymmetrised version of $\genaF{m}{2j}\psi(k_{1:n+2j})$ is given by 
%\begin{equs}\label{e:A+2junsym}
%\genaF{m}{2j}\psi(k_{1:n+2j})=&\iota c_{p,m}(n)\eps^{\frac{d}{2}(m-1)-1}\tindN{k_{1:m-p}}(w\cdot \overline{k_{1:m-p}})\times\\
%&\times\sum_{\overline{\ell_{1:p+1}}=-\overline{k_{1:m-p}}}\tindN{\ell_{1:p+1}} \hat\psi(\ell_{1:p+1}, k_{m-p+1:n+2j})
%\end{equs}
%where
%\begin{equ}[e:ChaosCoeff]
%c_{p,m}(n)\eqdef \frac{1}{(2\pi)^{\frac{d}{2}(m-1)}}\frac{n}{p+1} p! \binom{n-1}{p}\binom{m}{p}
%\end{equ}
%so that we only need to symmetrise it. Note that the expression above is invariant 
%with respect to permutations of either its first $m-p$ or its last $n-1-p$ variables. 
%Hence, the symmetrised version is 
%\begin{equs}[e:A+2j]
%\genaF{m}{2j}&\psi(k_{1:n+2j})\\
%=&\iota c_{p,m}(n)
%\binom{n+2j}{m-p}^{-1}\eps^{\frac{d}{2}(m-1)-1}\sum_{i_{1:m-p}\in\Delta_{m-p}} \tindN{k_{i_1:i_{m-p}}}(w\cdot \overline{k_{i_1:i_{m-p}}})\\
%&\qquad\qquad\times\sum_{\overline{\ell_{1:p+1}}=-\overline{k_{i_1:i_{m-p}}}}\tindN{\ell_{1:p+1}} \hat\psi(\ell_{1:p+1}, k_{1:n+2j\setminus\{i_1:i_{p+1}\}})\\
%=&\iota c_{p,m}(n)
%\binom{n+2j}{m-p}^{-1}\eps^{\frac{d}{2}(m-1)-1}\sum_{i_{1:m-p}\in\Delta_{m-p}}\tindN{k_{i_1:i_{m-p}}} (w\cdot \overline{k_{i_1:i_{m-p}}})\\
%&\qquad\qquad\times\sum_{\overline{\ell_{1:p+1}}=-\overline{k_{i_1:i_{m-p}}}}\tindN{\ell_{1:p+1}} \hat\psi(\ell_{1:p+1}, k_{1:n+2j\setminus\{i_1:i_{p+1}\}})
%\end{equs}
%where $\Delta_{m-p}$ is the discrete $(n+2j)$-dimensional simplex $\{1\leq i_1<\dots<i_{m-p}\leq n+2j\}$. 
%
%
%\subsection*{$m$ even}
%
%If $m$ is even, we write 
%\begin{equs}
%\genaF{m}{} I_n(K)=&\Big(\sum_{p=0}^{m/2-1}+\sum_{p=m/2+1}^{m-1}\Big) I_{n-1+m-2p}\left(\genaF{m}{m-1-2p} K\right)
%\end{equs}
%where $\genaF{m}{m-1-2p} K\in\fock_{n+m-1-2p}$ if $p\leq n-1$ and $0$ otherwise. 
%The previous expression means that, if $\psi\in\fock$, the $\gena\psi$ can be written as 
%\begin{equ}
%\genaF{m}{}  \psi=\genaF{m}{+} \psi +\genaF{m}{-} \psi  
%\end{equ}
%where further
%\begin{equ}
%\genaF{m}{+} \psi \eqdef \sum_{j=1}^{m/2} \genaF{m}{2j-1}   \psi\,,\qquad\text{and}\qquad\genaF{m}{-} \psi \eqdef \sum_{j=1}^{m/2} \genaF{m}{-(2j-1)}  \psi\,.
%\end{equ}
%Then, the operators above can immediately be seen to satisfy the following properties
%\begin{itemize}[noitemsep]
%\item for every $j=1,\dots, m/2$ and $n\geq (m-2j)/2$, 
%	\begin{equ}
%	\genaF{m}{2j-1} \colon \fock_n\to\fock_{n+2j-1},,\qquad\genaF{m}{-(2j-1)} \colon \fock_n\to\fock_{n-(2j-1)}\,,
%	\end{equ}
%	so that in particular $\genaF{m}{+} $ increases the Fock space index, while
%	$\genaF{m}{-} $ decreases it, 
%\item by the previous, orthogonality of Fock spaces with different indices and since $\genaF{m}{}$ is antisymmetric, 
%	for every $j$ we have
%	\begin{equ}[e:AdjointsGenEven]
%	(\genaF{m}{2j-1} )^\ast =-\genaF{m}{-(2j-1)} 
%	\end{equ}
%	which in particular implies that $(\genaF{m}{+} )^\ast =-\genaF{m}{-} $
%\end{itemize}
%We now want to obtain an expression for the operators above. 
%In view of~\eqref{e:AdjointsGen}, it suffices to look at $\genaF{m}{2j-1}$, $j=1,\dots,m/2$. 
%
%Let $j\in\{1,\dots,m/2\}$ and $p\eqdef (m-2j)/2$. 
%Thanks to~\eqref{e:prodI}, the Fourier transform of 
%$\genaF{m}{2j-1} \psi$, for $\psi\in\fock_n$ with $n\geq (m-2j)/2$ is 
%\begin{equs}
%&\genaF{m}{2j-1} \psi(k_{1:n+2j-1})=\iota\frac{\eps^{\frac{d}{2}(m-1)-1}}{(2\pi)^{\frac{d}{2}(m-1)}} n p! \binom{n-1}{p}\binom{m}{p}\\
%&\times \sum_{\ell_{1:p}}\indN{\ell_{1:p},k_{1:m-p}}(w\cdot (\overline{\ell_{1:p}}+\overline{k_{1:m-p}})) \hat\psi(-\overline{\ell_{1:p}}-\overline{k_{1:m-p}}, \ell_{1:p}, k_{m-p+1:n+2j-1})
%\end{equs}
%which still needs to be symmetrised. Before doing so, we want to rewrite the sum in a way which 
%makes it more transparent. 
%Note that changing variables by setting $q\eqdef -(\overline{\ell_{1:p}}+\overline{k_{1:m-p}})$, 
%we get
%\begin{equs}
%&\sum_{\ell_{1:p}}\indN{\ell_{1:p},k_{1:m-p}}(w\cdot (\overline{\ell_{1:p}}+\overline{k_{1:m-p}})) \hat\psi(-\overline{\ell_{1:p}}-\overline{k_{1:m-p}}, \ell_{1:p}, k_{m-p+1:n+2j-1})\\
%&=-\sum_{q+\overline{\ell_{1:p}}=-\overline{k_{1:m-p}}}\indN{\ell_{1:p},k_{1:m-p}}(w\cdot q) \hat\psi(q, \ell_{1:p}, k_{m-p+1:n+2j-1})\\
%&=-\sum_{\overline{\ell_{1:p+1}}=-\overline{k_{1:m-p}}}\tindN{\ell_{1:p+1}}\tindN{k_{1:m-p}}(w\cdot \ell_1) \hat\psi(\ell_{1:p+1}, k_{m-p+1:n+2j-1})\\
%&=-\frac{1}{p+1}\sum_{i=1}^{p+1}\sum_{\overline{\ell_{1:p+1}}=-\overline{k_{1:m-p}}}\tindN{\ell_{1:p+1}}\tindN{k_{1:m-p}}(w\cdot \ell_i) \hat\psi(\ell_{1:p+1}, k_{m-p+1:n+2j-1})\\
%&=-\frac{1}{p+1}\sum_{\overline{\ell_{1:p+1}}=-\overline{k_{1:m-p}}}\tindN{\ell_{1:p+1}}\tindN{k_{1:m-p}}(w\cdot \overline{\ell_{1:p+1}}) \hat\psi(\ell_{1:p+1}, k_{m-p+1:n+2j-1})\\
%&=\tindN{k_{1:m-p}}\frac{(w\cdot \overline{k_{1:m-p}}) }{p+1}\sum_{\overline{\ell_{1:p+1}}=-\overline{k_{1:m-p}}}\tindN{\ell_{1:p+1}}\hat\psi(\ell_{1:p+1}, k_{m-p+1:n+2j-1})\,.
%\end{equs}
%As a consequence, the unsymmetrised version of $\genaF{m}{2j-1}\psi(k_{1:n+2j-1})$ is given by 
%\begin{equs}\label{e:A+2j-1unsym}
%\genaF{m}{2j-1}\psi(k_{1:n+2j-1})=&\iota c_{p,m}(n)\eps^{\frac{d}{2}(m-1)-1}\tindN{k_{1:m-p}}(w\cdot \overline{k_{1:m-p}})\times\\
%&\times\sum_{\overline{\ell_{1:p+1}}=-\overline{k_{1:m-p}}}\tindN{\ell_{1:p+1}} \hat\psi(\ell_{1:p+1}, k_{m-p+1:n+2j-1})
%\end{equs}
%where $c_{p,m}(n)$ is defined as in~\eqref{e:ChaosCoeff}, 
%so that we only need to symmetrise it. Note that the expression above is invariant 
%with respect to permutations of either its first $m-p$ or its last $n-1-p$ variables. 
%Hence, the symmetrised version is 
%\begin{equs}[e:A+2j-1]
%\genaF{m}{2j-1}&\psi(k_{1:n+2j-1})\\
%=&\iota c_{p,m}(n)
%\binom{n+2j-1}{m-p}^{-1}\eps^{\frac{d}{2}(m-1)-1}\sum_{i_{1:m-p}\in\Delta_{m-p}} \tindN{k_{i_1:i_{m-p}}}(w\cdot \overline{k_{i_1:i_{m-p}}})\\
%&\qquad\qquad\times\sum_{\overline{\ell_{1:p+1}}=-\overline{k_{i_1:i_{m-p}}}}\tindN{\ell_{1:p+1}} \hat\psi(\ell_{1:p+1}, k_{1:n+2j-1\setminus\{i_1:i_{p+1}\}})\\
%=&\iota c_{p,m}(n)
%\binom{n+2j-1}{m-p}^{-1}\eps^{\frac{d}{2}(m-1)-1}\sum_{i_{1:m-p}\in\Delta_{m-p}} \tindN{k_{i_1:i_{m-p}}}(w\cdot \overline{k_{i_1:i_{m-p}}})\\
%&\qquad\qquad\times\sum_{\overline{\ell_{1:p+1}}=-\overline{k_{i_1:i_{m-p}}}}\indN{\ell_{1:p+1}} \hat\psi(\ell_{1:p+1}, k_{1:n+2j-1\setminus\{i_1:i_{p+1}\}})
%\end{equs}
%where $\Delta_{m-p}$ is the discrete $(n+2j-1)$-dimensional simplex $\{1\leq i_1<\dots<i_{m-p}\leq n+2j\}$. 
%
%
%\subsection*{Crucial estimates}
%
%Based on the explicit expression in the previous section, we now want to obtain 
%those basic estimates which will be crucial in the analysis. 
%
%\begin{proposition}\label{p:MainEstimate}
%For $m$ odd, let $a=2j$, $j=-(m-1)/2,\dots, (m-1)/2$ and $n\geq (m+1-2j)/2$, while 
%for $m$ even, let $a=2j-1$ for $j=1,\dots,m/2$ and $a=2j+1$ for $j=-1,\dots,-m/2$, and $n\geq (m-2j)/2$. 
%Then, for every $\psi\in\fock_n$, we have 
%\begin{equ}[e:ImpBoundGen]
%\|(-\gensy)^{-1/2}\genaF{m}{a}\psi\|\lesssim \|\cN^{\frac{m-1}{2}}(-\gensy)^{1/2}\psi\|\,.
%\end{equ}
%\end{proposition}
%
%The proof of the previous proposition is based on the following lemma. 
%
%%\begin{lemma}
%%For any $j=-(m-1)/2,\dots,(m-1)/2$ and $\psi\in\fock$, the following bounds hold 
%%\begin{equ}
%%\|(-\gensy)^{-1/2}\genaF{m}{2j}\psi\|^2\lesssim \|\cN^{\frac{m}{2}}(-\gensy)^{1/2}\psi\|
%%\end{equ}
%%\end{lemma}
%%\begin{proof}
%%Let $\psi\in\fock_n$, set $p=(m-1-2j)/2$ and assume $n\geq p-1$ (if not then $\genaF{m}{2j}\psi=0$). 
%%We consider the following 
%%\begin{equs}
%%\|&(-\gensy)^{-1/2}\genaF{m}{2j}\psi\|^2=(n+2j)!\sum_{k_{1:n+2j}}\frac{1}{|k_{1:n+2j}|^2} |\genaF{m}{2j}\psi(k_{1:n+2j})|^2\\
%%&=\frac{1}{(2\pi)^{d(m-1)}}(n+2j)!\Big(\frac{n! m!}{(p+1)!(n+2j)!}\Big)^2\times\\
%%&\qquad\times \eps^{d(m-1)-2}\sum_{k_{1:n+2j}}\frac{1}{|k_{1:n+2j}|^2}\times\Big|\sum_{i_{1:m-p}\in\Delta_{m-p}}\frac{\binom{n+2j}{m-p}}{\binom{n+2j}{m-p}} (w\cdot \overline{k_{i_1:i_{m-p}}})\times\\
%%&\qquad\qquad\qquad\times 
%%\sum_{\overline{\ell_{1:p+1}}=-\overline{k_{i_1:i_{m-p}}}}\indN{\ell_{1:p+1},k_{i_{1:m-p}}} \hat\psi(\ell_{1:p+1}, k_{1:n+2j\setminus\{i_1:i_{p+1}\}})\Big|^2\\
%%&\leq (n+2j)!\Big(\frac{n! m!}{(p+1)!(n+2j)!}\Big)^2\binom{n+2j}{m-p}^2\times\\ 
%%&\,\times\eps^{d(m-1)-2}\sum_{k_{1:n+2j}}\frac{(w\cdot \overline{k_{1:m-p}})^2}{|k_{1:n+2j}|^2} \Big|\sum_{\overline{\ell_{1:p+1}}=-\overline{k_{i_1:i_{m-p}}}}\indN{\ell_{1:p+1},k_{1:m-p}} \hat\psi(\ell_{1:p+1}, k_{m-p+1:n+2j})\Big|^2\\
%%&\leq (n+2j)!\Big(\frac{n! m!}{(p+1)!(m-p)!(n-1-p)!}\Big)^2\times\label{e:A+GenBound}\\ 
%%&\,\times\eps^{d(m-1)-2}\sum_{k_{1:n+2j}}\frac{(w\cdot \overline{k_{1:m-p}})^2}{|k_{1:n+2j}|^2} \Big|\sum_{\overline{\ell_{1:p+1}}=-\overline{k_{i_1:i_{m-p}}}}\indN{\ell_{1:p+1},k_{1:m-p}} \hat\psi(\ell_{1:p+1}, k_{m-p+1:n+2j})\Big|^2\,.
%%\end{equs}
%%Let us have a closer look at the sum. By a simple change of variables it equals
%%\begin{equs}
%%&\eps^{dp} \sum_{q, k_{m-p+1:n+2j}}(w\cdot q)^2\Big|\sum_{\overline{\ell_{1:p+1}}=-q}\indN{\ell_{1:p+1}} \hat\psi(\ell_{1:p+1}, k_{m-p+1:n+2j})\Big|^2\times\\
%%&\qquad\qquad\qquad\qquad\qquad\qquad\qquad\times \eps^{d(m-p-1)-2} \sum_{\overline{k_{1:m-p}}=q}\frac{\indN{k_{1:m-p}} }{|k_{1:n+2j}|^2}\\
%%&\lesssim \eps^{dp} \sum_{q, k_{m-p+1:n+2j}}(w\cdot q)^2\Big|\sum_{\overline{\ell_{1:p+1}}=-q}\indN{\ell_{1:p+1}} \hat\psi(\ell_{1:p+1}, k_{m-p+1:n+2j})\Big|^2
%%\end{equs}
%%where we bounded the second sum appearing at the left hand side as 
%%\begin{equs}[e:MidBound]
%%\eps^{d(m-p-1)-2} \sum_{\overline{k_{1:m-p}}=q}\frac{\indN{k_{1:m-p}} }{|k_{1:n+2j}|^2}&=\eps^{d(m-p-1)}\sum_{\overline{\eps k_{1:m-p}}=\eps q}\frac{\indN{k_{1:m-p}} }{|\eps k_{1:n+2j}|^2}\\
%%&\lesssim \int_{[-1,1]^{d(m-p-1)}}\frac{\dd x_{1:m-p-1}}{|x_{1:m-p-1}|^2}
%%\end{equs}
%%and the latter integral is bounded uniformly over $q$, $\eps$ and $n$ (but clearly depends on $m$ and $p$).
%%Getting back to the remaining sum, we apply Cauchy-Schwarx and obtain 
%%\begin{equs}
%%\eps^{dp} &\sum_{q, k_{m-p+1:n+2j}}(w\cdot q)^2\Big|\sum_{\overline{\ell_{1:p+1}}=-q}\indN{\ell_{1:p+1}} \hat\psi(\ell_{1:p+1}, k_{m-p+1:n+2j})\Big|^2\\
%%&\lesssim \sum_{q, k_{m-p+1:n+2j}}(w\cdot q)^2 \sum_{\overline{\ell_{1:p+1}}=-q}|\hat\psi(\ell_{1:p+1}, k_{m-p+1:n+2j})|^2\times\\
%%&\qquad\qquad\qquad\qquad\qquad\qquad\qquad\qquad\qquad\times\Big(\eps^{dp}\sum_{\overline{\ell_{1:p+1}}=-q} \indN{\ell_{1:p+1}} \Big)\\
%%&\lesssim \sum_{k_{1:n+2j}}(w\cdot \overline{k_{1:p+1}})^2 |\hat\psi( k_{1:n+2j})|^2\lesssim \sum_{k_{1:n+2j}} |\overline{k_{1:p+1}}|^2 |\hat\psi( k_{1:n+2j})|^2\\
%%&\lesssim \frac{1}{n+2j}\sum_{k_{1:n+2j}} |k_{1:n+2j}|^2 |\hat\psi( k_{1:n+2j})|^2\leq \frac{1}{n}\sum_{k_{1:n+2j}} |k_{1:n+2j}|^2 |\hat\psi( k_{1:n+2j})|^2 \,,
%%\end{equs}
%%where in the last step we used the symmetry of $\psi$. 
%%Since further we have 
%%\begin{equs}
%%\frac{ n! m!}{(p+1)!(m-p)!(n-1-p)!}=\frac{1}{(p+1)}\binom{m}{p} \frac{n!}{(n-p-1)!}\lesssim_{m,p} n^{p+1}
%%\end{equs}
%%where we omitted the dependence on $m$ and $p$ since they are fixed, 
%%and estimating $(n+2j)!\lesssim n! n^{2j}$, 
%%we conclude that 
%%\begin{equs}
%%\|(-\gensy)^{-1/2}\gena_{+,2j}\psi\|^2&\lesssim n! \frac{n^{2j} n^{2p+2}}{n}\sum_{k_{1:n+2j}} \frac12 |k_{1:n+2j}|^2 |\hat\psi( k_{1:n+2j})|^2\\
%%&\lesssim n! n^m\sum_{k_{1:n+2j}} \frac12 |k_{1:n+2j}|^2 |\hat\psi( k_{1:n+2j})|^2=\|\cN^{\frac{m}{2}}(-\gensy)^{1/2}\psi\|^2
%%\end{equs}
%%Since the expression for all the other operators is similar, the same techniques as above provide 
%%the desired bounds.
%%\end{proof}
%
%
%\begin{lemma}
%In the same setting as Proposition~\ref{p:MainEstimate}, let $j\geq 0$, $\phi\in\fock_{n+a}$ and $\psi\in\fock_n$. 
%Then for every $\gamma>0$
%\begin{equ}
%|\langle\phi,\genaF{m}{a} \psi\rangle|\leq Cn^{\frac{m-1}{2}}\Big(\gamma \|(-\gensy)^{1/2}\psi\|^2 + \frac{1}{\gamma}\|(-\gensy)^{1/2}\phi\|^2\Big)\,.
%\end{equ}
%\end{lemma}
%\begin{proof}
%By definition of the scalar product in $\fock$ and of the operator $\genaF{m}{2j}$ in~\eqref{e:A+2j},~\eqref{e:A+2j-1}, we have 
%\begin{equs}
%|\langle\phi,\genaF{m}{a} \psi\rangle|=&(n+a)! c_{p,m}(n)\binom{n+a}{m-p}^{-1}  \eps^{\frac{d}{2}(m-1)-1}\Big|\sum_{k_{1:n+a}}\hat\phi(-k_{1:n+a})\times\\
%&\times\sum_{i_{1:m-p}\in\Delta_{m-p}} \tindN{k_{i_1:i_{m-p}}}(w\cdot \overline{k_{i_1:i_{m-p}}})\sum_{\overline{\ell_{1:p+1}}=-\overline{k_{i_1:i_{m-p}}}}\tindN{\ell_{1:p+1}} \hat\psi\Big|%(\ell_{1:p+1}, k_{1:n+2j\setminus\{i_1:i_{p+1}\}})\Big|\,.
%\\
%=&(n+a)! c_{p,m}(n)\eps^{\frac{d}{2}(m-1)-1}\Big|\sum_{k_{1:n+a}}\hat\phi(-k_{1:n+a})(w\cdot \overline{k_{1:m-p}})\tindN{k_{1:m-p}}\\
%&\times\sum_{\overline{\ell_{1:p+1}}=-\overline{k_{1:m-p}}}\tindN{\ell_{1:p+1}} \hat\psi(\ell_{1:p+1}, k_{m-p+1:n+a})\Big|
%\end{equs}
%Let us look at the last two lines and in particular, at the sum over $k_{1:m-p}$. This equals
%\begin{equs}
%&\Big|\sum_{k_{1:m-p}}\hat\phi(-k_{1:n+a})(w\cdot \overline{k_{1:m-p}})\tindN{k_{1:m-p}}\\
%&\qquad\times\sum_{\overline{\ell_{1:p+1}}=-\overline{k_{1:m-p}}}\tindN{\ell_{1:p+1}} \hat\psi(\ell_{1:p+1}, k_{m-p+1:n+a})\Big|\\
%&=\Big|\sum_{q} (w\cdot q)\sum_{\overline{\ell_{1:p+1}}=-q}\tindN{\ell_{1:p+1}} \hat\psi(\ell_{1:p+1}, k_{m-p+1:n+a})\\
%&\qquad\times\sum_{\overline{k_{1:m-p}}=q}\tindN{k_{1:m-p}}\hat\phi(-k_{1:n+a})\Big|\\
%&\lesssim \Big(\sum_q (w\cdot q)^2\Big|\sum_{\overline{\ell_{1:p+1}}=-q}\tindN{\ell_{1:p+1}} \hat\psi(\ell_{1:p+1}, k_{m-p+1:n+a})\Big|^2\Big)^{1/2}\times\\
%&\qquad\times\Big(\sum_q \Big|\sum_{\overline{k_{1:m-p}}=q}\tindN{k_{1:m-p}}\hat\phi(-k_{1:n+a})\Big|^2\Big)^{1/2}=:\one\times\two
%\end{equs}
%Now, notice that the two factors at the right hand side can be respectively bounded as 
%\begin{equs}[e:oneGen]
%\one&=\Big(\sum_q (w\cdot q)^2\Big|\sum_{\overline{\ell_{1:p+1}}=-q}\tindN{\ell_{1:p+1}} \hat\psi(\ell_{1:p+1}, k_{m-p+1:n+a})\Big|^2\Big)^{1/2}\\
%&\lesssim \Big(\sum_q |q|^2\sum_{\overline{\ell_{1:p+1}}=-q} |\hat\psi(\ell_{1:p+1}, k_{m-p+1:n+a})|^2\sum_{\overline{\ell_{1:p+1}}=-q}\tindN{\ell_{1:p+1}} \Big)^{1/2}\\
%&=\eps^{-\frac{pd}{2}}\Big(\sum_{\ell_{1:p+1}} |\overline{\ell_{1:p+1}}|^2 |\hat\psi(\ell_{1:p+1}, k_{m-p+1:n+a})|^2 \Big)^{1/2}=:\eps^{-\frac{pd}{2}} \sqrt{A}
%\end{equs}
%and 
%\begin{equs}[e:twoGen]
%\two&=\Big(\sum_q \Big|\sum_{\overline{k_{1:m-p}}=q}\tindN{k_{1:m-p}}\hat\phi(-k_{1:n+a})\Big|^2\Big)^{1/2}\\
%&\leq \Big(\sum_q \sum_{\overline{k_{1:m-p}}=q}|k_{1:m-p}|^2|\hat\phi(-k_{1:n+a})|^2\sum_{\overline{k_{1:m-p}}=q}\frac{\tindN{k_{1:m-p}}}{|k_{1:m-p}|^2}\Big)^{1/2}\\
%&\lesssim\eps^{-\frac{d(m-p-1)}{2}+1} \Big(\sum_{k_{1:m-p}}|k_{1:m-p}|^2|\hat\phi(-k_{1:n+a})|^2\Big)^{1/2}=:\eps^{-\frac{d(m-p-1)}{2}+1} B
%\end{equs}
%where we used that 
%\begin{equs}[e:MidBound]
%\eps^{d(m-p-1)-2}\sum_{\overline{k_{1:m-p}}=q}\frac{\tindN{k_{1:m-p}}}{|k_{1:m-p}|^2}&=\eps^{d(m-p-1)}\sum_{\overline{\eps k_{1:m-p}}=\eps q}\frac{\tindN{k_{1:m-p}} }{|\eps k_{1:m-p}|^2} \\
%&\lesssim \int_{[-1,1]^{d(m-p-1)}}\frac{\dd x_{1:m-p-1}}{|x_{1:m-p-1}|^2}\lesssim_{m,p} 1\,.
%\end{equs}
%As a consequence, 
%\begin{equs}[e:AlmostDoneBound]
%|\langle&\phi,\genaF{m}{a} \psi\rangle|\lesssim (n+a)! c_{p,m}(n)\sum_{k_{m-p+1:n+a}}\sqrt{AB}\\
%&\leq (n+a)! c_{p,m}(n)\left(\frac{\gamma}{n^\alpha} \sum_{k_{m-p+1:n+a}} A +\frac{n^\alpha}{\gamma}\sum_{k_{m-p+1:n+a}} B\right)\\
%&\leq (n+a)!c_{p,m}(n)\Big(\frac{1}{n!}\frac{\gamma}{n^{\alpha+1}}\|(-\gensy)^{1/2}\psi\|^2 +\frac{1}{(n+a)!}\frac{n^{\alpha-1}}{\gamma}\|(-\gensy)^{1/2}\phi\|^2\Big)
%\end{equs}
%where in the last step we used that, by the definition of $A$ and $B$ in~\eqref{e:oneGen} and~\eqref{e:twoGen} respectively, 
%we have 
%\begin{equs}
%&\sum_{k_{m-p+1:n+a}} A=\sum_{\ell_{1:p+1}, k_{m-p+1:n+a}}  |\overline{\ell_{1:p+1}}|^2 |\hat\psi(\ell_{1:p+1}, k_{m-p+1:n+a})|^2\\
%&=\sum_{k_{1:n}}|\overline{k_{1:p+1}}|^2 |\hat\psi(k_{1:n})|^2\lesssim\frac{1}{n}\sum_{k_{1:n}}\frac12 |k_{1:n} |\hat\psi(k_{1:n})|^2=\frac{n^{-1}}{n!} \|(-\gensy)^{1/2}\psi\|^2
%\end{equs}
%and similarly
%\begin{equs}
%\sum_{k_{m-p+1:n+a}} B=\sum_{k_{1:n+a}} |k_{1:m-p}|^2|\hat\phi(k_{1:n+a})|^2\lesssim \frac{(n+a)^{-1}}{(n+a)!} \|(-\gensy)^{1/2}\phi\|^2\,.
%\end{equs}
%To conclude, notice that by~\eqref{e:ChaosCoeff}, we have 
%\begin{equ}
%c_{p,m}(n)= \frac{1}{(2\pi)^{\frac{d}{2}(m-1)}}\frac{n}{p+1} p! \binom{n-1}{p}\binom{m}{p}\lesssim_{m,p} n^{p+1}
%\end{equ}
%and $(n+a)!/n!\lesssim n^{a}$, so that~\eqref{e:AlmostDoneBound} gives 
%\begin{equ}
%|\langle\phi,\genaF{m}{a} \psi\rangle|\lesssim n^{a+p-\alpha}\gamma\|(-\gensy)^{1/2}\psi\|^2+\frac{n^{p+\alpha}}{\gamma} \|(-\gensy)^{1/2}\phi\|^2
%\end{equ}
%which, by recalling that $p=(m-2j)/2=(m-1-a)/2$ if $m$ is odd and $p=(m-2j)/2=(m-1-a)/2$ if $m$ even, 
%and choosing $\alpha=a/2$ proves the lemma. 
%\end{proof}
%
%\begin{proof}[of Proposition~\ref{p:MainEstimate}]
%We begin by considering the case of $a$ non-negative. 
%The variational characterisation of the $(-\gensy)^{-1/2}\fock$-norm gives us that 
%\begin{equs}
%\|&(-\gensy)^{-1/2}\genaF{m}{a}\psi\|^2=\sup_{\phi\in \fock_{n+a}}\left(2\langle \phi,\genaF{m}{a}  \psi\rangle-
%    \| (-\cL_0)^{\frac12}\phi\|^2\right)\\
%&\leq \sup_{\phi\in \fock_{n+a}}\left(2Cn^{\frac{m-1}{2}}\Big(\gamma \|(-\gensy)^{1/2}\psi\|^2 + \frac{1}{\gamma}\|(-\gensy)^{1/2}\phi\|^2\Big)-
%    \| (-\cL_0)^{\frac12}\phi\|^2\right)\\
%    &\lesssim n^{m-1}\|(-\gensy)^{1/2}\psi\|^2=\|\cN^{\frac{m-1}{2}}(-\gensy)^{1/2}\psi\|^2
%\end{equs}
%where in the first step we used the previous lemma while in the last we chose $\gamma\eqdef 2Cn^{\frac{m-1}{2}}$. 
%
%If instead $a$ is negative, we use the skew-symmetric relations between $\genaF{m}{a}$ and $\genaF{m}{-a}$. 
%Arguing as above, we have 
%\begin{equs}
%\|&(-\gensy)^{-1/2}\genaF{m}{a}\psi\|^2=\sup_{\phi\in \fock_{n+a}}\left(2\langle \phi,\genaF{m}{a}  \psi\rangle-
%    \| (-\cL_0)^{\frac12}\phi\|^2\right)\\
% &=\sup_{\phi\in \fock_{n+a}}\left(-2\langle \genaF{m}{-a}\phi, \psi\rangle-
%    \| (-\cL_0)^{\frac12}\phi\|^2\right)\\
%&\leq \sup_{\phi\in \fock_{n+a}}\left(2Cn^{\frac{m-1}{2}}\Big(\gamma \|(-\gensy)^{1/2}\phi\|^2 + \frac{1}{\gamma}\|(-\gensy)^{1/2}\psi\|^2\Big)-
%    \| (-\cL_0)^{\frac12}\phi\|^2\right)\\
%    &\lesssim n^{m-1}\|(-\gensy)^{1/2}\psi\|^2=\|\cN^{\frac{m-1}{2}}(-\gensy)^{1/2}\psi\|^2
%\end{equs}
%where this time we chose $\gamma\eqdef 1/(2Cn^{\frac{m-1}{2}})$. 
%\end{proof}
%
%\subsection{A priori estimates: $m$ odd}
%
%Let $n>m$ and $P^n_2$ be the projection onto $n$-th inhomogeneous Fock space with the first 
%chaos removed, i.e. $\bigoplus_{j=2}^n\Gamma L^2_j$. 
%Let us introduce the truncated generator $\genn$, and the corresponding truncated operator $\genaF{m}{2,n}$, 
%which are respectively defined as 
%\begin{equ}[e:gennGen]
%\genF{m}{2,n}\eqdef \gensy +\genaF{m}{2,n}\qquad \text{and}\qquad \genaF{m}{2,n}\eqdef P_2^n\genaF{m}{} P_2^n\,.
%\end{equ}
%The goal of this section is to derive suitable a priori estimates {\color{red}(or at least see where they fail)} 
%on the solution of the {\it truncated generator equation} which is defined as 
%\begin{equation}\label{eq:TruncGenEqGen}
%-\genF{m}{2,n} u^{\eps,n}=\psi,
%\end{equation}
%for a generic $\psi\in\bigoplus_{j=2}^b\Gamma L^2_j$, $2\leq b\le n$.
%
%\begin{proposition}\label{P:AprioriGeneratorEqGen}
%Let $b\in\N$ and $\psi\in\bigoplus_{j=1}^b\Gamma L^2_j$. For $n\in\N$, $n\geq b$, let $u^{\eps,n}$ 
%be the solution of the truncated generator equation in~\eqref{eq:TruncGenEqGen}. 
%Then, there exists a 
%positive constant $C=C(b,m,k)$ independent of $n,\eps$ and $\psi$, such that
%\begin{equation}\label{e:AprioriMainGen}
%\|\mathcal N^k (-\gensy)^\frac12 u^{\eps,n}\|\leq C\|(-\gensy)^{-\frac12}\psi\|\,,
%\end{equation}
%where $\cN$ is the number operator in Definition~\ref{def:NoOp}. 
%\end{proposition}
%\begin{proof}
%As $n$ and $\eps$ are fixed throughout the proof, we will write $u$ and $u_i,\,i=2,\dots,n$ in place of 
%$u^{\eps,n}$ and $u^{\eps,n}_i$, respectively. 
%Also, by convention, we let $u_{n+a}=0=u_{-a}$ for all $a$ even and less than $(m-1)/2$ if $m$ 
%is odd, or $a$ odd and less than $m/2$ if $m$ even. 
%To denote constants which do not depend on $\eps,n$ or $\psi$ 
%we will use $C$, and such $C$ might change from line to line. 
%
%As a first observation, we notice that, by orthogonality of Fock spaces with different indices, 
%the first component of $u$ is independent of the others, in the sense that on the one hand it solves the equation 
%\begin{equ}[e:u1New]
%-\gensy u_1=\psi_1
%\end{equ}
%and on the other no other component of $u$ depends on $u_1$. 
%Further, a standard argument shows that~\eqref{e:u1} immediately implies 
%that~\eqref{e:AprioriMain} holds for $u_1$, hence we only need to consider $u_i$ for $i\geq 2$. 
%
%Let us test the $i$-th component, $i\geq 2$, with $\genF{m}{} u$, which gives
%  \begin{equs}
%    \langle u_i,\genF{m}{} u\rangle&= \langle u_i, \cL_0 u_i\rangle+\langle u_i,\genaF{m}{+} u\rangle+\langle u_i,\genaF{m}{-} u\rangle\\
%    &= \langle u_i, \cL_0 u_i\rangle+\sum_{a}\Big[\langle u_i,\genaF{m}{a} u_{i-a}\rangle+\langle u_i,\genaF{m}{-a} u_{i+a}\rangle\Big]\\
%    &=  \langle u_i, \cL_0 u_i\rangle+\sum_{a}\Big[\langle u_i,\genaF{m}{a} u_{i-a}\rangle-\langle u_{i+a}, \genaF{m}{a} u_i\rangle\Big]\,,
%  \end{equs}
%where, as above, the sum in $a$ ranges over positive even integers less than $(m-1)/2$ if $m$ 
%is odd, or positive odd integers less than $m/2$ if instead $m$ is even. 
%In the former case, we dropped the term containing $\genaF{m}{0}$, as the operator is antisymmetric, so that 
%$\langle u_i,\genaF{m}{0} u_i\rangle=0$. 
%
%Let $f\colon \R_+\to\R$ be a positive increasing function, whose first derivative is strictly increasing
%Via summation by parts, we have 
%\begin{equs}
%    \sum_{i=2}^n f(i)\langle u_i,(-\cL_0) u_i\rangle=  &\sum_{i=2}^n f(i)\langle u_i,(-\genF{m}{}) u_i\rangle\\
%    &+\sum_{i=2}^n \sum_{a} (f(i)-f(i-a))\langle u_i,\genaF{m}{a} u_{i-a}\rangle.
%  \end{equs}
%Next note  by orthogonality of Fock spaces with different indices,
%  \begin{equ}
%    |\langle u_i,(-\genF{m}{}) u_i\rangle|= |\langle u_i,\psi_i\rangle|\le \frac12 \langle u_i,(-\cL_0)u_i\rangle+\frac12\langle \psi_i,(-\cL_0)^{-1}\psi_i\rangle.
%  \end{equ}
%As a consequence, we obtain
%  \begin{equs}[eq:freddoNew]
%\sum_{i=2}^n& f(i)\langle u_i,(-\cL_0) u_i\rangle\\
%&\le \|f(\cN)^{\frac12}(-\cL_0)^{-\frac12}\psi\|^2+2\sum_{i=2}^n \sum_{a} (f(i)-f(i-a))\langle u_i,\genaF{m}{a} u_{i-a}\rangle\,.
%\end{equs}
%At this point, we apply~\eqref{e:ImpBoundGen} to the last summand, which gives
%\begin{equs}
%\langle u_i,\genaF{m}{a} u_{i-2j}\rangle&\leq \frac{1}{2\gamma}\|(-\gensy)^{1/2} u_i\|^2+\frac{\gamma}2\|(-\gensy)^{-1/2} \genaF{m}{a} u_{i-a}\|^2\\
%&\leq \frac{1}{2\gamma}\|(-\gensy)^{1/2} u_i\|^2+\frac{\gamma}2 C(i-a)^{m-1}\|(-\gensy)^{1/2} u_{i-2j}\|^2
%\end{equs}
%then we choose $\gamma=C^{-1/2}(i-a)^{\frac{1-m}{2}}$ so that 
%\begin{equs}
%\langle u_i,\genaF{m}{a} u_{i-a}\rangle&\leq \frac{\sqrt{C}}{2}(i-a)^{\frac{m-1}{2}}\Big(\|(-\gensy)^{1/2} u_i\|^2+\|(-\gensy)^{1/2} u_{i-a}\|^2\Big)\,.
%\end{equs}
%From the previous, we see that 
%\begin{equs}
%2&\sum_{i=2}^n \sum_{a} (f(i)-f(i-a))\langle u_i,\genaF{m}{a} u_{i-a}\rangle\\
%&\leq \sqrt{C}\sum_{i=2}^n \sum_{a}(f(i)-f(i-a))(i-a)^{\frac{m-1}{2}}\Big(\|(-\gensy)^{1/2} u_i\|^2+\|(-\gensy)^{1/2} u_{i-a}\|^2\Big)\\
%&\leq \sqrt{C}\sum_{i=2}^n \sum_{a} f'(i) a (i-a)^{\frac{m-1}{2}}\Big(\|(-\gensy)^{1/2} u_i\|^2+\|(-\gensy)^{1/2} u_{i-a}\|^2\Big)\\
%&\leq \sqrt{C}m \sum_{i=2}^n \sum_{a} f'(i)  (i-a)^{\frac{m-1}{2}}\Big(\|(-\gensy)^{1/2} u_i\|^2+\|(-\gensy)^{1/2} u_{i-a}\|^2\Big)\\
%&=\sqrt{C}m \sum_{i=2}^n \sum_{a} \Big[f'(i)  (i-a)^{\frac{m-1}{2}} +f'(i+a)  i^{\frac{m-1}{2}}\Big] \|(-\gensy)^{1/2} u_i\|^2
%\end{equs}
%
%
%
%
%\end{proof}
%
%
%
%
%
%
%
%
%
%
%
%
%
%\newpage
%
%\section{OLD STUFF (still useful for comparison)}
%\subsubsection{$F$ Hermite polynomial of degree $3$}
%
%To try to understand what is going on, we take $F$ to be a Hermite polynomial of 
%degree $3$. In particular,~\eqref{e:Fexpeta} becomes
%\begin{equ}
%F(\eps^{d/2}\Pi_{1/\eps}\eta^\eps(x))= \eps^{\frac{3}{2}d}\sum_{k_{1:3}\colon |k_i|\leq \eps^{-1}}\frac{e_{\overline{k_{1:3}}}(x)}{(2\pi)^{d}}\Wick{\hat\eta(k_1)\hat\eta(k_2)\hat\eta(k_3)}\,,
%\end{equ}
%where, once again, $\overline{k_{1:j}}\eqdef \sum_{i=1}^jk_i$.  
%The non-linearity then reads
%\begin{equs}
%\NN_k(\eta)&= \eps^{-\frac{d}{2}-1} w\cdot\Pi_{1/\eps}\nabla F(\eps^{d/2}\Pi_{1/\eps}\eta^\eps) (e_{-k})\\
%&=\frac{\iota}{(2\pi)^{d}}\eps^{d-1}\sum_{\overline{k_{1:3}}=k} \indN{k_1,\dots,k_3} w\cdot \overline{k_{1:3}}\Wick{\hat\eta(k_1)\hat\eta(k_2)\hat\eta(k_3)}\\
%&=\frac{\iota}{(2\pi)^{d}}\eps^{d-1}I_3\left(\sum_{\overline{k_{1:3}}=k} \indN{k_1,\dots,k_3} (w\cdot \overline{k_{1:3}})e_{-k_1}\otimes e_{-k_2}\otimes e_{-k_3}\right)
%\end{equs}
%
%By polarisation, it suffices to study the action of $\gena$ on $I_n(K)$ for $K=\otimes^n h$, and $h\in L^2(\T^d)$. 
%Recall that 
%\begin{equ}
%D_{-k} I_n(K)=n I_{n-1}(\otimes^{n-1} h) \hat h(-k)\,,
%\end{equ}
%and the product rule for the integral operators, namely for $f,g$ symmetric functions in $L^2(\T^{md})$ and $L^2(\T^{nd})$ respectively, 
%we have 
%\begin{equ}[e:Mult]
%\sint_n(f) \sint_m(g)  = \sum_{p=0}^{m\land n} p! \binom{n}{p}\binom{m}{p} 
%\sint_{m+n-2p} (f \otimes_p g),
%\end{equ} 
%where
%\[
%(f \otimes_p g)(x_{1: m+n-2p})\eqdef 
%\int_{\T^{pd}}\dd y_{1:p} f(x_{1:n-p}, y_{1:p})
%g(x_{n-p+1:m+n-2p}, y_{1:p})\,.
%\]
%\medskip
%
%We have now collected all the elements we need in order to compute $\gena I_n(K)$, which is 
%\begin{equs}
%\,&\gena I_n(K)=\sum_k \NN_k(\eta) D_{-k} I_n(K)=n I_{n-1}(\otimes^{n-1} h)\sum_k \NN_k(\eta) \hat h(-k)\\
%&=\frac{\iota \eps^{d-1}}{(2\pi)^{d}} n I_{n-1}(\otimes^{n-1} h)I_3\left(\sum_{\overline{k_{1:3}}=k} \indN{k_1,\dots,k_3} (w\cdot \overline{k_{1:3}})e_{-k_1}\otimes e_{-k_2}\otimes e_{-k_3}\right)\\
%&=\frac{\iota \eps^{d-1}}{(2\pi)^{d}} n\sum_{p=0}^{ 3} p! \binom{n-1}{p}\binom{3}{p} \1_{p\leq n-1}\times\\
%&\qquad\times I_{n+2-2p}\bigg(\sum_{k_{1:3}} \indN{k_1,\dots,k_3} w\cdot \overline{k_{1:3}} \hat h(-\overline{k_{1:3}}) \left(\prod_{i=0}^p \hat h(k_i)\right) \times\\
%&\qquad\qquad\qquad\times e_{-k_{p+1}}\otimes\dots\otimes e_{-k_3}\otimes \otimes^{n-1-p} h\bigg)\\
%&=: \1_{0\leq n-1} I_{n+2}\left(\gena_{++}K\right) +\1_{1\leq n-1} I_{n}\left(\gena_{0}K\right) + \1_{2\leq n-1} I_{n-2}\left(\gena_{--}K\right)
%%&\qquad\qquad+\sum_{p=\lfloor\frac{j-1}{2}\rfloor+1}^{j-1}\1_{p\leq n-1} I_{n-1+j-2p}\left(\gena_{-, j-1-2p}K\Big)
%\end{equs}
%where we adopted the convention that $\hat h(k_0)\eqdef 1$ and $\gena_{\cdot, j-1-2p} K$ is the symmetrised 
%kernels of $I_{n-1+j-2p}$. Further, we used the fact that the summand corresponding to $p=3$ is $0$. 
%Indeed, 
%\begin{equs}
%I_{n-4}&\left(\sum_{k_{1:3}} \indN{k_1,\dots,k_3} w\cdot \overline{k_{1:3}} \hat h(-\overline{k_{1:3}}) \left(\prod_{i=1}^3 \hat h(k_i)\right)  \otimes^{n-4} h\right)\\
%&=\left[\sum_{k_{1:3}} \indN{k_1,\dots,k_3} w\cdot \overline{k_{1:3}} \hat h(-\overline{k_{1:3}}) \left(\prod_{i=1}^3 \hat h(k_i)\right) \right] I_{n-4}\left(\otimes^{n-4}h\right)\\
%&=-\langle w\cdot \nabla h, h^3\rangle I_{n-4}\left(\otimes^{n-4}h\right)=\frac{1}{4}\langle w\cdot \nabla (h^4), 1\rangle I_{n-4}\left(\otimes^{n-4}h\right)=0\,.
%\end{equs}
%We can now take a closed look at the operators $\gena_{++},\gena_0,\gena_{--}$. 
%First, since $\gena$ is antisymmetric, we must have 
%\begin{equ}
%(\gena_{++})^\ast=-\gena_{--}\qquad\text{and}\qquad (\gena_0)^\ast=-\gena_0\,.
%\end{equ}
%Hence, it is sufficient to study $\gena_{++}$ and $\gena_0$. Notice that, by the computation above, 
%the unsymmetrised version of $\gena_{++}$ acts on $\psi\in\fock_n$ as 
%\begin{equ}
%\gena_{++} \psi (k_{1:n+2}) = \frac{\iota \eps^{d-1}}{(2\pi)^{d}} n \indN{k_1,\dots,k_3} (w\cdot \overline{k_{1:3}}) \hat \psi(\overline{k_{1:3}}, k_{4:n+2}) 
%\end{equ}
%and its symmetrised version is 
%\begin{equs}
%&\gena_{++} \psi (k_{1:n+2}) = \frac{6\iota }{(2\pi)^{d}}\times \\
%&\times\frac{\eps^{d-1}}{(n+2)(n+1)}\sum_{1\leq i_1<i_2<i_3\leq n+2} \indN{k_{i_1},\dots,k_{i_3}} (w\cdot \overline{k_{{i_1}:{i_3}}}) \hat \psi(\overline{k_{i_1:i_3}}, k_{\{1:n+2\}\setminus\{i_1,i_2,i_3\}})\,.
%\end{equs}
%This is the operator corresponding to the $\genap$ in the previous sections and probably satisfies similar 
%bounds. 
%\medskip
%
%We now turn to $\gena_0$ which only acts on $\psi\in\fock_n$ for $n\geq 2$. 
%The unsymmetrised version of $\gena_0$ is 
%\begin{equs}
%\gena_0 \psi (k_{1:n})&=\frac{\iota \eps^{d-1}}{(2\pi)^{d}} 3 n (n-1) \sum_{k} \indN{k,k_1,k_2} w\cdot \big(k+\overline{k_{1:2}}\big) \, \hat \psi(-k-\overline{k_{1:2}},k, k_{3:n}) \\
%&=-\frac{\iota \eps^{d-1}}{(2\pi)^{d}} 3 n (n-1)\sum_\ell\indN{\ell-\overline{k_{1:2}},k_1,k_2} w\cdot \ell \, \hat \psi(\ell, \overline{k_{1:2}}-\ell, k_{3:n})\\
%&=-\frac{\iota \eps^{d-1}}{(2\pi)^{d}} \frac32 n (n-1)w\cdot \overline{k_{1:2}} \sum_{\ell+m=\overline{k_{1:2}}}\indN{-m,k_1,k_2} \hat \psi(\ell, m, k_{3:n})\\
%&=-\frac{\iota \eps^{d-1}}{(2\pi)^{d}} \frac32 n (n-1)w\cdot \overline{k_{1:2}} \sum_{\ell+m=\overline{k_{1:2}}}\indN{\ell,m} \hat \psi(\ell, m, k_{3:n})
%\end{equs}
%\giuseppeText{Oh look, it looks very much like the $\genam$ we had before...}
%\begin{remark}
%This term is quite annoying. Let us forget about the dependence on $n$ and let us study the $(-\gensy)^{-1/2}\fock$- norm. 
%We have 
%\begin{equs}
%\|&(-\gensy)^{-1/2}\gena_0\psi\|^2\\
%&\lesssim \sum_{k_1:n}\frac{1}{|k_1|^2+|k_2|^2} \left|\eps^{d-1}\sum_{k} \indN{k,k_1,k_2} w\cdot \big(k+\overline{k_{1:2}}\big) \, \hat \psi(k+\overline{k_{1:2}},k, k_{3:n}) \right|^2\\
%&=\sum_{j, k_{3:n}}\left|\eps^{\frac{d}{2}}\sum_{k} \indN{k,j} w\cdot \big(k+j\big) \, \hat \psi(k+j,k, k_{3:n}) \right|^2\left(\eps^{d-2}\sum_{\overline{k_{1:2}}=j}\frac{\indN{k_1,k_2}}{|k_1|^2+|k_2|^2}\right)\\
%&\lesssim \sum_{j, k_{3:n}}\left|\eps^{\frac{d}{2}}\sum_{k} \indN{k,j} w\cdot \big(k+j\big) \, \hat \psi(k+j,k, k_{3:n}) \right|^2\\
%&\lesssim \sum_{j, k_{3:n}}\left(\sum_k \big(w\cdot \big(k+j\big)\big)^2 \, |\hat \psi(k+j,k, k_{3:n})|^2\right)\times\left(\eps^d\sum_k\indN{k,j}\right)\\
%&\lesssim  \sum_{j,k, k_{3:n}} \big(w\cdot \big(k+j\big)\big)^2 \, |\hat \psi(k+j,k, k_{3:n})|^2=\frac{1}{n} \|(-\gensy)^{1/2}\psi\|^2
%\end{equs}
%In particular, this term does NOT go to $0$ and would in principle need to be considered in the estimates. 
%\end{remark}
%
%\subsubsection{$F$ Hermite polynomial of degree $4$}
%
%Before turning to the general case, we consider the case of $F$ being a Hermite polynomial of degree $4$. 
%The goal is to see the relation with that of degree $2$ we have worked with so far. 
%
%In this case,~\eqref{e:Fexpeta} becomes
%\begin{equ}
%F(\eps^{d/2}\Pi_{1/\eps}\eta^\eps(x))= \eps^{2d}\sum_{k_{1:4}\colon |k_i|\leq \eps^{-1}}\frac{e_{\overline{k_{1:4}}}(x)}{(2\pi)^\frac{3d}{2}}\Wick{\prod_{i=1}^4\hat\eta(k_i)}\,,
%\end{equ}
%where, once again, $\overline{k_{1:j}}\eqdef \sum_{i=1}^jk_i$.  
%The non-linearity then reads
%\begin{equs}
%\NN_k(\eta)&=\frac{\iota}{(2\pi)^\frac{3d}{2}}\eps^{\frac{3d}{2}-1}I_4\left(\sum_{\overline{k_{1:4}}=k} \indN{k_1,\dots,k_4} (w\cdot \overline{k_{1:4}})\bigotimes_{i=1}^4e_{-k_i}\right)
%\end{equs}
%Adopting the same conventions as in the previous section, we see that 
%\begin{equs}
%\,&\gena I_n(K)=\frac{\iota \eps^{\frac{3d}{2}-1}}{(2\pi)^\frac{3d}{2}} n I_{n-1}(\otimes^{n-1} h)I_4\left(\sum_{\overline{k_{1:4}}=k} \indN{k_1,\dots,k_4} (w\cdot \overline{k_{1:4}})\bigotimes_{i=1}^4e_{-k_i}\right)\\
%&=\frac{\iota \eps^{\frac{3d}{2}-1}}{(2\pi)^\frac{3d}{2}} n\sum_{p=0}^{ 4} p! \binom{n-1}{p}\binom{4}{p} \1_{p\leq n-1}\times\\
%&\times I_{n+3-2p}\bigg(\sum_{k_{1:4}} \indN{k_1,\dots,k_4} w\cdot \overline{k_{1:4}} \hat h(-\overline{k_{1:4}}) \left(\prod_{i=0}^p \hat h(k_i)\right) \bigotimes_{i=p+1}^{4}e_{-k_{i}}\otimes \bigotimes^{n-1-p} h\bigg)\\
%&=: \1_{0\leq n-1} I_{n+3}\left(\gena_{+,3}K\right) +\1_{1\leq n-1} I_{n+1}\left(\gena_{+,1}K\right)\\
%&\qquad\qquad\qquad+\1_{2\leq n-1} I_{n-1}\left(\gena_{-,1}K\right) + \1_{3\leq n-1} I_{n-2}\left(\gena_{-,3}K\right)
%%&\qquad\qquad+\sum_{p=\lfloor\frac{j-1}{2}\rfloor+1}^{j-1}\1_{p\leq n-1} I_{n-1+j-2p}\left(\gena_{-, j-1-2p}K\Big)
%\end{equs}
%where once again, we omitted the summand corresponding to $p=4$ as it is $0$. Indeed, 
%\begin{equs}
%I_{n-5}&\left(\sum_{k_{1:4}} \indN{k_1,\dots,k_4} w\cdot \overline{k_{1:4}} \hat h(-\overline{k_{1:4}}) \left(\prod_{i=1}^4 \hat h(k_i)\right)  \otimes^{n-5} h\right)\\
%&=\left[\sum_{k_{1:4}} \indN{k_1,\dots,k_4} w\cdot \overline{k_{1:4}} \hat h(-\overline{k_{1:4}}) \left(\prod_{i=1}^4 \hat h(k_i)\right) \right] I_{n-5}\left(\otimes^{n-5}h\right)\\
%&=-\langle w\cdot \nabla h, h^4\rangle I_{n-5}\left(\otimes^{n-5}h\right)=\frac{1}{5}\langle w\cdot \nabla (h^5), 1\rangle I_{n-5}\left(\otimes^{n-5}h\right)=0\,.
%\end{equs}
%As above, note that since $\gena$ is antisymmetric, we must have 
%\begin{equ}
%(\gena_{+,3})^\ast=-\gena_{-,3}\qquad\text{and}\qquad (\gena_{+,1})^\ast=-\gena_{-,1}\,.
%\end{equ}
%Hence, it is sufficient to study $\gena_{+,3}$ and $\gena_{+,1}$. 
%The former is the operator we are essentially used to. As a matter of fact its unsymmetrised version 
%applied to $\psi\in\fock_n$ reads
%\begin{equ}
%\gena_{+,3} \psi (k_{1:n+3}) = \frac{\iota \eps^{\frac{3d}{2}-1}}{(2\pi)^\frac{3d}{2}} n \indN{k_1,\dots,k_4} (w\cdot \overline{k_{1:4}}) \hat \psi(\overline{k_{1:4}}, k_{5:n+3}) 
%\end{equ}
%and its symmetrised version is 
%\begin{equs}
%&\gena_{+,3} \psi (k_{1:n+2}) = \frac{\iota }{(2\pi)^\frac{3d}{2}}\frac{4!}{(n+3)(n+2)(n+1)}\times \\
%&\times\eps^{\frac{3d}{2}-1}\sum_{1\leq i_1<i_2<i_3<i_4\leq n+2} \indN{k_{i_1},\dots,k_{i_4}} (w\cdot \overline{k_{{i_1}:{i_4}}}) \hat \psi(\overline{k_{i_1:i_4}}, k_{\{1:n+2\}\setminus\{i_1,\dots,i_4\}})\,.
%\end{equs}
%
%Let us turn to $\gena_{+,1}$. its unsymmetrised version acting on $\psi\in\fock_n$, $n\geq 2$ is 
%\begin{equs}
%\gena_{+,1}\psi(k_{1:n+1})&=\frac{\iota \eps^{\frac{3d}{2}-1}}{(2\pi)^\frac{3d}{2}} 4n(n-1)\sum_{k} \indN{k,k_{1:3}} w\cdot \big(k+\overline{k_{1:3}}\big) \, \hat \psi(k+\overline{k_{1:3}},k, k_{4:n}) \\
%&=-\frac{\iota \eps^{\frac{3d}{2}-1}}{(2\pi)^\frac{3d}{2}} 2n(n-1)w\cdot\overline{k_{1:3}}\sum_{\ell+m=\overline{k_{1:3}}} \indN{\ell,m}  \, \hat \psi(\ell,m, k_{4:n})
%\end{equs} 
%so in particular, it has the same structure as the operator $\cA_0$ of the previous section and arguably it will admit 
%the same bounds. 
%
%
%\subsubsection{$F$ general Hermite polynomial}
%
%We now try to generalise the previous by taking $F$ to be a Hermite polynomial of 
%order $j$, i.e. 
%\begin{equ}[e:FHermite]
%F(x)=H_j(x)
%\end{equ}
%so that~\eqref{e:Fexpeta} becomes
%\begin{equ}
%F(\eps^{d/2}\Pi_{1/\eps}\eta^\eps(x))= \eps^{\frac{d}{2}j}\sum_{k_{1:j}\colon |k_i|\leq \eps^{-1}}\frac{e_{\overline{k_{1:j}}}(x)}{(2\pi)^{\frac{d}{2}(j-1)}}\Wick{\prod_{i=1}^j\hat\eta(k_i)}\,,
%\end{equ}
%where, once again, $\overline{k_{1:j}}\eqdef \sum_{i=1}^jk_i$ and 
%the non-linearity then reads
%\begin{equs}
%\NN_k(\eta)&=\frac{\iota}{(2\pi)^{\frac{d}{2}(j-1)}}\eps^{\frac{d}{2}(j-1)-1}I_j\left(\sum_{\overline{k_{1:j}}=k} \indN{k_1,\dots,k_j} w\cdot \overline{k_{1:j}}\bigotimes_{i=1}^je_{-k_i}\right)
%\end{equs}
%\medskip
%
%Again, by polarisation, it suffices to study the action of $\gena$ on $I_n(K)$ for $K=\otimes^n h$, and $h\in L^2(\T^d)$. 
%\begin{equs}
%\,&\gena I_n(K)=\sum_k \NN_k(\eta) D_{-k} I_n(K)=n I_{n-1}(\otimes^{n-1} h)\sum_k \NN_k(\eta) \hat h(-k)\\
%=&\frac{\iota \eps^{\frac{d}{2}(j-1)-1}}{(2\pi)^{\frac{d}{2}(j-1)}} n I_{n-1}(\otimes^{n-1} h)I_j\left(\sum_{k_{1:j}} \indN{k_1,\dots,k_j} w\cdot \overline{k_{1:j}} \hat h(-\overline{k_{1:j}}) \bigotimes_{i=1}^je_{-k_i}\right)\\
%=&\frac{\iota \eps^{\frac{d}{2}(j-1)-1}}{(2\pi)^{\frac{d}{2}(j-1)}} n\sum_{p=0}^{ j} p! \binom{n-1}{p}\binom{j}{p} \1_{p\leq n-1}\times\\
%&\times I_{n-1+j-2p}\left(\sum_{k_{1:j}} \indN{k_1,\dots,k_j} w\cdot \overline{k_{1:j}} \hat h(-\overline{k_{1:j}}) \left(\prod_{i=0}^p \hat h(k_i)\right) \bigotimes_{i=p+1}^j e_{-k_{i}}\otimes \bigotimes^{n-1-p} h\right)\\
%=&: \sum_{p=0}^{\lfloor\frac{j-1}{2}\rfloor-1} \1_{p\leq n-1} I_{n-1+j-2p}\left(\gena_{+, j-1-2p}K\right) +\1_{\lfloor\frac{j-1}{2}\rfloor=\frac{j-1}{2}, \frac{j-1}{2}\leq n-1} I_{n}\left(\gena_{0}K\right) \\
%&\qquad\qquad+\sum_{p=\lfloor\frac{j-1}{2}\rfloor+1}^{j-1}\1_{p\leq n-1} I_{n-1+j-2p}\left(\gena_{-,2p- j+1}K\right)
%\end{equs}
%where we adopted the convention that $\hat h(k_0)\eqdef 1$ and $\gena_{\cdot, j-1-2p} K$ is the symmetrised 
%kernels of $I_{n-1+j-2p}$. 
%Note that for $p=j$, we have 
%\begin{equs}
%I_{n-1-j}&\left(\sum_{k_{1:j}} \indN{k_1,\dots,k_j} w\cdot \overline{k_{1:j}} \hat h(-\overline{k_{1:j}}) \left(\prod_{i=0}^j \hat h(k_i)\right) \otimes^{n-1-j} h\right)\\
%&=\left[\sum_{k_{1:j}} \indN{k_1,\dots,k_j} w\cdot \overline{k_{1:j}} \hat h(-\overline{k_{1:j}}) \left(\prod_{i=0}^j \hat h(k_i)\right)\right]I_{n-1-j}\left(\otimes^{n-1-j} h\right)\\
%&= \NN_h(h) I_{n-1-j}\left(\otimes^{n-1-j} h\right)=0
%\end{equs}
%where the latter is a consequence of the fact that $\NN_h(h)=0$, so that the $j$-th summand does not contribute.
%\medskip
%
%We now spell out the operators $\gena_{\cdot,\cdot} K$. 
%Since $\gena$ is antisymmetric, for every $a\in\{1,\dots,j-1\colon i\text{ even}\}$ if $j$ odd, and 
%$a\in\{1,\dots,j-1\colon i\text{ odd}\}$ if $j$ even,
%we must have 
%\begin{equ}
%\gena_{+, a}=-(\gena_{-,a})^\ast\qquad\text{and}\qquad \gena_{0}=-(\gena_{0})^\ast\,,
%\end{equ}
%so that it suffices to consider the $\gena_{+,\cdot}$ and $\gena_0$. 
%Let $a$ as above, $p=(j-1-a)/2$ and $\psi\in\fock_n$, then 
%\begin{equs}
%&\gena_{+, a}\psi(k_{1:n+a})\\
%&=\frac{\iota \eps^{\frac{d}{2}(j-1)-1}}{(2\pi)^{\frac{d}{2}(j-1)}} n p! \binom{n-1}{p}\binom{j}{p}\times\\
%&\quad\times\sum_{\ell_{1:p}} \indN{k_1,\dots,k_j} w\cdot \big(\overline{\ell_{1:p}}+\overline{k_{1:j-p}}\big) \hat \psi(-\overline{\ell_{1:p}}-\overline{k_{1:j-p}}, \ell_{1:p}, k_{j-p:n+a})\\
%&=\frac{\iota \eps^{\frac{d}{2}(j-1)-1}}{(2\pi)^{\frac{d}{2}(j-1)}} n (p-1)! \binom{n-1}{p}\binom{j}{p}\times\\
%&\quad\times w\cdot \overline{k_{1:j-p}}\sum_{\overline{\ell_{1:p+1}}=\overline{k_{1:j-p}}} \indN{\ell_1,\dots,\ell_{p+1}}  \hat \psi(\ell_{1:p+1}, k_{j-p:n+a})\\
%\end{equs}


\begin{appendix}

%\section{A heuristic derivation of the Ansatz}\label{a:Heu}

\section{Some technical results}\label{a:Approx}

In this appendix, we show the technical part of the Replacement Lemma~\ref{l:Replacement}. 

\begin{proposition}\label{p:Approx}
 For $\eps>0$, $n\in\N$ and $k_{1:n}\in\Z^{2n}_0$, define  $P^\eps$ as in~\eqref{e:P1}, i.e. 
\begin{equ}[e:P]
P^\eps(k_{1:n})\eqdef\frac{\lambda_\eps^2}{\pi^2} \sum_{\ell + m=k_1} \frac{\indN{\ell,m}}{\tilde\Gamma + \tilde\Gamma^wG(\Ll(\tilde\Gamma))}
\end{equ}
for
\begin{equ}
\tilde\Gamma\eqdef\tfrac12(|\ell|^2+|m|^2+|k_{2:n}|^2)\,,\qquad\tilde\Gamma^w\eqdef \tfrac12((\fw\cdot\ell)^2+(\fw\cdot m)^2+(\fw\cdot k)^2_{2:n})
\end{equ}
and $G$ is as in \eqref{eq:defG}.
Then, there exists a constant $C>0$ independent of $\eps,n$ such that 
\begin{equ}[e:Approx]
\sup_{n\in \N,k_{1:n}\in\Z^{2n}_0}\Big|P^\eps(k_{1:n})-G\big(\Ll(\tfrac12|k_{1:n}|^2)\big)   \Big| \leq C \lambda_\eps^2\,.
\end{equ}
\end{proposition}

\begin{proof}
  Note first of all that we can assume that $|k_1|\le 1/(2\eps)$. Indeed,
  in the complementary case both $P^\eps(k_{1:n})$ and 
  $G(\Ll(\frac12|k_{1:n}|^2))$ would be $o(\lambda_\eps^2)$ (uniformly in
 $n$ and  $k_{1:n}$). For $G(\Ll(\frac12|k_{1:n}|^2))$, this easily follows from the definition of $G$ and $\Ll$. 
 As for $P^\eps(k_{1:n})$, if  $|k_1|\ge 1/(2\eps)$, then either $|\ell|$ or $|m|$ are larger than $1/(4\eps)$. 
 Assume, for instance, that the former is the case. Then,
  \begin{equ}
    \label{eq:isthecase}
P^\eps(k_{1:n})\lesssim \lambda_\eps^2\sum_{\frac1{4\eps}\le |\ell|\le\frac1\eps}\frac1{|\ell|^2}\lesssim \lambda_\eps^2.
  \end{equ}
  For  $|k_1|\le 1/(2\eps)$,  we  rewrite \eqref{e:P} as
  \begin{equ}
    \label{eq:PP}
    P^\eps(k_{1:n})\eqdef\frac{\lambda_\eps^2}{\pi^2} \sum_{\substack{\ell + m=k_1\\ |\ell|,|m|\le \eps^{-1}}} \frac{1}{\tilde\Gamma + \tilde\Gamma^wG(\Ll(\tilde\Gamma))}
  \end{equ}
Note  that, when $\ell+m=k_1$,
\begin{equs}
  \tilde\Gamma&=\Gamma-\ell\cdot k_1,
  \qquad\qquad\qquad \Gamma\eqdef
  |\ell|^2+\frac12|k_{1:n}|^2,\\ 
  \tilde \Gamma^w&=\Gamma^w-(\fw\cdot k_1)(\fw\cdot \ell), \qquad\Gamma^w\eqdef
  (\fw\cdot \ell)^2+\frac12(\fw\cdot k)^2_{1:n}.
\end{equs}
For future reference, we point out that $\Gamma$ and $\tilde\Gamma$ are comparable in that, 
by the triangular inequality, we have
\begin{equ}[e:lm]
\Gamma\lesssim \tilde\Gamma\lesssim \Gamma\,
\end{equ}
uniformly in $\ell,m=k_1-\ell$ and  $k_{1:n}\in\Z^{2n}_0$.
\medskip

The proof of Proposition~\ref{p:Approx} follows from a sequence of technical steps in which the summand
in the definition of $P^\eps$ is progressively simplified and 
the Riemann sum is replaced by an integral. More precisely, we will first replace $P_0^\eps\eqdef P^\eps$ 
with $P^\eps_1$ and then $P^\eps_i$ with $P^\eps_{i+1}$, $i=1,\dots, 4$, in such a way that 
for all $i=0,\dots,5$, 
\begin{equs}
  \label{eq:pallosissimi}
  \sup_{n\in \N,k_{1:n}\in\Z^{2n}_0}\Big|P^\eps_i(k_{1:n})-P_{i+1}^\eps(k_{1:n})  \Big| \lesssim \lambda_\eps^2\,. 
\end{equs}
At last, we will compute $P^\eps_5$ and show that it equals $G(\Ll)$. 
\medskip

\noindent {\bf Step 1} First of all, we define $P_1^\eps(k_{1:n})$ as $P^{\eps}(k_{1:n})$ in \eqref{eq:PP}, 
except that $\tilde\Gamma$ and $\tilde\Gamma^w$ are replaced by $\Gamma$ and $\Gamma^w$, respectively. 
The absolute value of the difference is upper bounded
by a constant times
\begin{equ}
   \lambda_\eps^2\sum_{\substack{\ell+m=k_1\\|\ell|,|m|\le 1/\eps}}\frac{A+B}
{(|\ell|^2+|k_1|^2)^2}
\end{equ}
where in the denominator we used $\Gamma\wedge \tilde \Gamma\gtrsim |\ell|^2+|k_1|^2$ and 
$\Gamma^w\wedge\tilde \Gamma^w\geq 0$, while in the numerator $A$ and $B$ are defined as
\begin{equs}
A&\eqdef  |\ell\cdot k_1+(\fw\cdot \ell)(\fw\cdot k_1)G(\Ll(\tilde\Gamma))|,\\
B&\eqdef |G(\Ll(\tilde\Gamma))-G(\Ll(\Gamma)) |\Big(|\ell|^2+\frac12|k_{1:n}|^2\Big)\,.
\end{equs}
Now, $A\lesssim |\ell||k_1|$ so that
\begin{equs}
\lambda_\eps^2\sum_{\substack{\ell+m=k_1\\|\ell|,|m|\le 1/\eps}}\frac A{(|\ell|^2+|k_1|^2)^2} \lesssim \lambda_\eps^2,
\end{equs} 
as can be immediately seen by splitting the sum according to  whether $|\ell|\le |k_1|$ or the converse.
The term containing $B$ is dealt with similarly. Indeed, since $|\Gamma-\tilde\Gamma|\le |\ell|| k_1|$,  one has 
\begin{equs}
  |G(\Ll(\tilde\Gamma))-G(\Ll(\Gamma)) |\lesssim |\ell||k_1|\sup_{y\in [\Gamma,\tilde \Gamma]}\left|\frac1{|\log\eps|}\frac{G'(\Ll(y))
    }{y(\eps y+1)}
  \right|\lesssim \frac{ |\ell||k_1|}{|\ell|^2+\frac12|k_{1:n}|^2}
\end{equs}
where we further used that the derivative of $G$ is bounded and \eqref{e:lm} holds. 
Hence, $B\lesssim |\ell||k_1|$ and one concludes as above. Inequality \eqref{eq:pallosissimi} for $i=0$ is then proven.
\medskip

\noindent {\bf Step 2}
Secondly, we define $P_2^\eps(k_{1:n})$ as $P_1^{\eps}(k_{1:n})$, except that $\Gamma^w$ is replaced by 
$(\fw\cdot \ell)^2$, that is,
\begin{equ}
  P_2^\eps(k_{1:n})\eqdef \frac{\lambda_\eps^2}{\pi^2}\sum_{\substack{\ell+m=k_1\\|\ell|,|m|\le \eps^{-1}}}\frac{1}{\Gamma+(\fw\cdot \ell)^2G(\Ll(\Gamma))}\,.
\end{equ}
The difference between $P^\eps_1$ and $P^\eps_2$ is bounded by a constant times
\begin{equs}
 \lambda_\eps^2\sum_{\substack{\ell+m=k_1\\|\ell|,|m|\le \eps^{-1}}}\frac{(\fw\cdot k)_{1:n}^2}{(|\ell|^2+|k_{1:n}|^2)^2}
\lesssim \lambda_\eps^2\sum_{\substack{\ell+m=k_1\\|\ell|,|m|\le \eps^{-1}}}\frac{|k_{1:n}|^2}{(|\ell|^2+|k_{1:n}|^2)^2}\lesssim \lambda_\eps^2\,,
\end{equs}
where we used that $G(\Ll(\Gamma))\lesssim 1$. 
\medskip

\noindent {\bf Step 3}
As a third step, we  replace $P_2^\eps(k_{1:n})$ with
\begin{equ}
  P_3^\eps(k_{1:n})\eqdef \frac{\lambda_\eps^2}{\pi^2}\sum_{|\ell|\le \eps^{-1}}\frac{1}{\Gamma+(\fw\cdot \ell)^2G(\Ll(\Gamma))}.
\end{equ}
With respect to $P^\eps_2$, we have only removed the indicator function of $|m|=|k_1-\ell|\leq \eps^{-1}$. 
Hence, the difference between $P_2^\eps$ and $P_3^\eps$ only involves the sum over those values of 
$\ell$ such that $|k_1-\ell|\ge 1/\eps$ while $|\ell|\le 1/\eps$. 
But, since $|k_1|\le 1/(2\eps)$, this corresponds to $1/(2\eps)\le |\ell|\le 1/\eps$ and
the resulting sum can be bounded as in \eqref{eq:isthecase}.
\medskip

\noindent {\bf Step 4}  Next, we turn the sum into an integral. Setting
$\alpha_\eps\eqdef \frac12|\eps k_{1:n}|^2$, we define $P_4^\eps$ as
\begin{equ}[e:P4]
  P_4^\eps(k_{1:n})\eqdef \frac{\lambda_\eps^2}{\pi^2}\int_{|x|\le 1}\frac{\dd x}{|x|^2+\alpha_\eps+(x\cdot \fw)^2G(\Ll(\eps^{-2}(|x|^2+\alpha_\eps)))}\,.
\end{equ}
To obtain~\eqref{eq:pallosissimi} for $i=3$, we argue as in~\cite[Lemma C.7]{CET}. 
Denoting by $I$ the integrand in~\eqref{e:P4} and by $Q_\ell^\eps$ the square of side-length $\eps$ centred at 
$\ell\in\Z^2_0$,
the difference between $P_3^\eps$ and $P_4^\eps$ is bounded above by
\begin{equs}
\lambda_\eps^2\sum_{1\leq |\ell|\leq \eps^{-1}}\int_{Q_\ell^\eps} |I(\eps\ell)-I(x)|\dd x&\leq \lambda_\eps^2\eps\sum_{1\leq |\ell|\leq \eps^{-1}}\int_{Q_\ell^\eps} \sup_{y\in Q^\eps_\ell} |\nabla I(y)|\dd x\\
&\leq\lambda_\eps^2\eps^3\sum_{1\leq |\ell|\leq \eps^{-1}}\sup_{y\in Q^\eps_\ell} |\nabla I(y)|\,.
\end{equs} 
Now, a simple computation yields  
\begin{equ}
\sup_{y\in Q^\eps_\ell} |\nabla I(y)|\lesssim \lambda_\eps^2\frac{1}{(|\eps\ell|^2 +\alpha_\eps)^{3/2}}
\end{equ}
so that~\eqref{eq:pallosissimi} follows at once. 
\medskip 

Before moving to the next step, note that, by passing to polar coordinates, 
i.e. setting $x=r v_\theta$ for $v_\theta=(\cos\theta,\sin\theta)$, we have 
\begin{equs}
P_4^\eps(k_{1:n})&= \frac{\lambda_\eps^2}{\pi^2}\int_0^{2\pi}\dd \theta\int_{0}^1\frac{r\dd r}{r^2+\alpha_\eps+r^2(\fw\cdot v_\theta)^2G(\Ll(\eps^{-2}(r^2+\alpha_\eps)))}\\
&=\frac{\lambda_\eps^2}{2\pi^2}\int_0^{2\pi}\dd \theta\int_{0}^1\frac{\dd r}{r+\alpha_\eps+r(\fw\cdot v_\theta)^2G(\Ll(\eps^{-2}(r+\alpha_\eps)))}\\
&=\frac{\lambda_\eps^2}{2\pi^2}\int_0^{2\pi}\dd \theta\int_{0}^1\frac{\dd r}{r+\alpha_\eps+r|\fw|^2 \cos^2(\theta-\theta_w)G(\Ll(\eps^{-2}(r+\alpha_\eps)))}\\
&=\frac{\lambda_\eps^2}{\pi^2}\int_{0}^1\dd r\int_0^{\pi}\frac{\dd \theta}{r+\alpha_\eps+r|\fw|^2 \cos^2(\theta)G(\Ll(\eps^{-2}(r+\alpha_\eps)))}
\end{equs}
where we set $w=|w|v_{\theta_w}$ and used periodicity of the integrand in $\theta$. 
\medskip

\noindent{\bf Step 5} In this step, we insert two elements at the denominator of the integrand of $P_4^\eps$, 
that will allow us to make a simple change of variables. 
More precisely, we replace $P_4^\eps$ with $P_5^\eps$, the latter being defined as  
\begin{equ}
P_5^\eps(k_{1:n})\eqdef \frac{\lambda_\eps^2}{\pi^2}\int_0^{\pi}\int_{0}^1\frac{\dd \theta\dd r}{(r+\alpha_\eps)(r+\alpha_\eps+1)[1+|\fw|^2 \cos^2(\theta)G(\Ll(\eps^{-2}(r+\alpha_\eps)))]}\,.
\end{equ}
Now,~\eqref{eq:pallosissimi} for $i=4$ 
can be easily seen to hold by arguing as in the proof of~\cite[Eq. (A.9)]{CETWeak}.
\medskip

At this point, it remains to analyse the last integral. With the change of variables
\begin{equs}
y= \Ll(\eps^{-2}(r+\alpha_\eps))&=\lambda_\eps^2\log\Big(1+\frac1{r+\alpha_\eps}\Big), \\
\frac{\dd r}{(r+\alpha_\eps)(1+r+\alpha_\eps)}&=- \lambda_\eps^{-2} \dd y,
\end{equs}
one finds
\begin{equs}
P_5^\eps(k_{1:n})=\frac{1}{\pi^2}\int_0^{\pi}\dd \theta \int_{\Ll(\eps^{-2}+\frac12|k_{1:n}|^2)}^{\Ll(\frac12|k_{1:n}|^2)}\frac{\dd y}{1+|\fw|^2 \cos^2(\theta)G(y)}
\end{equs}
It is then immediate to verify that, again at the price of an error of order $\lambda_\eps^2$, we can replace the lower 
integration index by $0$, and obtain
\begin{equs}[e:G?]
\frac{1}{\pi^2}\int_{0}^{\Ll(\frac12|k_{1:n}|^2)}&\dd y\int_0^{\pi} \frac{\dd \theta}{1+|\fw|^2 \cos^2(\theta)G(y)}\\
&=\frac{1}{\pi}\int_{0}^{\Ll(\frac12|k_{1:n}|^2)}\frac{\dd y}{\sqrt{1+|\fw|^2G(y)}}=G(\Ll(\frac12|k_{1:n}|^2))
\end{equs}
where the last step is a consequence of our choice of $G$ (actually, {\it we have chosen $G$ in such a way that 
it holds}). Hence,~\eqref{e:Approx} follows at once and 
the proof of the proposition is concluded. 
\end{proof}


%%   In order to lighten the exposition, let us introduce some notations. 
%% Set
%% \begin{equs}[e:conv]
%% \beta_\eps\eqdef |\eps k_{2:n}|^2\,,&\qquad \alpha_\eps\eqdef \tfrac12 |\eps k_{1:n}|^2\\
%% \beta_\eps^w\eqdef (w\cdot (\eps k))_{2:n}|^2\,,&\qquad \alpha_\eps^w\eqdef \tfrac12 ((w\cdot(\eps k))_{1:n}|^2
%% \end{equs}
%% and, for $x\in\R^2$ such that $|x|\leq 1$, 
%% \begin{equs}[e:Gammas]
%% \tilde\Gamma_\eps(x)\eqdef \tfrac12 (|x|^2+|\eps k_1-x|^2+\beta_\eps)\,,&\qquad \Gamma_\eps(x)\eqdef |x|^2 +\alpha_\eps\,,\\
%% \tilde\Gamma_\eps^w(x)\eqdef \tfrac12 ((w\cdot x)^2+(w\cdot(\eps k_1-x)^2+\beta_\eps^w)\,,&\qquad \Gamma^w_\eps(x)\eqdef (w\cdot x)^2 +\alpha_\eps^w\,.
%% \end{equs}
%% For future reference, we point out that $\Gamma_\eps$ and $\tilde\Gamma_\eps$ are comparable in that, 
%% by the triangular inequality, we have
%% \begin{equ}[e:lm]
%% \Gamma_\eps(x)\lesssim \tilde\Gamma_\eps(x)\lesssim \Gamma_\eps(x)\,
%% \end{equ}
%% for all $x\in\R^2$ and $k_{1:n}\in\Z^{2n}_0$.

\end{appendix}











%
%
%
%\newpage
%
%Let us single out some special situations for the operators we introduced above
%\begin{itemize}
%\item[$p=0$] this is the well-known 
%\begin{equs}
%\cF(\gena_{+,j-1} K)(k_{1:n+j-1})=\frac{\iota \eps^{\frac{d}{2}(j-1)-1}}{(2\pi)^{\frac{d}{2}(j-1)}} n w\cdot \overline{k_{1:j}}\hat h(-\overline{k_{1:j}}) \prod_{i\neq 1,\dots,j}\hat h(k_i)
%\end{equs}
%\item[$p=j-1$] this is the dual of the previous.
%\item[$p=j$] this is $0$. Indeed, in this case we have 
%\begin{equs}
%I_{n-1-j}&\left(\sum_{k_{1:j}} \indN{k_1,\dots,k_j} w\cdot \overline{k_{1:j}} \hat h(-\overline{k_{1:j}}) \left(\prod_{i=0}^j \hat h(k_i)\right) \otimes^{n-1-j} h\right)\\
%&=\left[\sum_{k_{1:j}} \indN{k_1,\dots,k_j} w\cdot \overline{k_{1:j}} \hat h(-\overline{k_{1:j}}) \left(\prod_{i=0}^j \hat h(k_i)\right)\right]I_{n-1-j}\left(\otimes^{n-1-j} h\right)\\
%&= \NN_h(h) I_{n-1-j}\left(\otimes^{n-1-j} h\right)=0
%\end{equs}
%where the latter is a consequence of the fact that $\NN_h(h)=0$. 
%%whose symmetrised version, acting on a generic symmetric function $\psi\in\fock_n$ is 
%
%\end{itemize}
%\newpage

\section*{Acknowledgments} 
The authors would like to thank L. Haunschmid-Sibitz for many
discussions in an early stage of this project, C.~Landim for
explaining some of the ideas in~\cite{Komorowski2012}, and R.~Bauerschmidt for useful discussions about RG.  A
special thank goes to L.~Gr\"afner, who presented us the strategy
in~\cite{GrPer} which allowed us to perform the steps
in~\eqref{e:defDn} and lead to the proof of the main theorem in the
super-critical case.  G.~C. gratefully acknowledges financial support
via the UKRI FL fellowship ``Large-scale universal behaviour of Random
Interfaces and Stochastic Operators'' MR/W008246/1.  F.~T.  was
supported by the Austrian Science Fund (FWF): P35428-N.

\bibliography{ASEPSPDE}
\bibliographystyle{Martin}


\end{document}
